% !TEX program = xelatex
\documentclass[11pt,openany,fontset=fandol]{ctexbook}
\usepackage[paperwidth=176mm,paperheight=250mm,inner=21mm,outer=18mm,top=20mm,bottom=20mm,headheight=16pt,headsep=6mm,footskip=10mm]{geometry}
\usepackage{fontspec,amsmath,amssymb,amsthm,bm,mathtools,booktabs,tabularx,array,enumitem}
\usepackage{tikz,fancyhdr,listings,etoolbox,float}
\usepackage{hyperref}
\usetikzlibrary{arrows.meta,calc,decorations.pathmorphing,angles,quotes}
\setmonofont{lmmono10-regular.otf}[Scale=0.85]
\hypersetup{hidelinks,pdftitle={大学物理中的数学方法},pdfauthor={}}
\setlength{\parindent}{2em}\setlength{\parskip}{3pt}
\setlength{\emergencystretch}{2em}\linespread{1.18}\raggedbottom
\setcounter{tocdepth}{1}
\setlist[enumerate]{leftmargin=2em,itemsep=5pt,topsep=5pt}
\setlist[itemize]{leftmargin=2em,itemsep=3pt,topsep=4pt}
\ctexset{chapter={name={第,讲},number=\chinese{chapter},format=\LARGE\bfseries\centering,beforeskip=16pt,afterskip=22pt},section={format=\Large\bfseries},subsection={format=\large\bfseries}}
\pagestyle{fancy}\fancyhf{}\fancyhead[LE,RO]{\small\thepage}
\fancyhead[LO]{\small 大学物理中的数学方法}\fancyhead[RE]{\small\leftmark}
\renewcommand{\headrulewidth}{0pt}
\fancypagestyle{plain}{\fancyhf{}\fancyfoot[C]{\small\thepage}\renewcommand{\headrulewidth}{0pt}}
\theoremstyle{plain}\newtheorem{theorem}{定理}[chapter]\newtheorem{proposition}[theorem]{命题}
\theoremstyle{definition}\newtheorem{definition}{定义}[chapter]\newtheorem{example}{例}[chapter]
\theoremstyle{remark}\newtheorem*{remark}{说明}
\newcommand{\SolutionStart}{\par\nobreak\medskip\nobreak\noindent\textbf{解}\quad\ignorespaces}
\newenvironment{solution}{\SolutionStart}{\par\medskip}
\newcommand{\dd}{\mathop{}\!\mathrm d}\newcommand{\ii}{\mathrm i}\newcommand{\ee}{\mathrm e}
\newcommand{\vct}[1]{\bm{#1}}\newcommand{\Rea}{\operatorname{Re}}\newcommand{\Ima}{\operatorname{Im}}
\newcommand{\Arg}{\operatorname{Arg}}\newcommand{\rank}{\operatorname{rank}}
\newcommand{\tr}{\operatorname{tr}}\newcommand{\diag}{\operatorname{diag}}
\newcommand{\atan}{\operatorname{atan2}}\newcommand{\Res}{\operatorname{Res}}
\lstset{language=Python,basicstyle=\footnotesize\ttfamily,keywordstyle=\ttfamily,commentstyle=\ttfamily,stringstyle=\ttfamily,showstringspaces=false,breaklines=true,columns=fullflexible,frame=none,aboveskip=8pt,belowskip=8pt}
\BeforeBeginEnvironment{tabularx}{\par\noindent}
\tikzset{exspring/.style={decorate,decoration={coil,segment length=4pt,amplitude=2.5pt}},bookaxis/.style={->,>=Stealth,line width=.4pt},bookcurve/.style={line width=.7pt},bookguide/.style={dashed,line width=.35pt},every node/.style={font=\small}}
\providecommand{\ExpansionDiagram}[3]{\begin{figure}[H]\centering\begin{tikzpicture}#3\end{tikzpicture}\caption{#2}\label{#1}\end{figure}}
\providecommand{\kb}{k_{\mathrm B}}

% 分册与十编合订本共用的定位页码；局部章、图、题号保持原有对应。
\newcommand{\CoursePartNumber}{1}
\geometry{inner=19.5mm,outer=19.5mm,headheight=16pt}
\AtBeginDocument{%
  \fancyhead{}\fancyhead[C]{\small 大学物理中的数学方法}%
  \fancyfoot[C]{\small\CoursePartNumber-\thepage}%
  \fancypagestyle{plain}{\fancyhf{}\fancyfoot[C]{\small\CoursePartNumber-\thepage}\renewcommand{\headrulewidth}{0pt}}%
}

\begin{document}
\frontmatter
\begin{titlepage}
\centering\vspace*{30mm}
{\Large 物理竞赛二轮讲义\quad 第一编\par}\vspace{8mm}

{\Huge\bfseries 大学物理中的数学方法\par}
\vspace{12mm}{\Large 物理竞赛二轮讲义\par}
\vspace{10mm}{\large 复数、几何、线性代数与AI计算实验\par}
\vfill{\large 2026\par}
\end{titlepage}

\chapter*{前言}\addcontentsline{toc}{chapter}{前言}
计算做到一半仍看不清物理关系，往往需要换一种表示。相量把振动的幅值与相位放在同一个式子里；合适的坐标能消去约束；正常模把相互牵连的运动拆开。本编围绕这些用法展开，面向已经完成一轮物理学习、具备一元微积分基础的学生。

四讲分别讨论复数与复变函数、几何与运动、线性代数，以及 AI 辅助计算实验。前三讲从基本关系推导公式，在共振、曲轨运动和弱耦合问题中讨论近似的使用。第四讲用可运行的实验核对这些推导，并研究单摆、阻力抛体及测量数据。程序可以由 AI 协助实现，模型、初值和检验依据仍须写清。

例题与习题均附解答和配图。章末题按内容与难度标注，书末另列综合习题。可选题标为“选做”，拓展内容标为“选读”；章末题的简短提示收入附录。

本编采用 SI 单位，三角函数的自变量用弧度。相量统一采用 $\ee^{\ii\omega t}$，幅值为峰值；有效值另作说明。向量默认写作列，球坐标极角从正 $z$ 轴量起。配套实验页可独立打开；书中的模拟数据和理论参考值均明确标注。

\tableofcontents
\let\cleardoublepage\clearpage
\mainmatter
\chapter{复数、相位与复变函数}
两个不同相位的振动叠加后，最大振幅怎样判断？受迫振子的位移峰和功率峰为何不同？本章先用复数组织相位，再研究这些问题。

\section{从二维几何到复数运算}
\subsection{为什么需要一个新的数系}
平面中的长度与方向可用两个实数记录。把这两个数写成复数后，还能定义一种乘法，使乘法对应平面的伸缩与旋转。相位的合成正适合用这种表示。

定义 $\ii^2=-1$，写 $z=x+\ii y$，其中 $x,y$ 为实数。两个复数相等当且仅当实部、虚部分别相等。因此复数可视为有序数对 $(x,y)$，要求实数的运算规则保持成立，再用 $\ii^2=-1$ 化简，就得到
\begin{align}
(x+\ii y)+(a+\ii b)&=(x+a)+\ii(y+b),\\
(x+\ii y)(a+\ii b)&=(xa-yb)+\ii(xb+ya).
\end{align}
第一式对应向量加法。第二式则同时混合了两个坐标，不能理解成对应坐标各自相乘。例如 $(1+\ii)^2=2\ii$，不是 $1+\ii$。

复数 $\bar z=x-\ii y$ 称为 $z$ 的共轭。几何上，共轭是关于实轴的反射。它满足 $\overline{z+w}=\bar z+\bar w$ 和 $\overline{zw}=\bar z\bar w$。复数的模定义为
\begin{equation}
|z|=\sqrt{x^2+y^2},\qquad z\bar z=|z|^2.
\end{equation}
当 $w\ne0$ 时，乘以分母的共轭便得到除法：
\begin{equation}
\frac zw=\frac{z\bar w}{|w|^2}.
\end{equation}
这种写法既能完成代数计算，也能把倒数的几何意义留待极形式解释。

\begin{example}
计算 $(2+\ii)/(1-2\ii)$。
\end{example}
% 题设配图；与紧随其后的解答分别编号。
\ExpansionDiagram{fig:audit-m-q-001}{被除数与除数在同一复平面中的方向}{
\draw[->](-.4,0)--(3.5,0)node[right]{$ \Rea z $};\draw[->](0,-2.7)--(0,2.3)node[above]{$ \Ima z $};\draw[->,thick](0,0)--(2,1)node[above right]{$ 2+\ii $};\draw[->,thick](0,0)--(1,-2)node[above right]{$ 1-2\ii $};
}

\begin{solution}
题设图\ref{fig:audit-m-q-001}的两箭头分别表示分子与分母。先乘分母的共轭求值，再用解图\ref{fig:eg-1-1}核对商的模与辐角。

分母含虚部，先乘它的共轭，使分母成为实数：
\[
\frac{2+\ii}{1-2\ii}=
\frac{(2+\ii)(1+2\ii)}{(1-2\ii)(1+2\ii)}.
\]
分母为 $1-(2\ii)^2=5$。分子逐项展开为
\[
2+4\ii+\ii+2\ii^2=2+5\ii-2=5\ii,
\]
所以原式等于 $\ii$。由于 $\ii^2=-1$，分子的实部相消。将结果乘回原分母，$\ii(1-2\ii)=\ii+2$，确实恢复原分子。

\begin{figure}[H]\centering
\begin{tikzpicture}[scale=.95]
\draw[bookaxis](-.4,0)--(3,0)node[right]{$\Rea$};\draw[bookaxis](0,-2.5)--(0,1.9)node[above]{$\Ima$};
\draw[bookaxis,thick](0,0)--(2,1)node[above right]{$2+\ii$};
\draw[bookaxis,thick](0,0)--(1,-2)node[below right]{$1-2\ii$};
\draw[bookaxis](0,0)--(0,1)node[left]{$\ii$};
\draw[->](-63.435:1.1)arc(-63.435:26.565:1.1);\node[right]at(1.2,-.65){相差 $\pi/2$};
\node[align=left,anchor=west]at(3.3,-1.15){两数的模均为 $\sqrt5$\\商的模为1，辐角为 $\pi/2$};
\end{tikzpicture}
\caption{例1.1：除法使模相除、辐角相减；分子比分母逆时针超前 $\pi/2$，所以商为 $\ii$。}\label{fig:eg-1-1}
\end{figure}

商的模为 $\sqrt5/\sqrt5=1$；分子的辐角减去分母的辐角为 $\pi/2$，与图中的 $\ii$ 方向一致。
\end{solution}

\subsection{模长与三角不等式}
由乘法展开可直接证明 $|zw|=|z||w|$。加法则满足
\begin{equation}
\bigl||z|-|w|\bigr|\le |z+w|\le |z|+|w|.
\end{equation}
右侧的不等式就是平面三角形的两边之和不小于第三边，也可用代数证明：
\[
|z+w|^2=|z|^2+|w|^2+2\Rea(z\bar w)
\le(|z|+|w|)^2.
\]
这里使用了 $\Rea\xi\le|\xi|$。左侧由 $|z|=|(z+w)-w|\le|z+w|+|w|$ 及其交换式得到。它将在振动合成中给出合幅值的上、下界：两振动既不可能合成比幅值和更大的振动，也不可能合成比幅值差的绝对值更小的振动。

\section{极形式、旋转与开方}
\subsection{辐角与象限}
非零复数可写成 $z=r(\cos\theta+\ii\sin\theta)$，其中 $r=|z|>0$。$\theta$ 称为辐角，同一复数的辐角相差 $2k\pi$。为给出单个数值，常取主辐角 $\Arg z\in(-\pi,\pi]$。零的辐角未定义，因为原点没有方向。

由 $\tan\theta=y/x$ 求辐角时必须判断象限。例如 $z=-1+\ii$ 的主辐角为 $3\pi/4$，而 $\arctan(-1)=-\pi/4$ 落在错误的象限。计算机的 $\atan(y,x)$ 依据两个坐标判断象限，适合相位计算；它仍需配合所采用的辐角范围理解跳变。

将两复数的极形式相乘，并使用三角加法公式，得到
\begin{equation}
z_1z_2=r_1r_2[\cos(\theta_1+\theta_2)+\ii\sin(\theta_1+\theta_2)].
\end{equation}
模长相乘而方向角相加。因此乘以 $\ii$ 是逆时针旋转 $\pi/2$，乘以 $-1$ 是旋转 $\pi$，乘以模为1的复数是纯旋转。除法则将模长相除、辐角相减。

\begin{figure}[htbp]
\centering
\begin{tikzpicture}[scale=.85,>=Stealth]
\draw[->] (-.2,0)--(4.2,0) node[right] {$\Rea z$};\draw[->] (0,-.2)--(0,3.7) node[above] {$\Ima z$};
\draw[->,thick] (0,0)--(25:3) node[right] {$z$};
\draw[->,thick] (0,0)--(70:3) node[above right] {$z\ee^{\ii\alpha}$};
\draw (25:1.1) arc(25:70:1.1);\node at (48:1.4) {$\alpha$};
\draw[dashed] (0,0) circle(3);
\end{tikzpicture}
\caption{复数乘以单位复数所表示的旋转}\label{fig:complex-rotation}
\end{figure}

\subsection{整次幂与方根}
反复应用乘法规则得到棣莫弗公式
\begin{equation}
(\cos\theta+\ii\sin\theta)^n=\cos n\theta+\ii\sin n\theta.
\end{equation}
当 $n$ 是正整数且 $z^n=r\ee^{\ii\theta}$、$r>0$ 时，根必须满足模长为 $r^{1/n}$，而辐角乘以 $n$ 后等于 $\theta$ 加上某个 $2\pi$ 的整数倍。因此全部根是
\begin{equation}
z_k=r^{1/n}\ee^{\ii(\theta+2k\pi)/n},\qquad k=0,\ldots,n-1.
\end{equation}
多取一个 $k=n$ 会重复 $k=0$ 的根。这些根是圆周上正 $n$ 边形的顶点。方根是一个多值关系；指定某个主值，并不意味着其他根不存在。

\begin{example}
求 $z^3=-8$ 的全部根。
\end{example}
% 题设配图；与紧随其后的解答分别编号。
\ExpansionDiagram{fig:audit-m-q-002}{用模与辐角描述三次方程的候选根}{
\draw[->](-.4,0)--(3.2,0)node[right]{$ \Rea z $};\draw[->](0,-0.5)--(0,2.1)node[above]{$ \Ima z $};\draw[dashed](0,0)circle(1.5);\draw[->,thick](0,0)--(1.15,.96)node[above right]{$z=r\ee^{\ii\theta}$};\draw[->](.5,0)arc(0:40:.5);\node[right]at(.65,.27){$\theta$};\node[right]at(3.6,1){$z^3=-8$};
}

\begin{solution}
图\ref{fig:audit-m-q-002}中的$r,\theta$是未知根的极坐标；三次乘幂分别作用于模与辐角。
三次方程的根应在复平面中全部求出。图\ref{fig:eg-1-2}按模和辐角标出它们的位置。

设 $z=r\ee^{\ii\theta}$，$r\ge0$。由模长得 $r^3=8$，故 $r=2$；由辐角得
\[
3\theta=\pi+2k\pi,\qquad
\theta_k=\frac\pi3+\frac{2k\pi}3.
\]
取 $k=0,1,2$ 可得互不相同的三个根；再取其他整数只会重复这三点。因此
\begin{align*}
z_0&=2(\cos\tfrac\pi3+\ii\sin\tfrac\pi3)=1+\sqrt3\ii,\\
z_1&=2(\cos\pi+\ii\sin\pi)=-2,\\
z_2&=2(\cos\tfrac{5\pi}3+\ii\sin\tfrac{5\pi}3)=1-\sqrt3\ii.
\end{align*}
三点在模为2的圆上等角间隔。根的和为0，积为 $-8$，分别符合 $z^3+8$ 的二次项系数及常数项给出的韦达关系。

\begin{figure}[H]\centering
\begin{tikzpicture}[scale=1.15]
\draw[bookaxis](-2.5,0)--(2.4,0)node[right]{$\Rea z$};\draw[bookaxis](0,-2.2)--(0,2.3)node[above]{$\Ima z$};
\draw[bookguide](0,0)circle(2);\foreach\ang/\name in{60/z_0,180/z_1,300/z_2}{\draw[bookguide](0,0)--(\ang:2);\fill(\ang:2)circle(1.7pt)node[shift={(\ang:.33)}]{$\name$};}
\draw[->](.65,0)arc(0:60:.65);\node at(30:1){$\pi/3$};\node[below right]at(2,0){$2$};
\end{tikzpicture}
\caption{例1.2：$z^3=-8$ 的三个根。起始辐角为 $\pi/3$，相邻根相差 $2\pi/3$，三个根的模均为2。}\label{fig:eg-1-2}
\end{figure}

沿图中圆周再转 $2\pi/3$ 便循环到下一根；三次转动回到起点，说明根没有重复也没有遗漏。
\end{solution}

\section{欧拉公式与复指数}
复指数把正弦振动求导产生的相位变化写成代数乘法。欧拉公式由级数建立，随后用于旋转和振动的计算。

实指数的幂级数在所有有限实数处收敛，同一幂级数也可定义复指数：
\begin{equation}
\ee^z=\sum_{n=0}^{\infty}\frac{z^n}{n!}.
\end{equation}
级数绝对收敛，因此可以对 $z=\ii\theta$ 分开偶次与奇次项。利用
$\ii^{2j}=(-1)^j$、$\ii^{2j+1}=\ii(-1)^j$，有
\begin{align*}
\ee^{\ii\theta}
&=\sum_{j=0}^{\infty}\frac{(-1)^j\theta^{2j}}{(2j)!}
\\&\quad+\ii\sum_{j=0}^{\infty}\frac{(-1)^j\theta^{2j+1}}{(2j+1)!}.
\end{align*}
两项分别是余弦和正弦的实幂级数，因而得到
\begin{equation}
\ee^{\ii\theta}=\cos\theta+\ii\sin\theta.
\end{equation}
平面旋转因而可以写成指数，连续转过两个角度对应两个复指数相乘。

对一般 $z=x+\ii y$，有 $\ee^z=\ee^x(\cos y+\ii\sin y)$。于是 $\ee^{z+2\pi\ii}=\ee^z$，复指数关于虚部具有周期性。由欧拉公式还得到
\begin{align}
\cos\theta&=\frac{\ee^{\ii\theta}+\ee^{-\ii\theta}}2,\\
\sin\theta&=\frac{\ee^{\ii\theta}-\ee^{-\ii\theta}}{2\ii}.
\end{align}
这些关系能统一处理三角函数的求导、积分与合成。对实参数 $t$，$\dd\ee^{\ii\omega t}/\dd t=\ii\omega\ee^{\ii\omega t}$，每次求导只增加一个因子 $\ii\omega$。

\subsection{复对数为何必须规定分支}
求 $\ee^w=z\ne0$ 时，设 $z=r\ee^{\ii\theta}$，则 $w=\ln r+\ii(\theta+2k\pi)$。因而复对数具有无穷多个值。局部选定连续的辐角，才能得到单值对数。

如果沿环绕原点一周的闭合路径连续增加辐角，终点的辐角比起点多 $2\pi$。因此在整个去心平面 $\mathbb C\setminus\{0\}$ 上，不可能同时得到连续、单值的对数。可去掉一条从原点出发的射线，选择一个分支，例如取 $-\pi<\theta<\pi$ 的割线平面。在这种规定下，不能不检查辐角就套用 $\log(z_1z_2)=\log z_1+\log z_2$：两侧可能相差 $2\pi\ii$。

\section{同频振动与相量}
\subsection{相量的定义}
设 $x(t)=A\cos(\omega t+\varphi)$，其中 $A\ge0$。用欧拉公式可写成
\begin{equation}
x(t)=\Rea[\widetilde X\ee^{\ii\omega t}],\qquad \widetilde X=A\ee^{\ii\varphi}.
\end{equation}
$\widetilde X$ 称为相量，记录峰值与初相位。它不包含公共时间因子，因此对于同一角频率的线性叠加，先加相量、后取实部即可：
\begin{equation}
\widetilde X=\sum_j\widetilde X_j,\quad A=|\widetilde X|,\quad\varphi=\arg\widetilde X.
\end{equation}
若总相量为零，合振动为零，总相位没有定义。直接相量相加的条件是同频；不同频率的时间因子不能抽成同一个因子。

求导和积分同样简单：谐波求导对应乘以 $\ii\omega$，谐波积分对应除以 $\ii\omega$，后者要求 $\omega\ne0$，且还应保留积分常数。复数方法适用于实系数线性方程，因为取实部与线性运算可交换。非线性运算一般不满足这个规则，例如真实信号的平方不能由相量平方后直接取实部得到。

\textbf{算例：同频振动的合成。} 合成 $\cos\omega t+\sin\omega t$，并求其初值和初始导数。
% 题设配图；与紧随其后的解答分别编号。
\ExpansionDiagram{fig:audit-m-q-003}{余弦与正弦采用同一个复指数约定}{
\draw[->](-.4,0)--(3.5,0)node[right]{$ \Rea z $};\draw[->](0,-1.7)--(0,2.3)node[above]{$ \Ima z $};\draw[->,thick](0,0)--(1,0)node[above right]{$ 1 $};\draw[->,thick](0,0)--(0,-1)node[above right]{$ -\ii $};
}

\begin{solution}
按图\ref{fig:audit-m-q-003}采用的余弦实部约定，$\cos\omega t=\Rea(\ee^{\ii\omega t})$，而
$\sin\omega t=\Rea(-\ii\ee^{\ii\omega t})$。因此图\ref{fig:warm-1}中的两相量为$1$和$-\ii$，合相量满足
\[
 |\widetilde X|=\sqrt{1^2+(-1)^2}=\sqrt2,\qquad
 \arg\widetilde X=-\frac\pi4.
\]
于是$x(t)=\sqrt2\cos(\omega t-\pi/4)$。分别代入函数及其导数，
\[
 x(0)=\sqrt2\cos(-\pi/4)=1,\qquad
 \dot x(0)=-\sqrt2\,\omega\sin(-\pi/4)=\omega.
\]
合相量从第四象限逆时针转动时，实轴投影初始增大，与正的初始导数一致。

\begin{figure}[H]\centering
\begin{tikzpicture}[scale=1.3]
\draw[bookaxis](-.3,0)--(2.2,0)node[right]{$\Rea$};\draw[bookaxis](0,-1.6)--(0,.8)node[above]{$\Ima$};
\draw[bookaxis](0,0)--(1,0)node[above]{$1$};\draw[bookaxis](1,0)--(1,-1)node[right]{$-\ii$};\draw[bookaxis,thick](0,0)--(1,-1)node[below right]{$1-\ii$};
\draw(.5,0)arc(0:-45:.5);\node[right]at(.7,-.2){$-\pi/4$};
\end{tikzpicture}
\caption{同频振动的相量合成：合相量 $1-\ii$ 的模为 $\sqrt2$，初始投影为1；逆时针转动时投影初始增大。}\label{fig:warm-1}
\end{figure}
\end{solution}

\begin{example}
合成 $x=3\cos\omega t+4\sin\omega t$。
\end{example}
% 题设配图；与紧随其后的解答分别编号。
\ExpansionDiagram{fig:audit-m-q-004}{两个相量的长度分别为3和4}{
\draw[->](-.4,0)--(3.7,0)node[right]{$ \Rea z $};\draw[->](0,-4.7)--(0,2.3)node[above]{$ \Ima z $};\draw[->,thick](0,0)--(3,0)node[above right]{$ 3 $};\draw[->,thick](0,0)--(0,-4)node[above right]{$ -4\ii $};
}

\begin{solution}
图\ref{fig:audit-m-q-004}中两项相差四分之一周期，先在同一相位约定下相加。
按图\ref{fig:orig-003}作首尾相接的相量加法，水平边来自余弦，向下的竖边来自正弦。

统一使用余弦和时间因子 $\ee^{\ii\omega t}$。因
\[
\sin\omega t=\cos(\omega t-\pi/2)
=\Rea[-\ii\ee^{\ii\omega t}],
\]
两个振动的相量分别为3和 $-4\ii$，合相量为 $X=3-4\ii$。设 $X=A\ee^{-\ii\delta}$，则
\[
A=\sqrt{3^2+4^2}=5,\qquad
\cos\delta=\frac35,\quad\sin\delta=\frac45,
\]
故 $\delta=\arctan(4/3)$，运动为
\[
x=5\cos(\omega t-\delta).
\]
直接展开也得到 $5\cos\delta\cos\omega t+5\sin\delta\sin\omega t$，与题式逐项相同。初值检查为 $x(0)=3$、$\dot x(0)=4\omega$；这两个量还能检验相位的正负号。

\begin{figure}[H]\centering
\begin{tikzpicture}[scale=.63]
\draw[bookaxis](-.4,0)--(4.5,0)node[right]{$\Rea$};\draw[bookaxis](0,-4.8)--(0,1)node[above]{$\Ima$};
\draw[bookaxis,thick](0,0)--(3,0)node[midway,above]{$3$};\draw[bookaxis](3,0)--(3,-4)node[midway,right]{$-4\ii$};
\draw[bookaxis,thick](0,0)--(3,-4)node[midway,left]{$3-4\ii$};\draw(.8,0)arc(0:-53.13:.8);
\node at(1.4,-.6){$\varphi$};
\end{tikzpicture}
\caption{余弦约定下的正弦相量与合相量。角度从正实轴逆时针为正。}
\label{fig:orig-003}
\end{figure}

图中合相量的两条直角边为3和4，故模为5；它位于第四象限，正好对应解中的 $-\delta$。
\end{solution}

\subsection{合幅值与拍}
两同频振动的相量平方模为
\begin{equation}
A^2=A_1^2+A_2^2+2A_1A_2\cos(\varphi_2-\varphi_1).
\end{equation}
相位差为0时幅值为和，相位差为 $\pi$ 时为差的绝对值。正交相位时 $A=\sqrt{A_1^2+A_2^2}$，这是勾股定理的相位形式。

两频率接近而不相同时，应保留时间因子。等幅情形可直接用和差化积：
\begin{equation}
A\cos\omega_1t+A\cos\omega_2t
=2A\cos\frac{\omega_1-\omega_2}{2}t\cos\frac{\omega_1+\omega_2}{2}t.
\end{equation}
快速振动的角频率接近平均值，而振幅大小由 $2A|\cos[(\omega_1-\omega_2)t/2]|$ 决定。注意绝对值使相邻幅度极大值的间隔为 $2\pi/|\omega_1-\omega_2|$，所以通常定义的拍频为 $|f_1-f_2|$，不能把余弦因子中的半频差直接称为拍频。

\section{复指数方法的运算边界}
复数表示直接适用于线性叠加和求导。乘积、取模和非线性方程须依据实际的实物理量处理，不能任意交换运算与取实部的次序。

\subsection{复三角函数与双曲函数}
将指数形式作为定义，可以把三角函数延伸到复变量。先计算复余弦：
\begin{align*}
\cos(x+\ii y)
&=\frac{\ee^{-y}(\cos x+\ii\sin x)
+\ee^y(\cos x-\ii\sin x)}2\\
&=\cos x\frac{\ee^y+\ee^{-y}}2
-\ii\sin x\frac{\ee^y-\ee^{-y}}2.
\end{align*}
复正弦则把分子中的两指数相减，再除以$2\ii$。分别合并实部和虚部，得到
\begin{align}
\cos(x+\ii y)&=\cos x\cosh y-\ii\sin x\sinh y,\\
\sin(x+\ii y)&=\sin x\cosh y+\ii\cos x\sinh y.
\end{align}
其中 $\cosh y=(\ee^y+\ee^{-y})/2$，$\sinh y=(\ee^y-\ee^{-y})/2$。所以复三角函数可以指数增长，实变量上 $|\sin x|\le1$ 的性质不能直接用于复变量。这正是闭合路径求积分时必须检查圆弧项的原因。

复指数也能统一阻尼与振荡。$\ee^{(-\gamma+\ii\omega)t}$ 的实部是 $\ee^{-\gamma t}\cos\omega t$。实部控制包络的增减，虚部控制转动速率。特征根的实部、虚部分别对应哪种物理行为，在矩阵微分方程中仍然成立。

\subsection{相量图必须明确参考方向}
将电流或驱动力相位设为零，并不改变物理信号，只是规定共同的相位参考。对 $3\cos\omega t+4\sin\omega t$，正弦对应向下的虚轴分量。相量和位于第四象限，相位为负。

不同频率不能合成一个固定相量，但可以考察缓慢变化的包络。对 $\omega_1=4.2,\omega_2=3.8\,\mathrm{rad/s}$ 的等幅信号，合成式是 $2A\cos(0.2t)\cos(4t)$。相邻包络大小极大值间隔为 $\pi/0.2$，对应拍频 $0.4/(2\pi)$。
\begin{figure}[htbp]\centering
\begin{tikzpicture}[x=.29cm,y=.72cm]
\draw[bookaxis](0,0)--(20.7,0)node[right]{$t/\mathrm s$};\draw[bookaxis](0,-2.4)--(0,2.6)node[above]{$x/A$};
\draw[domain=0:20,samples=600]plot(\x,{2*cos(deg(.2*\x))*cos(deg(4*\x))});
\draw[dashed,domain=0:20,samples=130]plot(\x,{2*cos(deg(.2*\x))});\draw[dashed,domain=0:20,samples=130]plot(\x,{-2*cos(deg(.2*\x))});
\node[left]at(0,2){$2$};\node[left]at(0,-2){$-2$};\draw(10,.08)--(10,-.08)node[below]{$10$};\draw(20,.08)--(20,-.08)node[below]{$20$};
\end{tikzpicture}
\caption{两个近频信号的拍。虚线是两条有符号包络，物理振幅大小取其绝对值。}
\label{fig:orig-004}
\end{figure}

\section{受迫阻尼振动：从瞬态到稳态}
\subsection{方程、参数与齐次解}
考虑质量为 $m$、刚度为 $k$、线性阻尼系数为 $b$ 的单自由度系统，$m,k>0,b\ge0$：
\begin{equation}
m\ddot x+b\dot x+kx=F_0\cos\omega t.
\end{equation}
定义 $\omega_0=\sqrt{k/m}$ 与 $\gamma=b/(2m)$。齐次方程的指数试探解为 $x_h=\ee^{st}$，特征方程 $s^2+2\gamma s+\omega_0^2=0$。因此阻尼改变的不只是幅值，它也改变瞬态的形式。

当 $0\le\gamma<\omega_0$ 时，$s=-\gamma\pm\ii\omega_d$，$\omega_d=\sqrt{\omega_0^2-\gamma^2}$，实齐次解为
\begin{equation}
x_h=\ee^{-\gamma t}(C\cos\omega_dt+D\sin\omega_dt).
\end{equation}
临界阻尼 $\gamma=\omega_0$ 时，$x_h=(C+Dt)\ee^{-\gamma t}$。过阻尼时，两个负实根给出两个衰减指数。对于 $b>0$，这些瞬态最终衰减；对于 $b=0$，自由振动不衰减，不能一般地说总解最终只剩稳态。

\subsection{相量法求稳态}
将右侧扩展为 $F_0\ee^{\ii\omega t}$，取复特解 $x_c=\widetilde X\ee^{\ii\omega t}$，得到
\begin{equation}
\widetilde X=\frac{F_0}{k-m\omega^2+\ii b\omega}.
\end{equation}
如果分母不为零，取实部得 $x_s=A\cos(\omega t-\delta)$，其中
\begin{align}
A&=\frac{F_0}{\sqrt{(k-m\omega^2)^2+b^2\omega^2}},\\
\delta&=\atan(b\omega,k-m\omega^2).
\end{align}
对 $b>0,\omega>0$，相位滞后在0与 $\pi$ 之间。低频时弹性力主导，$A\to F_0/k$，位移几乎同相；高频时惯性项主导，$A\sim F_0/(m\omega^2)$，位移几乎反相。这里幅值分母量纲为刚度，结果量纲为长度。

\begin{example}
取 $m=1\,\mathrm{kg},b=2\,\mathrm{N\,s/m},k=25\,\mathrm{N/m}$，驱动力 $10\cos3t\,\mathrm N$，且 $x(0)=\dot x(0)=0$。求完整运动。
\end{example}
% 题设配图；与紧随其后的解答分别编号。
\ExpansionDiagram{fig:audit-m-q-005}{质量、弹簧和阻尼器的受迫运动}{
\draw[very thick](0,-.8)--(0,.9);\draw[exspring](0,.35)--(2.4,.35);\node[above]at(1.2,.45){$ 25 $};\draw[thick](2.4,-.65)rectangle(3.5,.65);\node at(2.95,0){$ 1 $};\draw(0,-.35)--(.8,-.35);\draw(.8,-.6)--(.8,-.1)--(1.7,-.1);\draw(.8,-.6)--(1.7,-.6);\draw(1.35,-.55)--(1.35,-.15);\draw(1.35,-.35)--(2.4,-.35);\node[below]at(1.2,-.7){$ 2 $};\draw[->,thick](3.5,0)--(5,0)node[right]{$ 10\cos3t $};\draw[->](2.4,-1.2)--(4,-1.2)node[right]{$x$};
}

\begin{solution}
图\ref{fig:audit-m-q-005}中$x$向右为正，弹簧力和阻尼力分别反对位移与速度。
图\ref{fig:eg-1-4}用同一时间轴比较同频特解与满足初值的完整解。两者之差就是待求的齐次项。

运动方程与初值为
\[
\ddot x+2\dot x+25x=10\cos3t,\qquad x(0)=\dot x(0)=0.
\]
先求同频特解。设$x_s=\Rea[X\ee^{3\ii t}]$，则
\[\dot x_s=\Rea[3\ii X\ee^{3\ii t}],\qquad
\ddot x_s=\Rea[(3\ii)^2X\ee^{3\ii t}]
=\Rea[-9X\ee^{3\ii t}].\]
代入方程，约去公共时间因子，得
\[
(25-9+6\ii)X=10,\qquad
X=\frac{10(16-6\ii)}{16^2+6^2}=\frac{40-15\ii}{73}.
\]
因此 $x_s=(40\cos3t+15\sin3t)/73$，幅值 $5/\sqrt{73}\,\mathrm m$，滞后角为 $\arctan(3/8)$。这个特解的初始位移和速度分别为 $40/73$ 和 $45/73$，尚不满足给定初值。

\begin{figure}[H]\centering
\begin{tikzpicture}[x=1.15cm,y=1.5cm]
\draw[bookaxis](0,0)--(6.35,0)node[right]{$t/\mathrm s$};\draw[bookaxis](0,-.8)--(0,.9)node[above]{$x/\mathrm m$};
\draw[dashed,domain=0:6,samples=180]plot(\x,{(40*cos(3*\x r)+15*sin(3*\x r))/73});
\draw[bookcurve,domain=0:6,samples=200]plot(\x,{(40*cos(3*\x r)+15*sin(3*\x r)-exp(-\x)*(40*cos(sqrt(24)*\x r)+85/sqrt(24)*sin(sqrt(24)*\x r)))/73});
\foreach\t in{2,4,6}{\draw(\t,0)--(\t,-.05)node[below]{$\t$};}
\fill(0,0)circle(1.5pt);\fill(0,{40/73})circle(1.5pt)node[left]{$40/73$};
\draw[thick](.5,.87)--(1.05,.87)node[right]{完整解};\draw[dashed](3.2,.87)--(3.75,.87)node[right]{同频特解};
\end{tikzpicture}
\caption{例1.4：齐次项把初始位移与初始速度修正为零，随后指数衰减。虚线特解在 $t=0$ 不满足初值。}\label{fig:eg-1-4}
\end{figure}

齐次方程取$x_h=\ee^{\lambda t}$，得到$\lambda^2+2\lambda+25=0$，即$(\lambda+1)^2=-24$，两根为$-1\pm\ii\sqrt{24}$。故完整解应写成
\[
x=x_s+\ee^{-t}(C\cos\sqrt{24}t+D\sin\sqrt{24}t).
\]
代入位移初值得$C=-40/73$。记$\omega_d=\sqrt{24}$，按乘积法则，
\[\dot x_h=\ee^{-t}\bigl[-C\cos\omega_dt-D\sin\omega_dt
-C\omega_d\sin\omega_dt+D\omega_d\cos\omega_dt\bigr],\]
而$\dot x_s=(-120\sin3t+45\cos3t)/73$。令$t=0$：
\[
0=\frac{45}{73}-C+\sqrt{24}D,
\qquad D=-\frac{85}{73\sqrt{24}}.
\]
所求运动为
\[
x(t)=\frac{40\cos3t+15\sin3t}{73}
-\frac{\ee^{-t}}{73}\left(40\cos\sqrt{24}t+
\frac{85}{\sqrt{24}}\sin\sqrt{24}t\right)\ \mathrm m.
\]
其中指数项在长时间后消失。将 $t=0$ 代回，位移与速度均为零；由原方程还应有 $\ddot x(0)=10\,\mathrm{m/s^2}$，可作第二项独立检查。

虚线的初始位移为 $40/73$；实线从零出发，初始切线水平，满足给定的位移和速度初值。
\end{solution}

\subsection{共振频率依赖于观察量}
求位移幅值最大值，等价于极小化 $D(\omega)=(k-m\omega^2)^2+b^2\omega^2$。求导得
\[
D'=2(k-m\omega^2)(-2m\omega)+2b^2\omega
=2\omega[-2m(k-m\omega^2)+b^2].
\]
若 $\omega_0^2>2\gamma^2$，非零峰频率为 $\omega_r=\sqrt{\omega_0^2-2\gamma^2}$；否则幅值在正频率上没有非零局部峰值。$\omega_0$ 则始终是相位滞后为 $\pi/2$ 的频率。两者一般不同。

稳态速度峰值为$V=\omega A$。对$V^2=F_0^2\omega^2/D$用商法则，
\[\frac{\dd V^2}{\dd\omega}
=F_0^2\frac{2\omega D-\omega^2D'}{D^2}
=\frac{2F_0^2\omega(k^2-m^2\omega^4)}{D^2}.\]
在$b>0,\omega>0$时，$D>0$，导数于$\omega=\sqrt{k/m}$两侧由正变负，故速度峰在$\omega_0$。平均耗散功率$\overline P=bV^2/2$也在此频率最大。谈论共振时须说明所观察的是位移、速度还是功率。

\subsection{能量平衡与无阻尼共振}
令 $E=mv^2/2+kx^2/2$，方程乘以 $v$ 得
\begin{equation}
\frac{\dd E}{\dd t}=F_0v\cos\omega t-bv^2.
\end{equation}
在稳态一个周期的平均能量变化为零，输入平均功率等于耗散平均功率。由$v=-\omega A\sin(\omega t-\delta)$展开瞬时输入功率，
\[F_0v\cos\omega t=F_0\omega A
\bigl[\sin\delta\cos^2\omega t-\cos\delta\sin\omega t\cos\omega t\bigr].\]
取$T=2\pi/\omega$，两个平均积分分别为
\[\frac1T\left[\frac t2+\frac{\sin2\omega t}{4\omega}\right]_0^T=\frac12,
\qquad
\frac1T\left[\frac{\sin^2\omega t}{2\omega}\right]_0^T=0.\]
故输入平均功率为$F_0\omega A\sin\delta/2$。再代入
$\sin\delta=b\omega/\sqrt D$和$A=F_0/\sqrt D$，得到$b\omega^2A^2/2$。

当 $b=0,\omega=\omega_0$ 时，稳态相量分母为零，原公式失效。试探共振特解为
\begin{equation}
x_p=\frac{F_0}{2m\omega_0}t\sin\omega_0t.
\end{equation}
系数由下面的恒等式确定：
\begin{equation}\label{eq:resonant-secular-identity}
\frac{\dd^2}{\dd t^2}(t\sin\omega_0t)
+\omega_0^2t\sin\omega_0t=2\omega_0\cos\omega_0t.
\end{equation}
左侧再乘$F_0/(2m\omega_0)$和$m$，恰好恢复外力。两次乘积求导及初值的逐项验证见本章图\ref{fig:hw-1-5}所配习题解答。振幅随时间增长，说明理想无阻尼系统没有有限幅值的谐波稳态。在真实系统中，阻尼、非线性或其他限制会阻止这种无限增长。

\section{阻尼模型的完整比较}
\subsection{相同初值下的三类自由响应}
令$\tau=\omega_0t$，自由方程为$q''+2\zeta q'+q=0$。特征方程配方为$(s+\zeta)^2=\zeta^2-1$，对同一初态$q(0)=1,q'(0)=0$分别求常数。

欠阻尼$0\le\zeta<1$时，记$w=\sqrt{1-\zeta^2}$，通解为$q=\ee^{-\zeta\tau}(A\cos w\tau+B\sin w\tau)$。初值给出$A=1$、$-\zeta A+wB=0$，所以$B=\zeta/w$，即
\[
q=\ee^{-\zeta\tau}\left[\cos(\sqrt{1-\zeta^2}\tau)+\frac{\zeta}{\sqrt{1-\zeta^2}}\sin(\sqrt{1-\zeta^2}\tau)\right].
\]
临界阻尼时，设$q=\ee^{-\tau}u(\tau)$。两次乘积求导得到
\[q'=\ee^{-\tau}(u'-u),\qquad q''=\ee^{-\tau}(u''-2u'+u).\]
代入$q''+2q'+q=0$，括号内$-2u'+2u'=0$、$u-2u+u=0$，剩下$u''=0$。两次积分得$u=A+B\tau$，初值给出$A=1,B-A=0$，故$q=(1+\tau)\ee^{-\tau}$。

过阻尼时令$d=\sqrt{\zeta^2-1}$、$s_\pm=-\zeta\pm d$，通解为$q=A\ee^{s_+\tau}+B\ee^{s_-\tau}$。两初值为$A+B=1$、$s_+A+s_-B=0$，解得$A=-s_-/(s_+-s_-)=(\zeta+d)/(2d)$、$B=s_+/(s_+-s_-)=(d-\zeta)/(2d)$。因此
\[
q=\frac{\zeta+d}{2d}\ee^{(-\zeta+d)\tau}+\frac{d-\zeta}{2d}\ee^{(-\zeta-d)\tau}.
\]
过阻尼的慢衰减率可有理化为$\zeta-d=(\zeta^2-d^2)/(\zeta+d)=1/(\zeta+d)$，在大$\zeta$时约为$1/(2\zeta)$。因此阻尼过大时，回到平衡点的过程反而变慢。
\begin{figure}[htbp]\centering
\begin{tikzpicture}[x=.64cm,y=1.65cm]
\draw[bookaxis](0,0)--(8.7,0)node[right]{$\tau$};\draw[bookaxis](0,-.7)--(0,1.25)node[above]{$q$};
\draw[thick,domain=0:8.4,samples=180]plot(\x,{exp(-.2*\x)*(cos(deg(sqrt(.96)*\x))+.2/sqrt(.96)*sin(deg(sqrt(.96)*\x)))});
\draw[dashed,domain=0:8.4,samples=130]plot(\x,{(1+\x)*exp(-\x)});
\draw[dotted,thick,domain=0:8.4,samples=130]plot(\x,{(1.5+sqrt(1.25))/(2*sqrt(1.25))*exp((-1.5+sqrt(1.25))*\x)+(sqrt(1.25)-1.5)/(2*sqrt(1.25))*exp((-1.5-sqrt(1.25))*\x)});
\node[left]at(0,1){$1$};\foreach\t in{2,4,6,8}{\draw(\t,.025)--(\t,-.025)node[below]{$\t$};}
\end{tikzpicture}
\caption{同一初值的自由响应：实线 $\zeta=0.2$，虚线 $\zeta=1$，点线 $\zeta=1.5$。}
\label{fig:orig-005}
\end{figure}

\subsection{位移响应、相位响应与品质因数}
相位滞后在 $\Omega=1$ 时为 $\pi/2$，与阻尼比无关；但相位变化的陡峭程度与阻尼有关。小阻尼系统在固有频率附近对频率变化更敏感。这既体现在幅值峰，也体现在相位曲线。
\begin{figure}[htbp]\centering
\begin{tikzpicture}[x=2.1cm,y=2.7cm]
\draw[bookaxis](0,0)--(2.75,0)node[right]{$\Omega$};\draw[bookaxis](0,0)--(0,1.2)node[above]{$\delta$};
\draw[thick,domain=0:2.6,samples=150]plot(\x,{atan2(.2*\x,1-\x*\x)/180});
\draw[dashed,domain=0:2.6,samples=150]plot(\x,{atan2(.6*\x,1-\x*\x)/180});
\draw[dotted,thick,domain=0:2.6,samples=150]plot(\x,{atan2(1.4*\x,1-\x*\x)/180});
\draw[bookguide](1,0)--(1,.5)--(0,.5);\node[left]at(0,.5){$\pi/2$};\node[left]at(0,1){$\pi$};\foreach\t in{1,2}{\draw(\t,.02)--(\t,-.02)node[below]{$\t$};}
\end{tikzpicture}
\caption{不同阻尼的相位响应：实线 $\zeta=0.1$，虚线 $0.3$，点线 $0.7$。}
\label{fig:orig-006}
\end{figure}

对有阻尼振子，常定义品质因数$Q=\omega_0/(2\gamma)=1/(2\zeta)$。自由振幅包络按$\ee^{-\gamma t}$衰减，相同振动相位的能量按$\ee^{-2\gamma t}$衰减。弱阻尼时$T_d\simeq2\pi/\omega_0$，一周期损失比例为
\[\frac{\Delta E}{E}=1-\ee^{-2\gamma T_d}
\simeq2\gamma T_d\simeq\frac{4\pi\gamma}{\omega_0}.\]
因此$Q\simeq2\pi E/\Delta E$。这种能量解释使用了弱阻尼下周期与指数展开的近似，不能延用于过阻尼系统。

\begin{example}
相邻同方向位移峰值之比为 $x_n/x_{n+1}=\ee^{0.10}$，测得峰值间隔 $T_d$。求衰减率与固有角频率。
\end{example}
% 题设配图；与紧随其后的解答分别编号。
\ExpansionDiagram{fig:audit-m-q-006}{同向峰的时间间隔与峰值比}{
\draw[->](-.4,0)--(7,0)node[right]{$ t $};\draw[->](0,-1.5)--(0,2.5)node[above]{$ x $};\draw[thick,domain=0:7,samples=140]plot(\x,{1.8*exp(-.1*\x/(2*pi))*(cos(\x r)+.1/(2*pi)*sin(\x r))});\draw[dashed](0,1.8)--(6.283,1.8);\fill(0,1.8)circle(1.5pt)node[left]{$x_n$};\fill(6.283,1.6287)circle(1.5pt)node[above]{$x_{n+1}$};\draw[<->](0,-1.8)--(6.283,-1.8)node[midway,below]{$T_d$};
}

\begin{solution}
图\ref{fig:audit-m-q-006}中两次同向峰相隔$T_d$，峰值的比值只包含包络的衰减。
在图\ref{fig:eg-1-5}中标出两次同向峰，横向间隔为 $T_d$，纵向峰比由题给数据确定。

欠阻尼自由运动写成 $x=A\ee^{-\gamma t}\cos(\omega_dt+\varphi)$。极值由
\[
\gamma\cos(\omega_dt+\varphi)+\omega_d\sin(\omega_dt+\varphi)=0
\]
确定。同向相邻极大值的相位相差 $2\pi$，所以 $T_d=2\pi/\omega_d$；两处三角因子相同，峰值之比便是 $\ee^{\gamma T_d}$。

\begin{figure}[H]\centering
\begin{tikzpicture}[x=2.3cm,y=1.4cm]
\draw[bookaxis](0,0)--(3.3,0)node[right]{$t/T_d$};\draw[bookaxis](0,-1.15)--(0,1.4)node[above]{$x/x_n$};
\draw[bookcurve,domain=0:3.05,samples=180]plot(\x,{exp(-.1*\x)*cos(360*\x-atan(.1/(2*pi)))/cos(atan(.1/(2*pi)))});
\draw[bookguide,domain=0:3.05]plot(\x,{exp(-.1*\x)/cos(atan(.1/(2*pi)))});
\foreach\t in{1,2,3}{\draw[bookguide](\t,0)--(\t,{exp(-.1*\t)});\fill(\t,{exp(-.1*\t)})circle(1.5pt);}
\node[above right]at(0,1){$x_n$};\node[above right]at(1,{exp(-.1)}){$x_{n+1}$};
\draw[<->](0,-1.27)--(1,-1.27)node[midway,below]{$T_d$};
\node at(2.1,1.3){$x_{n+1}/x_n=\ee^{-0.10}$};
\end{tikzpicture}
\caption{例1.5：时间以实测峰间隔 $T_d$ 归一化。同向峰的三角因子相同，峰比只留下指数衰减。}\label{fig:eg-1-5}
\end{figure}

由题给对数比 $\gamma T_d=0.10$，得到
\[
\gamma=\frac{0.10}{T_d},\qquad
\omega_d=\frac{2\pi}{T_d},\qquad
\omega_0=\sqrt{\omega_d^2+\gamma^2}
=\frac{\sqrt{4\pi^2+0.01}}{T_d}.
\]
两频率之比为 $\omega_0/\omega_d=\sqrt{1+(0.10/2\pi)^2}$，比1稍大。测得的峰间隔直接对应阻尼频率，求无阻尼固有频率还需用衰减量修正。

峰值时刻的余弦因子未必等于1。相邻同向峰的相位相差整周，三角因子相同，所以峰比公式仍然精确。
\end{solution}

\section{近似的尺度：共振附近该保留什么}
\subsection{先找小参数，再展开}
对 $\omega_0=\sqrt{k/m}$，令 $\Omega=\omega/\omega_0$、$\zeta=b/(2m\omega_0)$。位移幅值除以静态位移 $F_0/k$，得到
\[
\mathcal A(\Omega)=\bigl[(1-\Omega^2)^2+4\zeta^2\Omega^2\bigr]^{-1/2}.
\]
设 $\zeta\ll1$。如果 $\Omega$ 与1相隔一个固定的量，阻尼项是小修正；在 $\Omega=1$ 附近，第一项也很小，阻尼却决定了振幅的尺度。展开之前必须先说明研究的是哪一段频率。

写 $\Omega=1+\varepsilon$，则 $1-\Omega^2=-2\varepsilon-\varepsilon^2$。在 $\varepsilon=O(\zeta)$ 的窄频带内，分母平方的主要项是 $4\varepsilon^2+4\zeta^2$，故
\begin{equation}\label{eq:near-resonance}
\mathcal A\simeq\frac1{2\sqrt{\varepsilon^2+\zeta^2}}.
\end{equation}
这里舍去的项比保留项小一个相对量级 $O(\zeta)$。即使阻尼很弱，也不能在该频带中先令 $\zeta=0$：在 $\varepsilon=0$，这样会把有限的峰抹成无穷大。

\begin{figure}[htbp]\centering
\begin{tikzpicture}[x=.55cm,y=2cm]
\draw[bookaxis](-5.3,0)--(5.4,0)node[right]{$\varepsilon/\zeta$};\draw[bookaxis](0,0)--(0,1.3)node[above]{$2\zeta\mathcal A$};
\draw[bookcurve,domain=-5:5,samples=150]plot(\x,{1/sqrt(1+\x*\x)});
\draw[bookguide](-1,0)--(-1,.707)--(1,.707)--(1,0);
\node[below]at(-1,0){$-1$};\node[below]at(1,0){$1$};\node[left]at(0,1){$1$};\node[right]at(1.4,.72){$1/\sqrt2$};
\end{tikzpicture}
\caption{弱阻尼共振带的缩放图。横轴改用失谐量与阻尼比之比后，不同弱阻尼系统具有同一首阶形状。}
\label{fig:orig-007}
\end{figure}

\subsection{半功率宽度与近似精度}
平均耗散功率为 $b\omega^2A^2/2$。在上述窄带内，$\omega^2$ 的相对变化仅为 $O(\zeta)$，因此首阶半功率点满足 $|\varepsilon|=\zeta$，全宽约为 $2\zeta\omega_0$，相应 $Q\simeq\omega_0/\Delta\omega=1/(2\zeta)$。

需要更高精度时，直接处理功率式。以 $r=\omega/\omega_0$ 表示，取$\zeta>0,r>0$，半功率条件为
$(1-r^2)^2=4\zeta^2r^2$。开方后须按峰的两侧分支：
\[
\begin{aligned}
0<r<1:&\quad 1-r^2=2\zeta r,
&& (r+\zeta)^2=1+\zeta^2,\\
r>1:&\quad r^2-1=2\zeta r,
&& (r-\zeta)^2=1+\zeta^2.
\end{aligned}
\]
各支保留符合相应区间的正根，得到
\[
r_\pm=\sqrt{1+\zeta^2}\pm\zeta.
\]
因此恒幅力驱动、黏性阻尼模型的功率半高全宽恰好为 $2\zeta\omega_0$。两半高点的算术平均不等于 $\omega_0$；它们的几何平均才等于 $\omega_0$。首阶缩放图的左右对称不能用来断言精确曲线也左右对称。

\begin{example}
某振子的功率峰出现在 $100\,\mathrm{rad/s}$，半功率点为 $\omega_-,\omega_+$，测得宽度 $4\,\mathrm{rad/s}$。估计阻尼比、品质因数和位移峰的位置。
\end{example}
% 题设配图；与紧随其后的解答分别编号。
\ExpansionDiagram{fig:audit-m-q-007}{功率峰和半高交点对应不同测量量}{
\draw[->](-.4,0)--(5,0)node[right]{$ \omega $};\draw[->](0,-0.3)--(0,2.8)node[above]{$ \overline P $};\draw[thick,domain=.2:4.5,samples=100]plot(\x,{2.2/(1+4*(\x-2.4)^2)});\draw[dashed](.2,1.1)--(4.5,1.1);\node[left]at(0,1.1){$P_{\max}/2$};\draw[dotted](2.4,0)--(2.4,2.2);\node[above]at(2.4,2.4){$\omega_0=100$};\node[below left]at(1.9,0){$\omega_-$};\node[below right]at(2.9,0){$\omega_+$};
}

\begin{solution}
图\ref{fig:audit-m-q-007}中给定的是功率曲线，半高条件应作用于平均功率。
图\ref{fig:eg-1-6}中的水平半高线把测量量转为方程：交点处 $\overline P=P_{\max}/2$。

恒幅外力驱动的黏性阻尼振子，其平均功率为
\[
\overline P(\omega)=\frac{F_0^2}{2b}
\frac{b^2\omega^2}{(k-m\omega^2)^2+b^2\omega^2}.
\]
分式在 $k-m\omega^2=0$ 时为1，所以功率峰给出 $\omega_0=100\,\mathrm{rad/s}$。半功率条件是
\[
|k-m\omega^2|=b\omega.
\]
令 $\gamma=b/(2m)$，两正根分别为
\[
\omega_-=\sqrt{\omega_0^2+\gamma^2}-\gamma,
\qquad \omega_+=\sqrt{\omega_0^2+\gamma^2}+\gamma.
\]
由此峰宽恰为 $2\gamma=4\,\mathrm{rad/s}$，从而
\[
\gamma=2\,\mathrm{s^{-1}},\quad
\zeta=\gamma/\omega_0=0.02,\quad Q=1/(2\zeta)=25.
\]
位移幅值的分母没有功率式分子中的 $\omega^2$。对它另行求极值得
\[
\omega_{x,r}=\sqrt{\omega_0^2-2\gamma^2}
=\sqrt{9992}\simeq99.960\,\mathrm{rad/s}.
\]
峰宽关系在本题模型中是精确的；最后的十进制数作了舍入。位移峰与功率峰约差 $0.040\,\mathrm{rad/s}$，测量分辨率需小于这个差值才有望区分。

\begin{figure}[H]\centering
\begin{tikzpicture}[x=1.1cm,y=2.3cm]
\draw[bookaxis](-3.4,0)--(3.6,0)node[right]{$\omega$};\draw[bookaxis](-3.2,0)--(-3.2,1.22)node[above]{$\overline P/P_{\max}$};
\draw[bookcurve,domain=-3:3,samples=180]plot(\x,{.0016*(1+.02*\x)^2/((1-(1+.02*\x)^2)^2+.0016*(1+.02*\x)^2)});
\draw[bookguide](-3.2,.5)--(3,.5);\node[left]at(-3.2,.5){$1/2$};
\draw[bookguide](0,0)--(0,1);\node[below]at(0,0){$100$};
\foreach\v/\lab in{-.9900009998/\omega_-,1.0099990002/\omega_+}{\fill(\v,.5)circle(1.6pt);\draw[bookguide](\v,0)--(\v,.5);\node[below]at(\v,0){$\lab$};}
\draw[<->](-.990001,-.23)--(1.009999,-.23)node[midway,below]{$4\,\mathrm{rad/s}$};
\end{tikzpicture}
\caption{例1.6：用两条半高交线读功率峰宽。半功率点不必关于峰频率严格对称，但其频差恰为 $2\gamma$。}\label{fig:eg-1-6}
\end{figure}

图上的宽度由两个半功率交点确定；位移峰与功率峰的距离则由两种响应各自的极值条件确定。
\end{solution}

\subsection{共振近似的应用}
\textbf{分母的尺度展开。} 在 $\Omega=1+\varepsilon$ 下展开分母平方，保留三阶项；当 $\varepsilon=O(\zeta)$ 时，哪些项同阶？

% 题设配图；与紧随其后的解答分别编号。
\ExpansionDiagram{fig:audit-m-q-008}{共振附近同时比较失谐与阻尼的尺度}{
\draw[->](-.4,0)--(5,0)node[right]{$ \Omega $};\draw[->](0,-0.3)--(0,1.8)node[above]{$  $};\draw[very thick](.5,0)--(4.5,0);\draw(2.5,-.15)--(2.5,.15);\node[below]at(2.5,0){$1$};\draw[<->](2.5,.7)--(3.5,.7)node[midway,above]{$\varepsilon$};\node[above]at(2.5,1.4){$\Omega=1+\varepsilon,\quad\varepsilon=O(\zeta)$};
}

\begin{solution}
图\ref{fig:audit-m-q-008}给定$\varepsilon=O(\zeta)$，因此展开时须把失谐与阻尼放在同一阶数下比较。解图\ref{fig:class-1-1}采用幅值分母平方
$D=(1-\Omega^2)^2+4\zeta^2\Omega^2$。代入$\Omega=1+\varepsilon$，先分开展开两项：
\[
 (1-\Omega^2)^2=(-2\varepsilon-\varepsilon^2)^2
 =4\varepsilon^2+4\varepsilon^3+\varepsilon^4,
\]
\[
 4\zeta^2\Omega^2=4\zeta^2+8\zeta^2\varepsilon
 +4\zeta^2\varepsilon^2.
\]
若$\varepsilon=O(\zeta)$，则二阶项为$4(\varepsilon^2+\zeta^2)$，
三阶项为$4\varepsilon^3+8\zeta^2\varepsilon$；剩余两项同为四阶。
令$u=\varepsilon/\zeta$，在$u=O(1)$、$\zeta\to0$时，
\[
 \frac{D}{4\zeta^2}
 =1+u^2+\zeta(u^3+2u)
 +\zeta^2\left(\frac{u^4}{4}+u^2\right).
\]
图中虚线$1+u^2$保留二阶分母；它与精确曲线的首个差别来自奇函数$u^3+2u$，故共振两侧的修正并不对称。

\begin{figure}[H]\centering
\begin{tikzpicture}[x=1.15cm,y=.48cm]
\draw[bookaxis](-2.3,0)--(2.6,0)node[right]{$u=\varepsilon/\zeta$};\draw[bookaxis](-2.2,0)--(-2.2,6.8)node[above]{$D/(4\zeta^2)$};
\draw[bookcurve,domain=-2:2,samples=100]plot(\x,{((1-(1+.02*\x)^2)^2+.0016*(1+.02*\x)^2)/.0016});
\draw[dashed,domain=-2:2,samples=100]plot(\x,{1+(\x)^2});\node at(.3,6.5){实线：精确；虚线：二阶项};
\end{tikzpicture}
\caption{随堂练习1解图：示取 $\zeta=0.02$。当 $u=O(1)$，二阶项形成 $1+u^2$，三阶项使曲线出现首个不对称修正。}\label{fig:class-1-1}
\end{figure}
\end{solution}

\textbf{位移半幅宽度。} 位移幅值降为峰值的一半时，能否仍令 $|\varepsilon|=\zeta$？

% 题设配图；与紧随其后的解答分别编号。
\ExpansionDiagram{fig:audit-m-q-009}{半幅条件取在位移幅值的纵轴上}{
\draw[->](-.4,0)--(4.5,0)node[right]{$ |\varepsilon|/\zeta $};\draw[->](0,-0.3)--(0,2.8)node[above]{$ A/A_{\max} $};\draw[thick,domain=0:4,samples=80]plot(\x,{2.2/sqrt(1+\x*\x)});\draw[dashed](0,1.1)--(4,1.1);\node[left]at(0,1.1){$1/2$};\node[left]at(0,2.2){$1$};
}

\begin{solution}
题设图\ref{fig:audit-m-q-009}指定幅值降至峰值的一半。解图\ref{fig:class-1-2}同时画出半幅线和半功率对应的幅值线。在弱阻尼、$|\varepsilon|=O(\zeta)$的首阶近似中，半幅条件为
\[
 \frac{A}{A_{\max}}\simeq
 \frac{\zeta}{\sqrt{\varepsilon^2+\zeta^2}}=\frac12.
\]
两边平方，再移项得$\varepsilon^2+\zeta^2=4\zeta^2$，因此
$|\varepsilon|=\sqrt3\zeta$，两交点间的无量纲半幅全宽为$2\sqrt3\zeta$。
若取半功率，在这一近共振近似下功率正比于幅值平方，
所以相应幅值比为$1/\sqrt2$，代入同一公式才得到$|\varepsilon|=\zeta$。
图中的两条水平线代表不同条件。这里以$\Omega=1$为首阶峰位；
精确位移峰位为$\sqrt{1-2\zeta^2}$，其偏移属于下一阶，精确宽度应从完整幅值式求解。

\begin{figure}[H]\centering
\begin{tikzpicture}[x=1.7cm,y=2.8cm]
\draw[bookaxis](0,0)--(3.4,0)node[right]{$|\varepsilon|/\zeta$};\draw[bookaxis](0,0)--(0,1.2)node[above]{$A/A_{\max}$};
\draw[bookcurve,domain=0:3.1,samples=90]plot(\x,{1/sqrt(1+(\x)^2)});\draw[bookguide](0,.5)--(3.1,.5);\draw[dotted](0,{1/sqrt(2)})--(3.1,{1/sqrt(2)});
\draw[bookguide]({sqrt(3)},0)--({sqrt(3)},.5);\draw[dotted](1,0)--(1,{1/sqrt(2)});
\node[below]at(1,0){$1$};\node[below]at({sqrt(3)},0){$\sqrt3$};\node[left]at(0,.5){$1/2$};\node[left]at(0,{1/sqrt(2)}){$1/\sqrt2$};
\end{tikzpicture}
\caption{随堂练习2解图：半幅线与半功率对应的幅值线位于不同高度，交点的失谐量分别为 $\sqrt3\zeta$ 与 $\zeta$。}\label{fig:class-1-2}
\end{figure}
\end{solution}

\textbf{无阻尼近似的适用范围。} 在全部频率都用 $A=F_0/|k-m\omega^2|$ 近似弱阻尼响应，哪里会失效？

% 题设配图；与紧随其后的解答分别编号。
\ExpansionDiagram{fig:audit-m-q-010}{弱阻尼模型仍保留阻尼元件}{
\draw[very thick](0,-.8)--(0,.9);\draw[exspring](0,.35)--(2.4,.35);\node[above]at(1.2,.45){$ k $};\draw[thick](2.4,-.65)rectangle(3.5,.65);\node at(2.95,0){$ m $};\draw(0,-.35)--(.8,-.35);\draw(.8,-.6)--(.8,-.1)--(1.7,-.1);\draw(.8,-.6)--(1.7,-.6);\draw(1.35,-.55)--(1.35,-.15);\draw(1.35,-.35)--(2.4,-.35);\node[below]at(1.2,-.7){$ b\ll\sqrt{km} $};\draw[->,thick](3.5,0)--(5,0)node[right]{$ F_0\cos\omega t $};\draw[->](2.4,-1.2)--(4,-1.2)node[right]{$x$};
}

\begin{solution}
图\ref{fig:audit-m-q-010}的阻尼器给出力$-b\dot x$。判断它能否忽略，需比较其相量与惯性、弹性合力的相量。解图\ref{fig:class-1-3}所用的精确稳态幅值为
\[
 A=\frac{F_0/k}{\sqrt{(1-\Omega^2)^2+4\zeta^2\Omega^2}}.
\]
远离$\Omega=1$时，可提取$|1-\Omega^2|$，写成
\[
 A=\frac{F_0/k}{|1-\Omega^2|}
 \left[1+\left(\frac{2\zeta\Omega}{1-\Omega^2}\right)^2\right]^{-1/2}.
\]
因此无阻尼近似要求$2\zeta\Omega/|1-\Omega^2|\ll1$，被忽略的相对修正以这个比值的平方起始。
在$\Omega=1+\varepsilon$附近，$|1-\Omega^2|\simeq2|\varepsilon|$；
若$|\varepsilon|=O(\zeta)$，阻尼项与失谐项同阶，近似就失效。
特别在$\Omega=1$，精确幅值为$F_0/(b\omega_0)$，而图中无阻尼虚线发散。

\begin{figure}[H]\centering
\begin{tikzpicture}[x=3.5cm,y=.48cm]
\fill[black!8](.9,0)rectangle(1.1,6);\begin{scope}\clip(0,0)rectangle(2,6);
\draw[bookcurve,domain=0:2,samples=180]plot(\x,{1/sqrt((1-(\x)^2)^2+.0016*(\x)^2)});
\draw[dashed,domain=0:.98,samples=110]plot(\x,{1/abs(1-(\x)^2)});\draw[dashed,domain=1.02:2,samples=110]plot(\x,{1/abs(1-(\x)^2)});\end{scope}
\draw[bookaxis](0,0)--(2.2,0)node[right]{$\Omega$};\draw[bookaxis](0,0)--(0,6.8)node[above]{$A/(F_0/k)$};\node[below]at(1,0){$1$};
\node[above]at(1,6){峰顶截去};\node at(1.75,4.3){实线：有阻尼};\node at(1.75,3.3){虚线：忽略阻尼};
\end{tikzpicture}
\caption{随堂练习3解图：示取 $\zeta=0.02$。共振附近的浅灰频带需保留阻尼，远离该频带时两曲线接近；阴影边界仅示意频带量级。}\label{fig:class-1-3}
\end{figure}
\end{solution}

\section{交流电路中的复阻抗}
\subsection{从元件方程推导阻抗}
仍使用 $\ee^{\ii\omega t}$ 的峰值相量约定。电阻满足 $u=Ri$，所以 $Z_R=R$；电感满足 $u=L\dot i$，所以 $Z_L=\ii\omega L$；电容满足 $i=C\dot u$，所以 $Z_C=1/(\ii\omega C)$。对 $\omega>0$，后者等于 $-\ii/(\omega C)$。

电感电压超前电流，对应阻抗的正虚部；电容电流超前电压，对应阻抗的负虚部。判断方向时应从 $\widetilde U=Z\widetilde I$ 出发，避免仅凭一个相位角的正负猜测。

串联电路阻抗相加，串联 RLC 满足
\begin{equation}
Z=R+\ii\left(\omega L-\frac1{\omega C}\right),\qquad \widetilde I=\frac{\widetilde U}{Z}.
\end{equation}
并联电路则电压相同，电流相加，因此\textbf{导纳} $Y=1/Z$ 相加，而不是阻抗直接相加。两类连接必须从公共变量判断。

\begin{example}
$R=30\,\Omega,L=0.20\,\mathrm H,C=100\,\mu\mathrm F$ 串联，电压 $u=10\cos100t\,\mathrm V$。求稳态电流与平均功率。
\end{example}
% 题设配图；与紧随其后的解答分别编号。
\ExpansionDiagram{fig:audit-m-q-011}{串联RLC回路中三个元件共用电流}{
\draw(0,0)--(0,.55);\draw(0,1.45)--(0,2)--(.7,2);\draw(0,1)circle(.45);\node[left]at(-.5,1){$u(t)$};\draw(1.7,2)--(2.4,2);\draw(0.7,1.87)rectangle(1.7,2.13);\node[above]at(1.2,2.45){$ R $};\draw(3.4,2)--(4.1,2);\draw[exspring](2.4,2)--(3.4,2);\node[above]at(2.9,2.45){$ L $};\draw(5.1,2)--(5.8,2);\draw(4.1,2)--(4.4799999999999995,2);\draw(4.72,2)--(5.1,2);\draw(4.4799999999999995,1.65)--(4.4799999999999995,2.35);\draw(4.72,1.65)--(4.72,2.35);\node[above]at(4.6,2.45){$ C $};\draw(5.8,2)--(5.8,0)--(0,0);\draw[->](2,0)--(1,0)node[midway,below]{$I$};
}

\begin{solution}
图\ref{fig:audit-m-q-011}与图\ref{fig:orig-008}中的三个元件串联，共用电流，因此总阻抗等于各元件阻抗之和。图\ref{fig:eg-1-7}将总阻抗的实部和虚部画成两条直角边。驱动角频率为 $100\,\mathrm{rad/s}$，电感、电容电抗分别为
\[
\omega L=20\,\Omega,\qquad (\omega C)^{-1}=100\,\Omega.
\]
故 $Z=30+\ii(20-100)=30-80\ii\,\Omega$。以电压为实相量 $U=10\,\mathrm V$，
\[
I=\frac{U}{Z}=\frac{10(30+80\ii)}{7300}
=\frac{3+8\ii}{73}\,\mathrm A.
\]
电流峰值为 $|I|=1/\sqrt{73}\,\mathrm A$，相位为 $\phi=\arctan(8/3)>0$，所以
\[
i(t)=\frac1{\sqrt{73}}\cos(100t+\phi)\ \mathrm A.
\]
正相位说明电流超前电压，与总电抗为负相符。因题给10 V为峰值，平均电阻功率为
\[
\overline P=R I_{\rm rms}^2=\frac12R|I|^2=\frac{15}{73}\,\mathrm W.
\]
也可用 $U|I|\cos\phi/2$ 得到同值；若把峰值直接当有效值，会多算一倍。

\begin{figure}[H]\centering
\begin{tikzpicture}[scale=.85]
\draw[bookaxis](-.5,0)--(4.4,0)node[right]{$\Rea Z/\Omega$};\draw[bookaxis](0,-3.4)--(0,.8)node[above]{$\Ima Z/\Omega$};
\draw[bookaxis,thick](0,0)--(1.2,0)node[midway,above]{$30$};
\draw[bookaxis](1.2,0)--(1.2,-3.2)node[midway,right]{$-80$};
\draw[bookaxis,thick](0,0)--(1.2,-3.2)node[below right]{$Z=30-80\ii$};
\draw(.6,0)arc(0:-69.444:.6);\node[right]at(2.3,-1){$\Arg I=-\Arg Z>0$};
\end{tikzpicture}
\caption{例1.7：阻抗三角形的竖边是电感与电容电抗之差。电压相位取零后，电流相位等于阻抗相位的相反数。}\label{fig:eg-1-7}
\end{figure}

阻抗箭头指向第四象限，因 $I=U/Z$，电流相对实电压向相反方向转角，故是超前而非滞后。
\end{solution}

\subsection{电流共振与电容电压共振}
恒幅电压驱动下，电流幅值 $I_0=U_0/|Z|$ 在 $\omega L=1/(\omega C)$ 时最大，角频率为 $1/\sqrt{LC}$。但电容电压幅值为 $U_{C0}=I_0/(\omega C)$，多出的 $1/\omega$ 会改变峰值位置。

以电荷 $q$ 为变量，串联电路满足 $L\ddot q+R\dot q+q/C=u$。这与受迫振动严格对应：$L\leftrightarrow m,R\leftrightarrow b,1/C\leftrightarrow k$。电容电压 $q/C$ 与电荷同频同相，其非零峰频率存在时为
\begin{equation}
\omega_{C,r}^2=\frac1{LC}-\frac{R^2}{2L^2}>0.
\end{equation}
电流对应速度，因此电流峰频率又与电荷、电容电压的峰频率不同。两种系统具有相同形式的微分方程，因而频率响应也可逐项对应。

\section{并联电路、复功率与相位三角形}
\subsection{串联相量图与并联导纳}
串联 RLC 以电流为共同参考，电阻电压沿实轴，电感电压沿正虚轴，电容电压沿负虚轴。总电压是三者的矢量和，电抗电压可能相互抵消，但每个元件的电压不必为零。
\begin{figure}[htbp]\centering
\begin{tikzpicture}[scale=.9]
\draw(0,0)--(0,1.1);\draw(0,1.5)circle(.4);\node at(0,1.5){$\sim$};\draw(0,1.9)--(0,2.5)--(1.1,2.5);
\draw(1.1,2.25)rectangle(2.1,2.75);\node[above]at(1.6,2.8){$R$};\draw(2.1,2.5)--(2.6,2.5);
\draw[decorate,decoration={coil,segment length=4pt,amplitude=3pt}](2.6,2.5)--(3.8,2.5);\node[above]at(3.2,2.9){$L$};
\draw(3.8,2.5)--(4.4,2.5);\draw(4.4,2.1)--(4.4,2.9) (4.65,2.1)--(4.65,2.9);\node[above]at(4.52,2.95){$C$};
\draw(4.65,2.5)--(5.2,2.5)--(5.2,0)--(0,0);\draw[bookaxis](.4,2.5)--(.9,2.5)node[midway,above]{$i$};
\node[left]at(-.45,1.5){$u(t)$};
\end{tikzpicture}
\caption{串联 RLC 的共同电流。元件理想化，导线电阻忽略。}
\label{fig:orig-008}
\end{figure}

\begin{figure}[htbp]\centering
\begin{tikzpicture}[scale=.72]
\draw[bookaxis](0,0)--(4.8,0)node[right]{$\widetilde I\text{ 的方向}$};\draw[bookaxis](0,-1.9)--(0,3.3);
\draw[bookaxis,thick](0,0)--(3,0)node[midway,below]{$U_R$};\draw[bookaxis](3,0)--(3,2.7)node[midway,left]{$U_L$};
\draw[bookguide](3,2.7)--(3.65,2.7) (3,1.1)--(3.65,1.1);\draw[bookaxis](3.65,2.7)--(3.65,1.1)node[midway,right]{$U_C$};\draw[bookaxis,thick](0,0)--(3,1.1)node[midway,above]{$U$};
\end{tikzpicture}
\caption{感性串联电路的电压相量；电容电压箭头向下，总电压超前电流。}
\label{fig:orig-009}
\end{figure}

并联电路以电压为共同参考，电流相加。理想并联 RC 导纳为 $Y=1/R+\ii\omega C$，并联 RLC 导纳为 $Y=1/R+\ii(\omega C-1/\omega L)$。导纳虚部是正时，电流超前电压；它与串联阻抗虚部的相位判据方向不同。

\subsection{复功率的因子与符号}
令电压、电流峰值相量为 $U,I$，定义复功率 $S=UI^*/2=P+\ii Q_{\rm re}$。其中 $P$ 是平均有功功率，$Q_{\rm re}$ 是无功功率，避免与品质因数混淆。若电压相位减去电流相位为 $\phi$，则
\[
P=\tfrac12U_0I_0\cos\phi,\qquad Q_{\rm re}=\tfrac12U_0I_0\sin\phi.
\]
采用有效值相量时不再有因子 $1/2$。共轭也不可省去，否则结果会依赖共同的时间相位参考，而物理平均功率不应依赖参考时刻。

\begin{example}
并联 RC 取 $R=100\,\Omega,C=20\,\mu\mathrm F,\omega=1000\,\mathrm{rad/s}$，电压峰值 $10\,\mathrm V$，相位为零。求稳态总电流及功率。
\end{example}
% 题设配图；与紧随其后的解答分别编号。
\ExpansionDiagram{fig:audit-m-q-012}{并联RC回路中两支路共用电压}{
\draw(-.8,0)--(-.8,.55);\draw(-.8,1.45)--(-.8,2)--(5,2);\draw(-.8,0)--(5,0);\draw(-.8,1)circle(.45);\node[left]at(-1.3,1){$u(t)$};\draw(1.0,0)--(1.0,.6);\draw(1.0,1.4)--(1.0,2);\draw(0.87,.6)rectangle(1.13,1.4);\node[right]at(1.3,1){$ R $};\draw(2.7,0)--(2.7,.6);\draw(2.7,1.4)--(2.7,2);\draw(2.7,.6)--(2.7,.88);\draw(2.7,1.12)--(2.7,1.4);\draw(2.4000000000000004,.88)--(3.0,.88);\draw(2.4000000000000004,1.12)--(3.0,1.12);\node[right]at(3.0,1){$ C $};
}

\begin{solution}
图\ref{fig:audit-m-q-012}中电阻、电容两支路共用电压，须分别求电流再相加。
图\ref{fig:eg-1-8}以共同电压为实轴，把电阻电流与电容电流画成互相垂直的两边。

并联支路共用电压，应先相加电流或导纳。电阻导纳为 $1/R=0.01\,\mathrm S$，电容导纳为 $\ii\omega C=0.02\ii\,\mathrm S$，因此
\[
Y=0.01+0.02\ii,\qquad I=YU=0.1+0.2\ii\ \mathrm A.
\]
分别算得电阻电流 $0.1\cos\omega t$，电容电流 $-0.2\sin\omega t$；合成后为
\[
i(t)=\sqrt{0.05}\cos[1000t+\arctan(2)]\ \mathrm A.
\]
总电流峰值约0.2236 A。使用峰值相量时，复功率需带因子 $1/2$：
\[
S=\frac12UI^*=\frac12\times10(0.1-0.2\ii)=0.5-\ii\ \mathrm{VA}.
\]
所以有功功率为0.5 W，无功功率为 $-1\,\mathrm{var}$。另算电阻的 $U_0^2/(2R)=0.5\,\mathrm W$，与复功率的实部相符。电容的周期平均净吸收功率为零。

\begin{figure}[H]\centering
\begin{tikzpicture}[scale=1.25]
\draw[bookaxis](-.3,0)--(3.1,0)node[right]{与 $U$ 同相};\draw[bookaxis](0,-.2)--(0,2.65)node[above]{超前 $\pi/2$};
\draw[bookaxis,thick](0,0)--(1,0)node[midway,below]{$I_R=0.1$};
\draw[bookaxis](1,0)--(1,2)node[midway,right]{$I_C=0.2\ii$};
\draw[bookaxis,thick](0,0)--(1,2)node[above right]{$I=0.1+0.2\ii$};
\draw(.6,0)arc(0:63.435:.6);\node at(32:.95){$\phi$};
\node[right]at(1.7,.8){电流单位：A};
\end{tikzpicture}
\caption{例1.8：总电流是两支路电流的矢量和。水平投影 $I_R$ 决定有功功率，竖直分量来自电容。}\label{fig:eg-1-8}
\end{figure}

总电流的水平投影与电压同相。这一分量等于电阻电流，所以有功功率也等于电阻支路功率。
\end{solution}

\section{傅里叶展开与多频线性响应}
单一相量记录一个频率。周期外力的各频率分量由正交性确定，线性系统逐频响应后再叠加；它们一般不能合成一个固定相量。

\subsection{系数由正交性确定}
周期 $T$ 的信号，设 $\omega_f=2\pi/T$，可在适当条件下展开为
\[
f(t)=\frac{a_0}{2}+\sum_{n=1}^{\infty}[a_n\cos n\omega_ft+b_n\sin n\omega_ft].
\]
例如，周期函数分段连续、每段具有连续的一阶导数时，傅里叶级数在连续点收敛到函数值，在跳跃点收敛到左右极限的平均值。以下展开按这一常用的充分条件理解；平方可积函数则可在均方意义下展开，但不必在每一点都逐点收敛。

不同整数频率的正弦、余弦在一个完整周期上正交。其依据是对任意非零整数$\ell$，
\[\int_0^T\ee^{\ii\ell\omega_ft}\dd t
=\left[\frac{\ee^{\ii\ell\omega_ft}}{\ii\ell\omega_f}\right]_0^T=0,\]
而$\ell=0$时积分为$T$。将两三角函数写成复指数或用积化和差，就得到异频乘积积分为零、同频正弦和余弦平方积分为$T/2$。将有限展开式乘以$\cos m\omega_ft$并积分，只有对应余弦项留下$Ta_m/2$；正弦项同理。因此系数为
\[
a_n=\frac2T\int_0^Tf(t)\cos n\omega_ft\dd t,\qquad b_n=\frac2T\int_0^Tf(t)\sin n\omega_ft\dd t.
\]
常量项也必须单独积分：$a_0=2T^{-1}\int_0^Tf(t)\dd t$。正交性只有在规定的完整周期或相应内积区间上成立，不能对任意截取的一小段信号直接认为交叉项为零。

\begin{example}
对以 $2\pi$ 为周期、在 $(0,\pi)$ 为1、在 $(-\pi,0)$ 为 $-1$ 的方波，求傅里叶系数。
\end{example}
% 题设配图；与紧随其后的解答分别编号。
\ExpansionDiagram{fig:audit-m-q-013}{一周期方波的正、负平台}{
\draw[->](-.4,0)--(4,0)node[right]{$ t $};\draw[->](0,-1.8)--(0,1.8)node[above]{$ f $};\draw[thick](-3.14,-1)--(0,-1);\draw[thick](0,1)--(3.14,1);\draw[dashed](-3.14,-1)--(-3.14,1);\draw[dashed](0,-1)--(0,1);\node[below]at(-3.14,0){$-\pi$};\node[below]at(3.14,0){$\pi$};\node[left]at(0,1){1};\node[right]at(0,-1){$-1$};
}

\begin{solution}
图\ref{fig:audit-m-q-013}中的方波关于原点反对称，傅里叶系数可先按奇偶性筛选。
图\ref{fig:orig-010}中原方波关于原点反对称，先据此删去常量与余弦系数，再积分求正弦系数。

按 $f=a_0/2+\sum_{n\ge1}(a_n\cos nt+b_n\sin nt)$ 定义系数。方波为奇函数，乘偶函数 $1,\cos nt$ 后仍为奇函数，因此 $a_0=a_n=0$。

\begin{figure}[H]\centering
\begin{tikzpicture}[x=.85cm,y=1.2cm]
\draw[bookaxis](0,0)--(6.65,0)node[right]{$t$};\draw[bookaxis](0,-1.45)--(0,1.55)node[above]{$f$};
\draw[dashed,thick](0,1)--(3.14159,1) (3.14159,-1)--(6.28318,-1);
\draw[domain=0:6.28318,samples=350]plot(\x,{4/pi*(sin(deg(\x))+sin(deg(3*\x))/3+sin(deg(5*\x))/5)});
\node[left]at(0,1){$1$};\node[left]at(0,-1){$-1$};\node[below]at(3.14159,0){$\pi$};\node[below]at(6.28318,0){$2\pi$};
\end{tikzpicture}
\caption{方波与前三个非零谐波的部分和。虚线表示方波的两段，不把跳跃处连接成有限斜率。}
\label{fig:orig-010}
\end{figure}

对正弦系数利用奇偶性，将积分区间减半：
\begin{align*}
b_n&=\frac1\pi\int_{-\pi}^{\pi}f(t)\sin nt\dd t
=\frac2\pi\int_0^\pi\sin nt\dd t\\
&=\frac2\pi\left[-\frac{\cos nt}{n}\right]_0^\pi
=\frac2{n\pi}[1-\cos n\pi]
=\frac{2[1-(-1)^n]}{n\pi}.
\end{align*}
偶数 $n$ 时系数为零，奇数 $n$ 时为 $4/(n\pi)$，所以
\[
f(t)\sim\frac4\pi\left(\sin t+\frac{\sin3t}3+\frac{\sin5t}5+\cdots\right).
\]
在每个连续点，级数和等于原函数；在 $t=k\pi$ 的跳跃处，和为左右极限的平均值0。图\ref{fig:orig-010}的粗曲线仅保留前三个非零谐波，在跳跃点处过零，平顶附近仍有振荡。跳跃附近的过冲属于截断近似，不能当作原方波的取值。
\end{solution}

\subsection{把非正弦驱动拆成各频率}
线性振子对每个谐波独立响应。若驱动力的第 $n$ 次谐波相量为 $F_n$，稳态位移对应
\[
X_n=\frac{F_n}{k-mn^2\omega_f^2+\ii bn\omega_f}.
\]
将各频率响应取实部后相加，便得到总稳态。对于方波驱动，高次系数已经按 $1/n$ 下降；位移响应在高频又按 $1/n^2$ 降低，所以位移通常比驱动力平滑。若某一谐波接近系统共振，它会被选择性放大，不能只依据原驱动的幅值大小判断最终最显著的频率。

\subsection{频率分解与平均能量}
在一个完整周期上，正交性使交叉项积分为零，因此
\[
\frac1T\int_0^Tf^2\dd t=\frac{a_0^2}{4}+\frac12\sum_{n=1}^{\infty}(a_n^2+b_n^2)
\]
（在平方可积条件下理解）。这把时间域均方值与各谐波幅值联系起来。对于有限和可以直接逐项证明；对无穷级数交换求和与积分还需相应收敛条件。由均方值求有效值时，不能把所有谐波峰值直接相加。

\section{复可导、解析性与调和函数}
二维场的两个分量可合并成复函数，但沿不同方向求导必须得到同一结果，才具有复导数。以下先建立这一条件，再说明它与调和势的联系。

\subsection{复导数为何比实导数严格}
写 $f(z)=u(x,y)+\ii v(x,y)$。复导数定义为
\begin{equation}
f'(z_0)=\lim_{\Delta z\to0}\frac{f(z_0+\Delta z)-f(z_0)}{\Delta z}.
\end{equation}
实导数只从实轴两侧靠近，复导数却要求沿任意平面路径靠近都得到同一极限。沿实方向得到 $u_x+\ii v_x$，沿虚方向得到 $v_y-\ii u_y$，于是必须满足柯西--黎曼条件
\begin{equation}
u_x=v_y,\qquad u_y=-v_x.
\end{equation}
必要条件不能不加说明地当成充分条件。若 $u,v$ 的一阶偏导在邻域连续，实可微给出
\[
\Delta f=(u_x+\ii v_x)\Delta x+(u_y+\ii v_y)\Delta y+o(|\Delta z|).
\]
代入柯西--黎曼条件，右侧变为 $(u_x+\ii v_x)(\Delta x+\ii\Delta y)+o(|\Delta z|)$，再除以 $\Delta z$ 就证明复导数存在。这个推导同时说明了连续偏导等正则性条件的作用。

\begin{definition}
函数在某开集的每一点复可导，称其在该开集解析。
\end{definition}
解析是一个区域性概念。一个函数可能只在孤立的一点复可导，而不在任何邻域解析。

\begin{example}
比较 $z^2,\bar z,|z|^2$ 的复可导性。
\end{example}
% 题设配图；与紧随其后的解答分别编号。
\ExpansionDiagram{fig:audit-m-q-014}{同一点的实方向与虚方向增量}{
\draw[->](-.4,0)--(3,0)node[right]{$ x $};\draw[->](0,-0.3)--(0,2.5)node[above]{$ y $};\fill(.8,.7)circle(2pt)node[below right]{$z_0$};\draw[->,thick](.8,.7)--(2,.7)node[right]{$h\in\mathbb R$};\draw[->,thick](.8,.7)--(.8,1.9)node[above]{$\ii h$};
}

\begin{solution}
三个函数的实部、虚部均有连续偏导，可先用柯西--黎曼条件逐点判断。图\ref{fig:audit-m-q-014}给出沿实、虚两个方向的差商比较；对共轭函数，解答图\ref{fig:eg-1-10}将这两条趋近路径单独画出。孤立的候选可导点还可直接代入差商确认。

对 $z^2=(x^2-y^2)+2xy\ii$，有 $u_x=2x=v_y$、$u_y=-2y=-v_x$，处处满足条件。由 $u_x+\ii v_x$ 得导数 $2x+2y\ii=2z$，故它在整个复平面解析。

\begin{figure}[H]\centering
\begin{tikzpicture}[scale=1.25]
\draw[bookaxis](-.4,0)--(3,0)node[right]{$\Rea h$};\draw[bookaxis](0,-.3)--(0,2.1)node[above]{$\Ima h$};
\draw[bookaxis,thick](2.4,0)--(.5,0);\draw[bookaxis,thick](0,1.8)--(0,.45);
\node[above]at(1.65,.15){$h=\xi\to0$};\node[left]at(-.05,1.2){$h=\ii\eta\to0$};
\node[align=left,right]at(3.7,.8){沿实轴：$\bar h/h=1$\\沿虚轴：$\bar h/h=-1$};
\end{tikzpicture}
\caption{例1.10：共轭函数的差商随趋近方向而变。判断复可导时，所有复方向必须得到同一个极限。}\label{fig:eg-1-10}
\end{figure}

对 $\bar z=x-\ii y$，$u_x=1$ 而 $v_y=-1$，两者处处不相等。也可从差商 $\overline h/h$ 看出：沿实轴趋近为1，沿虚轴趋近为 $-1$，极限不存在，任何点都不复可导。

对 $|z|^2=x^2+y^2$，$v=0$，条件要求 $2x=2y=0$，只有原点可能可导。原点差商为
\[
\frac{|h|^2-0}{h}=\bar h\longrightarrow0,
\]
故 $f'(0)=0$。它在任何非空开区域内都不解析，因为开区域还含有非原点，而这些点不可导。

\end{solution}

\subsection{解析函数与二维无源场}
若 $u,v$ 有连续二阶偏导并满足柯西--黎曼条件，则
\[
u_{xx}+u_{yy}=v_{yx}-v_{xy}=0,
\]
虚部也有$v_{xx}+v_{yy}=-u_{yx}+u_{xy}=0$。两式中混合偏导的抵消使用了二阶偏导连续。满足拉普拉斯方程的函数称为调和函数，二维无电荷区域的静电势$\Phi$正满足这一方程，因此解析函数的实部可用来构造电势。

例如 $f(z)=z^2$ 的实部 $x^2-y^2$ 是调和函数。若 $\Phi=C(x^2-y^2)$，电场 $\vct E=-\nabla\Phi=(-2Cx,2Cy)$，散度为零，符合局部无电荷条件；要成为题目的完整物理解，还须满足给定边界条件。

\section{复平面上的轨迹与映射}
\subsection{用模和辐角描述几何约束}
把复数当作平面上的点，许多几何条件会变得紧凑。$|z-z_0|=R$ 表示圆心为 $z_0$ 的圆；$|z-z_1|=|z-z_2|$ 表示线段 $z_1z_2$ 的垂直平分线。由模平方展开，后一条件成为
\[
2\Rea[z(\bar z_2-\bar z_1)]=|z_2|^2-|z_1|^2,
\]
这是一个关于 $x,y$ 的一次方程。式子右侧不应在平移坐标时任意丢掉，它记录两端点到原点距离的差异。

加法的几何意义可以用平行四边形看得更清楚。若 $z_1,z_2$ 分别沿两个方向，合数的模不是一般的模长之和。只有同向时，三角不等式取等号；若两向量反向，其模差给出合数的大小。
\begin{figure}[htbp]\centering
\begin{tikzpicture}[scale=.85]
\draw[bookaxis](-.2,0)--(4.5,0)node[right]{$\Rea$};\draw[bookaxis](0,-.2)--(0,3.1)node[above]{$\Ima$};
\coordinate(A)at(2.6,.5);\coordinate(B)at(.9,1.7);\coordinate(C)at(3.5,2.2);
\draw[bookaxis,thick](0,0)--(A)node[midway,below]{$z_1$};\draw[bookaxis,thick](0,0)--(B)node[midway,left]{$z_2$};
\draw[bookguide](A)--(C)--(B);\draw[bookaxis,thick](0,0)--(C)node[above right]{$z_1+z_2$};
\end{tikzpicture}
\caption{复数加法的平行四边形。虚部是实数坐标，不能把虚部的符号误当成模长的符号。}
\label{fig:orig-011}
\end{figure}

\begin{example}
求满足 $|z-1|=2|z+1|$ 的点集。
\end{example}
% 题设配图；与紧随其后的解答分别编号。
\ExpansionDiagram{fig:audit-m-q-015}{两定点到动点的距离比}{
\draw[->](-3.2,0)--(3.2,0)node[right]{$x$};\fill(-1,0)circle(2pt)node[below]{$-1$};\fill(1,0)circle(2pt)node[below]{$ 1 $};\fill(-1.66667,1.33333)circle(2pt)node[above]{$P(z)$};\draw(-1,0)--(-1.66667,1.33333)node[midway,left]{$|z+1|$};\draw(1,0)--(-1.66667,1.33333)node[midway,right]{$|z-1|$};
}

\begin{solution}
图\ref{fig:audit-m-q-015}中两线段分别表示$|z-1|$与$|z+1|$，平方后可转为直角坐标方程。
图\ref{fig:eg-1-11}用两条线段表示题给距离比；配方后应得到同一圆心和半径。

两边都是距离，可以平方而不改变解集。令 $z=x+\ii y$，有
\[
(x-1)^2+y^2=4[(x+1)^2+y^2].
\]
展开并把同类项移到右侧，再除以3，得
\[
3x^2+10x+3y^2+3=0,
\qquad x^2+\frac{10}3x+y^2+1=0.
\]
对 $x$ 配方：
\[
\left(x+\frac53\right)^2+y^2=\frac{25}9-1=\frac{16}9.
\]
所以点集为圆心 $(-5/3,0)$、半径 $4/3$ 的圆。它与实轴的交点为 $-3$、$-1/3$；分别代入距离比，可得 $4/2=2$、$(4/3)/(2/3)=2$，两点均符合原条件。固定点 $z=-1$ 不在这个圆上，原式中距离为零的情形也没有被遗漏。

\begin{figure}[H]\centering
\begin{tikzpicture}[scale=1.12]
\draw[bookaxis](-3.4,0)--(1.65,0)node[right]{$x$};\draw[bookaxis](0,-1.65)--(0,2)node[above]{$y$};
\draw[bookcurve](-1.666667,0)circle(1.333333);\fill(-1.666667,0)circle(1.2pt)node[below left]{$C$};
\coordinate(P)at(-1.666667,1.333333);\fill(P)circle(1.5pt)node[above left]{$P$};
\fill(-1,0)circle(1.5pt)node[below right]{$-1$};\fill(1,0)circle(1.5pt)node[below]{$1$};
\draw(P)--(-1,0)node[midway,right]{$d$};\draw(P)--(1,0)node[midway,above right]{$2d$};
\draw[bookguide](-1.666667,0)--(P)node[midway,left]{$4/3$};
\end{tikzpicture}
\caption{例1.11：圆上的点到 $1,-1$ 的距离分别为 $2d,d$。标出圆心与半径后，可用两实轴交点核对配方结果。}\label{fig:eg-1-11}
\end{figure}

图中的圆心不是两个定点的中点。本题中动点到 $-1$ 的距离较小，所得圆也位于负实轴一侧。
\end{solution}

\subsection{分支、连续相位与相位展开}
主辐角方便输出单个数值，却不一定连续。例如一个点从负实轴上方绕到下方，主辐角会从接近 $\pi$ 跳到接近 $-\pi$。点的运动很平滑，这个跳变来自角度代表元的选择。研究连续相位时，可以每次补上或减去 $2\pi$，使相邻角度连续，这称为相位展开。

对数的单值解析分支必须去掉一条割线。取负实轴为割线时，$\log z=\ln|z|+\ii\theta$，$-\pi<\theta<\pi$。在割线两侧趋近同一个负实数，两个极限相差 $2\pi\ii$。不能把这两个侧极限误当成函数连续延拓后的同一个值。
\begin{figure}[htbp]\centering
\begin{tikzpicture}[scale=1]
\draw[bookaxis](-2.3,0)--(1.5,0)node[right]{$\Rea$};\draw[bookaxis](0,-1)--(0,1)node[above]{$\Ima$};
\draw[thick,dashed](-2,0)--(0,0);\fill(-1.4,.18)circle(1.5pt);\fill(-1.4,-.18)circle(1.5pt);
\node[above]at(-1.4,.32){$\theta\to\pi$};\node[below]at(-1.4,-.32){$\theta\to-\pi$};\node[below right]at(0,0){$0$};
\end{tikzpicture}
\caption{负实轴割线的两侧。主值跳变是分支选择的结果。}
\label{fig:orig-012}
\end{figure}

\subsection{解析映射怎样改变局部几何}
在 $z_0$ 附近，解析函数满足 $f(z_0+\Delta z)=f(z_0)+f'(z_0)\Delta z+o(|\Delta z|)$。若 $f'(z_0)\ne0$，主要作用是乘以一个非零复数，所以局部长度放大 $|f'(z_0)|$ 倍，方向旋转 $\arg f'(z_0)$，两条曲线的夹角保持不变。这就是保角性的局部来源。

以$w=z^2$为例，$u=x^2-y^2,v=2xy$。对非零常数$x=c$，消去$y$得到$u=c^2-v^2/(4c^2)$；对非零常数$y=c$，消去$x$得到$u=v^2/(4c^2)-c^2$，故网格映成两组抛物线。两条坐标轴则分别退化为实轴上的半直线。在非零点，映射后的网格仍正交；原点$f'(0)=0$，局部保角结论不适用。从原点出发的两射线经平方后辐角差加倍，也不违反非零点的局部结论。
\begin{figure}[htbp]\centering
\begin{tikzpicture}[scale=.95]
\begin{scope}
\draw[bookaxis](0,0)--(1.45,0)node[right]{$x$};\draw[bookaxis](0,0)--(0,1.45)node[above]{$y$};
\foreach\c in{.25,.5,.75,1}{\draw(\c,0)--(\c,1);\draw[dashed](0,\c)--(1,\c);}
\node at(.6,-.45){$z$ 平面};\end{scope}
\draw[bookaxis](2,.55)--(3,.55)node[midway,above]{$w=z^2$};
\begin{scope}[shift={(4.5,0)},x=.95cm,y=.95cm]
\draw[bookaxis](-1.2,0)--(1.4,0)node[right]{$u$};\draw[bookaxis](0,0)--(0,2.35)node[above]{$v$};
\foreach\c in{.25,.5,.75,1}{\draw[domain=0:1,samples=45]plot({\c*\c-\x*\x},{2*\c*\x});\draw[dashed,domain=0:1,samples=45]plot({\x*\x-\c*\c},{2*\c*\x});}
\node at(0,-.45){$w$ 平面};\end{scope}
\end{tikzpicture}
\caption{平方映射下的坐标网格。实线对应 $x$ 为常量，虚线对应 $y$ 为常量。}
\label{fig:orig-013}
\end{figure}

\section{复积分、奇点与留数（选读）}
围道积分把路径、方向和函数的奇点联系起来。用于实积分之前，先确定物理或几何允许的路径，再判断哪些奇点被包围及无穷远弧段能否略去。

\subsection{曲线积分从参数化开始}
对有向分段光滑曲线 $C:z=z(t)$，$a\le t\le b$，定义
\begin{equation}
\int_C f(z)\dd z=\int_a^b f(z(t))z'(t)\dd t.
\end{equation}
例如单位圆逆时针参数化为$z=\ee^{\ii t}$、$0\le t\le2\pi$，则$\dd z=\ii\ee^{\ii t}\dd t$，因此
\[\oint_C\frac{\dd z}{z}
=\int_0^{2\pi}\frac{\ii\ee^{\ii t}}{\ee^{\ii t}}\dd t
=\ii[t]_0^{2\pi}=2\pi\ii.\]
闭路积分并不自动为零，$1/z$在内部原点有奇点。

若 $f=u+\ii v$ 且 $\dd z=\dd x+\ii\dd y$，则
\[
f\dd z=(u\dd x-v\dd y)+\ii(v\dd x+u\dd y).
\]
在偏导连续、函数于简单闭曲线及内部解析的条件下，对实、虚部用格林公式，面积积分分别为 $-v_x-u_y$ 与 $u_x-v_y$，由柯西--黎曼条件均为零。这给出柯西积分定理的一种直接解释：解析性导致不同微元贡献相互抵消。更一般的定理不要求预先假设连续偏导，但这里的版本已经足以处理本书的例题。

\subsection{柯西公式与简单极点}
若 $f$ 在包含简单闭曲线及其内部的开集解析，$a$ 在内部，曲线逆时针，则
\begin{equation}
f(a)=\frac1{2\pi\ii}\oint_C\frac{f(z)}{z-a}\dd z.
\end{equation}
取内部小圆$C_\rho:z=a+\rho\ee^{\ii t}$，两条边界均约定逆时针。对外路径与小圆之间的无奇点区域用积分定理，得到外路径积分等于小圆积分；小圆若按该区域内边界定向，本应顺时针，移到等式右侧后方向反转。因此
\[\oint_C\frac{f(z)}{z-a}\dd z
=\oint_{C_\rho}\frac{f(z)}{z-a}\dd z
=\ii\int_0^{2\pi}f(a+\rho\ee^{\ii t})\dd t.\]
将$f(a)$分出，常数项积分为$2\pi\ii f(a)$，余项模不超过
$2\pi\max_{|z-a|=\rho}|f(z)-f(a)|$，由连续性在$\rho\to0$时趋零，即得公式。

若某孤立奇点附近有洛朗展开 $f(z)=\sum_{n=-\infty}^{\infty}c_n(z-a)^n$，其中 $c_{-1}$ 称为留数。圆周积分中，只有 $n=-1$ 项保留非零贡献。对于简单极点，留数可直接算为
\begin{equation}
\Res(f,a)=\lim_{z\to a}(z-a)f(z).
\end{equation}
若路径内部只有有限个孤立奇点，路径上无奇点，且其他地方解析，则留数定理给出
\begin{equation}
\oint_C f(z)\dd z=2\pi\ii\sum_j\Res(f,a_j).
\end{equation}
曲线反向会改变符号，外部奇点不计入求和。

\begin{example}
求 $I=\int_{-\infty}^{\infty}\dd x/(x^2+a^2)$，其中 $a>0$。
\end{example}
% 题设配图；与紧随其后的解答分别编号。
\ExpansionDiagram{fig:audit-m-q-016}{上半平面的极点与闭合积分路径}{
\draw[->](-.4,0)--(3.2,0)node[right]{$ \Rea z $};\draw[->](0,-1.6)--(0,2.7)node[above]{$ \Ima z $};\draw[very thick,->](-2.5,0)--(2.5,0);\draw[->,thick](2.5,0)arc(0:175:2.5);\draw(175:2.5)arc(175:180:2.5);\fill(0,1)circle(2pt)node[right]{$\ii a$};\fill(0,-1)circle(2pt)node[right]{$-\ii a$};\node[below]at(2.5,0){$R$};\node[below]at(-2.5,0){$-R$};
}

\begin{solution}
图\ref{fig:audit-m-q-016}与图\ref{fig:orig-014}采用同一逆时针上半圆闭路。实线段对应目标积分，内部极点为$\ii a$；补入的圆弧需另行估计。

被积函数在实轴上无奇点，且在无穷远按 $x^{-2}$ 衰减。取从 $-R$ 到 $R$ 的实线段和上半圆弧组成逆时针闭路，$R>a$。复函数
\[
f(z)=\frac1{(z-\ii a)(z+\ii a)}
\]
在闭路内只有 $z=\ii a$ 一个极点，留数为
\[
\lim_{z\to\ii a}(z-\ii a)f(z)=\frac1{2\ii a}.
\]
留数定理给出实线段积分与弧积分之和等于 $2\pi\ii/(2\ii a)=\pi/a$。还需证明补入的弧不贡献极限。在 $|z|=R$ 上，
\[
|z^2+a^2|\ge R^2-a^2,
\qquad\left|\int_{\rm arc}f(z)\dd z\right|
\le\frac{\pi R}{R^2-a^2}\longrightarrow0.
\]
因此$I=\pi/a$。实变量方法令$x=au$、$\dd x=a\dd u$，全轴边界不变，也可得到
\[I=\frac1a\int_{-\infty}^{\infty}\frac{\dd u}{1+u^2}
=\frac1a[\arctan u]_{-\infty}^{\infty}
=\frac1a\left(\frac\pi2+\frac\pi2\right)=\frac\pi a.\]
因子$1/a$同时给出积分随$a$变化的尺度。

\begin{figure}[H]\centering
\begin{tikzpicture}[scale=1,>=Stealth]
\draw[bookaxis](-2.5,0)--(2.7,0)node[right]{$\Rea z$};\draw[bookaxis](0,-1)--(0,2.5)node[above]{$\Ima z$};
\draw[thick](2,0)arc(0:180:2);\draw[bookaxis](.35,0)--(.95,0);\draw[bookaxis](60:2)arc(60:80:2);
\fill(0,.9)circle(1.7pt)node[right]{$\ii a$};\fill(0,-.9)circle(1.7pt)node[right]{$-\ii a$};
\node[below]at(-2,0){$-R$};\node[below]at(2,0){$R$};
\end{tikzpicture}
\caption{实轴积分的上半圆闭合路径。内部极点、方向和圆弧估计缺一不可。}
\label{fig:orig-014}
\end{figure}

图上 $\ii a$ 在内部、$-\ii a$ 在外部，因此留数和只取上方极点；弧长因子 $\pi R$ 则来自同一条圆弧。
\end{solution}

\begin{example}
求 $J(b)=\int_{-\infty}^{\infty}\cos(bx)\dd x/(x^2+a^2)$，$a>0,b\in\mathbb R$。
\end{example}
% 题设配图；与紧随其后的解答分别编号。
\ExpansionDiagram{fig:audit-m-q-017}{上半平面的极点与闭合积分路径}{
\draw[->](-.4,0)--(3.2,0)node[right]{$ \Rea z $};\draw[->](0,-1.6)--(0,2.7)node[above]{$ \Ima z $};\draw[very thick,->](-2.5,0)--(2.5,0);\draw[->,thick](2.5,0)arc(0:175:2.5);\draw(175:2.5)arc(175:180:2.5);\fill(0,1)circle(2pt)node[right]{$\ii a$};\fill(0,-1)circle(2pt)node[right]{$-\ii a$};\node[below]at(2.5,0){$R$};\node[below]at(-2.5,0){$-R$};
}

\begin{solution}
图\ref{fig:audit-m-q-017}标出两个极点；选哪一个计入闭路，还取决于指数因子在哪个半平面衰减。
沿用图\ref{fig:eg-1-13}的路径。先按 $b$ 的符号选择指数衰减的半平面，才能控制圆弧项。

先取 $b\ge0$，将余弦换成复指数的实部，考虑
\[
f(z)=\frac{\ee^{\ii bz}}{z^2+a^2}.
\]
在上半圆中，$|\ee^{\ii bz}|=\ee^{-b\Ima z}\le1$，因此例1.12的弧积分上界仍可使用。唯一内部极点 $\ii a$ 处的留数为
\[
\Res(f,\ii a)=\frac{\ee^{\ii b(\ii a)}}{2\ii a}
=\frac{\ee^{-ab}}{2\ii a}.
\]
令半径趋于无穷，得到复实轴积分 $\pi\ee^{-ab}/a$。其虚部来自 $\sin(bx)/(x^2+a^2)$，是奇函数的全轴积分，确实为零。取实部便得 $b\ge0$ 时的 $J(b)=\pi\ee^{-ab}/a$。

\begin{figure}[H]\centering
\begin{tikzpicture}[scale=1,>=Stealth]
\draw[bookaxis](-2.5,0)--(2.7,0)node[right]{$\Rea z$};\draw[bookaxis](0,-1)--(0,2.5)node[above]{$\Ima z$};
\draw[thick](2,0)arc(0:180:2);\draw[bookaxis](.35,0)--(.95,0);\draw[bookaxis](60:2)arc(60:80:2);
\fill(0,.9)circle(1.7pt)node[right]{$\ii a$};\fill(0,-.9)circle(1.7pt)node[right]{$-\ii a$};
\node[below]at(-2,0){$-R$};\node[below]at(2,0){$R$};
\end{tikzpicture}
\caption{例1.13解图：$b\ge0$ 时在上半平面闭合，指数模为 $\ee^{-b\Ima z}\le1$；只取 $\ii a$ 处的留数。}
\label{fig:eg-1-13}
\end{figure}

原积分关于 $b$ 为偶函数，故全部实数 $b$ 的结果为
\[
J(b)=\frac\pi a\ee^{-a|b|}.
\]
在 $b=0$ 时恢复例1.12。不能对 $\cos(bz)$ 直接使用模不超过1的估计，因为复余弦含有增长的指数项；这里先选衰减的复指数，才使所选闭路有效。

在所选上半圆上 $\Ima z\ge0$，这使 $b\ge0$ 时的指数模不超过1；图中的几何位置正是闭路选择的依据。
\end{solution}

\section{复积分的路径选择与估计（选读）}
\subsection{方向、奇点与大圆弧}
使用留数法前应先画路径。对 $1/(z^2+a^2)$，上半平面只有 $\ii a$，下半平面只有 $-\ii a$。从 $-R$ 沿实轴走向 $R$，再走上半圆回去是逆时针；若走下半圆回去则是顺时针，留数定理的符号也随之改变。

对包含 $\ee^{\ii bz}$ 的积分，$b>0$ 选上半平面，因为模为 $\ee^{-b\Ima z}$；$b<0$ 则选下半平面。选择路径是为了同时包含容易计算的奇点，并使补加路径的贡献趋于零。只算留数而没有证明补加部分消失，尚未完成实积分的推导。

\subsection{高阶极点与留数提取}
若 $f(z)=g(z)/(z-a)^m$，其中 $g$ 在 $a$ 解析且 $g(a)\ne0$，则将 $g$ 在 $a$ 泰勒展开，只有 $m-1$ 次项与分母结合成 $(z-a)^{-1}$，因此
\[
\Res(f,a)=\frac{g^{(m-1)}(a)}{(m-1)!}.
\]
这一公式取出了展开中负一次幂的系数，与留数的定义一致。

\begin{example}
求 $\int_{-\infty}^{\infty}\dd x/(x^2+a^2)^2$，$a>0$。
\end{example}
% 题设配图；与紧随其后的解答分别编号。
\ExpansionDiagram{fig:audit-m-q-018}{上半平面的极点与闭合积分路径}{
\draw[->](-.4,0)--(3.2,0)node[right]{$ \Rea z $};\draw[->](0,-1.6)--(0,2.7)node[above]{$ \Ima z $};\draw[very thick,->](-2.5,0)--(2.5,0);\draw[->,thick](2.5,0)arc(0:175:2.5);\draw(175:2.5)arc(175:180:2.5);\fill(0,1)circle(2pt)node[right]{$\ii a$};\fill(0,-1)circle(2pt)node[right]{$-\ii a$};\node[below]at(2.5,0){$R$};\node[below]at(-2.5,0){$-R$};
}

\begin{solution}
图\ref{fig:audit-m-q-018}仍采用上半圆闭路，但分母平方使内部极点$\ii a$变为二阶极点。
路径仍见图\ref{fig:eg-1-14}，极点位置没有改变，改变的是上方极点的阶数。

仍用上半圆闭路。分解为
\[
f(z)=\frac1{(z-\ii a)^2}\,g(z),\qquad g(z)=(z+\ii a)^{-2}.
\]
内部极点 $\ii a$ 的阶数为2。把解析因子展开为 $g(\ii a)+g'(\ii a)(z-\ii a)+\cdots$，可见乘上二阶主因子后，负一次幂的系数正是 $g'(\ii a)$。计算得
\[
g'(z)=-2(z+\ii a)^{-3},\qquad
\Res(f,\ii a)=g'(\ii a)
=-\frac2{(2\ii a)^3}=\frac1{4\ii a^3}.
\]
圆弧上 $|f(z)|\le(R^2-a^2)^{-2}$，故弧积分模不超过 $\pi R/(R^2-a^2)^2$，趋于零。于是
\[
\int_{-\infty}^{\infty}\frac{\dd x}{(x^2+a^2)^2}
=2\pi\ii\frac1{4\ii a^3}=\frac\pi{2a^3}.
\]
另对例1.12的$I(a)=\pi/a$求参数导数。链式法则给出
\[\frac{\partial}{\partial a}(x^2+a^2)^{-1}
=-(x^2+a^2)^{-2}(2a),\qquad
I'(a)=\pi(-a^{-2})=-\frac\pi{a^2}.\]
所以
\[-2a\int_{-\infty}^{\infty}\frac{\dd x}{(x^2+a^2)^2}
=-\frac\pi{a^2},\]
除以$-2a$再次得到$\pi/(2a^3)$。交换求导与积分须在固定$a>0$的邻域验证可积控制；这里没有移动的实轴奇点。

\begin{figure}[H]\centering
\begin{tikzpicture}[scale=1,>=Stealth]
\draw[bookaxis](-2.5,0)--(2.7,0)node[right]{$\Rea z$};\draw[bookaxis](0,-1)--(0,2.5)node[above]{$\Ima z$};
\draw[thick](2,0)arc(0:180:2);\draw[bookaxis](.35,0)--(.95,0);\draw[bookaxis](60:2)arc(60:80:2);
\fill(0,.9)circle(1.7pt)node[right]{$\ii a$};\fill(0,-.9)circle(1.7pt)node[right]{$-\ii a$};
\node[below]at(-2,0){$-R$};\node[below]at(2,0){$R$};
\end{tikzpicture}
\caption{例1.14解图：相同上半圆只包住 $\ii a$，此处已是二阶极点；圆弧上被积函数按 $R^{-4}$ 衰减。}
\label{fig:eg-1-14}
\end{figure}

闭路图决定取哪个极点，局部展开决定取哪个系数。二阶极点要提取负一次幂项，不能照搬简单极点的乘极限公式。
\end{solution}

\clearpage
\section*{习题与解答}\addcontentsline{toc}{section}{习题与解答}
\begin{enumerate}[label=1.\arabic*]
\item \begin{minipage}[t]{\linewidth}
【基础】求 $z^4=-16$ 的全部根，并求四根的和与积。
\end{minipage}
% 题设配图；与紧随其后的解答分别编号。
\ExpansionDiagram{fig:audit-m-q-019}{四次方程中未知根的模与辐角}{
\draw[->](-.4,0)--(3.2,0)node[right]{$ \Rea z $};\draw[->](0,-0.5)--(0,2.1)node[above]{$ \Ima z $};\draw[dashed](0,0)circle(1.5);\draw[->,thick](0,0)--(1.15,.96)node[above right]{$z=r\ee^{\ii\theta}$};\draw[->](.5,0)arc(0:40:.5);\node[right]at(.65,.27){$\theta$};\node[right]at(3.6,1){$z^4=-16$};
}

\begin{solution}
图\ref{fig:audit-m-q-019}中未知根写成$r\ee^{\ii\theta}$，模方程和相位方程应分别求解。
按图\ref{fig:hw-1-1}先确定根的模与起始辐角，再检查四个根的对称关系。

设 $z=r\ee^{\ii\theta}$。模长方程为 $r^4=16$，故 $r=2$；相位方程为 $4\theta=\pi+2k\pi$。取 $k=0,1,2,3$，得到
\[
z_k=2\ee^{\ii(\pi/4+k\pi/2)}.
\]
换成直角形式即 $\sqrt2(1+\ii)$、$\sqrt2(-1+\ii)$、$\sqrt2(-1-\ii)$、$\sqrt2(1-\ii)$。两两相对，故根的和为0。四个辐角之和为 $4\pi$，所以乘积为 $2^4\ee^{4\pi\ii}=16$。

\begin{figure}[H]\centering
\begin{tikzpicture}[scale=1.1]
\draw[bookaxis](-2.2,0)--(2.5,0)node[right]{$\Rea z$};\draw[bookaxis](0,-2.2)--(0,2.5)node[above]{$\Ima z$};
\draw[bookguide](0,0)circle(2);\foreach\a/\lab in{45/z_0,135/z_1,225/z_2,315/z_3}{\draw[bookguide](0,0)--(\a:2);\fill(\a:2)circle(1.7pt)node[shift={(\a:.3)}]{$\lab$};}
\draw[->](.55,0)arc(0:45:.55);\node at(22.5:.9){$\pi/4$};\node[below right]at(2,0){$2$};
\end{tikzpicture}
\caption{题1.1解图：四根在半径2的圆上每隔 $\pi/2$ 排列；相对的根相消，解释了根和为零。}\label{fig:hw-1-1}
\end{figure}

多项式 $z^4+16$ 的三次项系数为零，根的和应为零；四次为偶数，根的积应为常数项16。这两项检查同时成立。

图上每个点都有一个关于原点对称的根，故根和为零；乘积仍须按全部四个辐角计算，不能只看模。
\end{solution}

\item \begin{minipage}[t]{\linewidth}
【基础】合成 $2\cos\omega t+2\cos(\omega t+2\pi/3)$，并同时检查初始值和初始导数。

\end{minipage}
% 题设配图；与紧随其后的解答分别编号。
\ExpansionDiagram{fig:audit-m-q-020}{两个等长相量的夹角为$2\pi/3$}{
\draw[->](-.4,0)--(3.5,0)node[right]{$ \Rea z $};\draw[->](0,-0.7)--(0,2.432)node[above]{$ \Ima z $};\draw[->,thick](0,0)--(2,0)node[above right]{$ 2 $};\draw[->,thick](0,0)--(-1,1.732)node[above right]{$ 2\ee^{2\pi\ii/3} $};
}

\begin{solution}
图\ref{fig:audit-m-q-020}中的两相量等长，第二项相对第一项转过$2\pi/3$。
在图\ref{fig:ex-1-2}中补成平行四边形，粗对角线就是要计算的合相量。

以第一项为相位参考，合相量为
\[
X=2+2\ee^{2\pi\ii/3}
=2+2(-1/2+\sqrt3\ii/2)=1+\sqrt3\ii.
\]
其模为2，辐角为 $\pi/3$，因此 $x=2\cos(\omega t+\pi/3)$。原式初值为 $2+2\cos(2\pi/3)=1$，合成式为 $2\cos(\pi/3)=1$。

\begin{figure}[H]\centering
\begin{tikzpicture}[scale=1.15]
\draw[bookaxis](-1.7,0)--(3,0)node[right]{$\Rea$};\draw[bookaxis](0,-.3)--(0,2.7)node[above]{$\Ima$};
\draw[bookaxis](0,0)--(2,0)node[below]{$X_1=2$};\draw[bookaxis](0,0)--(120:2)node[above left]{$X_2$};
\draw[bookguide](120:2)--(1,1.732051)--(2,0);\draw[bookaxis,thick](0,0)--(1,1.732051)node[above right]{$X=1+\sqrt3\ii$};
\draw[bookguide](1,0)--(1,1.732051);\node[below]at(1,0){$1$};
\draw[->](.55,0)arc(0:60:.55);\node at(30:.9){$\pi/3$};
\end{tikzpicture}
\caption{题1.2解图：虚线补成平行四边形，粗箭头为合相量。合相位为 $+\pi/3$，实轴投影给出 $x(0)=1$。}\label{fig:ex-1-2}
\end{figure}

求导后，原式在0处为 $-2\omega\sin(2\pi/3)=-\sqrt3\omega$；合成式给出 $-2\omega\sin(\pi/3)$，两者相等。若把合相位误写成 $-\pi/3$，初始位移检查仍会通过，速度检查则会发现错误。

合相量指向第一象限，实轴投影为1。其随时间逆时针转动时，实投影初始减小，对应 $\dot x(0)<0$。
\end{solution}

\item \begin{minipage}[t]{\linewidth}
【基础】设 $x=A\cos\omega t$。写出 $x^2$ 的常量部分与二倍频部分，说明为何不能把相量平方后直接取实部当作 $x^2$。
\end{minipage}
% 题设配图；与紧随其后的解答分别编号。
\ExpansionDiagram{fig:audit-m-q-021}{平方运算作用于完整实信号}{
\draw[->](-.4,0)--(6.7,0)node[right]{$ \omega t $};\draw[->](0,-1.5)--(0,1.8)node[above]{$ x/A $};\draw[thick,domain=0:6.283,samples=100]plot(\x,{cos(\x r)});\draw[dashed](0,1)--(6.3,1);\draw[dashed](0,-1)--(6.3,-1);
}

\begin{solution}
图\ref{fig:audit-m-q-021}中实信号每半周期变号，平方后这些正、负部分对平均值均有贡献。
图\ref{fig:hw-1-3}把平方后的平均高度与重复周期分开显示，可据此核对恒等式。

先对实函数使用半角恒等式：
\[
x^2=A^2\cos^2\omega t=\frac{A^2}{2}+\frac{A^2}{2}\cos2\omega t.
\]
常量部分和二倍频部分的幅值都是 $A^2/2$。直接平方相量 $A$ 并取二倍频实部，会得到 $A^2\cos2\omega t$，与上式不符。

\begin{figure}[H]\centering
\begin{tikzpicture}[x=.8cm,y=1.8cm]
\draw[bookaxis](0,0)--(6.7,0)node[right]{$\omega t$};\draw[bookaxis](0,-.15)--(0,1.3)node[above]{$x^2/A^2$};
\draw[bookcurve,domain=0:6.283185,samples=160]plot(\x,{cos(\x r)^2});\draw[bookguide](0,.5)--(6.3,.5)node[right]{$1/2$};
\foreach\t/\lab in{0/0,3.141593/\pi,6.283185/2\pi}{\node[below]at(\t,0){$\lab$};}
\node at(3.1,1.2){一个原周期内出现两个峰};
\end{tikzpicture}
\caption{题1.3解图：平方后曲线始终非负，平均值为 $A^2/2$，起伏角频率为 $2\omega$。}\label{fig:hw-1-3}
\end{figure}

对一般复相量，令 $u=X\ee^{\ii\omega t}$，用 $\Rea u=(u+\bar u)/2$ 展开：
\[
(\Rea u)^2=\frac{u^2+2u\bar u+\bar u^2}{4}
=\frac{|X|^2}{2}+\frac12\Rea(X^2\ee^{2\ii\omega t}).
\]
中间的 $u\bar u$ 项正是直流部分的来源。平方是非线性运算，不能按单个同频相量的线性规则处理。

曲线的平均高度为 $1/2$，对应展开中的常量项；二倍频项描述它相对这条中线的起伏。
\end{solution}

\item \begin{minipage}[t]{\linewidth}
【基础】对 $m=1,b=2,k=10,F_0=6,\omega=2$（SI 单位），求稳态幅值、滞后相位、位移峰频率和平均耗散功率。

\end{minipage}
% 题设配图；与紧随其后的解答分别编号。
\ExpansionDiagram{fig:audit-m-q-022}{给定驱动下的位移和阻尼速度方向}{
\draw[very thick](0,-.8)--(0,.9);\draw[exspring](0,.35)--(2.4,.35);\node[above]at(1.2,.45){$ 10 $};\draw[thick](2.4,-.65)rectangle(3.5,.65);\node at(2.95,0){$ 1 $};\draw(0,-.35)--(.8,-.35);\draw(.8,-.6)--(.8,-.1)--(1.7,-.1);\draw(.8,-.6)--(1.7,-.6);\draw(1.35,-.55)--(1.35,-.15);\draw(1.35,-.35)--(2.4,-.35);\node[below]at(1.2,-.7){$ 2 $};\draw[->,thick](3.5,0)--(5,0)node[right]{$ 6\cos2t $};\draw[->](2.4,-1.2)--(4,-1.2)node[right]{$x$};
}

\begin{solution}
图\ref{fig:audit-m-q-022}中弹簧和阻尼器接同一固定端，复刚度分别包含惯性、弹性和阻尼项。
按图\ref{fig:ex-1-4}的同一位移正方向列方程，弹簧形变为 $x$，阻尼器相对速度为 $\dot x$。

从方程 $\ddot x+2\dot x+10x=6\cos2t$ 出发，设 $x_s=\Rea(X\ee^{2\ii t})$，得到
\[
(10-4+4\ii)X=6,\qquad X=\frac6{6+4\ii}.
\]
因此幅值 $A=6/\sqrt{52}=3/\sqrt{13}\,\mathrm m$，滞后角 $\delta=\arctan(2/3)$，即 $x_s=A\cos(2t-\delta)$。

\begin{figure}[H]\centering
\begin{tikzpicture}[scale=.9]
\draw[thick](0,-.8)--(0,.9);
\draw[exspring](0,.35)--(3,.35);\node[above]at(1.5,.55){$k$};
\draw(0,-.35)--(1.15,-.35);\draw(1.15,-.62)--(1.15,-.08)--(2.05,-.08);
\draw(1.15,-.62)--(2.05,-.62);\draw(1.7,-.56)--(1.7,-.14);\draw(1.7,-.35)--(3,-.35);
\node[below]at(1.55,-.65){$b$};\draw[thick](3,-.65)rectangle(4,.65);\node at(3.5,0){$m$};
\draw[bookaxis](3,-1.1)--(4.5,-1.1)node[right]{$x$};
\draw[bookaxis,thick](4,0)--(5.6,0)node[right]{$F_0\cos\omega t$};
\end{tikzpicture}
\caption{题1.4、1.5、1.17：弹簧与黏性阻尼器并联连接物块，位移从平衡位置量起。题1.5取 $b=0$。}\label{fig:ex-1-4}
\end{figure}

恒幅驱动力下，幅值分母平方为$D(\omega)=(10-\omega^2)^2+4\omega^2$。链式求导得
\[D'=2(10-\omega^2)(-2\omega)+8\omega
=4\omega(\omega^2-8).\]
在正频率区间，$D'$于$\sqrt8$两侧由负变正，故分母在此最小，位移峰为$\omega_r=\sqrt8\,\mathrm{rad/s}$，小于固有频率$\sqrt{10}$。

阻尼器的瞬时耗散功率为$2\dot x_s^2$。由$\dot x_s=-\omega A\sin(\omega t-\delta)$及$T=2\pi/\omega$，
\[\frac1T\int_0^T\sin^2(\omega t-\delta)\dd t
=\frac1T\left[\frac t2-\frac{\sin(2\omega t-2\delta)}{4\omega}\right]_0^T
=\frac12.\]
故周期平均给出
\[
\overline P=\frac{b\omega^2A^2}{2}
=\frac{2\times4\times9}{2\times13}=\frac{36}{13}\,\mathrm W.
\]
另以输入功率 $F_0\omega A\sin\delta/2$ 计算，使用 $\sin\delta=2/\sqrt{13}$，也得同值。

两元件对向右运动的物块分别给出 $-kx,-b\dot x$，它们进入复刚度的不同部分；平均耗散只由阻尼器计算。
\end{solution}

\item \begin{minipage}[t]{\linewidth}
【提高】对无阻尼共振，验证 $x_p=F_0t\sin\omega_0t/(2m\omega_0)$，并加上齐次解使 $x(0)=\dot x(0)=0$。 装置见图\ref{fig:ex-1-4}。
\end{minipage}
% 题设配图；与紧随其后的解答分别编号。
\ExpansionDiagram{fig:audit-m-q-023}{无阻尼振子受到固有频率的外力}{
\draw[very thick](0,-.8)--(0,.9);\draw[exspring](0,.35)--(2.4,.35);\node[above]at(1.2,.45){$ k $};\draw[thick](2.4,-.65)rectangle(3.5,.65);\node at(2.95,0){$ m $};\draw[->,thick](3.5,0)--(5,0)node[right]{$ F_0\cos\omega_0t $};\draw[->](2.4,-1.2)--(4,-1.2)node[right]{$x$};
}

\begin{solution}
沿图\ref{fig:ex-1-4}的正方向取位移，并按图\ref{fig:audit-m-q-023}去掉阻尼器作用，即$b=0$。弹簧决定固有频率$\omega_0=\sqrt{k/m}$，驱动力与它同频；图\ref{fig:hw-1-5}画出待验证的共振响应。

记 $C=F_0/(2m\omega_0)$，对候选解逐次求导：
\[
\dot x_p=C\sin\omega_0t+C\omega_0t\cos\omega_0t,
\]
\[
\ddot x_p=2C\omega_0\cos\omega_0t-C\omega_0^2t\sin\omega_0t.
\]
代入 $m\ddot x_p+m\omega_0^2x_p$，两个含 $t\sin\omega_0t$ 的项相消，余下 $2mC\omega_0\cos\omega_0t=F_0\cos\omega_0t$，验证成立。

\begin{figure}[H]\centering
\begin{tikzpicture}[x=.62cm,y=.3cm]
\draw[bookaxis](0,0)--(13.8,0)node[right]{$\omega_0t$};\draw[bookaxis](0,-13)--(0,14)node[above]{$x_p/[F_0/(2m\omega_0^2)]$};
\draw[bookguide](0,0)--(12.6,12.6) (0,0)--(12.6,-12.6);
\draw[bookcurve,domain=0:12.56637,samples=200]plot(\x,{\x*sin(\x r)});\node at(8,12.5){包络正比于 $t$};
\end{tikzpicture}
\caption{题1.5解图：无阻尼共振解为 $t\sin\omega_0t$ 型，位移从零开始且初始切线水平，包络随时间线性增长。}\label{fig:hw-1-5}
\end{figure}

完整解还可加 $A\cos\omega_0t+B\sin\omega_0t$。本题 $x_p(0)=\dot x_p(0)=0$，故初值要求 $A=B=0$，候选解本身就是完整解。随时间增长的因子是无阻尼共振的结果，不能用有限常幅的单频稳态公式代替。

相量图中固定长度的响应箭头对应常幅振动，因而不适用于本题随时间增长的振幅。
\end{solution}

\item \begin{minipage}[t]{\linewidth}
【提高】对串联 RLC，推导电容电压幅值；求 $R=10\,\Omega,L=0.1\,\mathrm H,C=100\,\mu\mathrm F$ 时电流峰频率和电容电压峰频率。

\end{minipage}
% 题设配图；与紧随其后的解答分别编号。
\ExpansionDiagram{fig:audit-m-q-024}{串联RLC中电容两端的测量位置}{
\draw(0,0)--(0,.55);\draw(0,1.45)--(0,2)--(.7,2);\draw(0,1)circle(.45);\node[left]at(-.5,1){$u(t)$};\draw(1.7,2)--(2.4,2);\draw(0.7,1.87)rectangle(1.7,2.13);\node[above]at(1.2,2.45){$ R $};\draw(3.4,2)--(4.1,2);\draw[exspring](2.4,2)--(3.4,2);\node[above]at(2.9,2.45){$ L $};\draw(5.1,2)--(5.8,2);\draw(4.1,2)--(4.4799999999999995,2);\draw(4.72,2)--(5.1,2);\draw(4.4799999999999995,1.65)--(4.4799999999999995,2.35);\draw(4.72,1.65)--(4.72,2.35);\node[above]at(4.6,2.45){$ C $};\draw(5.8,2)--(5.8,0)--(0,0);\draw[->](2,0)--(1,0)node[midway,below]{$I$};\node[below]at(4.6,1.3){$U_C$};
}

\begin{solution}
图\ref{fig:audit-m-q-024}中电容电压由共同电荷$q$决定，先写整个串联回路方程。
图\ref{fig:ex-1-6}规定共同电流和电容正端。电荷的二阶方程由这一路径上的电压降相加得到。

设电容电荷为$q$，基尔霍夫方程为$L\ddot q+R\dot q+q/C=U_0\cos\omega t$。代入$q=\Rea(\widetilde q\ee^{\ii\omega t})$，一次与二次导数分别产生$\ii\omega$、$-\omega^2$，故
\[\left(\frac1C-L\omega^2+\ii R\omega\right)\widetilde q=U_0,
\qquad
\widetilde U_C=\frac{\widetilde q}{C}
=\frac{U_0}{1-LC\omega^2+\ii RC\omega}.\]
取模得到
\[
U_{C0}=\frac{U_0}{\sqrt{(1-LC\omega^2)^2+R^2C^2\omega^2}}.
\]
求电容电压峰，就是使分母平方 $D_C$ 最小。由
\[
D_C'=2(1-LC\omega^2)(-2LC\omega)+2R^2C^2\omega
=2\omega[-2LC+2L^2C^2\omega^2+R^2C^2]=0
\]
得正频率驻点 $\omega_{C,r}^2=1/(LC)-R^2/(2L^2)$，只有右侧为正时才存在该峰。

\begin{figure}[H]\centering
\begin{tikzpicture}[scale=1.05]
\draw(0,0)--(0,.45) (0,1.25)--(0,1.8)--(1,1.8);
\draw(0,.85)circle(.4);\node at(0,.85){$\sim$};\node[left]at(-.4,.85){$u(t)$};
\draw(1,1.6)rectangle(2,2);\node[above]at(1.5,2.05){$R$};\draw(2,1.8)--(2.5,1.8);
\draw(2.5,1.8) arc(180:0:.1875) arc(180:0:.1875) arc(180:0:.1875) arc(180:0:.1875);\node[above]at(3.25,2.05){$L$};\draw(4,1.8)--(4.65,1.8);
\draw(4.65,1.4)--(4.65,2.2) (4.9,1.4)--(4.9,2.2);\node[above]at(4.78,2.25){$C$};
\draw(4.9,1.8)--(5.7,1.8)--(5.7,0)--(0,0);
\draw[bookaxis](5.7,1.35)--(5.7,.4)node[midway,right]{$i$};
\node at(4.65,1.1){$+$};\node at(5.05,1.1){$-$};\node[below]at(4.8,.9){$u_C$};
\end{tikzpicture}
\caption{题1.6：串联 RLC 电路。电容电压按电流流入端为正。}\label{fig:ex-1-6}
\end{figure}

电流幅值为 $U_0/\sqrt{R^2+(\omega L-1/\omega C)^2}$，分母在电抗抵消时最小，故 $\omega_{I,r}=1/\sqrt{LC}$。代入题给参数，
\[
\omega_{C,r}=\sqrt{95000}\simeq308.221\,\mathrm{rad/s},\qquad
\omega_{I,r}=\sqrt{100000}\simeq316.228\,\mathrm{rad/s}.
\]
电容电压与电流是不同的响应量，它们的峰值条件应分别求出。

电容电压对应电荷，电流对应电荷的导数；图中的测量端不同，因此两个幅频峰的位置也不同。
\end{solution}

\item \begin{minipage}[t]{\linewidth}
【选做】对 $f(z)=x^2+\ii y^2$，判断在哪些点复可导，在哪些开区域解析，并求存在时的导数。
\end{minipage}
% 题设配图；与紧随其后的解答分别编号。
\ExpansionDiagram{fig:audit-m-q-025}{在实、虚两个方向上检验差商}{
\draw[->](-.4,0)--(3,0)node[right]{$ x $};\draw[->](0,-0.3)--(0,2.5)node[above]{$ y $};\fill(.8,.7)circle(2pt)node[below right]{$z_0$};\draw[->,thick](.8,.7)--(2,.7)node[right]{$h\in\mathbb R$};\draw[->,thick](.8,.7)--(.8,1.9)node[above]{$\ii h$};
}

\begin{solution}
图\ref{fig:audit-m-q-025}中增量可沿实轴或虚轴趋零；本题的实部、虚部都具有连续偏导。
图\ref{fig:hw-1-7}用于区分“在某点可导”和“在某开区域解析”。先求全部可导点。

令 $u=x^2,v=y^2$。四个偏导为 $u_x=2x,u_y=0,v_x=0,v_y=2y$，柯西--黎曼条件是
\[
u_x=v_y,\quad u_y=-v_x
\quad\Longleftrightarrow\quad x=y.
\]
偏导处处连续，故这条直线上的每一点都复可导，且 $f'=u_x+\ii v_x=2x$；直线外不复可导。

\begin{figure}[H]\centering
\begin{tikzpicture}[scale=1.15]
\draw[bookaxis](-1.8,0)--(2.1,0)node[right]{$x$};\draw[bookaxis](0,-1.8)--(0,2.1)node[above]{$y$};
\draw[bookcurve](-1.5,-1.5)--(1.65,1.65)node[above]{$x=y$};\fill(.8,.8)circle(1.7pt)node[below right]{$z_0$};
\draw[bookguide](.8,.8)circle(.45);\node[right,align=left]at(2.3,.4){直线上：复可导\\任意小圆盘仍含不可导点};
\end{tikzpicture}
\caption{题1.7解图：可导点集是一条直线，不包含任何开圆盘，因此不能构成解析区域。}\label{fig:hw-1-7}
\end{figure}

还可在 $z_0=x_0+\ii x_0$ 处直接展开。增量 $h=\xi+\ii\eta$ 满足
\[
f(z_0+h)-f(z_0)=2x_0h+\xi^2+\ii\eta^2.
\]
最后一项的模不超过 $|h|^2$，除以 $h$ 后趋零，因此导数确为 $2x_0$。可导点集是一条线，没有非空开内部，所以不存在解析的非空开区域。

图中以任意直线点为圆心画小圆，总有圆内点不在直线上；因此找不到全部可导的开邻域。
\end{solution}

\item \begin{minipage}[t]{\linewidth}
【选做】逆时针沿 $|z|=2$，求 $\oint z\dd z/(z^2+1)$；另求 $\int_{-\infty}^{\infty}\dd x/(x^2+a^2)^2$，$a>0$。

\end{minipage}
% 题设配图；与紧随其后的解答分别编号。
\ExpansionDiagram{fig:audit-m-q-026}{圆周围道同时包围正、负虚轴上的两个极点}{
\draw[->](-.4,0)--(3.5,0)node[right]{$ \Rea z $};\draw[->](0,-0.7)--(0,2.3)node[above]{$ \Ima z $};\draw[thick,->](2,0)arc(0:330:2);\draw(330:2)arc(330:360:2);\fill(0,1)circle(2pt)node[left]{$\ii$};\fill(0,-1)circle(2pt)node[left]{$-\ii$};\node[above right]at(1.3,1.5){$|z|=2$};
}

\begin{solution}
图\ref{fig:audit-m-q-026}中圆周取逆时针方向，$\pm\ii$都在内部，两个留数均须计入。
图\ref{fig:ex-1-8}的圆周把 $\pm\ii$ 都包入内部，且走向为逆时针。

第一积分的两个简单极点 $\ii,-\ii$ 都位于 $|z|=2$ 内，且
\[
\Res\left(\frac z{z^2+1},\pm\ii\right)
=\left.\frac z{2z}\right|_{z=\pm\ii}=\frac12.
\]
闭路逆时针，故积分为 $2\pi\ii(1/2+1/2)=2\pi\ii$。

\begin{figure}[H]\centering
\begin{tikzpicture}[scale=.75]
\draw[bookaxis](-2.6,0)--(2.8,0)node[right]{$\Rea z$};\draw[bookaxis](0,-2.5)--(0,2.6)node[above]{$\Ima z$};
\draw[bookcurve](0,0)circle(2);\draw[bookaxis](25:2)arc(25:65:2);
\fill(0,1)circle(2pt)node[right]{$\ii$};\fill(0,-1)circle(2pt)node[right]{$-\ii$};
\node[below left]at(0,0){$O$};\node[below right]at(2,0){$2$};\node[above left]at(-1.4,1.5){$C$};
\end{tikzpicture}
\caption{题1.8：闭曲线 $C:|z|=2$ 的方向及被积函数的奇点位置。}\label{fig:ex-1-8}
\end{figure}

第二积分取上半圆闭路，内部只有 $\ii a$ 这一二阶极点。写 $f=(z-\ii a)^{-2}g(z)$、$g=(z+\ii a)^{-2}$，得
\[
g'(z)=-2(z+\ii a)^{-3},\qquad
\Res(f,\ii a)=g'(\ii a)=-\frac2{(2\ii a)^3}=\frac1{4\ii a^3}.
\]
半径 $R>a$ 的圆弧上，积分模不超过 $\pi R/(R^2-a^2)^2$，随 $R\to\infty$ 趋零。于是实积分等于 $2\pi\ii/(4\ii a^3)=\pi/(2a^3)$。正被积函数的结果为正，量纲为 $a^{-3}$，均符合检查。

第一问沿圆闭合，需计入两个简单极点；第二问沿图\ref{fig:orig-014}的上半圆闭合，只计入 $\ii a$ 处的二阶极点。
\end{solution}

\item \begin{minipage}[t]{\linewidth}
【提高】设 $d>0,\kappa>0$。求满足 $|z-d|=\kappa|z+d|$ 的轨迹，分别讨论 $\kappa=1$ 和 $\kappa\ne1$，给出圆心与半径。

\end{minipage}
% 题设配图；与紧随其后的解答分别编号。
\ExpansionDiagram{fig:audit-m-q-027}{距离比模型的两个固定点与候选动点}{
\draw[->](-3.2,0)--(3.2,0)node[right]{$x$};\fill(-1,0)circle(2pt)node[below]{$-d$};\fill(1,0)circle(2pt)node[below]{$ d $};\fill(0.5,1.7)circle(2pt)node[above]{$P(z)$};\draw(-1,0)--(0.5,1.7)node[midway,left]{$|z+d|$};\draw(1,0)--(0.5,1.7)node[midway,right]{$|z-d|$};
}

\begin{solution}
图\ref{fig:audit-m-q-027}中的$d$固定而$\kappa$可变；化为二次方程后还要单独处理二次项消失的情形。
图\ref{fig:ex-1-9}中的两线段分别对应 $|z-d|$ 与 $|z+d|$，先平方消去距离根号。

两边平方并展开，得
\[
(1-\kappa^2)(x^2+y^2+d^2)-2d(1+\kappa^2)x=0.
\]
当 $\kappa=1$，方程化为 $-4dx=0$，故轨迹为 $x=0$，即两定点的垂直平分线。

\begin{figure}[H]\centering
\begin{tikzpicture}[scale=.94]
\begin{scope}
\draw[bookaxis](-1.6,0)--(1.7,0)node[right]{$x/d$};\draw[bookaxis](0,-1.5)--(0,1.8)node[above]{$y/d$};
\draw[bookcurve](0,-1.3)--(0,1.4);\fill(-1,0)circle(1.5pt)node[below]{$-1$};\fill(1,0)circle(1.5pt)node[below]{$1$};
\fill(0,1)circle(1.5pt)node[above right]{$P$};\draw(-1,0)--(0,1)--(1,0);\node at(0,-2){$\kappa=1$：$x=0$};
\end{scope}
\begin{scope}[shift={(6,0)}]
\draw[bookaxis](-3.3,0)--(1.4,0)node[right]{$x/d$};\draw[bookaxis](0,-1.5)--(0,1.8)node[above]{$y/d$};
\draw[bookcurve](-1.666667,0)circle(1.333333);\fill(-1.666667,0)circle(1.2pt)node[below left]{$C$};
\fill(-1,0)circle(1.5pt)node[below]{$-1$};\fill(1,0)circle(1.5pt)node[below]{$1$};
\fill(-1.666667,1.333333)circle(1.5pt)node[above]{$P$};\draw(-1,0)--(-1.666667,1.333333)--(1,0);
\node at(-1,-2){示取 $\kappa=2$：圆};
\end{scope}
\end{tikzpicture}
\caption{题1.9解图：两距离相等时为垂直平分线；右图取 $\kappa=2$，得到 $C=(-5d/3,0)$、半径 $4d/3$。两图按同一 $d$ 归一化。}\label{fig:ex-1-9}
\end{figure}

当 $\kappa\ne1$，除以 $1-\kappa^2$，记
\[
x_0=\frac{d(1+\kappa^2)}{1-\kappa^2},
\]
配方得到 $(x-x_0)^2+y^2=x_0^2-d^2$。右侧整理为 $4d^2\kappa^2/(1-\kappa^2)^2$，所以圆心为 $(x_0,0)$，半径为 $2d\kappa/|1-\kappa^2|$。圆心坐标保留符号，半径必须为正。

因原式两侧非负，平方没有引入伪解。$\kappa\to1$ 时圆心和半径都变大，有限区域内的轨迹趋近直线，不能把两者分别当成有限极限。

当两距离相等时，图中动点在两定点的垂直平分线上；距离比改变后，圆心沿两定点所在实轴移动。
\end{solution}

\item \begin{minipage}[t]{\linewidth}
【提高】变换 $w=1/z$，$z\ne0$，把圆 $|z-1|=1$ 变成什么轨迹？用代数方程求解，并解释为何圆上的 $z=0$ 要先排除。

\end{minipage}
% 题设配图；与紧随其后的解答分别编号。
\ExpansionDiagram{fig:audit-m-q-028}{过原点的圆在倒数映射中必须删去原点}{
\draw[->](-.4,0)--(3,0)node[right]{$ x $};\draw[->](0,-1.8)--(0,2.6)node[above]{$ y $};\draw[thick](1,0)circle(1);\draw(0,0)circle(2.5pt);\node[below left]at(0,0){删去};\fill(1,0)circle(1.5pt)node[below]{$1$};\draw[->](3.7,.4)--(5,.4)node[midway,above]{$w=1/z$};
}

\begin{solution}
图\ref{fig:audit-m-q-028}中的空心点$z=0$不属于倒数映射的定义域，须在代入前排除。
图\ref{fig:ex-1-10}在原平面排除了奇点，并在像平面补上下面求出的轨迹。

由 $w=1/z$ 得 $z=1/w$，有限像点必有 $w\ne0$。代入给定圆，
\[
\left|\frac1w-1\right|=1
\quad\Longleftrightarrow\quad |1-w|=|w|.
\]
令 $w=u+\ii v$，平方得 $(1-u)^2+v^2=u^2+v^2$，故 $u=1/2$。这条直线不含 $w=0$，其任一点的逆像 $z=1/w$ 都满足原圆方程，因而没有漏掉有限像点。

\begin{figure}[H]\centering
\begin{tikzpicture}[scale=.9]
\draw[bookaxis](-.6,0)--(2.8,0)node[right]{$\Rea z$};\draw[bookaxis](0,-1.6)--(0,1.8)node[above]{$\Ima z$};
\draw[bookcurve](1.2,0)circle(1.2);\fill[white](0,0)circle(2.5pt);\draw(0,0)circle(2.5pt);
\fill(2.4,0)circle(1.5pt)node[below]{$2$};\fill(1.2,1.2)circle(1.5pt)node[above]{$1+\ii$};
\node[below left]at(0,-.1){$0$（排除）};\draw[bookaxis](3.1,.45)--(4.3,.45)node[midway,above]{$1/z$};
\begin{scope}[shift={(5.6,0)}]
\draw[bookaxis](-.5,0)--(2.1,0)node[right]{$u$};\draw[bookaxis](0,-1.6)--(0,1.8)node[above]{$v$};
\draw[bookcurve](.65,-1.45)--(.65,1.45)node[above]{$u=1/2$};\fill(.65,0)circle(1.5pt)node[below right]{$1/2$};
\fill(.65,-.65)circle(1.5pt)node[right]{$(1-\ii)/2$};
\end{scope}
\end{tikzpicture}
\caption{题1.10解图：圆过变换的奇点，有限像为直线 $u=1/2$。两对已标出的点可以直接用 $w=1/z$ 核对。}\label{fig:ex-1-10}
\end{figure}

原圆上的 $z=0$ 不在变换定义域中；沿圆趋近该点时 $|w|=1/|z|\to\infty$。所以像为整条直线 $\Rea w=1/2$。

图中 $2\mapsto1/2$、$1+\ii\mapsto(1-\ii)/2$，两像都落在所得直线上，是代数结果的直接检验。
\end{solution}

\item \begin{minipage}[t]{\linewidth}
【提高】欠阻尼自由振动的相邻同向位移峰值对数比为 $\Delta=\ln(x_n/x_{n+1})$。不用弱阻尼近似，求品质因数 $Q$ 与 $\Delta$ 的关系，再给出小 $\Delta$ 极限。

\end{minipage}
% 题设配图；与紧随其后的解答分别编号。
\ExpansionDiagram{fig:audit-m-q-029}{同向峰的对数比用于测量阻尼}{
\draw[->](-.4,0)--(7,0)node[right]{$ t $};\draw[->](0,-1.5)--(0,2.5)node[above]{$ x $};\draw[thick,domain=0:7,samples=140]plot(\x,{1.8*exp(-.1*\x/(2*pi))*(cos(\x r)+.1/(2*pi)*sin(\x r))});\draw[dashed](0,1.8)--(6.283,1.8);\fill(0,1.8)circle(1.5pt)node[left]{$x_n$};\fill(6.283,1.6287)circle(1.5pt)node[above]{$x_{n+1}$};\draw[<->](0,-1.8)--(6.283,-1.8)node[midway,below]{$T_d$};
}

\begin{solution}
图\ref{fig:audit-m-q-029}与图\ref{fig:ex-1-11}中的相邻同向峰由两条竖线定位。横向间隔给出阻尼周期，纵向峰比给出包络衰减，无须预先采用弱阻尼近似。

设欠阻尼运动为 $x=A\ee^{-\gamma t}\cos(\omega_dt+\varphi)$。同向相邻峰的相位相差 $2\pi$，所以 $T_d=2\pi/\omega_d$，峰值的共同三角因子相消，给出
\[
\Delta=\ln(x_n/x_{n+1})=\gamma T_d
=\frac{2\pi\gamma}{\omega_d}.
\]
再用 $\omega_0^2=\omega_d^2+\gamma^2$ 和 $Q=\omega_0/(2\gamma)$，得
\[
Q=\frac12\sqrt{1+\frac{\omega_d^2}{\gamma^2}}
=\frac{\sqrt{\Delta^2+4\pi^2}}{2\Delta}.
\]
对$0<\Delta\ll1$，先把根号中的$4\pi^2$提取出来，再用
$(1+u)^{1/2}=1+u/2+O(u^2)$，其中$u=\Delta^2/(4\pi^2)$：
\[
Q=\frac\pi\Delta\sqrt{1+\frac{\Delta^2}{4\pi^2}}
=\frac\pi\Delta\left[1+\frac{\Delta^2}{8\pi^2}
+O(\Delta^4)\right].
\]
首项为$\pi/\Delta$。若测得$\Delta=0$，则$\gamma=0$，对应无阻尼极限$Q\to\infty$，不能在含$1/\Delta$的式子中直接代零。两相邻峰处速度均为零，机械能正比位移平方，所以能量的对数衰减量是$2\Delta$。

\begin{figure}[H]\centering
\begin{tikzpicture}[scale=.9]
\draw[bookaxis](0,0)--(6.6,0)node[right]{$t$};\draw[bookaxis](0,-1.5)--(0,1.7)node[above]{$x$};
\draw[bookcurve,domain=0:6.2,samples=190]plot(\x,{1.35*exp(-.12*\x)*cos(180*\x)});
\draw[bookguide,domain=0:6.2,samples=70]plot(\x,{1.35*exp(-.12*\x)});
\foreach\n/\lab in{1/x_n,2/x_{n+1}}{
\pgfmathsetmacro{\tpeak}{2*\n-atan(.12/pi)/180}
\pgfmathsetmacro{\xpeak}{1.35*exp(-.12*\tpeak)*cos(180*\tpeak)}
\draw[bookguide](\tpeak,0)--(\tpeak,\xpeak);
\fill(\tpeak,\xpeak)circle(1.7pt)node[above right]{$\lab$};}
\end{tikzpicture}
\caption{题1.11：相邻同向位移峰值；虚线示意衰减包络。}\label{fig:ex-1-11}
\end{figure}

纵向峰比只由指数衰减产生；因此用同一图读出的峰间距与峰比，可以分别求阻尼频率和衰减率。
\end{solution}

\item \begin{minipage}[t]{\linewidth}
【提高】并联 RLC 取 $R=100\,\Omega,L=0.1\,\mathrm H,C=10\,\mu\mathrm F,\omega=1000\,\mathrm{rad/s}$，电压峰值10 V。求正弦稳态下各支路电流、总电流、复功率，说明无功抵消是否意味着两储能支路没有电流。

\end{minipage}
% 题设配图；与紧随其后的解答分别编号。
\ExpansionDiagram{fig:audit-m-q-030}{并联RLC的三支路电流分别计入}{
\draw(-.8,0)--(-.8,.55);\draw(-.8,1.45)--(-.8,2)--(5,2);\draw(-.8,0)--(5,0);\draw(-.8,1)circle(.45);\node[left]at(-1.3,1){$u(t)$};\draw(1.0,0)--(1.0,.6);\draw(1.0,1.4)--(1.0,2);\draw(0.87,.6)rectangle(1.13,1.4);\node[right]at(1.3,1){$ R $};\draw(2.7,0)--(2.7,.6);\draw(2.7,1.4)--(2.7,2);\draw[decorate,decoration={coil,aspect=.5,segment length=4pt,amplitude=2pt}](2.7,.6)--(2.7,1.4);\node[right]at(3.0,1){$ L $};\draw(4.4,0)--(4.4,.6);\draw(4.4,1.4)--(4.4,2);\draw(4.4,.6)--(4.4,.88);\draw(4.4,1.12)--(4.4,1.4);\draw(4.1000000000000005,.88)--(4.7,.88);\draw(4.1000000000000005,1.12)--(4.7,1.12);\node[right]at(4.7,1){$ C $};
}

\begin{solution}
图\ref{fig:audit-m-q-030}中三个支路电压相同，储能支路的电流虽可相互抵消，却未必各自为零。
图\ref{fig:ex-1-12}各支路两端电压相同，按箭头正方向计算三支路电流后再相加。

以电压相量 $U=10\,\mathrm V$ 为参考。三个支路的峰值电流为
\[
I_R=U/R=0.1\,\mathrm A,
\qquad I_L=U/(\ii\omega L)=-0.1\ii\,\mathrm A,
\qquad I_C=\ii\omega CU=0.1\ii\,\mathrm A.
\]
因此总电流 $I=I_R+I_L+I_C=0.1\,\mathrm A$，对应 $i=0.1\cos1000t\,\mathrm A$。电感电流滞后电压 $\pi/2$，电容电流超前 $\pi/2$，两者在总电流中逐时相消。

\begin{figure}[H]\centering
\begin{tikzpicture}[scale=.9]
\draw(0,0)--(6,0) (0,2.2)--(6,2.2);\draw(0,0)--(0,.7) (0,1.5)--(0,2.2);
\draw(0,1.1)circle(.4);\node at(0,1.1){$\sim$};\node[left]at(-.4,1.1){$u$};
\draw(2,0)--(2,.65) (2,1.55)--(2,2.2);\draw(1.8,.65)rectangle(2.2,1.55);\node[left]at(1.8,1.1){$R$};
\draw(4,0)--(4,.6) (4,1.6)--(4,2.2);\draw(4,.6) arc(-90:90:.125) arc(-90:90:.125) arc(-90:90:.125) arc(-90:90:.125);\node[left]at(3.7,1.1){$L$};
\draw(6,0)--(6,.95) (6,1.2)--(6,2.2);\draw(5.55,.95)--(6.45,.95) (5.55,1.2)--(6.45,1.2);\node[right]at(6.45,1.1){$C$};
\foreach\x/\j in{2/R,4/L,6/C}{\draw[bookaxis](\x+.35,2)--(\x+.35,1.55)node[right]{$i_\j$};}
\node[left]at(0,2.2){$+$};\node[left]at(0,0){$-$};
\end{tikzpicture}
\caption{题1.12：并联 RLC，各支路电流正方向由电压正端流向负端。}\label{fig:ex-1-12}
\end{figure}

复功率为 $S=UI^*/2=0.5\,\mathrm{VA}$，实部为0.5 W，虚部为零。分别看支路，有 $S_L=+0.5\ii\,\mathrm{VA}$、$S_C=-0.5\ii\,\mathrm{VA}$，它们各自都非零。无功抵消指对电源的合成结果，不表示储能支路没有电流。

电感和电容相量分别沿负、正虚轴，长度相等时相消；电阻电流沿实轴留下，图中的两储能支路仍各有电流。
\end{solution}

\item \begin{minipage}[t]{\linewidth}
【提高】周期 $2\pi$ 的函数在 $[-\pi,\pi]$ 上为 $f(t)=|t|$。求傅里叶系数，并用 $t=0$ 处的展开推出 $\sum_{n=1}^{\infty}n^{-2}$。

\end{minipage}
% 题设配图；与紧随其后的解答分别编号。
\ExpansionDiagram{fig:audit-m-q-031}{偶对称三角波在一个周期中的形状}{
\draw[->](-.4,0)--(4,0)node[right]{$ t $};\draw[->](0,-0.3)--(0,3.6)node[above]{$ f $};\draw[thick](-3.14,3.14)--(0,0)--(3.14,3.14);\node[below]at(-3.14,0){$-\pi$};\node[below]at(3.14,0){$\pi$};
}

\begin{solution}
图\ref{fig:audit-m-q-031}中$f(t)=|t|$为偶函数，因此展开只需常量项和余弦项。
图\ref{fig:ex-1-13}是偶对称三角波，因此先令全部正弦系数为零，再算余弦项。

函数为偶函数，所以 $b_n=0$。按常量项为 $a_0/2$ 的约定，
\[
a_0=\frac2\pi\int_0^\pi t\dd t
=\frac2\pi\left[\frac{t^2}{2}\right]_0^\pi=\pi.
\]
对余弦系数取$u=t$、$\dd v=\cos nt\dd t$，则$\dd u=\dd t$、$v=\sin nt/n$。分部积分得
\begin{align*}
a_n&=\frac2\pi\int_0^\pi t\cos nt\dd t\\
&=\frac2\pi\left\{\left[\frac{t\sin nt}{n}\right]_0^\pi
-\frac1n\int_0^\pi\sin nt\dd t\right\}\\
&=\frac2\pi\left[\frac{t\sin nt}{n}+\frac{\cos nt}{n^2}\right]_0^\pi
=\frac{2[(-1)^n-1]}{\pi n^2}.
\end{align*}
所以展开为$|t|=\pi/2-(4/\pi)\sum_{j\ge0}\cos[(2j+1)t]/(2j+1)^2$。系数绝对可和，且$t=0$是连续点，在此求值得
\[0=\frac\pi2-\frac4\pi\sum_{j=0}^{\infty}\frac1{(2j+1)^2}
\quad\Longrightarrow\quad
\sum_{j=0}^{\infty}\frac1{(2j+1)^2}=\frac{\pi^2}{8}.\]

\begin{figure}[H]\centering
\begin{tikzpicture}[scale=.9]
\draw[bookaxis](-4,0)--(4.1,0)node[right]{$t$};\draw[bookaxis](0,-.3)--(0,2)node[above]{$f(t)$};
\draw[bookcurve](-3.6,0)--(-1.8,1.5)--(0,0)--(1.8,1.5)--(3.6,0);
\draw[bookguide](-1.8,0)--(-1.8,1.5)--(1.8,1.5)--(1.8,0);
\node[below]at(-1.8,0){$-\pi$};\node[below]at(1.8,0){$\pi$};\node[below left]at(0,0){$0$};\node[above left]at(0,1.5){$\pi$};
\end{tikzpicture}
\caption{题1.13：$f(t)=|t|$ 在一个周期外按周期 $2\pi$ 延拓。}\label{fig:ex-1-13}
\end{figure}

设全和为 $S$，偶数项之和为 $S/4$，从而 $S-S/4=\pi^2/8$，解得 $S=\pi^2/6$。

在图中的谷底 $t=0$，函数连续且取零，故可把级数在此求值；这一步给出奇数平方倒数和。
\end{solution}

\item \begin{minipage}[t]{\linewidth}
【选做】求 $\int_{-\infty}^{\infty}\dd x/[(x^2+a^2)(x^2+b^2)]$，$a,b>0$。分别处理 $a\ne b$ 与 $a=b$，检查两结果的连续衔接。
\end{minipage}
% 题设配图；与紧随其后的解答分别编号。
\ExpansionDiagram{fig:audit-m-q-032}{两对虚轴极点及参数合并的可能性}{
\draw[->](-.4,0)--(2.8,0)node[right]{$ \Rea z $};\draw[->](0,-2.8)--(0,3)node[above]{$ \Ima z $};\foreach \y/\lab in {1/\ii a,2/\ii b,-1/-\ii a,-2/-\ii b}{\fill(0,\y)circle(2pt)node[right]{$\lab$};}\node[right]at(1.4,.4){示意$a<b$};
}

\begin{solution}
图\ref{fig:audit-m-q-032}中四个极点在$a=b$时两两合并，部分分式公式应先在$a\ne b$下使用。
图\ref{fig:hw-1-14}把两个上半平面极点及其合并过程画出。分式拆开后，每项还可以分别使用图\ref{fig:orig-014}的路径，将参数换成相应的 $a$ 或 $b$。

先设 $a\ne b$，令部分分式系数满足
\[
\frac1{(x^2+a^2)(x^2+b^2)}
=\frac A{x^2+a^2}+\frac B{x^2+b^2}.
\]
比较分子得 $A+B=0$、$Ab^2+Ba^2=1$，故 $A=1/(b^2-a^2)$、$B=-A$。使用 $\int_{-\infty}^{\infty}\dd x/(x^2+a^2)=\pi/a$，积分为
\[
\frac{\pi/a-\pi/b}{b^2-a^2}=\frac\pi{ab(a+b)}.
\]
当 $a=b$ 时，部分分式中出现的零分母只表示该分解不再适用；原积分仍存在。由二阶极点结果得 $\pi/(2a^3)$，也等于上述最终式令 $b\to a$ 的极限。参数取固定正值附近时有可积控制，因此该连续衔接合法。

\begin{figure}[H]\centering
\begin{tikzpicture}[scale=1.05]
\draw[bookaxis](-3,0)--(3.2,0)node[right]{$\Rea z$};\draw[bookaxis](0,-1)--(0,3)node[above]{$\Ima z$};
\draw[bookcurve](-2.6,0)--(2.6,0);\draw[bookcurve](2.6,0)arc(0:180:2.6);
\draw[->](-.9,0)--(-.3,0);\draw[->](55:2.6)arc(55:75:2.6);
\fill(0,.9)circle(1.7pt)node[right]{$\ii a$};\fill(0,1.8)circle(1.7pt)node[right]{$\ii b$};
\node[left]at(0,1.35){$a\ne b$};\node[right,align=left]at(3.8,1.1){$b\to a$：\\两个简单极点\\合并为二阶极点};
\end{tikzpicture}
\caption{题1.14解图：示意取 $a<b$；反序时两标记交换，闭路结论不变。分式拆分与参数合并极限对应两种局部极点结构。}\label{fig:hw-1-14}
\end{figure}

当 $b\to a$，两简单分式的系数发散而积分组合仍有有限极限；原函数的极点合并成二阶极点，与例1.14对应。
\end{solution}

\item \begin{minipage}[t]{\linewidth}
【综合】质量 $m$ 通过弹簧 $k$ 和黏性阻尼器 $b$ 与同一运动底座连接，底座位移为 $y=Y\cos\omega t$，物块绝对位移为 $x$。自行列方程，求 $x$ 与相对位移 $x-y$ 的稳态相量。检查低频、高频极限，并说明为何不能把底座位移直接当成外力。

\end{minipage}
% 题设配图；与紧随其后的解答分别编号。
\ExpansionDiagram{fig:audit-m-q-033}{弹簧与阻尼器共同连接在运动底座上}{
\draw[very thick](0,-.8)--(0,.9);\draw[exspring](0,.35)--(2.4,.35);\node[above]at(1.2,.45){$ k $};\draw[thick](2.4,-.65)rectangle(3.5,.65);\node at(2.95,0){$ m $};\draw(0,-.35)--(.8,-.35);\draw(.8,-.6)--(.8,-.1)--(1.7,-.1);\draw(.8,-.6)--(1.7,-.6);\draw(1.35,-.55)--(1.35,-.15);\draw(1.35,-.35)--(2.4,-.35);\node[below]at(1.2,-.7){$ b $};\draw[->](2.4,-1.2)--(4,-1.2)node[right]{$x$};\draw[->](-.5,-1.2)--(.7,-1.2)node[below]{$y(t)$};
}

\begin{solution}
图\ref{fig:audit-m-q-033}中底座位移为$y$，两元件的变形和相对速度分别是$x-y$与$\dot x-\dot y$。
图\ref{fig:ex-1-15}把底座端和物块端分别标为 $y,x$。弹簧、阻尼器都连接这两个端点。

弹簧伸长是 $x-y$，阻尼器两端相对速度是 $\dot x-\dot y$。只对物块列牛顿方程，
\[
m\ddot x=-k(x-y)-b(\dot x-\dot y).
\]
取底座相量为实数 $Y$，将谐波形式代入，得
\[
(k-m\omega^2+\ii b\omega)X=(k+\ii b\omega)Y.
\]
记分母为 $D$，则 $X=(k+\ii b\omega)Y/D$，进一步相减得
\[
X-Y=\frac{(k+\ii b\omega)-D}{D}Y
=\frac{m\omega^2Y}{D}.
\]
低频极限$X/Y\to1$、$(X-Y)/Y\to0$；高频极限$X/Y\to0$、$(X-Y)/Y\to-1$。若改用$r=x-y$，则$x=r+y$、$\ddot x=\ddot r+\ddot y$，代入物块方程得到
\[m(\ddot r+\ddot y)=-kr-b\dot r,
\qquad m\ddot r+b\dot r+kr=-m\ddot y.\]
由$\dot y=-\omega Y\sin\omega t$、$\ddot y=-\omega^2Y\cos\omega t$，等效外力幅值为$m\omega^2Y$。位移$Y$本身没有力的量纲。

\begin{figure}[H]\centering
\begin{tikzpicture}[scale=.95]
\draw[thick](0,-.8)--(0,.9);
\draw[exspring](0,.35)--(3,.35);\node[above]at(1.5,.55){$k$};
\draw(0,-.35)--(1.15,-.35);\draw(1.15,-.62)--(1.15,-.08)--(2.05,-.08);
\draw(1.15,-.62)--(2.05,-.62);\draw(1.7,-.56)--(1.7,-.14);\draw(1.7,-.35)--(3,-.35);
\node[below]at(1.55,-.65){$b$};\draw[thick](3,-.65)rectangle(4,.65);\node at(3.5,0){$m$};
\draw[bookaxis](3,-1.1)--(4.5,-1.1)node[right]{$x$};
\draw[bookaxis](-.7,-1.1)--(.6,-1.1)node[below right]{$y(t)$};\node[above]at(0,1){运动底座};
\end{tikzpicture}
\caption{题1.15：弹簧和阻尼器的左端随同一底座运动；$x,y$ 均为绝对位移。}\label{fig:ex-1-15}
\end{figure}

图上的形变量是两端位移之差，相对速度同理；把底座加速度移到相对坐标方程右侧，才得到等效惯性驱动力。
\end{solution}

\item \begin{minipage}[t]{\linewidth}
【提高】令 $\zeta=0.02$，分别用精确公式和共振附近首阶近似计算 $\Omega=0.98,1,1.02$ 处的无量纲位移幅值。比较相对误差；保持 $\Omega=0.5$ 不变而令 $\zeta\to0$ 时，应采用哪种近似？
\end{minipage}
% 题设配图；与紧随其后的解答分别编号。
\ExpansionDiagram{fig:audit-m-q-034}{同一小阻尼下的三个近共振采样频率}{
\draw[->](-.4,0)--(5,0)node[right]{$ \Omega $};\draw[->](0,-0.3)--(0,1.8)node[above]{$  $};\foreach \x/\lab in {1/.98,2.5/1,4/1.02}{\draw(\x,-.15)--(\x,.15);\node[below]at(\x,0){$\lab$};}\node[above]at(2.5,1){$\zeta=0.02$};
}

\begin{solution}
图\ref{fig:audit-m-q-034}中的频率偏离1的量与$\zeta$同阶，因此可比较近共振首阶式与精确式。
图\ref{fig:hw-1-16}放大近共振频带，分别画精确曲线与首阶式，三条竖线对应题给频率。

精确无量纲幅值为
\[
\mathcal A(\Omega)=\frac1{\sqrt{(1-\Omega^2)^2+4\zeta^2\Omega^2}}.
\]
在近共振缩放 $\varepsilon=\Omega-1=O(\zeta)$ 下，保留首阶后分母为 $2\sqrt{\varepsilon^2+\zeta^2}$。代入 $\zeta=0.02$，得到
\[
\begin{array}{c|rr}
\Omega&\text{精确值}&\text{首阶近似}\\
0.98&17.9466&17.6777\\
1&25.0000&25.0000\\
1.02&17.4162&17.6777
\end{array}
\]
以精确值为分母定义相对误差，两侧分别约为 $-1.50\%$ 和 $+1.50\%$，中心为零。近似将两侧写成对称，略去了高阶的不对称项。

\begin{figure}[H]\centering
\begin{tikzpicture}[x=1.1cm,y=2.2cm]
\draw[bookaxis](-3.2,0)--(3.5,0)node[right]{$\Omega$};\draw[bookaxis](-3,0)--(-3,1.2)node[above]{$\mathcal A/25$};
\draw[bookcurve,domain=-2.8:2.8,samples=170]plot(\x,{.04/sqrt((1-(1+.02*\x)^2)^2+.0016*(1+.02*\x)^2)});
\draw[dashed,domain=-2.8:2.8,samples=170]plot(\x,{1/sqrt(1+(\x)^2)});
\foreach\u/\lab in{-1/.98,0/1,1/1.02}{\draw[bookguide](\u,0)--(\u,{1/sqrt(1+(\u)^2)});\node[below]at(\u,0){$\lab$};}
\node at(1.8,1.05){实线：精确};\node at(1.8,.84){虚线：首阶};
\end{tikzpicture}
\caption{题1.16解图：近共振曲线的首阶式左右对称，精确式保留轻微不对称。图中仅放大 $\Omega\approx1$ 的频带。}\label{fig:hw-1-16}
\end{figure}

固定 $\Omega=0.5$ 时，$1-\Omega^2=3/4$ 并不小，应在非共振分母上展开，首项为 $4/3$，阻尼修正为 $O(\zeta^2)$。此时不满足 $\Omega-1=O(\zeta)$，不能沿用近共振的缩放式。

左右误差符号相反，与图中精确曲线略偏斜的形状对应；$\Omega=0.5$ 不在这个放大频带内。
\end{solution}

\item \begin{minipage}[t]{\linewidth}
【综合】质量为1 kg的受迫振子在角频率2 rad/s时位移滞后恒幅驱动力 $\pi/2$，在1 rad/s时滞后 $\pi/4$。求刚度和黏性阻尼系数，判断位移曲线是否有非零频率峰，并求平均功率峰频率。 装置见图\ref{fig:ex-1-4}。
\end{minipage}
% 题设配图；与紧随其后的解答分别编号。
\ExpansionDiagram{fig:audit-m-q-035}{由两个不同频率的相位记录反求系统参数}{
\draw[very thick](0,-.8)--(0,.9);\draw[exspring](0,.35)--(2.4,.35);\node[above]at(1.2,.45){$ k $};\draw[thick](2.4,-.65)rectangle(3.5,.65);\node at(2.95,0){$ 1 $};\draw(0,-.35)--(.8,-.35);\draw(.8,-.6)--(.8,-.1)--(1.7,-.1);\draw(.8,-.6)--(1.7,-.6);\draw(1.35,-.55)--(1.35,-.15);\draw(1.35,-.35)--(2.4,-.35);\node[below]at(1.2,-.7){$ b $};\draw[->,thick](3.5,0)--(5,0)node[right]{$ F_0\cos\omega t $};\draw[->](2.4,-1.2)--(4,-1.2)node[right]{$x$};\node[below]at(2.5,-1.8){$\phi(2)=\pi/2,\quad\phi(1)=\pi/4$};
}

\begin{solution}
图\ref{fig:audit-m-q-035}中刚度与阻尼未知，两次测得的相位分别约束复刚度的实部和虚部。
装置见图\ref{fig:ex-1-4}，两次相位测量则画在图\ref{fig:hw-1-17}的复刚度平面中。

稳态相量为 $X=F_0/(k-m\omega^2+\ii b\omega)$。若位移滞后力的角度为 $\phi$，分母的辐角就是 $\phi$。在 $\omega=2\,\mathrm{rad/s}$ 时 $\phi=\pi/2$，分母实部必须为零。由 $k=m\omega^2$ 和 $m=1\,\mathrm{kg}$，得 $k=4\,\mathrm{N/m}$。

\begin{figure}[H]\centering
\begin{tikzpicture}[scale=.95]
\draw[bookaxis](-.4,0)--(4.5,0)node[right]{$\Rea D$};\draw[bookaxis](0,-.3)--(0,4)node[above]{$\Ima D$};
\draw[bookaxis,thick](0,0)--(1.65,1.65)node[right]{$D(1)=3+3\ii$};\draw[bookguide](1.65,1.65)--(1.65,0)node[below]{$3$};
\draw[bookaxis,thick](0,0)--(0,3.3)node[left]{$D(2)=6\ii$};
\draw(.7,0)arc(0:45:.7);\node at(23:1.05){$\pi/4$};\node[right]at(.2,2.7){$\Arg D(2)=\pi/2$};
\end{tikzpicture}
\caption{题1.17解图：以 SI 数值表示复刚度 $D=k-m\omega^2+\ii b\omega$。两条射线分别落实题给的 $\pi/4,\pi/2$ 滞后角。}\label{fig:hw-1-17}
\end{figure}

在 $\omega=1\,\mathrm{rad/s}$ 时，$\phi=\pi/4$ 且实、虚部都为正，所以 $b=(k-m\omega^2)/\omega=3\,\mathrm{N\,s/m}$。非零位移峰要求 $k/m-b^2/(2m^2)>0$，但本题
\[
\frac km-\frac{b^2}{2m^2}=(4-4.5)\,\mathrm{s^{-2}}=-0.5\,\mathrm{s^{-2}}<0,
\]
该条件不成立，故位移幅值从零频起单调下降。功率的峰仍满足 $k-m\omega^2=0$，所以为 $2\,\mathrm{rad/s}$。相位经过 $\pi/2$ 的事实，并不能单独证明位移幅频曲线有非零频率峰。

第二个测量点落在正虚轴，先给出 $k=m\omega^2$；第一个点的两直角边相等，再确定 $b$。
\end{solution}

\end{enumerate}

\chapter{向量、曲线与空间几何}
粒子沿曲线运动时，位置可以只用一个参数描述，加速度却通常不沿切线。积分连续电荷的场时，每个微元的场方向也不同。本章从这些几何关系出发，练习选择坐标、分解运动和构造积分元。

\section{向量的分量与几何意义}
\subsection{向量与坐标分量}
位移、速度、力是有大小与方向的物理对象。写成 $(a_x,a_y,a_z)$，只是它们在某组基矢下的分量。若采用右手正交单位基 $\vct e_x,\vct e_y,\vct e_z$，则 $\vct a=a_x\vct e_x+a_y\vct e_y+a_z\vct e_z$。改变基矢会改变分量，却不改变物理向量。

直角坐标基矢通常固定，求导时只需对分量求导。极坐标基矢随方向改变，即使某个分量为常量，整个向量仍可能变化；求导时要把基矢的变化一并计入。

向量加法与数乘满足分配律。长度 $|\vct a|=\sqrt{a_x^2+a_y^2+a_z^2}$，非零向量对应单位向量 $\hat{\vct a}=\vct a/|\vct a|$。单位向量只记录方向，模长为1，不是给每个分量分别取1。

\subsection{点积、投影与功}
点积定义为 $\vct a\cdot\vct b=|\vct a||\vct b|\cos\theta$。在正交单位基下，
\begin{equation}
\vct a\cdot\vct b=a_xb_x+a_yb_y+a_zb_z.
\end{equation}
取单位向量 $\vct n$，$\vct a\cdot\vct n$ 是沿 $\vct n$ 的带符号标量分量，而 $(\vct a\cdot\vct n)\vct n$ 才是投影向量。负号说明投影与所选正方向相反，不说明长度为负。

恒力在位移 $\Delta\vct r$ 上做功 $W=\vct F\cdot\Delta\vct r$。若力与位移垂直，功为零，即使力很大；若方向反向，功为负。曲线上的变力做功为 $\int\vct F\cdot\dd\vct r$，这里每个微元都必须取相应方向的投影。

\subsection{叉积、力矩与定向}
三维叉积的方向由右手规则确定，模长为 $|\vct a||\vct b|\sin\theta$。分量式为
\begin{equation}
\vct a\times\vct b=
\begin{vmatrix}\vct e_x&\vct e_y&\vct e_z\\a_x&a_y&a_z\\b_x&b_y&b_z\end{vmatrix}.
\end{equation}
交换次序会改变符号。其模长是两向量张成的平行四边形面积，三角形面积则为一半。力矩 $\vct\tau=\vct r\times\vct F$ 的模长等于力乘力臂；角动量 $\vct L=\vct r\times\vct p$ 则记录相对指定参考点的转动状态。

\begin{example}
力 $\vct F=(F,2F,0)$ 作用于点 $\vct r=(a,0,0)$，求相对原点的力矩，并判断该力沿 $x$ 方向位移 $\Delta x$ 时做的功。
\end{example}
% 题设配图；与紧随其后的解答分别编号。
\ExpansionDiagram{fig:audit-m-q-036}{位置向量、作用力及给定的水平位移}{
\draw[->](-.4,0)--(4.8,0)node[right]{$ x $};\draw[->](0,-0.3)--(0,3)node[above]{$ y $};\fill(2,0)circle(2pt)node[below]{$(a,0)$};\draw[->,thick](2,0)--(3,2)node[above]{$\vct F=(F,2F)$};\draw[->](0,-.55)--(2,-.55)node[midway,below]{$\vct r$};\draw[->](2,-.55)--(3.5,-.55)node[midway,below]{$\Delta x$};
}

\begin{solution}
图\ref{fig:audit-m-q-036}中力矩取相对原点的位置向量，做功则取实际位移；两者使用不同乘积。
按图\ref{fig:eg-2-1}先分解力：竖直分量与位置垂直，水平分量与给定位移平行。

相对原点的力矩由位置与力的叉积决定。按行列式展开，
\[
\vct\tau=\begin{vmatrix}\vct e_x&\vct e_y&\vct e_z\\a&0&0\\F&2F&0\end{vmatrix}
=(0,0,2aF).
\]
假设 $aF>0$，方向沿正 $z$，与右手规则一致。沿 $x$ 的力分量平行位置矢量，力臂为零；沿 $y$ 的分量为 $2F$，力臂为 $a$。

\begin{figure}[H]\centering
\begin{tikzpicture}[scale=.9]
\draw[bookaxis](-.4,0)--(6.2,0)node[right]{$x$};\draw[bookaxis](0,-.3)--(0,3.4)node[above]{$y$};
\coordinate(P)at(2,0);\draw[bookaxis,thick](0,0)--(P)node[midway,below]{$a$};\fill(P)circle(2pt)node[below]{$P$};
\draw[bookaxis,thick](P)--(3.3,2.6)node[above right]{$\vct F$};\draw[bookguide](P)--(3.3,0)--(3.3,2.6);
\node[below]at(2.65,0){$F$};\node[right]at(3.3,1.25){$2F$};
\draw[bookaxis](2,-.75)--(4,-.75)node[midway,below]{$\Delta x$};\node[right]at(4.3,2.2){$\tau_z=a(2F)$};
\end{tikzpicture}
\caption{例2.1：位置矢量沿 $x$，只有竖直力分量贡献力矩；给定水平位移上的功则由水平力分量贡献。图中的力分解采用另一尺度。}\label{fig:eg-2-1}
\end{figure}

若位移为 $\Delta\vct r=(\Delta x,0,0)$，恒力做功为
\[
W=\vct F\cdot\Delta\vct r=F\Delta x.
\]
沿 $y$ 的分量虽然贡献力矩，却与这次位移垂直，不贡献功。结果中的符号由 $\Delta x$ 决定，向负 $x$ 位移时该功为负。

图上两条投影分别进入叉积与点积；同一力的两个分量在这两个问题中承担不同作用。
\end{solution}

\subsection{混合积与常用恒等式}
混合积 $(\vct a\times\vct b)\cdot\vct c$ 是有向体积，其绝对值为平行六面体体积。循环置换三向量不改变它，交换任意两个会变号。零混合积意味着三向量共面或至少有一个退化为零。

由分量展开或叉积的几何定义可得
\begin{align}
\vct a\times(\vct b\times\vct c)
&=\vct b(\vct a\cdot\vct c)-\vct c(\vct a\cdot\vct b),\\
|\vct a\times\vct b|^2
&=|\vct a|^2|\vct b|^2-(\vct a\cdot\vct b)^2.
\end{align}
第一式可先核对$x$分量。将内层叉积展开，有
\begin{align*}
[\vct a\times(\vct b\times\vct c)]_x
&=a_y(b_xc_y-b_yc_x)-a_z(b_zc_x-b_xc_z)\\
&=b_x(a_yc_y+a_zc_z)-c_x(a_yb_y+a_zb_z).
\end{align*}
补入相互抵消的$b_xa_xc_x-c_xa_xb_x$，便得到
$b_x(\vct a\cdot\vct c)-c_x(\vct a\cdot\vct b)$；
$y,z$分量按坐标循环置换得到。这个结果也表明叉积不满足结合律。
第二式可由$\sin^2\theta+\cos^2\theta=1$得到，可据此用点积计算叉积的模。

\section{直线、平面与空间距离}
直线和平面约束分别限制可达位置或运动方向。方向向量与法向量把几何条件化为方程，距离和投影的计算再计入向量的长度。

\subsection{先识别方向与法线}
过 $\vct r_0$、方向为 $\vct s\ne0$ 的直线写成 $\vct r=\vct r_0+t\vct s$。参数 $t$ 一般只是定位参数，除非明确规定，不能把它当作物理时间。过 $\vct r_1$、法向量为 $\vct n\ne0$ 的平面满足 $\vct n\cdot(\vct r-\vct r_1)=0$。

平面写成 $Ax+By+Cz=D$ 时，法向量为 $(A,B,C)$。若直线方向与法向量点积非零，直线与平面唯一相交；若点积为零，直线平行于平面，进一步检查直线上一个点是否在平面内，可区分整条直线在平面内和直线与平面无交点。

\subsection{用投影推导距离}
从平面上一点到 $P$ 的向量在单位法向量上的投影，就是点到平面的带符号距离。取绝对值得
\begin{equation}
d=\frac{|\vct n\cdot\vct P-D|}{|\vct n|}.
\end{equation}
从 $P$ 减去这段法向投影，得到垂足
\begin{equation}
\vct H=\vct P-\frac{\vct n\cdot\vct P-D}{|\vct n|^2}\vct n.
\end{equation}
代入可验证 $\vct n\cdot\vct H=D$，而 $\vct P-\vct H$ 与法向量平行。距离公式与垂足公式是同一投影关系的两个结果。

\begin{figure}[htbp]
\centering
\begin{tikzpicture}[scale=.85,>=Stealth]
\draw (-2,-.7)--(1.4,-.7)--(2.4,.5)--(-1,.5)--cycle;
\node at (1.8,.15) {$\Pi$};\coordinate(H)at(0,0);\coordinate(P)at(0,2.2);
\draw[dashed] (H)--(P);\fill(H)circle(1.5pt)node[below]{$H$};\fill(P)circle(1.5pt)node[right]{$P$};
\draw[->] (.6,.2)--(.6,1.6)node[right]{$\vct n$};
\draw (0,.2)--(.2,.2)--(.2,0);\node[left] at(0,1.1){$d$};
\end{tikzpicture}
\caption{点到平面的法向投影}\label{fig:plane-distance}
\end{figure}

点到直线的距离可由面积求高：
\begin{equation}
d_{P,l}=\frac{|(\vct P-\vct r_0)\times\vct s|}{|\vct s|}.
\end{equation}
也可将 $\vct P-\vct r_0$ 减去沿直线的投影，再求模。两种方法等价，前者紧凑，后者能同时给出垂足参数 $t=(\vct P-\vct r_0)\cdot\vct s/|\vct s|^2$。

\subsection{异面直线与夹角}
两非平行直线的公垂线方向为 $\vct s_1\times\vct s_2$，因而最短距离为
\begin{equation}
d=\frac{|(\vct r_2-\vct r_1)\cdot(\vct s_1\times\vct s_2)|}{|\vct s_1\times\vct s_2|}.
\end{equation}
这也可理解为平行六面体体积除以底面积。若两方向平行，分母为零，公式不适用；此时选一条直线上的点，改用点到另一条直线的距离。

两直线的锐角满足 $\cos\alpha=|\vct s_1\cdot\vct s_2|/(|\vct s_1||\vct s_2|)$。两平面的锐角用法向量的同一余弦公式。直线与平面的角是方向与法向量夹角的余角，因此应使用
\begin{equation}
\sin\alpha=\frac{|\vct s\cdot\vct n|}{|\vct s||\vct n|}.
\end{equation}
把这三个角度的对象分别画清，通常比记三个公式更可靠。

\section{圆锥曲线的几何定义}
轨道方程中的参数应能对应到焦点、准线和距离。先从几何定义推导方程，再由实际允许区域选择曲线分支。

\subsection{抛物线：焦点与准线}
抛物线是到固定点和固定直线距离相等的点集。取焦点 $(p,0)$、准线 $x=-p$，$p>0$，距离关系为 $\sqrt{(x-p)^2+y^2}=|x+p|$。平方并化简得到 $y^2=4px$。曲线顶点在原点，开口沿正 $x$，焦点到顶点距离为 $p$。

物理中常出现形式 $y=ax^2+bx+c$，它与标准形式之间还存在平移、旋转和轴名称变换。辨认曲线时要区分几何类型与当前坐标方向，不能因为方程中的平方在 $x$ 上就断言曲线不属于抛物线。

\subsection{椭圆：两个焦距之和}
椭圆定义为到两焦点距离之和为定值$2a$。取焦点$(\pm c,0)$，$a>c\ge0$，并记$r_\pm=\sqrt{(x\mp c)^2+y^2}$。由$r_++r_-=2a$及$r_-^2-r_+^2=4cx$，得$r_--r_+=2cx/a$，进而$r_-=a+cx/a$。再次平方，
\[(x+c)^2+y^2=\left(a+\frac{cx}{a}\right)^2.\]
两边的$2cx$抵消，移项成为$(1-c^2/a^2)x^2+y^2=a^2-c^2$。除以正数$a^2-c^2$，得到
\begin{equation}
\frac{x^2}{a^2}+\frac{y^2}{b^2}=1,\qquad b^2=a^2-c^2.
\end{equation}
离心率 $e=c/a$ 在 $[0,1)$ 内。$e=0$ 退化为圆，但仍是正常的椭圆特例；$e\to1$ 时短轴趋于零，必须小心极限与几何退化。
\begin{figure}[htbp]
\centering
\begin{tikzpicture}[scale=1,>=Stealth]
\draw[->](-2.6,0)--(2.7,0)node[right]{$x$};\draw[->](0,-1.6)--(0,1.8)node[above]{$y$};
\draw[thick](0,0)ellipse(2.2 and 1.35);
\pgfmathsetmacro{\fdist}{sqrt(2.2*2.2-1.35*1.35)}
\coordinate(F1)at(-\fdist,0);\coordinate(F2)at(\fdist,0);\coordinate(P)at({2.2*cos(110)},{1.35*sin(110)});
\fill(F1)circle(1.5pt)node[below]{$F_1$};\fill(F2)circle(1.5pt)node[below]{$F_2$};\fill(P)circle(1.5pt)node[above left]{$P$};
\draw[dashed](F1)--(P)node[midway,left]{$r_1$};\draw[dashed](F2)--(P)node[midway,above]{$r_2$};
\node[below right]at(2.2,0){$a$};\node[above right]at(0,1.35){$b$};
\end{tikzpicture}
\caption{椭圆的焦距和与标准坐标}\label{fig:ellipse}
\end{figure}

椭圆可参数化为 $x=a\cos u,y=b\sin u$。参数 $u$ 称为一种角参数，却一般不是由焦点量起的极角，更不自动等于 $\omega t$。同一几何曲线可以对应无穷多种运动速度，必须另外给出时间规律。

\subsection{双曲线与渐近线}
若到两焦点距离之差的绝对值为$2a$，其中$c>a>0$，取同样的$r_\pm$。在右支有$r_--r_+=2a$，再由平方差$4cx$得到$r_-+r_+=2cx/a$，所以$r_-=a+cx/a$。平方后仍有
\[(1-c^2/a^2)x^2+y^2=a^2-c^2,\]
但现在两侧的常数因子为负。移项并除以$c^2-a^2$，得
\begin{equation}
\frac{x^2}{a^2}-\frac{y^2}{b^2}=1,\qquad b^2=c^2-a^2.
\end{equation}
右支取$x\ge a$，另一分支由$x\mapsto-x$得到。焦距参数满足$c^2=a^2+b^2$，与椭圆不同。远离原点时常量$1$相对二次项可忽略，故渐近线为$y=\pm(b/a)x$；它们描述远处方向，不是双曲线上的有限线段。

\subsection{切线与光学反射}
对椭圆 $x^2/a^2+y^2/b^2=1$，在 $(x_0,y_0)$ 处的法向量与 $(x_0/a^2,y_0/b^2)$ 平行，因此切线为 $xx_0/a^2+yy_0/b^2=1$。

更一般地，椭圆的焦距和函数$S(P)=|P-F_1|+|P-F_2|$沿曲线恒定。记$r_j=|\vct r-\vct F_j|>0$、$\vct u_j=(\vct r-\vct F_j)/r_j$。沿任意微小位移，距离的链式求导给出
\[
\dd r_j=\frac{2(\vct r-\vct F_j)\cdot\dd\vct r}{2r_j}
=\vct u_j\cdot\dd\vct r,
\qquad \dd S=(\vct u_1+\vct u_2)\cdot\dd\vct r.
\]
沿椭圆切线有$\dd S=0$，故法线平行于两个单位向量之和，平分两焦点方向的夹角。以单位法向$\vct n=(\vct u_1+\vct u_2)/|\vct u_1+\vct u_2|$写反射定律，出射方向为
\[
\vct u_{\rm out}=\vct u_1-2(\vct u_1\cdot\vct n)\vct n
=-\vct u_2,
\]
这里用到$|\vct u_1+\vct u_2|^2=2(1+\vct u_1\cdot\vct u_2)$。所以从一焦点发出的光经椭圆反射后指向另一焦点。实际三维装置还需考虑反射面的空间形状。

\section{抛体轨迹、可达区域与包络}
\subsection{消去时间得到轨迹}
忽略空气阻力，在恒定重力场中取$x$水平、$y$向上。以速率$v_0$、仰角$\alpha$从原点抛出，方程$\ddot x=0,\ddot y=-g$先积分为
$\dot x=v_0\cos\alpha$、$\dot y=v_0\sin\alpha-gt$，其中常数由初速度确定。再用初位置为零积分，
\begin{align*}
x&=\int_0^tv_0\cos\alpha\dd s=[v_0s\cos\alpha]_0^t,\\
y&=\int_0^t(v_0\sin\alpha-gs)\dd s
=\left[v_0s\sin\alpha-\frac12gs^2\right]_0^t.
\end{align*}
因此
\begin{equation}
x=v_0\cos\alpha\,t,\qquad y=v_0\sin\alpha\,t-\frac12gt^2.
\end{equation}
若 $\cos\alpha\ne0$，消去时间得到 $y=x\tan\alpha-gx^2/(2v_0^2\cos^2\alpha)$。当起落等高且 $0<\alpha<\pi/2$ 时，飞行时间为 $2v_0\sin\alpha/g$，射程为 $v_0^2\sin2\alpha/g$，最大射程在 $\alpha=\pi/4$。

上述最优角适用于无阻力、起落等高的模型。改变起落高度或加入阻力后，应重新确定飞行时间与射程。

\subsection{固定初速时哪些点能到达}
给定 $v_0$，而仰角可变。令 $u=\tan\alpha$，用 $1/\cos^2\alpha=1+u^2$，轨迹写成
\begin{equation}
y=xu-\frac{gx^2}{2v_0^2}(1+u^2).
\end{equation}
对于固定$x>0$，对发射参数求导，
\[\frac{\partial y}{\partial u}=x-\frac{gx^2}{v_0^2}u,\qquad
\frac{\partial^2y}{\partial u^2}=-\frac{gx^2}{v_0^2}<0.\]
故唯一极大点为$u=v_0^2/(gx)$。代回原式，$xu=v_0^2/g$，而二次项为$gx^2/(2v_0^2)+v_0^2/(2g)$，相减得最大可达高度
\begin{equation}
y_{\rm env}(x)=\frac{v_0^2}{2g}-\frac{gx^2}{2v_0^2}.
\end{equation}
这条曲线是各条抛物线轨迹的包络。目标点位于它下方时，一般有两个发射角；在包络上时二者合并；在包络上方时没有实发射角。在 $x=0$ 处可取竖直发射的极限，但前面消去时间的式子本身不直接包含竖直轨迹。
\begin{figure}[htbp]
\centering
\begin{tikzpicture}[x=5.5cm,y=5cm,>=Stealth]
\draw[->](0,0)--(1.18,0)node[right]{$X$};\draw[->](0,0)--(0,.64)node[above]{$Y$};
\foreach\ang in{15,30,45,60,75}{\pgfmathsetmacro{\endx}{sin(2*\ang)}\draw[domain=0:\endx,samples=70]plot(\x,{\x*tan(\ang)-\x*\x/(2*cos(\ang)*cos(\ang))});}
\draw[thick,domain=0:1,samples=70]plot(\x,{.5-.5*\x*\x});
\node[above right]at(.66,.31){包络};\node[below]at(1,0){$1$};\node[left]at(0,.5){$1/2$};
\end{tikzpicture}
\caption{固定初速的抛体轨迹与包络。横、纵坐标分别为 $X=gx/v_0^2$、$Y=gy/v_0^2$；细线对应 $15^\circ,30^\circ,45^\circ,60^\circ,75^\circ$ 发射角，粗线为包络。}\label{fig:projectile-envelope}
\end{figure}

\begin{example}
以固定初速 $v_0$ 射向目标 $(X,Y)$，$X>0$。求可达条件与发射角。
\end{example}
% 题设配图；与紧随其后的解答分别编号。
\ExpansionDiagram{fig:audit-m-q-037}{固定初速下向指定目标发射}{
\draw[->](-.4,0)--(5.8,0)node[right]{$ x $};\draw[->](0,-0.3)--(0,3.1)node[above]{$ y $};\fill(0,0)circle(2pt);\draw[->,thick](0,0)--(1.3,1)node[above]{$v_0$};\draw[->](.6,0)arc(0:37.6:.6);\node[right]at(.65,0.2){$\alpha$};\fill(4,1.7)circle(2pt)node[above right]{$P(X,Y)$};\draw[dashed](4,0)--(4,1.7)--(0,1.7);\draw[->](6.2,2.8)--(6.2,1.9)node[right]{$g$};
}

\begin{solution}
图\ref{fig:audit-m-q-037}中目标$P(X,Y)$固定，未知的是发射方向与到达时间。
图\ref{fig:eg-2-2}的粗线给出固定初速的可达边界。下面把“目标在边界之下”写成判别式条件。

令发射方向与正 $x$ 轴夹角为 $\alpha$，并取 $\cos\alpha>0$，使粒子能到达 $X>0$ 的目标。无阻力运动满足
\[
x=v_0\cos\alpha\,t,\qquad y=v_0\sin\alpha\,t-\frac12gt^2.
\]
由第一式消去 $t$，再令 $u=\tan\alpha$，得
\[
y=xu-\frac{gx^2}{2v_0^2}(1+u^2).
\]
在目标点代入 $x=X,y=Y$，整理为
\[
gX^2u^2-2v_0^2Xu+(gX^2+2v_0^2Y)=0.
\]
发射角存在的必要条件是二次方程有实根，其判别式为 $4X^2[v_0^4-g(gX^2+2v_0^2Y)]$，故
\[
Y\le\frac{v_0^2}{2g}-\frac{gX^2}{2v_0^2}.
\]
满足严格不等式时，一般有两个角度：
\[
\tan\alpha=\frac{v_0^2\pm\sqrt{v_0^4-g(gX^2+2v_0^2Y)}}{gX}.
\]
等号成立时两角重合，目标在包络上。每一根还要按允许的发射方向、地面和障碍物条件筛选；对应飞行时间为 $X/(v_0\cos\alpha)>0$。

\begin{figure}[H]
\centering
\begin{tikzpicture}[x=5.5cm,y=5cm,>=Stealth]
\draw[->](0,0)--(1.18,0)node[right]{$gx/v_0^2$};\draw[->](0,0)--(0,.64)node[above]{$gy/v_0^2$};
\foreach\ang in{15,30,45,60,75}{\pgfmathsetmacro{\endx}{sin(2*\ang)}\draw[domain=0:\endx,samples=70]plot(\x,{\x*tan(\ang)-\x*\x/(2*cos(\ang)*cos(\ang))});}
\draw[thick,domain=0:1,samples=70]plot(\x,{.5-.5*\x*\x});
\node[above right]at(.66,.31){包络};\node[below]at(1,0){$1$};\node[left]at(0,.5){$1/2$};
\end{tikzpicture}
\caption{例2.2解图：粗线为可达边界，细线为不同角度轨迹。坐标按 $v_0^2/g$ 缩放，目标在粗线下方才有实发射角。}\label{fig:eg-2-2}
\end{figure}

目标在包络内部时通常有两条轨迹，位于包络上时两发射角合并；这与角度二次方程的两个根、重根相对应。
\end{solution}

\section{极坐标运动学与中心力}
中心力的径向与绕心转动适合用极坐标表示。基矢随位置转动的贡献由位置矢量求导得到，再与受力及守恒关系连接。

\subsection{基矢的时间导数}
平面极坐标取 $x=r\cos\theta,y=r\sin\theta$。定义 $\vct e_r=(\cos\theta,\sin\theta)$，$\vct e_\theta=(-\sin\theta,\cos\theta)$。求导得 $\dot{\vct e}_r=\dot\theta\vct e_\theta$，$\dot{\vct e}_\theta=-\dot\theta\vct e_r$。位置为 $\vct r=r\vct e_r$，于是
\begin{align}
\vct v&=\dot r\vct e_r+r\dot\theta\vct e_\theta,\\
\vct a&=\ddot r\vct e_r+\dot r\dot\theta\vct e_\theta
+(\dot r\dot\theta+r\ddot\theta)\vct e_\theta-r\dot\theta^2\vct e_r\\
&=(\ddot r-r\dot\theta^2)\vct e_r+(r\ddot\theta+2\dot r\dot\theta)\vct e_\theta.
\end{align}
向心项来自切向基矢的改变，$2\dot r\dot\theta$ 则由两次求导中的两个同类项相加而来。漏掉基矢求导会把圆周运动的加速度错误地算成零。

\begin{figure}[htbp]
\centering
\begin{tikzpicture}[scale=.85,>=Stealth]
\draw[->] (0,0)--(4.2,0)node[right]{$x$};\draw[->](0,0)--(0,3.4)node[above]{$y$};
\coordinate(P)at(40:2.8);\draw[->,thick](0,0)--(P)node[midway,above]{$r$};
\draw[->](P)--++(40:1.1)node[right]{$\vct e_r$};\draw[->](P)--++(130:1.1)node[above]{$\vct e_\theta$};
\draw(.7,0)arc(0:40:.7);\node at(20:1){$\theta$};\fill(P)circle(1.5pt);
\end{tikzpicture}
\caption{极坐标的两个单位基矢}\label{fig:polar-basis}
\end{figure}

\begin{example}
粒子满足 $r=at,\theta=bt$，$m,a,b>0,t>0$。求作用力的分量，并判断它是否是中心力。
\end{example}
% 题设配图；与紧随其后的解答分别编号。
\ExpansionDiagram{fig:audit-m-q-038}{极坐标轨迹和随位置旋转的基矢}{
\draw[->](-.4,0)--(4,0)node[right]{$ x $};\draw[->](0,-0.4)--(0,3.6)node[above]{$ y $};\draw[thick,domain=0:75,samples=100]plot({.04*\x*cos(\x)},{.04*\x*sin(\x)});\draw[->,thick](0,0)--(1.2,2.07846)node[above right]{$P$};\node[left]at(.5,1){$r=at$};\draw[->](.65,0)arc(0:60:.65);\node[right]at(.7,.45){$\theta=bt$};\draw[->](1.2,2.07846)--(1.7,2.94449)node[right]{$\vct e_r$};\draw[->](1.2,2.07846)--(.33397,2.57846)node[left]{$\vct e_\theta$};
}

\begin{solution}
图\ref{fig:audit-m-q-038}中的$r$与$\theta$均随时间变化，径向和横向基矢也必须求导。
图\ref{fig:eg-2-3}把速度分成径向与横向两项；求导时还必须计入两根基矢随角度转动。

力应由完整加速度求出，不能只对径向距离作两次导数。极坐标速度为
\[
\vct v=\dot r\vct e_r+r\dot\theta\vct e_\theta
=a\vct e_r+abt\vct e_\theta.
\]
利用 $\dot{\vct e}_r=b\vct e_\theta$、$\dot{\vct e}_\theta=-b\vct e_r$ 再求导，有
\[
\vct a=-ab^2t\vct e_r+(ab+ab)\vct e_\theta.
\]
于是 $F_r=-mab^2t$、$F_\theta=2mab$。第二项的两个贡献分别来自径向速度的转向与横向速度分量的增长。

\begin{figure}[H]
\centering
\begin{tikzpicture}[scale=.85,>=Stealth]
\draw[->] (0,0)--(4.2,0)node[right]{$x$};\draw[->](0,0)--(0,3.4)node[above]{$y$};
\coordinate(P)at(40:2.8);\draw[->,thick](0,0)--(P)node[midway,above]{$r$};
\draw[->](P)--++(40:1.1)node[right]{$\vct e_r$};\draw[->](P)--++(130:1.1)node[above]{$\vct e_\theta$};
\draw(.7,0)arc(0:40:.7);\node at(20:1){$\theta$};\fill(P)circle(1.5pt);
\end{tikzpicture}
\caption{例2.3解图：径向与横向基矢随粒子转动。速度在两方向上的分量分别为 $a,abt$，横向加速度需计入基矢求导。}\label{fig:eg-2-3}
\end{figure}

也可检查角动量：$L_z=mr^2\dot\theta=ma^2bt^2$，其导数 $2ma^2bt$ 等于 $rF_\theta$。因 $F_\theta\ne0$，合力不始终指向或背离原点，不能视为相对原点的中心力。

图中的两分量在本题分别为 $a,abt$。非零横向力说明合力偏离径向，因而不是仅由中心力产生的运动。
\end{solution}

\subsection{角动量与面积速度}
中心力没有切向分量，故 $r\ddot\theta+2\dot r\dot\theta=0$。乘以 $r$ 得 $\dd(r^2\dot\theta)/\dd t=0$。角动量 $L_z=mr^2\dot\theta$ 因而守恒。

短时间扫过的扇形面积为 $\dd S=\tfrac12r^2\dd\theta$，故面积速度 $\dot S=L_z/(2m)$ 恒定。这个结论不要求轨迹是椭圆，也不要求力满足反平方律；它只依赖于关于该中心的力矩为零。反平方律决定轨迹形状，中心力决定面积定律，两种条件的作用应分开。

\subsection{反平方力与圆锥曲线轨道}
考虑单位质量的引力加速度 $-\mu\vct e_r/r^2$，其中 $\mu>0$。设单位质量角动量 $h=r^2\dot\theta\ne0$，令 $u=1/r$，用撇号表示对 $\theta$ 求导。由链式法则
\[
\dot r=-u^{-2}u'\dot\theta=-hu',\qquad
\ddot r=-hu''\dot\theta=-h^2u^2u'',\qquad
r\dot\theta^2=h^2u^3.
\]
代入径向方程$\ddot r-r\dot\theta^2=-\mu/r^2$，约去非零的$-h^2u^2$，得到$u''+u=\mu/h^2$。常数特解为$\mu/h^2$，齐次特征方程为$\lambda^2+1=0$，故
\[u=\frac\mu{h^2}+A\cos\theta+B\sin\theta.\]
对非圆轨道，选择近心点为$\theta=0$时$u'(0)=0$，得$B=0$；再以偏心率$e\ge0$记$A=e\mu/h^2$。圆轨道对应$A=B=e=0$。于是轨道为
\begin{equation}
r=\frac p{1+e\cos\theta},\qquad p=\frac{h^2}{\mu},\quad e\ge0.
\end{equation}
这正是焦点位于原点的圆锥曲线极坐标方程。$0\le e<1$ 为椭圆，$e=1$ 为抛物线，$e>1$ 为双曲线；只有分母为正的角域对应物理上的正半径。

径向与切向速度分别为 $\dot r=he\sin\theta/p$ 和 $r\dot\theta=h(1+e\cos\theta)/p$。将$p=h^2/\mu$代入并相加两分量的平方，
\[
\begin{aligned}
v^2&=\frac{\mu^2}{h^2}
\left[e^2\sin^2\theta+(1+e\cos\theta)^2\right]\\
&=\frac{\mu^2}{h^2}(1+2e\cos\theta+e^2),
\qquad \frac\mu r=\frac{\mu^2}{h^2}(1+e\cos\theta).
\end{aligned}
\]
动能中的$e\cos\theta$项与势能中的同项相消，因此单位质量机械能为
\begin{equation}
\mathcal E=\frac{v^2}{2}-\frac\mu r
=\frac{\mu^2(e^2-1)}{2h^2}.
\end{equation}
角度项相消，所以能量确为常量。对椭圆，$r_{\min}=p/(1+e),r_{\max}=p/(1-e)$，半长轴 $a=(r_{\min}+r_{\max})/2=p/(1-e^2)$，从而 $\mathcal E=-\mu/(2a)$。几何尺度、角动量和能量在同一推导中连起来。

\section{曲线的切向、法向与曲率}
\subsection{从参数速度到弧长参数}
设光滑曲线 $\vct r(u)$ 的导数非零。相邻两点间弧长微元为 $\dd s=|\vct r'(u)|\dd u$，因此单位切向量 $\vct T=\dd\vct r/\dd s=\vct r'(u)/|\vct r'(u)|$。参数速度 $\dd\vct r/\dd u$ 与物理速度不同，除非 $u$ 明确就是时间。

定义曲率 $\kappa=|\dd\vct T/\dd s|$。因为 $\vct T\cdot\vct T=1$，求导得到 $\vct T\cdot\dd\vct T/\dd s=0$，所以切向量的变化垂直于切线。曲率非零时，定义主法向量 $\vct N=(\dd\vct T/\dd s)/\kappa$，它指向曲线局部弯曲的一侧。曲率半径 $\rho_c=1/\kappa$ 描述该点密切圆的半径。

物理速度为 $\vct v=v\vct T$，继续求导，得到
\[
\vct a=\dot v\vct T+v^2\kappa\vct N.
\]
切向加速度改变速率，法向加速度改变方向。这与极坐标的径向、切向分解不同：只有在圆周等特定曲线上，主法向才与径向直接一致。

平面参数曲线令$w=|\vct r'|=\sqrt{x'^2+y'^2}>0$，则
\[w'=\frac{\vct r'\cdot\vct r''}{w},\qquad
\frac{\dd\vct T}{\dd s}
=\frac1w\frac{\dd}{\dd u}\left(\frac{\vct r'}w\right)
=\frac{\vct r''}{w^2}
-\frac{\vct r'(\vct r'\cdot\vct r'')}{w^4}.\]
取模平方，并用前面的叉积模恒等式，得到
\[\kappa^2=\frac{w^2|\vct r''|^2-(\vct r'\cdot\vct r'')^2}{w^6}
=\frac{|\vct r'\times\vct r''|^2}{w^6}.\]
二维叉积只剩垂直平面的分量$x'y''-y'x''$，因此
\[
\kappa=\frac{|x'y''-y'x''|}{(x'^2+y'^2)^{3/2}}.
\]
若直接写成 $y=f(x)$，则 $\kappa=|f''|/(1+f'^2)^{3/2}$。分母记录斜率对弧长的影响，不能把二阶导数本身当成任意位置的曲率。

\textbf{算例：匀速圆周运动。} 圆周运动的速率不变时，加速度是否为零？写出 $\vct v=v\vct e_t$ 并求导。
% 题设配图；与紧随其后的解答分别编号。
\ExpansionDiagram{fig:audit-m-q-039}{速率恒定时速度方向仍沿圆周变化}{
\draw[thick](0,0)circle(1.4);\fill(0,0)circle(1.3pt)node[below]{$O$};\fill(1.4,0)circle(2pt)node[below right]{$P$};\draw[->](0,0)--(1.4,0)node[midway,below]{$R$};\draw[->,thick](1.4,0)--(1.4,1.1)node[right]{$\vct v$};\draw[->](0,-1.9)--(2.5,-1.9)node[right]{$x$};
}

\begin{solution}
题设图\ref{fig:audit-m-q-039}中速度沿切线，解图\ref{fig:warm-2}补出指向圆心的加速度。设$OP=R$，取
$\vct e_r=(\cos\theta,\sin\theta)$、
$\vct e_t=(-\sin\theta,\cos\theta)$，则
\[
 \dot{\vct e}_t=-\dot\theta\,\vct e_r,\qquad v=R\dot\theta.
\]
对$\vct v=v\vct e_t$求导，得到
\[
 \vct a=\dot v\,\vct e_t-v\dot\theta\,\vct e_r
 =-\frac{v^2}{R}\vct e_r.
\]
速率恒定只使$\dot v=0$。基矢的导数仍非零，所以加速度沿图中从$P$指向$O$的箭头，大小为$v^2/R$。

\begin{figure}[H]\centering
\begin{tikzpicture}[scale=1.2]
\draw[bookguide](0,0)circle(1.3);\fill(1.3,0)circle(1.7pt)node[below right]{$P$};\fill(0,0)circle(1.3pt)node[left]{$O$};
\draw[bookaxis,thick](1.3,0)--(1.3,1.1)node[right]{$\vct v$};\draw[bookaxis,thick](1.3,0)--(.35,0)node[midway,above]{$\vct a$};
\node at(.2,-1.8){$|\vct a|=v^2/R$，指向圆心};
\end{tikzpicture}
\caption{匀速圆周运动的速度与加速度：速率恒定仍有切向单位矢量的变化，加速度垂直速度、指向圆心。}\label{fig:warm-2}
\end{figure}
\end{solution}

\begin{example}
粒子以恒定速率 $v$ 沿 $y=ax^2$ 运动，$a>0$。求顶点的加速度和密切圆。
\end{example}
% 题设配图；与紧随其后的解答分别编号。
\ExpansionDiagram{fig:audit-m-q-040}{顶点处的切线与待求的密切圆}{
\draw[->](-.4,0)--(3.5,0)node[right]{$ x $};\draw[->](0,-0.4)--(0,2.8)node[above]{$ y $};\draw[thick,domain=-2:2,samples=70]plot(\x,{.4*\x*\x});\draw[dashed](-1,0)--(1.8,0);\fill(0,0)circle(2pt);\draw[->](0,.15)--(1.2,.15)node[right]{$v$};\node[left]at(-1.8,1.8){$y=ax^2$};
}

\begin{solution}
图\ref{fig:audit-m-q-040}中顶点的切线水平。速率恒定时，切向加速度为零，加速度沿曲线凹侧的主法线；密切圆心也在这一方向上，见图\ref{fig:orig-029}。

对函数图像 $y=f(x)$，沿 $x$ 增大方向的单位切向为 $(1,f')/\sqrt{1+f'^2}$。本题 $f'=2ax,f''=2a$，在顶点 $x=0$，切线水平，弯曲侧向上。

\begin{figure}[H]\centering
\begin{tikzpicture}[scale=1]
\draw[bookaxis](-1.6,0)--(1.8,0)node[right]{$x$};\draw[bookaxis](0,-.2)--(0,2.5)node[above]{$y$};
\draw[dashed](0,1)circle(1);\draw[thick,domain=-1.5:1.5,samples=100]plot(\x,{.5*\x*\x});
\fill(0,1)circle(1.3pt)node[right]{$C$};\draw[bookaxis](0,0)--(0,.7)node[midway,left]{$\vct N$};\draw[bookaxis](0,0)--(.85,0)node[below]{$\vct T$};
\draw[<->](.35,0)--(.35,1)node[midway,right]{$\rho$};\node[below left]at(0,0){$O$};
\end{tikzpicture}
\caption{抛物线顶点的密切圆。图取 $a=1/2$，虚线圆与实线曲线只有局部接触。}
\label{fig:orig-029}
\end{figure}

由曲率公式，
\[
\kappa(0)=\left.\frac{|f''|}{(1+f'^2)^{3/2}}\right|_{x=0}=2a,
\qquad\rho=\frac1{2a}.
\]
速率恒定时 $\dot v=0$，加速度全部沿主法向，因此
\[
\vct a(0)=v^2\kappa(0)\vct e_y=(0,2av^2).
\]
密切圆圆心在顶点沿正 $y$ 走距离 $\rho$ 处，圆方程为
\[
x^2+(y-\rho)^2=\rho^2.
\]
取经过原点的下支，有
\[
y=\rho-\sqrt{\rho^2-x^2}
=\rho\left[1-\sqrt{1-\frac{x^2}{\rho^2}}\right].
\]
在 $|x|\ll\rho$ 时，用 $\sqrt{1-u}=1-u/2+O(u^2)$ 展开，得到
\[
y=\frac{x^2}{2\rho}+O(x^4)=ax^2+O(x^4).
\]
圆与原曲线的位置、切线和曲率相同，因而密切圆给出局部二阶近似。

\end{solution}

\section{旋转坐标与惯性力}
\subsection{主动旋转与被动换轴}
在固定直角坐标中，向量主动逆时针转角 $\alpha$，分量变为 $R(\alpha)\vct x$，其中
\begin{equation}
R(\alpha)=\begin{pmatrix}\cos\alpha&-\sin\alpha\\\sin\alpha&\cos\alpha\end{pmatrix}.
\end{equation}
若向量本身固定，仅把坐标轴逆时针转 $\alpha$，新分量则为 $R(-\alpha)\vct x$。两者互为逆变换。通过检查固定向量 $(1,0)$ 在新坐标下的 $y'$ 分量为负，可以快速检查旋转矩阵符号。

\subsection{旋转系中的求导法则}
设旋转基矢的角速度为 $\vct\Omega(t)$，则每个基矢满足 $\dot{\vct e}_j=\vct\Omega\times\vct e_j$。对任意向量 $\vct A=\sum A_j\vct e_j$，惯性系导数与旋转系分量导数关系为
\begin{equation}
\left(\frac{\dd\vct A}{\dd t}\right)_{\rm in}
=\left(\frac{\dd\vct A}{\dd t}\right)_{\rm rot}+\vct\Omega\times\vct A.
\end{equation}
先假设两坐标系原点重合且固定。该法则用于位置得到$\vct v_{\rm in}=\vct v_{\rm rel}+\vct\Omega\times\vct r$。再作惯性系导数，
\[\begin{aligned}
\vct a_{\rm in}
&=\vct a_{\rm rel}+\vct\Omega\times\vct v_{\rm rel}
+\dot{\vct\Omega}\times\vct r
+\vct\Omega\times\vct v_{\rm in}\\
&=\vct a_{\rm rel}+\vct\Omega\times\vct v_{\rm rel}
+\dot{\vct\Omega}\times\vct r
+\vct\Omega\times(\vct v_{\rm rel}+\vct\Omega\times\vct r).
\end{aligned}\]
合并两个同类项，得到
\begin{align}
\vct v_{\rm in}&=\vct v_{\rm rel}+\vct\Omega\times\vct r,\\
\vct a_{\rm in}&=\vct a_{\rm rel}+2\vct\Omega\times\vct v_{\rm rel}
+\dot{\vct\Omega}\times\vct r+\vct\Omega\times(\vct\Omega\times\vct r).
\end{align}
若旋转系原点还在平动，应另外加上原点的惯性系速度和加速度。

牛顿方程在旋转系中写成
\[
m\vct a_{\rm rel}=\vct F-2m\vct\Omega\times\vct v_{\rm rel}
-m\dot{\vct\Omega}\times\vct r-m\vct\Omega\times(\vct\Omega\times\vct r).
\]
右侧新增项分别是科里奥利力、欧拉力和离心力。这些项来自坐标系的运动，应与物体之间的实作用力分别记录。

\begin{example}
物体固定在匀速旋转圆盘上，距轴 $R$，角速度为 $\Omega$。解释两坐标系中受力方程。
\end{example}
% 题设配图；与紧随其后的解答分别编号。
\ExpansionDiagram{fig:audit-m-q-041}{随匀速转盘固定的质点}{
\draw[thick](0,0)circle(1.4);\fill(0,0)circle(1.3pt)node[below]{$O$};\fill(1.4,0)circle(2pt)node[below right]{$P$};\draw[->](0,0)--(1.4,0)node[midway,below]{$R$};\draw[->,thick](1.4,0)--(1.4,1.1)node[right]{$\vct v$};\draw[->](0,-1.9)--(2.5,-1.9)node[right]{$x$};\draw[->](-1.1,-.5)arc(204:260:1.2);\node[left]at(-1.2,-1){$\Omega$};
}

\begin{solution}
图\ref{fig:audit-m-q-041}中质点相对转盘静止。图\ref{fig:eg-2-5}分别画出两个参考系的受力，其中虚箭头表示只加入转动系方程的惯性力。

惯性系中，物体虽然与圆盘没有相对运动，却在空间作半径为 $R$ 的圆周运动。由加速度 $\vct a=-\Omega^2R\vct e_r$，牛顿方程要求真实合力为
\[
\vct F_{\rm real}=-m\Omega^2R\vct e_r.
\]
这部分力必须由实际约束提供，例如静摩擦或连接件。若约束不能提供所需向心力，物体便无法保持固定。

\begin{figure}[H]\centering
\begin{tikzpicture}[scale=.88]
\foreach\x/\title in{0/惯性系,5.6/随盘转动系}{
\begin{scope}[shift={(\x,0)}]
\draw[bookguide](0,0)circle(1.3);\fill(0,0)circle(1.3pt)node[left]{$O$};\fill(1.3,0)circle(2pt)node[below]{$m$};
\draw[bookaxis,thick](1.3,0)--(.35,0)node[midway,above]{$F_{\rm real}$};\node at(.6,-1.8){\title};
\end{scope}}
\draw[bookaxis](1.3,0)--(1.3,1.1)node[right]{$\vct v$};
\draw[bookaxis,dashed](6.9,0)--(8,0)node[right]{$m\Omega^2R$};
\node at(5.9,1.7){$\vct v_{\rm rel}=0$};
\end{tikzpicture}
\caption{例2.5：左图由真实向心力产生圆周加速度；右图在转动系中另加入离心惯性力（虚箭头），静止物体满足平衡。}\label{fig:eg-2-5}
\end{figure}

改用与圆盘一同转动的坐标系，$\vct v_{\rm rel}=\vct a_{\rm rel}=0$，且 $\dot{\vct\Omega}=0$。各项方程变为
\[
0=\vct F_{\rm real}-2m\vct\Omega\times0
-m\vct\Omega\times(\vct\Omega\times\vct r).
\]
利用向量三重积及 $\vct\Omega\cdot\vct r=0$，有
\[
\vct\Omega\times(\vct\Omega\times\vct r)
=\vct\Omega(\vct\Omega\cdot\vct r)-\Omega^2\vct r
=-\Omega^2R\vct e_r.
\]
所以离心惯性力为 $+m\Omega^2R\vct e_r$，与真实向心力平衡。科里奥利项为零，因为它使用的是相对速度；不能把惯性系圆周速度 $\Omega R$ 代入该项。

\end{solution}

\section{从约束选坐标：抛物线轨道上的小珠}
\subsection{把几何关系写进动能}
光滑细轨道位于竖直平面内，形状为 $y=x^2/(2R)$，$R>0$；$y$ 轴竖直向上。质量为 $m$ 的小珠从 $x=a>0$ 处由静止释放。求运动方程、最低点的支持力，并研究振幅较小时的周期。

轨道只允许一个独立位移，选 $x$ 为坐标即可。高度和竖直速度由约束确定：
\[
y=\frac{x^2}{2R},\qquad \dot y=\frac xR\dot x.
\]
于是
\begin{equation}\label{eq:bead-energy}
T=\frac m2\left(1+\frac{x^2}{R^2}\right)\dot x^2,\qquad
U=\frac{mgx^2}{2R}.
\end{equation}
虽然势能恰好是 $x$ 的二次函数，动能的系数仍随位置变化，所以有限振幅运动并非严格简谐。只看势能形状会漏掉约束的几何作用。

\begin{figure}[htbp]\centering
\begin{tikzpicture}[scale=1.15]
\draw[bookaxis](-2.3,0)--(2.5,0)node[right]{$x$};\draw[bookaxis](0,-.3)--(0,2.6)node[above]{$y$};
\draw[bookcurve,domain=-2.1:2.1,samples=80]plot(\x,{\x*\x/2});
\fill(1.2,.72)circle(2pt);\node[right]at(1.25,.75){$P$};
\draw[bookaxis,thick](1.2,.72)--(1.2,-.2)node[below]{$mg$};
\draw[bookaxis,thick](1.2,.72)--(.36,1.42)node[left]{$N$};
\draw[bookguide](.8,.24)--(1.8,1.44);\node[right]at(1.8,1.44){切线};
\draw[bookguide](2,0)--(2,2);\node[below]at(2,0){$a$};\node[left]at(-.1,.12){$O$};
\end{tikzpicture}
\caption{小珠的重力与支持力。支持力沿轨道法线，水平位移与沿轨道的路程不同。图中坐标按 $R$ 缩放。}
\label{fig:orig-030}
\end{figure}

\subsection{沿切线列方程，与能量法互相核对}
对总能量求时间导数，乘积及链式法则给出
\[\begin{aligned}
0&=\frac{\dd}{\dd t}\left[\frac m2\left(1+\frac{x^2}{R^2}\right)\dot x^2
+\frac{mgx^2}{2R}\right]\\
&=m\left(1+\frac{x^2}{R^2}\right)\dot x\ddot x
+\frac{mx\dot x^3}{R^2}+\frac{mgx\dot x}{R}.
\end{aligned}\]
在$\dot x\ne0$的位置约去$m\dot x$，得到
\begin{equation}\label{eq:bead-ode}
\left(1+\frac{x^2}{R^2}\right)\ddot x+
\frac{x\dot x^2}{R^2}+\frac gR x=0.
\end{equation}
方程在转折点由连续性成立。也可直接将牛顿方程投影到切线：切向单位矢量为 $(1,x/R)/\sqrt{1+x^2/R^2}$；加速度为 $(\ddot x,(\dot x^2+x\ddot x)/R)$。二者点积乘以 $m$，等于重力的切向分量 $-mgx/(R\sqrt{1+x^2/R^2})$，化简后得到同一方程。第二条路线说明支持力确实被正确消去。

最低点速度由能量得 $v_O^2=ga^2/R$。该点曲率半径为 $R$，法向向上，故
\[
N_O-mg=\frac{mv_O^2}{R},\qquad N_O=mg\left(1+\frac{a^2}{R^2}\right).
\]
当 $a\to0$，支持力趋于 $mg$；振幅增大时，轨道为小珠提供更大的向心加速度。

\subsection{周期修正来自哪一项}
由式\eqref{eq:bead-energy}，在从最低点走向右端的四分之一周期中，
\[
\dot x=\sqrt{\frac gR}\sqrt{\frac{a^2-x^2}{1+x^2/R^2}},\qquad
\frac T4=\sqrt{\frac Rg}\int_0^a\frac{\sqrt{1+x^2/R^2}}{\sqrt{a^2-x^2}}\dd x.
\]
令$x=a\sin u$，则$\dd x=a\cos u\dd u$、$\sqrt{a^2-x^2}=a\cos u$，$x=0,a$对应$u=0,\pi/2$，转折点根式因子因而消去。记$\epsilon=a/R$，展开被积函数，
\[\sqrt{1+\epsilon^2\sin^2u}
=1+\frac12\epsilon^2\sin^2u+O(\epsilon^4).\]
其余项在整个固定区间一致受控。利用
\[\int_0^{\pi/2}\sin^2u\dd u
=\left[\frac u2-\frac{\sin2u}{4}\right]_0^{\pi/2}=\frac\pi4,\]
得到
\begin{align}
T&=4\sqrt{R/g}\int_0^{\pi/2}\sqrt{1+\epsilon^2\sin^2u}\dd u\\
&=2\pi\sqrt{R/g}\left[1+\frac{\epsilon^2}{4}+O(\epsilon^4)\right].
\end{align}
周期随振幅增长。由于 $\dot x$ 的典型量级为 $a\sqrt{g/R}$，运动方程中的 $x^2\ddot x/R^2$ 和 $x\dot x^2/R^2$ 相对线性项都是 $O(\epsilon^2)$；若保留其中一个而删去另一个，得到的二阶周期修正通常会错。

\subsection{曲轨运动的应用}
\textbf{首阶周期。} 对轨道 $y=x^2/(2R)$，从 $a\ll R$ 出发求首阶周期。

% 题设配图；与紧随其后的解答分别编号。
\ExpansionDiagram{fig:audit-m-q-042}{抛物线轨道上的小振幅释放点}{
\draw[->](-.4,0)--(3,0)node[right]{$ x $};\draw[->](0,-0.4)--(0,2.4)node[above]{$ y $};\draw[thick,domain=-2:2,samples=80]plot(\x,{.35*\x*\x});\fill(1.7,{.35*1.7^2})circle(2pt)node[above right]{$A$};\draw[dashed](1.7,0)--(1.7,{.35*1.7^2});\node[below]at(1.7,0){$a$};\node[left]at(-1.7,1.8){$ y=x^2/(2R) $};\draw[->](2.6,1.7)--(2.6,.7)node[right]{$g$};
}

\begin{solution}
图\ref{fig:audit-m-q-042}的释放点横坐标为$a$。运动中$x/R=O(a/R)$，故小振幅极限下轨道斜率也趋零。结合解图\ref{fig:class-2-1}的约束关系，有
\[
 T=\frac m2(1+x^2/R^2)\dot x^2,\qquad U=\frac{mg}{2R}x^2.
\]
取线性运动所需的二次能量，动能保留$m\dot x^2/2$，势能保留$mgx^2/(2R)$。
等价地，在完整运动方程中删去相对量级为$O((a/R)^2)$的两项，得到
\[
 \ddot x+\frac gR x=0.
\]
特征方程为$\lambda^2+g/R=0$，记$\omega_0=\sqrt{g/R}$，通解为$x=A\cos\omega_0t+B\sin\omega_0t$。由$x(0)=A=a$、$\dot x(0)=B\omega_0=0$得$B=0$，所以$x=a\cos(\sqrt{g/R}\,t)$，故
$\omega_0=\sqrt{g/R}$、$T_0=2\pi\sqrt{R/g}$。
首阶周期不依赖质量；弧长与水平距离的差别进入下一阶振幅修正。

\begin{figure}[H]\centering
\begin{tikzpicture}[scale=1.15]
\draw[bookaxis](-2.3,0)--(2.5,0)node[right]{$x$};\draw[bookaxis](0,-.3)--(0,2.6)node[above]{$y$};
\draw[bookcurve,domain=-2.1:2.1,samples=80]plot(\x,{\x*\x/2});
\fill(1.2,.72)circle(2pt);\node[right]at(1.25,.75){$P$};
\draw[bookaxis,thick](1.2,.72)--(1.2,-.2)node[below]{$mg$};
\draw[bookaxis,thick](1.2,.72)--(.36,1.42)node[left]{$N$};
\draw[bookguide](.8,.24)--(1.8,1.44);\node[right]at(1.8,1.44){切线};
\draw[bookguide](2,0)--(2,2);\node[below]at(2,0){$a$};\node[left]at(-.1,.12){$O$};
\end{tikzpicture}
\caption{第二章首阶周期练习解图：轨道最低点附近切线近水平，高度为 $x^2/(2R)$，重力沿轨道的分量给出线性回复力。}
\label{fig:class-2-1}
\end{figure}
\end{solution}

\textbf{支持力。} 在同一轨道的任意位置，为什么仅用能量不能直接求支持力？补出计算式。

% 题设配图；与紧随其后的解答分别编号。
\ExpansionDiagram{fig:audit-m-q-043}{同一轨道上反力沿局部法线作用}{
\draw[->](-.4,0)--(3,0)node[right]{$ x $};\draw[->](0,-0.4)--(0,2.4)node[above]{$ y $};\draw[thick,domain=-2:2,samples=80]plot(\x,{.35*\x*\x});\fill(1.7,{.35*1.7^2})circle(2pt)node[above right]{$A$};\draw[dashed](1.7,0)--(1.7,{.35*1.7^2});\node[below]at(1.7,0){$a$};\node[left]at(-1.7,1.8){$ y=x^2/(2R) $};\draw[->](2.6,1.7)--(2.6,.7)node[right]{$g$};\draw[->](1.1,.4235)--(.4,1.3326)node[left]{$N$};
}

\begin{solution}
题设图\ref{fig:audit-m-q-043}把支持力画在局部法线上。取解图\ref{fig:class-2-2}中朝凹侧的主法向
$\vct n=(-x/R,1)/\sqrt{1+x^2/R^2}$，曲率为
$\kappa=1/[R(1+x^2/R^2)^{3/2}]$。支持力不做功，能量式中因而没有$N$，
它只能先确定速率。由$x=a$处静止释放，
\[
 \frac12mv^2+\frac{mgx^2}{2R}=\frac{mga^2}{2R},
 \qquad v^2=\frac gR(a^2-x^2).
\]
把重力投影到图中的法向，得$-mg/\sqrt{1+x^2/R^2}$。牛顿方程沿此方向为
\[
 N-\frac{mg}{\sqrt{1+x^2/R^2}}
 =\frac{mv^2}{R(1+x^2/R^2)^{3/2}}.
\]
代入速率并通分，最后得到
\[
 \boxed{N(x)=\frac{mg(1+a^2/R^2)}{(1+x^2/R^2)^{3/2}}}.
\]
最低点给出$N(0)=mg(1+a^2/R^2)$；转折点给出$N(a)=mg/\sqrt{1+a^2/R^2}$。
两处都与相应的法向加速度相符。

\begin{figure}[H]\centering
\begin{tikzpicture}[scale=1.15]
\draw[bookaxis](-2.3,0)--(2.5,0)node[right]{$x$};\draw[bookaxis](0,-.3)--(0,2.6)node[above]{$y$};
\draw[bookcurve,domain=-2.1:2.1,samples=80]plot(\x,{\x*\x/2});
\fill(1.2,.72)circle(2pt);\node[right]at(1.25,.75){$P$};
\draw[bookaxis,thick](1.2,.72)--(1.2,-.2)node[below]{$mg$};
\draw[bookaxis,thick](1.2,.72)--(.36,1.42)node[left]{$N$};
\draw[bookguide](.8,.24)--(1.8,1.44);\node[right]at(1.8,1.44){切线};
\draw[bookguide](2,0)--(2,2);\node[below]at(2,0){$a$};\node[left]at(-.1,.12){$O$};
\end{tikzpicture}
\caption{第二章支持力练习解图：法向指向弯曲侧，重力的法向投影与支持力合成向心加速度。}
\label{fig:class-2-2}
\end{figure}
\end{solution}

\textbf{改变轨道。} 将轨道改为半径 $R$ 的圆弧，比较两轨道的小振幅周期及首个振幅修正。

% 题设配图；与紧随其后的解答分别编号。
\ExpansionDiagram{fig:audit-m-q-044}{圆弧与抛物线在最低点具有相同曲率半径}{
\draw[->](-.4,0)--(3,0)node[right]{$ x $};\draw[->](0,-0.4)--(0,2.4)node[above]{$ y $};\draw[thick,domain=-2:2,samples=80]plot(\x,{.35*\x*\x});\fill(1.7,{.35*1.7^2})circle(2pt)node[above right]{$A$};\draw[dashed](1.7,0)--(1.7,{.35*1.7^2});\node[below]at(1.7,0){$a$};\node[left]at(-1.7,1.8){$ y=x^2/(2R) $};\draw[->](2.6,1.7)--(2.6,.7)node[right]{$g$};\draw[dashed,domain=-1.35:1.35,samples=60]plot(\x,{1.42857143-sqrt(1.42857143^2-\x*\x)});\node[right]at(2,2){实线抛物线};
}

\begin{solution}
图\ref{fig:audit-m-q-044}中比较两条曲线时，应先统一最低点的曲率半径，再比较各自的振幅修正。
图\ref{fig:class-2-3}中，圆弧与抛物线在最低点的曲率半径都为$R$，所以小振幅极限同为$T_0=2\pi\sqrt{R/g}$。比较首个修正时，两条轨道均以水平转折坐标$a$为振幅，取$a/R\ll1$。
\begin{figure}[H]\centering
\begin{tikzpicture}[scale=2.3]
\draw[bookaxis](-1.05,0)--(1.2,0)node[right]{$x/R$};\draw[bookaxis](0,-.1)--(0,1.05)node[above]{$y/R$};
\draw[bookcurve,domain=-.92:.92,samples=100]plot(\x,{(\x)^2/2});\draw[dashed,domain=-.92:.92,samples=100]plot(\x,{1-sqrt(1-(\x)^2)});
\fill(0,0)circle(.8pt);\node[right]at(.75,.18){抛物线};\node[above]at(.65,.73){圆弧};
\end{tikzpicture}
\caption{第二章随堂比较解图：两轨道在最低点的曲率半径同为 $R$，局部二阶形状相同，离开最低点后圆弧上升更快。}\label{fig:class-2-3}
\end{figure}
圆弧上取从最低点铅直向下半径量起的偏角$\theta$，则$x=R\sin\theta$、$y=R(1-\cos\theta)$、$\dd s=R\dd\theta$。设转折角为$\theta_0$，能量守恒给出
\[
 \frac12mR^2\dot\theta^2=mgR(\cos\theta-\cos\theta_0).
\]
由最低点到右转折点所需时间为四分之一周期，故
\[
 T_{\rm circle}=4\sqrt{\frac{R}{2g}}
 \int_0^{\theta_0}\frac{\dd\theta}{\sqrt{\cos\theta-\cos\theta_0}}.
\]
令$k=\sin(\theta_0/2)$、$\sin(\theta/2)=k\sin u$。利用
$\cos\theta-\cos\theta_0=2k^2\cos^2u$及
$\dd\theta=2k\cos u\,\dd u/\sqrt{1-k^2\sin^2u}$，消去转折点的根式奇性：
\[
 T_{\rm circle}=4\sqrt{\frac Rg}
 \int_0^{\pi/2}\frac{\dd u}{\sqrt{1-k^2\sin^2u}}
 .
\]
这里使用$(1-z)^{-1/2}=1+z/2+O(z^2)$和
$\int_0^{\pi/2}\sin^2u\,\dd u=[u/2-\sin2u/4]_0^{\pi/2}=\pi/4$，因此
$T_{\rm circle}=T_0[1+k^2/4+O(k^4)]$。又
$k^2=\theta_0^2/4+O(\theta_0^4)$、
$a/R=\sin\theta_0=\theta_0+O(\theta_0^3)$，所以
\[
 \boxed{\frac{T_{\rm circle}}{T_0}
 =1+\frac{a^2}{16R^2}+O((a/R)^4)}.
\]
前面对抛物线已由弧长微元和能量积分得到
$T_{\rm parabola}/T_0=1+a^2/(4R^2)+O((a/R)^4)$。
相同水平振幅下，抛物线的首个周期修正为圆弧的四倍。图中最低点的二阶接触只确定线性化频率；离开最低点后的高度与弧长变化共同决定振幅修正。
\end{solution}

\section{面积元、体积元与雅可比}
适合物体形状的坐标也会改变积分元。微小边的几何关系给出长度、面积和体积因子，实际区域的覆盖次数则决定积分范围与分支。

\subsection{从微小边长看坐标变换}
平面极坐标的微小区域沿径向边长为 $\dd r$，沿切向边长为 $r\dd\theta$，因此 $\dd A=r\dd r\dd\theta$。柱坐标取 $x=\rho\cos\phi,y=\rho\sin\phi,z=z$，体积元为 $\rho\dd\rho\dd\phi\dd z$。

对一般参数化 $\vct r=\vct r(q_1,q_2,q_3)$，三个小边向量近似为 $\partial\vct r/\partial q_j\,\dd q_j$。微小平行六面体体积是混合积的绝对值，所以
\begin{equation}
\dd V=\left|\det\frac{\partial(x,y,z)}{\partial(q_1,q_2,q_3)}\right|
\dd q_1\dd q_2\dd q_3.
\end{equation}
行列式的绝对值称为雅可比因子。它给出坐标变换前后微元的面积或体积之比。

\subsection{球坐标与面积元}
定义极角 $\theta$ 从正 $z$ 轴量起，方位角 $\phi$ 在 $xy$ 平面中绕 $z$ 轴量起：
\[
x=r\sin\theta\cos\phi,\quad y=r\sin\theta\sin\phi,\quad z=r\cos\theta.
\]
\begin{figure}[htbp]\centering
\begin{tikzpicture}[scale=.95]
\begin{scope}
\draw[bookaxis](0,0)--(3,0)node[right]{$\rho$};\draw[bookaxis](0,0)--(0,3.2)node[above]{$z$};
\draw[bookaxis,thick](0,0)--(2,2.5)node[midway,above left]{$r$};\fill(2,2.5)circle(1.5pt)node[above right]{$P$};
\draw[bookguide](0,2.5)--(2,2.5)--(2,0);\draw(0,.9)arc(90:51.34:.9);\node at(69:1.25){$\theta$};
\node at(1.3,-.6){子午平面};
\end{scope}
\begin{scope}[shift={(5,0)}]
\draw[bookaxis](0,0)--(3.2,0)node[right]{$x$};\draw[bookaxis](0,0)--(0,2.7)node[above]{$y$};
\draw[bookaxis,thick](0,0)--(1.6,1.2)node[midway,above left]{$\rho$};\fill(1.6,1.2)circle(1.5pt)node[above right]{$P'$};
\draw[bookguide](0,1.2)--(1.6,1.2)--(1.6,0);\draw(.85,0)arc(0:36.87:.85);\node at(18:1.18){$\phi$};
\node at(1.3,-.6){$xy$ 平面};
\end{scope}
\end{tikzpicture}
\caption{球坐标分解为两次投影。$\rho=r\sin\theta$ 为柱坐标半径，$P'$ 为 $P$ 在 $xy$ 平面上的投影。}
\label{fig:orig-031}
\end{figure}
局部三个正交边长分别为 $\dd r,r\dd\theta,r\sin\theta\dd\phi$，故
\begin{equation}
\dd V=r^2\sin\theta\dd r\dd\theta\dd\phi.
\end{equation}
在半径$R$的球面上，$\dd S=R^2\sin\theta\dd\theta\dd\phi$。全空间角域为$0\le\theta\le\pi,0\le\phi<2\pi$，故
\[S=R^2[-\cos\theta]_0^\pi[\phi]_0^{2\pi}=4\pi R^2,\]
\[V=[r^3/3]_0^R[-\cos\theta]_0^\pi[\phi]_0^{2\pi}
=\frac{4\pi R^3}{3}.\]

\begin{example}
推导均匀实心球相对过球心轴的转动惯量。
\end{example}
% 题设配图；与紧随其后的解答分别编号。
\ExpansionDiagram{fig:audit-m-q-045}{绕中心轴求实心球惯量的切片}{
\draw[thick](0,0)circle(1.6);\draw[dashed](0,0)ellipse(1.6 and .4);\draw[->](0,-2)--(0,2.3)node[above]{$z$};\draw[->](0,0)--(1.3,.93)node[midway,below]{$R$};\draw[dashed](-1.2,1.05)--(1.2,1.05);\node[left]at(-1.4,1.1){薄层};
}

\begin{solution}
图\ref{fig:audit-m-q-045}中各水平薄层到轴的距离不同，转动惯量需对距离平方加权积分。
图\ref{fig:eg-2-6}中的水平投影 $\rho=r\sin\theta$ 就是质元到 $z$ 轴的垂距，它应先平方再乘体积元。

均匀球的体密度为 $\rho_m=M/(4\pi R^3/3)$。取球坐标 $r,\theta,\phi$，极角从转轴 $z$ 量起。质元到转轴的距离是 $r\sin\theta$，而不是到球心的距离 $r$；质元质量为 $\dd m=\rho_mr^2\sin\theta\dd r\dd\theta\dd\phi$。

\begin{figure}[H]\centering
\begin{tikzpicture}[scale=.95]
\begin{scope}
\draw[bookaxis](0,0)--(3,0)node[right]{$\rho$};\draw[bookaxis](0,0)--(0,3.2)node[above]{$z$};
\draw[bookaxis,thick](0,0)--(2,2.5)node[midway,above left]{$r$};\fill(2,2.5)circle(1.5pt)node[above right]{$P$};
\draw[bookguide](0,2.5)--(2,2.5)--(2,0);\draw(0,.9)arc(90:51.34:.9);\node at(69:1.25){$\theta$};
\node at(1.3,-.6){子午平面};
\end{scope}
\begin{scope}[shift={(5,0)}]
\draw[bookaxis](0,0)--(3.2,0)node[right]{$x$};\draw[bookaxis](0,0)--(0,2.7)node[above]{$y$};
\draw[bookaxis,thick](0,0)--(1.6,1.2)node[midway,above left]{$\rho$};\fill(1.6,1.2)circle(1.5pt)node[above right]{$P'$};
\draw[bookguide](0,1.2)--(1.6,1.2)--(1.6,0);\draw(.85,0)arc(0:36.87:.85);\node at(18:1.18){$\phi$};
\node at(1.3,-.6){$xy$ 平面};
\end{scope}
\end{tikzpicture}
\caption{例2.6解图：左图的水平投影 $\rho=r\sin\theta$ 是到转轴的垂距；右图再分解方位角，两图共同确定球坐标积分元。}
\label{fig:eg-2-6}
\end{figure}

按 $I_z=\int d_\perp^2\dd m$，三个积分可分开：
\[
I_z=\rho_m\int_0^Rr^4\dd r
\int_0^\pi\sin^3\theta\dd\theta\int_0^{2\pi}\dd\phi.
\]
径向积分为$[r^5/5]_0^R=R^5/5$，方位角积分为$[\phi]_0^{2\pi}=2\pi$。在极角积分中令$u=\cos\theta$，则$\dd u=-\sin\theta\dd\theta$，两端$0,\pi$对应$1,-1$，故
\[
\int_0^\pi\sin^3\theta\dd\theta
=-\int_1^{-1}(1-u^2)\dd u
=\left[u-\frac{u^3}{3}\right]_{-1}^{1}=\frac43.
\]
故 $I_z=8\pi\rho_mR^5/15=2MR^2/5$。结果为质量乘长度平方；其系数小于薄球壳的 $2/3$，因为实心球的一部分质量更靠近转轴。

沿球心射线改变 $r$ 会同时改变垂距与壳层体积，因而径向被积因子是 $r^4$，不能只用表面半径 $R$。
\end{solution}

\section{立体角、通量与对称性}
立体角定义为观察点所见面元在垂直视线方向的投影面积除以距离平方：$\dd\Omega=\dd S_\perp/r^2$。在以观察点为球心的球面上，$\dd\Omega=\sin\theta\dd\theta\dd\phi$，整球立体角为 $4\pi$。

半顶角为 $\theta_0$ 的圆锥对应
\begin{equation}
\Omega=\int_0^{2\pi}\dd\phi\int_0^{\theta_0}\sin\theta\dd\theta
=[\phi]_0^{2\pi}[-\cos\theta]_0^{\theta_0}
=2\pi(1-\cos\theta_0).
\end{equation}
小角时 $1-\cos\theta_0\approx\theta_0^2/2$，所以 $\Omega\approx\pi\theta_0^2$，与单位球上小圆面积相符。不能用二维角度 $2\theta_0$ 代替三维立体角。

\begin{figure}[htbp]
\centering
\begin{tikzpicture}[scale=.85,>=Stealth]
\draw(0,0)circle(1.7);\draw[->](0,-2)--(0,2.3)node[above]{$z$};
\draw(0,0)--(55:1.7) (0,0)--(125:1.7);\draw[very thick](55:1.7)arc(55:125:1.7);
\draw(55:.6)arc(55:90:.6);\node at(72:1){$\theta_0$};\node[below right]at(0,0){$O$};
\node[right]at(2.3,.7){球冠的轴截面};
\end{tikzpicture}
\caption{圆锥对应的球冠与立体角}\label{fig:solid-angle}
\end{figure}

点电荷位于观察点时，电场沿径向，$\dd\Phi_E=\vct E\cdot\vct n\dd S=q\dd\Omega/(4\pi\varepsilon_0)$，其中 $\dd\Omega$ 应按所选法向赋予符号。对球冠径向向外的法向，$\Phi_E=q\Omega/(4\pi\varepsilon_0)$。这解释了为什么点电荷总通量与包围面的形状无关：它测量的是有向立体角总和。

\section{连续分布的积分：圆环与圆盘}
\subsection{先处理方向，再积分大小}
半径 $R$ 的均匀带电细圆环，总电荷为 $Q$。轴线上观察点距圆心为 $z$。每个电荷元到观察点距离相同，对称元的横向分量相消，轴向分量相加。因此
\begin{equation}
E_z=\frac{Qz}{4\pi\varepsilon_0(R^2+z^2)^{3/2}}.
\end{equation}
圆心处 $E_z=0$，远处 $E_z\sim Q\operatorname{sgn}(z)/(4\pi\varepsilon_0z^2)$。圆环半径通过微元距离与方向投影进入结果。只有在远场，首项才与同电荷量的点电荷相同。

对 $Q>0,z>0$ 求极值，导数的符号由 $R^2-2z^2$ 决定，所以峰值在 $z=R/\sqrt2$，场强为 $Q[2/(3\sqrt3R^2)]/(4\pi\varepsilon_0)$。场在远处减小，在圆心也为零，因此最大值必须在中间出现；求导前可先用这两个极限预判曲线的大致形状。

\subsection{电场分量的几何投影}
圆环上相对的两电荷元产生等大电场，它们在环面内的分量反向抵消，轴向分量同向相加。图\ref{fig:orig-033}中到观察点的距离为$s=\sqrt{R^2+z^2}$，投影因子为$z/s$。逐元叠加得到
\[
\dd E_z=\frac{\dd q}{4\pi\varepsilon_0s^2}\frac zs,
\qquad E_z=\frac{z}{4\pi\varepsilon_0s^3}\int_{\rm ring}\dd q
=\frac{Qz}{4\pi\varepsilon_0(R^2+z^2)^{3/2}}.
\]
距离和投影对同一圆环均为常数，才能移到积分号外。分母的三次幂同时来自库仑场的距离平方与方向投影；轴场极值的逐项求导见图\ref{fig:ex-2-7}所配习题解答。
\begin{figure}[H]\centering
\begin{tikzpicture}[scale=.9]
\draw(0,0)ellipse(1.8 and .65);\draw[bookaxis](0,-.3)--(0,3.4)node[above]{$z$};
\coordinate(P)at(0,2.5);\coordinate(A)at(-1.8,0);\coordinate(B)at(1.8,0);
\fill(P)circle(1.5pt)node[right]{$P$};\fill(A)circle(1.5pt)node[left]{$\dd q$};\fill(B)circle(1.5pt)node[right]{$\dd q$};
\draw[bookguide](A)--(P) (B)--(P);\draw[bookaxis](P)--(.65,3.4);\draw[bookaxis](P)--(-.65,3.4);
\draw[bookaxis,thick](P)--(0,3.65);\draw[bookguide](0,0)--(B)node[midway,below]{$R$};
\end{tikzpicture}
\caption{对称电荷元的电场分解。图为几何示意，环画成椭圆以表示空间透视。}
\label{fig:orig-033}
\end{figure}

\subsection{将圆盘分成同心圆环}
面电荷密度 $\sigma$ 的圆盘半径为 $R$。半径 $\rho$、宽 $\dd\rho$ 的圆环电荷 $\dd q=2\pi\sigma\rho\dd\rho$。对 $z>0$，用圆环结果积分：
\begin{align}
E_z&=\frac{\sigma z}{2\varepsilon_0}\int_0^R\frac{\rho\dd\rho}{(\rho^2+z^2)^{3/2}}\\
&=\frac{\sigma z}{4\varepsilon_0}\int_{z^2}^{z^2+R^2}u^{-3/2}\dd u
=\frac{\sigma z}{4\varepsilon_0}[-2u^{-1/2}]_{z^2}^{z^2+R^2}\\
&=\frac{\sigma}{2\varepsilon_0}\left(1-\frac z{\sqrt{R^2+z^2}}\right).
\end{align}
这里换元$u=\rho^2+z^2$、$\dd u=2\rho\dd\rho$，并用$z>0$处理下限平方根。下方电场方向相反。$z\to0^+$时回到盘中心近旁的无限平面极限；$z\gg R$时展开$z/\sqrt{R^2+z^2}$，得到$E_z\approx\sigma R^2/(4\varepsilon_0z^2)$，与总电荷$Q=\pi R^2\sigma$的点电荷远场一致。

电势叠加不需要方向投影。取无穷远为零，对 $z>0$，
\begin{equation}
\Phi(z)=\frac{\sigma}{2\varepsilon_0}\int_0^R\frac{\rho\dd\rho}{\sqrt{\rho^2+z^2}}
=\frac{\sigma}{2\varepsilon_0}[\sqrt{\rho^2+z^2}]_0^R
=\frac{\sigma}{2\varepsilon_0}(\sqrt{R^2+z^2}-z).
\end{equation}
再用 $E_z=-\dd\Phi/\dd z$ 得到相同场强。分别从电场和电势出发，是对积分方向与符号的一次独立检验。

\section{曲线积分与曲面积分的对象}
沿路径的功与穿过曲面的通量使用不同的方向信息。先确定累加的物理量，再选参数、切向或法向，最后把积分范围对应到实际路径和曲面。

\subsection{标量曲线积分和做功积分}
沿细线分布的质量可写成 $M=\int_C\lambda(s)\dd s$，其中线密度为标量。这种积分不依赖走向。力做功则是 $W=\int_C\vct F\cdot\dd\vct r$，反向走同一路径会改变符号。两者的微元和方向性质不同。

对参数化 $\vct r(t)$，标量积分为 $\int\lambda(\vct r(t))|\vct r'(t)|\dd t$，做功积分为 $\int\vct F(\vct r(t))\cdot\vct r'(t)\dd t$。后者没有额外再乘一个速度模，否则会重复计入路径微元。

\begin{example}
平面力场 $\vct F=C(-y,x)$，$C>0$。从 $(0,0)$ 到 $(a,a)$ 分别沿先横后纵和先纵后横路径运动，求做功。
\end{example}
% 题设配图；与紧随其后的解答分别编号。
\ExpansionDiagram{fig:audit-m-q-046}{连接同一对端点的两条折线路径}{
\draw[->](-.4,0)--(3.5,0)node[right]{$ x $};\draw[->](0,-0.3)--(0,3)node[above]{$ y $};\draw[very thick,->](0,0)--(2.5,0);\draw[very thick,->](2.5,0)--(2.5,2.5);\draw[dashed,->](0,0)--(0,2.5);\draw[dashed,->](0,2.5)--(2.5,2.5);\fill(2.5,2.5)circle(2pt)node[above right]{$(a,a)$};
}

\begin{solution}
图\ref{fig:audit-m-q-046}中的两条路径方向相同而经过区域不同，应逐段代入力场。
图\ref{fig:orig-034}给两条折线路径标了走向；每段积分都按这条走向确定微位移的符号。

做功积分的微元是
\[
\dd W=F_x\dd x+F_y\dd y=-Cy\dd x+Cx\dd y.
\]
先横后纵时，第一段 $y=0,\dd y=0$，积分为0；第二段 $x=a,\dd x=0$，所以
\[
W_A=0+\int_0^a Ca\dd y=[Cay]_0^a=Ca^2.
\]
先纵后横时，第一段 $x=0,\dd x=0$，积分为0；第二段 $y=a,\dd y=0$，所以
\[
W_B=0+\int_0^a(-Ca)\dd x=[-Cax]_0^a=-Ca^2.
\]
二者端点完全相同，功却不同，故该场不保守。若去程取路径A、回程反向走路径B，闭路功为 $W_A-W_B=2Ca^2$，正好等于标量旋度 $2C$ 乘所围面积 $a^2$。沿对角线 $x=y=s$ 时，微元为 $-Cs\dd s+Cs\dd s=0$，又得到第三种路径的零功结果。

\begin{figure}[H]\centering
\begin{tikzpicture}[scale=2,>=Stealth]
\draw[bookaxis](-.1,0)--(1.35,0)node[right]{$x$};\draw[bookaxis](0,-.1)--(0,1.35)node[above]{$y$};
\draw[bookaxis,thick](0,0)--(1,0);\draw[bookaxis,thick](1,0)--(1,1);
\draw[bookaxis,dashed](0,0)--(0,1);\draw[bookaxis,dashed](0,1)--(1,1);
\draw[dotted](0,0)--(1,1);\fill(1,1)circle(.8pt)node[right]{$(a,a)$};
\node[below]at(.55,-.12){路径 A};\node[above]at(.5,1.06){路径 B};
\end{tikzpicture}
\caption{端点相同的两条路径。实线路径做功 $Ca^2$，虚线路径做功 $-Ca^2$。}
\label{fig:orig-034}
\end{figure}

实线路径只有竖直段贡献正功，虚线路径只有水平段贡献负功。沿点线对角线，两项微元相消。
\end{solution}

\subsection{有向面积元从叉积得到}
曲面写成 $\vct r(u,v)$，两参数方向的微小边向量分别为 $\vct r_u\dd u,\vct r_v\dd v$。因此标量面积元为 $\dd S=|\vct r_u\times\vct r_v|\dd u\dd v$，有向面积元为
\[
\dd\vct S=(\vct r_u\times\vct r_v)\dd u\dd v.
\]
交换参数次序会反转法向。标量面积不变，通量的符号改变。对图形曲面 $z=f(x,y)$，取向上法向，有
\[
\dd\vct S=(-f_x,-f_y,1)\dd x\dd y,
\quad \dd S=\sqrt{1+f_x^2+f_y^2}\dd x\dd y.
\]
其中根号是斜面的真实面积与水平投影面积之比。

\begin{example}
场 $\vct A=(0,0,A_0)$ 穿过斜面 $z=px+qy$，投影区域为面积 $S_0$ 的矩形，法向向上。求通量。
\end{example}
% 题设配图；与紧随其后的解答分别编号。
\ExpansionDiagram{fig:audit-m-q-047}{斜面及其在水平面上的投影区域}{
\draw[thick,fill=gray!10](0,0)--(3,.5)--(4.5,2.2)--(1.5,1.7)--cycle;\draw[dashed](0,-1.1)--(3,-.6)--(4.5,.1)--(1.5,-.4)--cycle;\draw[dotted](0,0)--(0,-1.1);\draw[dotted](4.5,2.2)--(4.5,.1);\draw[->,thick](2.2,1)--(1.8,2.5)node[above]{$\vct n$};\node at(2.4,-.65){$S_0$};\node[right]at(4.5,2.2){$z=px+qy$};
}

\begin{solution}
图\ref{fig:audit-m-q-047}中场沿$z$方向，曲面的有向面积元须保留其水平投影分量。
图\ref{fig:orig-035}区分真实斜面与水平投影。通量中的竖直场与向上面积向量点乘，最后只留下投影面积。

将曲面参数化为 $\vct r(x,y)=(x,y,px+qy)$，其中 $(x,y)$ 遍历给定投影矩形。两个切向量为
\[
\vct r_x=(1,0,p),\qquad\vct r_y=(0,1,q),
\]
向上定向的面积向量为
\[
\dd\vct S=(\vct r_x\times\vct r_y)\dd x\dd y
=(-p,-q,1)\dd x\dd y.
\]
故 $\vct A\cdot\dd\vct S=A_0\dd x\dd y$，积分即为
\[
\Phi=A_0\iint_{D}\dd x\dd y=A_0S_0.
\]
从单位法向看，$\vct n=(-p,-q,1)/\sqrt{1+p^2+q^2}$，面积为 $S_0\sqrt{1+p^2+q^2}$，点乘带来的余弦因子恰好抵消面积的放大。若法向改向下，通量整体变号。

\begin{figure}[H]\centering
\begin{tikzpicture}[scale=.9]
\draw(0,0)--(3,0)--(4,1.1)--(1,1.1)--cycle;\draw[bookguide](0,1.8)--(0,0) (3,1.8)--(3,0) (4,2.9)--(4,1.1);
\draw[thick](0,1.8)--(3,1.8)--(4,2.9)--(1,2.9)--cycle;
\draw[bookaxis](2,2.35)--(2,3.7)node[above]{$\vct A$};\node at(2,.55){投影区域};\node[left]at(0,1.8){曲面};
\draw[bookaxis](1.4,2.35)--(.8,3.1)node[above left]{$\vct n$};\node at(1.7,.55){$S_0$};
\end{tikzpicture}
\caption{曲面通量与投影面积的关系。竖直虚线表示投影，不表示场线。}
\label{fig:orig-035}
\end{figure}

图上的投影关系为 $S\,n_z=S_0$，正好解释计算中根号因子为何抵消。反向定向只会改变通量的符号。
\end{solution}

\section{梯度、方向导数与等值面}
势能的空间变化决定力的方向和大小。先把微小位移引起的势变化写成线性形式，再讨论沿指定方向的变化率及等值面的几何。

标量函数 $f(x,y,z)$ 的全微分为 $\dd f=f_x\dd x+f_y\dd y+f_z\dd z=\nabla f\cdot\dd\vct r$，其中 $\nabla f=(f_x,f_y,f_z)$。沿单位方向 $\vct n$ 的变化率为 $D_{\vct n}f=\nabla f\cdot\vct n$，最大值为 $|\nabla f|$，发生在梯度方向。

沿等值面移动时 $\dd f=0$，所以梯度与所有切向位移垂直，是等值面的法线。若梯度为零，不能据此唯一确定法线，那里可能是极值点、鞍点或其他临界点。电场 $\vct E=-\nabla\Phi$ 指向电势下降最快的方向，负号具有明确的物理意义。
\begin{figure}[htbp]\centering
\begin{tikzpicture}[scale=.9]
\foreach\r in{.6,1,1.414}{\draw(0,0)circle(\r);}\fill(1,1)circle(1.5pt);
\draw[bookaxis,thick](1,1)--(1.7,1.7)node[above]{$\nabla U$};\draw[bookaxis,thick](1,1)--(.45,.45)node[left]{$\vct F$};
\draw[bookaxis](0,0)--(1.8,0)node[right]{$x$};\draw[bookaxis](0,0)--(0,1.8)node[above]{$y$};
\end{tikzpicture}
\caption{$U=k(x^2+y^2)/2$ 的等势线、梯度与恢复力。梯度向外，力 $-\nabla U$ 向内。}
\label{fig:orig-036}
\end{figure}

\begin{example}
温度场 $T=T_0+ax^2+by^2$，求点 $(x_0,y_0)$ 处沿方向 $(1,1)$ 的温度变化率。
\end{example}
% 题设配图；与紧随其后的解答分别编号。
\ExpansionDiagram{fig:audit-m-q-048}{温度场中的给定点及单位化前的方向}{
\draw[->](-.4,0)--(4.5,0)node[right]{$ x $};\draw[->](0,-0.3)--(0,3)node[above]{$ y $};\draw[thick](0,0)ellipse(2.2 and 1.4);\draw[dashed](0,0)ellipse(1.4 and .9);\fill(1,.8)circle(2pt)node[below right]{$(x_0,y_0)$};\draw[->,thick](1,.8)--(2.2,2)node[above]{$(1,1)$};
}

\begin{solution}
图\ref{fig:audit-m-q-048}中箭头只给出方向，计算方向导数前要先化为单位向量。
图\ref{fig:eg-2-9}同时画出梯度和题给的单位方向。沿方向走单位路程时，只取梯度在该方向上的投影。

题目求的是单位路程上的温度变化率，先把方向向量归一化为 $\vct n=(1,1)/\sqrt2$。从 $P=(x_0,y_0)$ 沿该方向走距离 $s$ 后，位置是
\[
x=x_0+s/\sqrt2,\qquad y=y_0+s/\sqrt2.
\]
代入温度函数，链式求导为
\[\frac{\dd T}{\dd s}
=2a\left(x_0+\frac s{\sqrt2}\right)\frac1{\sqrt2}
+2b\left(y_0+\frac s{\sqrt2}\right)\frac1{\sqrt2}.\]
在$s=0$得
\[
\left.\frac{\dd T}{\dd s}\right|_0
=\frac{2ax_0+2by_0}{\sqrt2}
=\sqrt2(ax_0+by_0).
\]
用梯度法也得到 $\nabla T(P)=(2ax_0,2by_0)$，点乘单位 $\vct n$ 得到相同结果。常量 $T_0$ 不影响变化率。如果把方向反转，方向导数变号；若该式为零，只能说明在此方向上首阶不变，不能推出整个温度场为常量。

\begin{figure}[H]\centering
\begin{tikzpicture}[scale=1.1]
\draw[bookaxis](-.3,0)--(4,0)node[right]{$x$};\draw[bookaxis](0,-.3)--(0,3.3)node[above]{$y$};
\draw[bookguide](0,0)circle(1.118);\draw[bookguide](0,0)circle(1.8);
\coordinate(P)at(1,.5);\fill(P)circle(1.7pt)node[below right]{$P$};
\draw[bookaxis,thick](P)--(2.6,1.3)node[right]{$\nabla T$};
\draw[bookaxis](P)--(2.1,1.6)node[above]{$\vct n=(1,1)/\sqrt2$};
\node[right,align=left]at(2.3,2.65){$D_{\vct n}T=\nabla T\cdot\vct n$};
\end{tikzpicture}
\caption{例2.9：方向导数取梯度在单位方向上的投影。图示 $a=b>0$ 的等温线，计算式仍适用于一般的 $a,b$。}\label{fig:eg-2-9}
\end{figure}

题给方向一般偏离梯度方向，方向导数的绝对值因而不超过梯度的模。方向向量须先单位化，才能得到每单位路程的变化率。
\end{solution}

\section{散度、旋度与局部测量}
\subsection{散度测量净流出}
向量场 $\vct A=(A_x,A_y,A_z)$ 的散度定义为
\[
\nabla\cdot\vct A=\frac{\partial A_x}{\partial x}+\frac{\partial A_y}{\partial y}+\frac{\partial A_z}{\partial z}.
\]
取以$(x,y,z)$为中心、边长为$\Delta x,\Delta y,\Delta z$的微小长方体。将两张$x$法向面的场值在面中心展开，首阶净通量为
\[
\begin{aligned}
\Delta\Phi_x
&=\left[A_x\left(x+\frac{\Delta x}{2},y,z\right)
-A_x\left(x-\frac{\Delta x}{2},y,z\right)\right]\Delta y\Delta z\\
&=(\partial_xA_x)\Delta x\Delta y\Delta z+o(\Delta V).
\end{aligned}
\]
这里取各边按固定比例趋零；面内变化也只影响更高阶项。另两对面分别给出$(\partial_yA_y)\Delta V$与$(\partial_zA_z)\Delta V$。三者相加并除以体积，极限就是散度。正散度表示局部源，负散度表示汇；零散度不意味着场为零，也不意味着场没有弯曲或旋转。

\subsection{旋度测量局部环流}
在$xy$平面中取以$(x,y)$为中心的小矩形，逆时针绕行。两条水平边走向相反，按边中点展开，其首阶环流之和为
\[
\left[A_x\left(x,y-\frac{\Delta y}{2}\right)
-A_x\left(x,y+\frac{\Delta y}{2}\right)\right]\Delta x
=-(\partial_yA_x)\Delta x\Delta y+o(\Delta S).
\]
两条竖直边则给出$(\partial_xA_y)\Delta x\Delta y+o(\Delta S)$。相加后除以面积并取极限，得到$\partial_xA_y-\partial_yA_x$，即旋度的$z$分量。对另外两个坐标平面作同样的定向，完整分量为
\[
\nabla\times\vct A=(\partial_yA_z-\partial_zA_y,\ \partial_zA_x-\partial_xA_z,\ \partial_xA_y-\partial_yA_x).
\]
旋度方向按右手规则对应局部环流方向。旋度由单位面积上的局部环流决定。仅看一条场线是否弯曲，无法判断旋度。
\begin{figure}[htbp]\centering
\begin{tikzpicture}[scale=.8]
\foreach\s/\kind/\title in{0/1/源场,3.7/2/旋转场,7.4/3/鞍形场}{
\begin{scope}[shift={(\s,0)}]
\foreach\x in{-1,0,1}{\foreach\y in{-1,0,1}{\pgfmathtruncatemacro{\nz}{\x*\x+\y*\y}\ifnum\nz>0
\ifnum\kind=1\draw[bookaxis](\x,\y)--({1.27*\x},{1.27*\y});\fi
\ifnum\kind=2\draw[bookaxis](\x,\y)--({\x-.27*\y},{\y+.27*\x});\fi
\ifnum\kind=3\draw[bookaxis](\x,\y)--({1.27*\x},{.73*\y});\fi\fi}}
\node at(0,-1.8){\title};\end{scope}}
\end{tikzpicture}
\caption{三种向量场：$(x,y)$ 的散度为2；$(-y,x)$ 的旋度 $z$ 分量为2；$(x,-y)$ 的散度和旋度都为0。}
\label{fig:orig-037}
\end{figure}

\subsection{保守场与区域条件}
若 $\vct F=-\nabla U$，混合偏导可交换时必有 $\nabla\times\vct F=0$。反过来，零旋度要推出全局势函数，还需要区域条件，例如在单连通区域中。去掉奇点的平面可能出现零旋度但绕奇点环流非零的场。

例如$\vct A=(-y/r^2,x/r^2)$，$r^2=x^2+y^2$，在$r>0$处的两个偏导为
\[\partial_xA_y=\frac1{r^2}-\frac{2x^2}{r^4}
=\frac{y^2-x^2}{r^4},\qquad
\partial_yA_x=-\frac1{r^2}+\frac{2y^2}{r^4}
=\frac{y^2-x^2}{r^4}.\]
两者相减，旋度为零。但沿图~\ref{fig:math-vortex-puncture}的逆时针单位圆，取$x=\cos t,y=\sin t$，有
\[\vct A=(-\sin t,\cos t),\qquad
\dd\vct r=(-\sin t,\cos t)\dd t,\]
因此
\[\oint_C\vct A\cdot\dd\vct r
=\int_0^{2\pi}(\sin^2t+\cos^2t)\dd t=[t]_0^{2\pi}=2\pi.\]
该场局部是辐角的梯度，辐角却不能在整个去心平面单值连续定义。若用圆盘套Stokes公式，圆盘就穿过了被剔除的原点，违反光滑性条件；这与第一章$\oint\dd z/z$的非零值对应。
\begin{figure}[htbp]\centering
\begin{tikzpicture}[>=Stealth,x=1.3cm,y=1.3cm]
\draw[->](-1.55,0)--(1.65,0)node[right]{$x$};
\draw[->](0,-1.5)--(0,1.55)node[above]{$y$};
\draw[thick](0,0)circle(1);
\draw[->,thick](1,0)arc(0:55:1);
\node[above left]at(-.707,.707){$C$};
\draw[fill=white](0,0)circle(2.5pt)node[below left]{$O$};
\fill(1,0)circle(1.5pt)node[below right]{$P=(1,0)$};
\draw[->,thick](1,0)--(1,.65)node[right]{$\vct A(P)$};
\end{tikzpicture}
\caption{去心平面中的单位圆与切向场。原点不属于场的定义域；边界环流不能由穿过该点的圆盘作无源推断。}
\label{fig:math-vortex-puncture}
\end{figure}

\section{格林、散度与斯托克斯定理}
边界积分可以换成区域内部的积分，但方向约定与光滑性条件须同时保留。以下从通量和环量出发，分别说明三种定理处理的对象。

\subsection{格林定理的边界抵消机制}
设平面区域 $D$ 的边界 $C$ 逆时针，场的偏导在区域邻域连续，则
\[
\oint_C(P\dd x+Q\dd y)=\iint_D(\partial_xQ-\partial_yP)\dd x\dd y.
\]
先对矩形应用微积分基本定理，再把一般区域分成许多小块。内部公共边在两个相邻小块中走向相反，线积分抵消，只留下外边界。这个抵消机制解释了局部旋度怎样积累成整体环流。
\begin{figure}[htbp]\centering
\begin{tikzpicture}[scale=1.2]
\draw(0,0)rectangle(3,2);\foreach\x in{1,2}{\draw[dashed](\x,0)--(\x,2);}\draw[dashed](0,1)--(3,1);
\draw[bookaxis,thick](.4,0)--(1,0);\draw[bookaxis,thick](3,.35)--(3,.95);\draw[bookaxis,thick](2.6,2)--(2,2);\draw[bookaxis,thick](0,1.65)--(0,1.05);
\draw[bookaxis](.94,.15)--(.94,.7);\draw[bookaxis](1.06,.85)--(1.06,.3);\node[right]at(3.2,1){$C$};
\end{tikzpicture}
\caption{小区域公共边的相反走向抵消，外边界的正向环流保留下来。图中内部两箭头表示两侧小块的贡献。}
\label{fig:orig-038}
\end{figure}

对 $P=-Cy,Q=Cx$，右侧为 $2C$ 乘面积。边长为 $a$ 的方形环流为 $2Ca^2$，正好等于前面两条路径做功之差。格林定理既可帮助计算，也可用来验证路径积分。

\subsection{散度定理与电场通量}
若体积区域 $V$ 的边界闭合，法向向外，且场在区域邻域具有连续一阶偏导，则
\[
\iint_{\partial V}\vct A\cdot\dd\vct S=\iiint_V\nabla\cdot\vct A\dd V.
\]
同样地，内部小块公共面的通量相消，局部净流出累积到外边界。对于点电荷的场，原点是奇点，不能把整个含原点区域当作光滑无源区域应用公式；应挖去小球，或使用包含源的适当广义形式。
\begin{figure}[htbp]\centering
\begin{tikzpicture}[scale=1]
\draw(0,0)rectangle(2.5,2);\draw(0,2)--(.8,2.7)--(3.3,2.7)--(3.3,.7)--(2.5,0);\draw(2.5,2)--(3.3,2.7);\draw[dashed](0,0)--(.8,.7)--(3.3,.7) (.8,.7)--(.8,2.7);
\draw[bookaxis](2.5,1)--(3.4,1)node[right]{$+\vct n_x$};\draw[bookaxis](0,1)--(-.9,1)node[left]{$-\vct n_x$};
\draw[bookaxis](1.7,2.4)--(1.7,3.4)node[above]{$+\vct n_z$};
\end{tikzpicture}
\caption{小体元各面的外法向。相反两面的通量差产生对应方向的偏导。}
\label{fig:orig-039}
\end{figure}

\begin{example}
场 $\vct A=(ax,by,cz)$ 穿过任意闭合曲面的外向通量是多少？曲面包围体积为 $V$。
\end{example}
% 题设配图；与紧随其后的解答分别编号。
\ExpansionDiagram{fig:audit-m-q-049}{闭合曲面的外法向与内部体积}{
\draw[thick](0,0)ellipse(2 and 1.2);\node at(0,0){$V$};\draw[->](2,0)--(3,0)node[right]{$\vct n$};\draw[->](0,1.2)--(0,2.1);\draw[->](-2,0)--(-2.8,0);\node[below]at(0,-1.6){$\vct A=(ax,by,cz)$};
}

\begin{solution}
图\ref{fig:audit-m-q-049}中的通量取外法向，体积分使用内部体积$V$。
图\ref{fig:eg-2-10}的小体元展示相对两面的通量相减。先把三对面的贡献相加，再用散度定理推广到任意闭合边界。

对光滑场和闭合边界，优先考虑散度定理。逐分量求偏导，
\[
\nabla\cdot\vct A=\partial_x(ax)+\partial_y(by)+\partial_z(cz)=a+b+c.
\]
因此外向通量为
\[
\Phi=\iiint_V(\nabla\cdot\vct A)\dd V=(a+b+c)V.
\]
可用边长 $L_x,L_y,L_z$ 的长方体核对：相对两张 $x$ 法向面的净通量为 $aL_xL_yL_z=aV$，其余两对面分别为 $bV,cV$。相加即恢复上式。

\begin{figure}[H]\centering
\begin{tikzpicture}[scale=1]
\draw(0,0)rectangle(2.5,2);\draw(0,2)--(.8,2.7)--(3.3,2.7)--(3.3,.7)--(2.5,0);\draw(2.5,2)--(3.3,2.7);\draw[dashed](0,0)--(.8,.7)--(3.3,.7) (.8,.7)--(.8,2.7);
\draw[bookaxis](2.5,1)--(3.4,1)node[right]{$+\vct n_x$};\draw[bookaxis](0,1)--(-.9,1)node[left]{$-\vct n_x$};
\draw[bookaxis](1.7,2.4)--(1.7,3.4)node[above]{$+\vct n_z$};
\end{tikzpicture}
\caption{例2.10解图：相对的两个 $x$ 法向面的通量差为 $a\dd V$，另两对面分别贡献 $b\dd V,c\dd V$。}
\label{fig:eg-2-10}
\end{figure}

公式与曲面位置、形状无关，只依赖它包围的体积。若 $a+b+c=0$，净通量为零，但各面仍可有流入和流出，场也不必为零。

沿三个坐标方向的净通量分别为 $aV,bV,cV$，对应图上的三对外法向；内部公共面在拼接体元时相消。
\end{solution}

\subsection{斯托克斯定理与曲面选择}
对有向曲面 $S$ 及其边界 $C$，方向按右手规则一致，
\[
\oint_C\vct A\cdot\dd\vct r=\iint_S(\nabla\times\vct A)\cdot\dd\vct S.
\]
给定同一边界，可选择容易计算的跨接曲面，只要场在需要的区域满足光滑条件。边界走向相同，曲面法向就不能任意反转。平面情形直接恢复格林定理。
\begin{figure}[htbp]\centering
\begin{tikzpicture}[scale=1,>=Stealth]
\draw(0,0)ellipse(2 and .75);\draw[dashed](-2,0)--(2,0);\draw[bookaxis](0,0)--(0,2.1)node[above]{$\vct n$};
\draw[bookaxis,thick](1.732,.375)arc(30:75:2 and .75);\node[right]at(2.15,0){$C$};\node at(-.8,.2){$S$};
\end{tikzpicture}
\caption{曲面法向与边界走向的右手关系。从法向一侧看，正向边界为逆时针。}
\label{fig:orig-040}
\end{figure}

\section{曲线坐标中的微分算子（选读）}
\subsection{柱坐标中的梯度和散度}
柱坐标三个局部边长为 $\dd r,r\dd\phi,\dd z$。沿单位切向每走一段长度，方位角改变的量为距离除以 $r$，因此
\[
\nabla f=\vct e_r f_r+\vct e_\phi\frac{f_\phi}{r}+\vct e_z f_z.
\]
对微小楔形体元，径向两面的面积分别与$r$和$r+\dd r$成正比。保留首阶通量差，
\[
\begin{aligned}
\dd\Phi_r
&=[(r+\dd r)A_r(r+\dd r,\phi,z)-rA_r(r,\phi,z)]\dd\phi\dd z\\
&=\partial_r(rA_r)\dd r\dd\phi\dd z.
\end{aligned}
\]
方位角两面的面积为$\dd r\dd z$，其净通量为
\[
\dd\Phi_\phi=\partial_\phi A_\phi\dd\phi\dd r\dd z.
\]
上下两面给出
\[
\dd\Phi_z=\partial_zA_z\,r\dd r\dd\phi\dd z.
\]
相加后除以体积$r\dd r\dd\phi\dd z$，得到
\[
\nabla\cdot\vct A=\frac1r\partial_r(rA_r)+\frac1r\partial_\phi A_\phi+\partial_zA_z.
\]
旋度的 $z$ 分量为 $(\nabla\times\vct A)_z=r^{-1}[\partial_r(rA_\phi)-\partial_\phi A_r]$。曲线坐标中的这些尺度因子来自微元几何，不能把直角坐标公式只换变量名。
\begin{figure}[htbp]\centering
\begin{tikzpicture}[scale=1]
\draw[bookaxis](0,0)--(3.7,0)node[right]{$x$};\draw[bookaxis](0,0)--(0,2.7)node[above]{$y$};
\draw[thick](25:2)arc(25:45:2)--(45:2.6)arc(45:25:2.6)--cycle;
\draw[bookguide](0,0)--(25:2.6) (0,0)--(45:2.6);\draw(.6,0)arc(0:25:.6);\node at(12:1){$\phi$};
\node[right]at(25:2.3){$\dd r$};\node[above]at(35:2.7){$r\dd\phi$};
\end{tikzpicture}
\caption{极坐标微元的两条实际边长。增加相同角度时，外侧边更长。}
\label{fig:orig-041}
\end{figure}

\subsection{球坐标的尺度与源项}
由三个局部边长，球坐标梯度为
\[
\nabla f=\vct e_r f_r+\vct e_\theta\frac{f_\theta}{r}+\vct e_\phi\frac{f_\phi}{r\sin\theta}.
\]
散度由小体元六个面的净通量求出。取边长为$\dd r,r\dd\theta,r\sin\theta\dd\phi$的球坐标体元，其体积为
\[
\dd V=r^2\sin\theta\,\dd r\dd\theta\dd\phi.
\]
径向两面的面积分别含因子$r^2$与$(r+\dd r)^2$，所以净径向通量的一阶项是
\[
\dd\Phi_r=\partial_r(r^2A_r)\sin\theta\,\dd r\dd\theta\dd\phi.
\]
在极角为$\theta$的一面，面积为$r\sin\theta\,\dd r\dd\phi$；对面的面积还须把$\theta$换成$\theta+\dd\theta$。因此
\[
\dd\Phi_\theta=r\partial_\theta(\sin\theta A_\theta)\,\dd r\dd\theta\dd\phi.
\]
两张方位角面的面积同为$r\dd r\dd\theta$，其通量差为
\[
\dd\Phi_\phi=r\partial_\phi A_\phi\,\dd r\dd\theta\dd\phi.
\]
三项相加后除以$\dd V$，并令体元尺寸趋零，得到
\[
\nabla\cdot\vct A=\frac1{r^2}\partial_r(r^2A_r)
+\frac1{r\sin\theta}\partial_\theta(\sin\theta A_\theta)
+\frac1{r\sin\theta}\partial_\phi A_\phi.
\]
对于纯径向场 $A_r=C/r^2$，在 $r>0$ 处散度为零，但任意以原点为球心的球面通量都是 $4\pi C$。原点存在奇点，含原点的区域不满足上述无源光滑条件。坐标轴上的 $\sin\theta=0$ 也可能只是坐标退化，需改用直角坐标检查场是否真的奇异。

\subsection{拉普拉斯算子来自梯度的散度}
对标量场，定义 $\nabla^2f=\nabla\cdot(\nabla f)$。把刚得到的梯度分量代入散度公式，柱坐标中为
\begin{equation}
\nabla^2f=\frac1r\partial_r(rf_r)+\frac1{r^2}f_{\phi\phi}+f_{zz},
\end{equation}
球坐标中则为
\begin{equation}
\nabla^2f=\frac1{r^2}\partial_r(r^2f_r)
+\frac1{r^2\sin\theta}\partial_\theta(\sin\theta f_\theta)
+\frac1{r^2\sin^2\theta}f_{\phi\phi}.
\end{equation}
这些公式使用沿单位基矢的物理分量。若把 $f_\phi$ 直接当成单位切向的变化率，就会遗漏梯度中的尺度因子，并在第二次求导时继续放大错误。

\begin{example}
分别求三维球对称与二维圆对称的调和势，讨论无源区域内的径向形式。
\end{example}
% 题设配图；与紧随其后的解答分别编号。
\ExpansionDiagram{fig:audit-m-q-050}{球对称与圆柱对称的径向几何不同}{
\begin{scope}[xshift=-2cm]\draw(0,0)circle(1.2);\draw[dashed](0,0)ellipse(1.2 and .3);\draw[->](0,0)--(1.2,0)node[midway,above]{$r$};\node[below]at(0,-1.6){三维球面对称};\end{scope}\begin{scope}[xshift=2cm]\draw(0,0)circle(1.2);\draw[->](0,0)--(1.2,0)node[midway,above]{$r$};\node[below]at(0,-1.6){二维圆对称};\end{scope}
}

\begin{solution}
图\ref{fig:audit-m-q-050}中两种对称性对应不同的面积随$r$变化规律，径向拉普拉斯算子也不同。
图\ref{fig:eg-2-11}对比球面与单位长度圆柱侧面的面积增长，先据此写径向通量守恒。

球对称函数只依赖半径。在 $r>0$ 的无源区域，三维拉普拉斯方程化为
\[
\frac1{r^2}\frac{\dd}{\dd r}(r^2f')=0.
\]
先乘$r^2$，对零导数积分得到$r^2f'=C_1$；再除以$r^2$积分，
\[
f(r)=C_0+C_1\int r^{-2}\dd r=C_0-\frac{C_1}{r}.
\]
二维圆对称时的方程为$r^{-1}(rf')'=0$。先乘$r$积分得$rf'=C_1$，再除以$r$积分，
\[
f(r)=C_0+C_1\int_{r_0}^{r}\frac{\dd s}{s}
=C_0+C_1\left[\ln(s/r_0)\right]_{r_0}^r
=C_0+C_1\ln(r/r_0),\qquad r_0>0.
\]
可逐次求导代回检查。三维 $r^2f'$ 与二维 $rf'$ 都为常量，分别表示球面和单位长度圆柱侧面的通量不随半径改变。

\begin{figure}[H]\centering
\begin{tikzpicture}[scale=.9]
\begin{scope}
\draw(0,0)circle(1.1);\draw[bookguide](0,0)circle(1.45);\draw[bookaxis](0,0)--(35:1.45)node[midway,above]{$r$};
\draw[bookaxis,thick](1.1,0)--(1.9,0)node[right]{$f'\vct e_r$};\node[align=center]at(0,-2){球面面积 $4\pi r^2$\\$r^2f'=C_1$};
\end{scope}
\begin{scope}[shift={(6,0)}]
\draw(0,0)circle(1.1);\draw[bookguide](0,0)circle(1.45);\draw[bookaxis](0,0)--(35:1.45)node[midway,above]{$r$};
\draw[bookaxis,thick](1.1,0)--(1.9,0)node[right]{$f'\vct e_r$};\node[align=center]at(0,-2){单位长度圆柱侧面 $2\pi r$\\$rf'=C_1$};
\end{scope}
\end{tikzpicture}
\caption{例2.11：无源壳层中内、外通量相同。两种几何的面积增长率不同，分别导致 $1/r$ 势与对数势。箭头示意所讨论的径向导数为正的情形。}\label{fig:eg-2-11}
\end{figure}

这些形式只适用于不含原点的无源区域。若要求势在原点附近有界，对应的 $1/r$ 或 $\ln r$ 项必须消失；若另给内外边界电势，则由边界条件确定积分常数。二维对数势一般不能再把无穷远选成有限零点。

相同的壳层通量分别要求 $r^2f'$、$rf'$ 为常量，积分后便得到不同的径向势。
\end{solution}

\section{复势与二维无源场（选读）}
若 $W(z)=\Phi(x,y)+\ii\Psi(x,y)$ 解析，柯西--黎曼条件给出 $\nabla\Phi\cdot\nabla\Psi=0$。因此 $\Phi$ 等值线与 $\Psi$ 等值线在梯度非零处正交。电场 $-\nabla\Phi$ 沿 $\Psi$ 等值线切向，$\Psi$ 可用于描画二维场线。

由 $W'=\Phi_x+\ii\Psi_x=\Phi_x-\ii\Phi_y$，得到 $E_x-\ii E_y=-W'(z)$。这个关系中的虚部符号依赖所选择的复场写法，若改写成 $E_x+\ii E_y$，就要相应取共轭。

\begin{example}
无限长线电荷在垂直截面中的电势为 $\Phi=-C\ln(r/r_0)$，$C=\lambda/(2\pi\varepsilon_0)$。求场线、等势线和穿过圆柱侧面的通量。
\end{example}
% 题设配图；与紧随其后的解答分别编号。
\ExpansionDiagram{fig:audit-m-q-051}{无限直线电荷的柱面与径向场线}{
\draw[very thick](0,-1.5)--(0,2.1)node[above]{$\lambda$};\draw(0,-1)ellipse(1.4 and .35);\draw(0,1.3)ellipse(1.4 and .35);\draw(-1.4,-1)--(-1.4,1.3);\draw(1.4,-1)--(1.4,1.3);\draw[->](.2,.1)--(2.2,.1)node[right]{$\vct e_r$};\draw[<->](-1.9,-1)--(-1.9,1.3)node[midway,left]{$L$};
}

\begin{solution}
图\ref{fig:audit-m-q-051}中圆柱侧面的场沿外法向，端面与径向场垂直。
图\ref{fig:orig-042}中虚线圆为 $r$ 固定，实箭头沿径向。电势的负径向导数因此直接给出场强。

电势仅与 $r$ 有关，因此
\[
\vct E=-\nabla\Phi=-\frac{\dd\Phi}{\dd r}\vct e_r
=\frac Cr\vct e_r.
\]
对 $\lambda>0$，电场沿径向向外；对负线电荷则向内。等势条件为 $r=\text{常量}$，截面中是同心圆。场线的切线沿 $\vct e_r$，故是径向射线，两族曲线处处正交。

\begin{figure}[H]\centering
\begin{tikzpicture}[scale=.95]
\foreach\r in{.55,1,1.45}{\draw[dashed](0,0)circle(\r);}
\foreach\a in{0,45,...,315}{\draw[bookaxis](\a:.2)--(\a:1.85);}\fill(0,0)circle(1.7pt)node[below right]{$\lambda>0$};
\end{tikzpicture}
\caption{正线电荷的截面：径向实箭头为电场线，虚线圆为等势线。}
\label{fig:orig-042}
\end{figure}

取半径 $r$、长度 $\ell$ 的同轴圆柱面。场在侧面与外法向平行且大小恒定，所以
\[
\Phi_E=\frac Cr(2\pi r\ell)=2\pi C\ell
=\frac{\lambda\ell}{\varepsilon_0}.
\]
圆柱两端面的法向沿轴向，与电场垂直，不贡献通量。这与所围线电荷为 $\lambda\ell$ 的高斯定理相符。改变 $r_0$ 只给电势加常数，既不改变电场，也不改变通量。

\end{solution}

\clearpage
\section*{习题与解答}\addcontentsline{toc}{section}{习题与解答}
\begin{enumerate}[label=2.\arabic*]
\item \begin{minipage}[t]{\linewidth}
【基础】对 $\vct a=(1,2,0),\vct b=(0,1,2),\vct c=(2,0,1)$，求 $\vct a\cdot\vct b$、$\vct a\times\vct b$、$\vct a$ 与 $\vct b$ 的夹角余弦，以及 $(\vct a\times\vct b)\cdot\vct c$，并解释叉积的模与混合积的几何意义。

\end{minipage}
% 题设配图；与紧随其后的解答分别编号。
\ExpansionDiagram{fig:audit-m-q-052}{三个给定向量在同一三维基底下表示}{
\begin{scope}[x={(1cm,0cm)},y={(-.3cm,.6cm)},z={(.25cm,1cm)}]\draw[->](0,0,0)--(2.8,0,0)node[right]{$x$};\draw[->](0,0,0)--(0,2.7,0)node[left]{$y$};\draw[->](0,0,0)--(0,0,2.5)node[above]{$z$};\draw[->,thick](0,0,0)--(1,2,0)node[right,xshift=3mm,yshift=-2mm]{$\vct a$};\draw[->,thick](0,0,0)--(0,1,2)node[left]{$\vct b$};\draw[->,thick](0,0,0)--(2,0,1)node[right]{$\vct c$};\end{scope}
}

\begin{solution}
图\ref{fig:audit-m-q-052}中三个向量共用右手坐标系，叉积和混合积均应按这套次序计算。
图\ref{fig:ex-2-1}只表示空间方向关系。长度、夹角及体积都用题给三维分量计算，不能量取投影图。

先对 $\vct a,\vct b$ 运算：$\vct a\cdot\vct b=1\times0+2\times1+0\times2=2$，两模都为 $\sqrt5$，故夹角余弦为 $2/5$。叉积按分量展开为
\[
\vct a\times\vct b=(2\times2-0\times1,\ 0\times0-1\times2,\ 1\times1-2\times0)
=(4,-2,1).
\]
其模 $\sqrt{21}$ 是两向量张成的平行四边形面积。与 $\vct c$ 点乘，得混合积 $4\times2+(-2)\times0+1\times1=9$，绝对值9为相应平行六面体体积，正号给出所选次序的定向。

\begin{figure}[H]\centering
\begin{tikzpicture}[scale=.95]
\draw[bookaxis](0,0)--(3,0)node[right]{$x$};\draw[bookaxis](0,0)--(-1.05,-1.8)node[below left]{$y$};\draw[bookaxis](0,0)--(0,2.8)node[above]{$z$};
\draw[bookaxis,thick](0,0)--(.3,-1.2)node[below right]{$\vct a=(1,2,0)$};
\draw[bookaxis,thick](0,0)--(-.35,1.6)node[above left]{$\vct b=(0,1,2)$};
\draw[bookaxis,thick](0,0)--(2,1.1)node[above right]{$\vct c=(2,0,1)$};\node[below right]at(0,0){$O$};
\end{tikzpicture}
\caption{题2.1：三个向量的空间投影示意。长度与夹角应按题给分量计算。}\label{fig:ex-2-1}
\end{figure}

用恒等式检查叉积模：$|\vct a|^2|\vct b|^2-(\vct a\cdot\vct b)^2=25-4=21$，与直接计算一致。

叉积垂直于前两向量张成的面；混合积再取第三向量在该法向上的投影，得到带定向的体积9。
\end{solution}

\item \begin{minipage}[t]{\linewidth}
【基础】求点 $(1,2,3)$ 到平面 $2x-y+2z=4$ 的距离与垂足；另求直线 $(1,0,0)+t(1,2,-1)$ 与平面 $x+y+z=3$ 的交点与锐角。

\end{minipage}
% 题设配图；与紧随其后的解答分别编号。
\ExpansionDiagram{fig:audit-m-q-053}{点到平面作法向投影；另以参数直线求交}{
\begin{scope}[x={(.8cm,0cm)},y={(-.35cm,.24cm)},z={(0cm,.6cm)}]\filldraw[fill=gray!10](0,0,2)--(3,0,-1)--(3,3,.5)--(0,3,3.5)--cycle;\fill(1,2,3)circle(2pt)node[above]{$P(1,2,3)$};\fill(.555556,2.222222,2.555556)circle(2pt)node[below left]{$H$};\draw[dashed,thick](1,2,3)--(.555556,2.222222,2.555556);\draw[->](1,1,1.5)--(2,.5,2.5)node[right]{$\vct n=(2,-1,2)$};\end{scope}\node[below]at(1,-1.3){$2x-y+2z=4$};\node[right,align=left]at(4,1.2){另求交点：\\$\vct r=(1,0,0)+t(1,2,-1)$\\$x+y+z=3$};
}

\begin{solution}
图\ref{fig:audit-m-q-053}中$PH$沿平面法向；另一个求交问题使用图右列出的直线与平面，二者的平面条件须分别代入。
图\ref{fig:ex-2-2}左图沿法向找垂足，右图先找交点，再用法向计算线面角。

第一平面的法向为 $\vct n=(2,-1,2)$，模为3。把 $P=(1,2,3)$ 代入左侧减右侧，得 $2-2+6-4=2$，故距离为 $2/3$。垂足沿法向从 $P$ 减去有符号投影：
\[
H=P-\frac{\vct n\cdot P-4}{|\vct n|^2}\vct n
=(1,2,3)-\frac29(2,-1,2)
=(5/9,20/9,23/9).
\]
此点满足平面式，且 $P-H$ 平行法向，符合垂足定义。

\begin{figure}[H]\centering
\begin{tikzpicture}[scale=.85]
\begin{scope}
\draw(-.8,-.3)--(2,-.3)--(2.8,.6)--(0,.6)--cycle;\node at(2,.1){$\Pi_1$};
\coordinate(P)at(.75,2);\coordinate(H)at(.75,.2);\fill(P)circle(2pt)node[above]{$P$};\fill(H)circle(1.5pt)node[below]{$H$};
\draw[bookguide](P)--(H);\draw(.75,.4)--(.95,.4)--(.95,.2);\draw[bookaxis](H)--(.75,1.25)node[left]{$\vct n$};
\node[below,align=center]at(1,-.7){$PH=2/3$};
\end{scope}
\begin{scope}[shift={(5.7,0)}]
\draw(-.8,-.3)--(2,-.3)--(2.8,.6)--(0,.6)--cycle;\node at(2,.1){$\Pi_2$};
\coordinate(C)at(.6,.2);\draw[bookaxis](-.3,-.9)--(1.95,1.85)node[right]{$\vct d$};\fill(C)circle(1.7pt)node[below right]{$C$};
\draw[bookaxis](C)--(.6,1.65)node[left]{$\vct n_2$};\draw[bookguide](C)--(2,.2);\draw(1,.2)arc(0:50.71:.4);\node[right]at(1.05,.55){$\alpha$};
\node[below,align=center]at(1,-1.1){$C=(2,2,-1)$};
\end{scope}
\end{tikzpicture}
\caption{题2.2解图：左图沿法向求垂足，右图用法向与线方向的余角求线面角。图是空间关系的投影，角度数值由点积计算。}\label{fig:ex-2-2}
\end{figure}

第二问中直线点为 $(1+t,2t,-t)$，代入 $x+y+z=3$ 得 $1+2t=3$，所以 $t=1$，交点为 $(2,2,-1)$。线方向为 $(1,2,-1)$，平面法向为 $(1,1,1)$，两者点积为2。线与平面的锐角 $\alpha$ 因而满足
\[
\sin\alpha=\frac2{\sqrt6\sqrt3}=\frac{\sqrt2}{3}.
\]
法向与直线方向的角是所求角的余角，因此这里取正弦。

图中线面角是直线与其平面投影的锐角，它和线法向角互余，因此最后取正弦。
\end{solution}

\item \begin{minipage}[t]{\linewidth}
【基础】从包络方程说明起落等高的最大射程；对于 $v_0=10\,\mathrm{m/s},g=10\,\mathrm{m/s^2}$，求到达 $(5,0)\,\mathrm m$ 的两个发射角。

\end{minipage}
% 题设配图；与紧随其后的解答分别编号。
\ExpansionDiagram{fig:audit-m-q-054}{等高目标与固定大小的初速度}{
\draw[->](-.4,0)--(5.8,0)node[right]{$ x $};\draw[->](0,-0.3)--(0,3.1)node[above]{$ y $};\fill(0,0)circle(2pt);\draw[->,thick](0,0)--(1.3,1)node[above]{$v_0$};\draw[->](.6,0)arc(0:37.6:.6);\node[right]at(.65,0.2){$\alpha$};\fill(4,0)circle(2pt)node[above right]{$P(5,0)$};\draw[dashed](4,0)--(4,0)--(0,0);\draw[->](6.2,2.8)--(6.2,1.9)node[right]{$g$};
}

\begin{solution}
图\ref{fig:audit-m-q-054}中发射点与目标等高，发射方向决定同一初速对应的射程。
图\ref{fig:ex-2-3}把同一目标的低抛、高抛轨迹画在一起；两条曲线使用同一位置与长度尺度。

固定初速的可达区域满足 $y\le v_0^2/(2g)-gx^2/(2v_0^2)$。在地面 $y=0$ 上，条件化为 $x^2\le v_0^4/g^2$，所以正向最大射程为 $v_0^2/g$，题给参数下是10 m。

\begin{figure}[H]\centering
\begin{tikzpicture}[scale=.7]
\draw[bookaxis](-.3,0)--(6,0)node[right]{$x$};\draw[bookaxis](0,-.2)--(0,5.4)node[above]{$y$};
\draw[bookcurve,domain=0:5,samples=80]plot(\x,{.2679491924*\x-.0535898385*\x*\x});
\draw[bookcurve,domain=0:5,samples=80]plot(\x,{3.732050808*\x-.7464101615*\x*\x});
\fill(5,0)circle(2pt)node[below]{$P(5,0)\,\mathrm m$};\node[below left]at(0,0){$O$};
\draw[bookaxis](0,0)--({1.2*cos(15)},{1.2*sin(15)});
\draw[bookaxis](0,0)--({1.2*cos(75)},{1.2*sin(75)})node[left]{$v_0$};
\node[right]at(3.8,3.7){同一初速};
\node[right]at(1.15,.21){$15^\circ$};\node[right]at(.36,1.43){$75^\circ$};
\end{tikzpicture}
\caption{题2.3解图：同一初速的两条真实抛物线分别以 $15^\circ,75^\circ$ 发射，均在 $(5,0)$ 落地。}\label{fig:ex-2-3}
\end{figure}

起落等高时，除去 $t=0$ 的落地根，飞行时间为 $2v_0\sin\alpha/g$。代回水平位置，得到 $R=(v_0^2/g)\sin2\alpha$。要求 $R=5\,\mathrm m$，故 $\sin2\alpha=1/2$，在 $0<\alpha<\pi/2$ 中得到 $\alpha=15^\circ$ 或 $75^\circ$。

两角分别对应低抛、高抛，均小于最大射程条件。两条轨迹的飞行时间随 $\sin\alpha$ 改变，并不相同。

两条初速度的长度相同，倾角分别为 $15^\circ,75^\circ$；它们在同一位置落地，却不在同一时刻落地。
\end{solution}

\item \begin{minipage}[t]{\linewidth}
【基础】对 $r=at,\theta=bt$，其中 $m,a,b>0,t>0$，求速率、角动量与面积速度，并由角动量变化求切向力，核对运动学所得结果。

\end{minipage}
% 题设配图；与紧随其后的解答分别编号。
\ExpansionDiagram{fig:audit-m-q-055}{极坐标中位置、基矢及转动方向}{
\draw[->](-.4,0)--(4,0)node[right]{$ x $};\draw[->](0,-0.4)--(0,3.6)node[above]{$ y $};\draw[thick,domain=0:75,samples=100]plot({.04*\x*cos(\x)},{.04*\x*sin(\x)});\draw[->,thick](0,0)--(1.2,2.07846)node[above right]{$P$};\node[left]at(.5,1){$r=at$};\draw[->](.65,0)arc(0:60:.65);\node[right]at(.7,.45){$\theta=bt$};\draw[->](1.2,2.07846)--(1.7,2.94449)node[right]{$\vct e_r$};\draw[->](1.2,2.07846)--(.33397,2.57846)node[left]{$\vct e_\theta$};
}

\begin{solution}
图\ref{fig:audit-m-q-055}中面积速度由径向距离和角速度共同决定；其变化与切向力矩对应。
图\ref{fig:ex-2-4}先确定径向与角度增大方向；面积速度用径向距离与横向速度相乘求出。

设 $a,b>0,t>0$。极坐标速度两分量为 $v_r=\dot r=a$、$v_\theta=r\dot\theta=abt$，因此
\[
v=\sqrt{v_r^2+v_\theta^2}=a\sqrt{1+b^2t^2}.
\]
相对原点的角动量 $L_z=mr^2\dot\theta=ma^2bt^2$，面积速度为 $\dot S=r^2\dot\theta/2=a^2bt^2/2$。由角动量定理，
\[
\dot L_z=2ma^2bt=rF_\theta,
\qquad F_\theta=2mab.
\]
直接用极坐标加速度也得到 $a_\theta=r\ddot\theta+2\dot r\dot\theta=2ab$，两法一致。面积速度不恒定，故该运动不只受相对原点的中心力控制。

\begin{figure}[H]\centering
\begin{tikzpicture}[scale=.85]
\draw[bookaxis](-1.8,0)--(3.3,0)node[right]{$x$};\draw[bookaxis](0,-.5)--(0,4.1)node[above]{$y$};
\draw[bookcurve,domain=0:105,samples=100]plot({.035*\x*cos(\x)},{.035*\x*sin(\x)});
\coordinate(P)at(1.113693,1.113693);\draw[bookguide](0,0)--(P);\fill(P)circle(2pt)node[below right]{$P(t)$};
\draw[bookaxis](P)--(1.75,1.75)node[above right]{$\vct e_r$};\draw[bookaxis](P)--(.4773,1.75)node[above left]{$\vct e_\theta$};
\draw[->](.5,0)arc(0:45:.5);\node at(22.5:.8){$\theta$};\node[above left]at(.5,.5){$r$};\node[below left]at(0,0){$O$};
\end{tikzpicture}
\caption{题2.4：极坐标基矢与螺线轨迹示意；$\vct e_r$ 沿径向，$\vct e_\theta$ 沿角度增大方向。}\label{fig:ex-2-4}
\end{figure}

图中的横向分量 $r\dot\theta$ 随时间增加，故角动量不恒定，必须有非零横向力矩。
\end{solution}

\item \begin{minipage}[t]{\linewidth}
【提高】用球面面积元推导均匀薄球壳的转动惯量 $2MR^2/3$，再与实心球比较，解释系数的差别。

\end{minipage}
% 题设配图；与紧随其后的解答分别编号。
\ExpansionDiagram{fig:audit-m-q-056}{薄球壳的纬度圆带}{
\draw[thick](0,0)circle(1.6);\draw[dashed](0,0)ellipse(1.6 and .4);\draw[->](0,-2)--(0,2.3)node[above]{$z$};\draw[->](0,0)--(1.3,.93)node[midway,below]{$R$};\draw[dashed](-1.2,1.05)--(1.2,1.05);\node[left]at(-1.4,1.1){圆带};
}

\begin{solution}
图\ref{fig:audit-m-q-056}中球壳微元是球面圆带，不能使用实心球的体积薄层。
图\ref{fig:ex-2-5}的薄带是本题积分元，其到轴距离为 $R\sin\theta$，带宽沿子午线量取。

均匀薄球壳面密度为 $\sigma_m=M/(4\pi R^2)$。极角为 $\theta$ 的薄带到轴距离为 $R\sin\theta$，面积微元为 $R^2\sin\theta\dd\theta\dd\phi$。于是
\[
I_z=\int(R\sin\theta)^2\sigma_mR^2\sin\theta\dd\theta\dd\phi
=2\pi\sigma_mR^4\int_0^\pi\sin^3\theta\dd\theta.
\]
令$u=\cos\theta$，则$\dd u=-\sin\theta\dd\theta$，极角两端对应$u=1,-1$，故
\[\int_0^\pi\sin^3\theta\dd\theta
=-\int_1^{-1}(1-u^2)\dd u
=\left[u-\frac{u^3}{3}\right]_{-1}^1=\frac43.\]
所以$I_z=8\pi\sigma_mR^4/3=2MR^2/3$。

\begin{figure}[H]\centering
\begin{tikzpicture}[scale=1.05]
\draw[bookcurve](0,0)circle(1.6);
\draw[bookaxis](0,-1.9)--(0,2.15)node[above]{$z$};
\draw[bookguide,domain=0:180,variable=\t,samples=65]plot({1.6*cos(\t)},{1.6*sqrt(1-.93*.93)*sin(\t)});
\draw[domain=180:360,variable=\t,samples=65]plot({1.6*cos(\t)},{1.6*sqrt(1-.93*.93)*sin(\t)});
\foreach\ang in{55,65}{
\pgfmathsetmacro{\latr}{1.6*sin(\ang)}
\pgfmathsetmacro{\laty}{1.6*.93*cos(\ang)}
\pgfmathsetmacro{\lath}{\latr*sqrt(1-.93*.93)}
\pgfmathsetmacro{\cutangle}{asin(sqrt(1-.93*.93)*cos(\ang)/(.93*sin(\ang)))}
\draw[bookguide,domain=\cutangle:{180-\cutangle},variable=\t,samples=65]plot({\latr*cos(\t)},{\laty+\lath*sin(\t)});
\draw[bookcurve,domain={180-\cutangle}:{360+\cutangle},variable=\t,samples=90]plot({\latr*cos(\t)},{\laty+\lath*sin(\t)});}
\draw(0,0)--({1.6*sin(55)},{1.6*.93*cos(55)})node[pos=.8,below right]{$R$};
\draw[bookaxis,domain=0:55,variable=\t,samples=35]plot({.65*sin(\t)},{.65*.93*cos(\t)});
\node at(.23,.48){$\theta$};
\draw[<->,domain=55:65,variable=\t,samples=25]plot({1.85*sin(\t)},{1.85*.93*cos(\t)});
\node[right]at(1.85,.87){$\dd\theta$};
\draw[bookguide](0,.235)--(2,-.3)node[right]{薄带};\node[below left]at(0,0){$O$};
\end{tikzpicture}
\caption{题2.5：绕 $z$ 轴转动的球壳。取极角 $\theta$ 至 $\theta+\dd\theta$ 的薄带作积分元，图中带宽作放大示意。}\label{fig:ex-2-5}
\end{figure}

同质量、同外半径的实心球为 $2MR^2/5$。还可把实心球分成半径 $r$ 的薄壳，逐壳加上 $(2/3)r^2\dd m$；内部壳的 $r<R$，因此单位质量的平均到轴距离平方较小，系数也较小。

薄带质量中已有一个 $\sin\theta$，垂距平方再给出两个，因此极角积分出现 $\sin^3\theta$。
\end{solution}

\item \begin{minipage}[t]{\linewidth}
【提高】反平方力轨道的近心距和远心距分别为 $R,3R$。求 $a,e,p,h^2$ 与单位质量能量，参数 $\mu$ 给定。

\end{minipage}
% 题设配图；与紧随其后的解答分别编号。
\ExpansionDiagram{fig:audit-m-q-057}{近心点、远心点与焦点距离}{
\draw[thick](0,0)ellipse(2.2 and 1.9053);\fill(-1.1,0)circle(2pt)node[below]{$F$};\fill(-2.2,0)circle(2pt)node[left]{$P$};\fill(2.2,0)circle(2pt)node[right]{$A$};\draw[<->](-2.2,-.5)--(-1.1,-.5)node[midway,below]{$R$};\draw[<->](-1.1,.4)--(2.2,.4)node[midway,above]{$3R$};
}

\begin{solution}
图\ref{fig:audit-m-q-057}中两个端点距离均从力心$F$量起，不是从椭圆中心量起。
图\ref{fig:ex-2-6}的 $R,3R$ 都从焦点量起。先把它们分别对应 $a(1-e),a(1+e)$。

椭圆两端的焦距关系为 $r_p=a(1-e)$、$r_a=a(1+e)$。将 $r_p=R,r_a=3R$ 相加、相减，得
\[
a=2R,\qquad e=\frac{r_a-r_p}{r_a+r_p}=\frac12.
\]
由 $p=a(1-e^2)$ 和反平方轨道的 $p=h^2/\mu$，得 $p=3R/2$、$h^2=3\mu R/2$。椭圆单位质量能量为 $\mathcal E=-\mu/(2a)=-\mu/(4R)$。

\begin{figure}[H]\centering
\begin{tikzpicture}[scale=.95]
\draw[bookcurve](1,0)ellipse(2 and 1.732);\draw[bookaxis](-1.6,0)--(3.7,0);\fill(0,0)circle(2pt)node[above]{$F$};
\fill(-1,0)circle(2pt)node[above left]{$P_{\rm near}$};\fill(3,0)circle(2pt)node[above right]{$P_{\rm far}$};
\draw[<->](-1,-.55)--(0,-.55)node[midway,below]{$R$};\draw[<->](0,-.55)--(3,-.55)node[midway,below]{$3R$};
\end{tikzpicture}
\caption{题2.6：引力中心 $F$ 位于焦点，标出的两个距离均从 $F$ 量起。}\label{fig:ex-2-6}
\end{figure}

用近心点检验：该点径向速度为零、横向速度 $h/R$，所以 $\mathcal E=h^2/(2R^2)-\mu/R=3\mu/(4R)-\mu/R=-\mu/(4R)$。远心点以速度 $h/(3R)$ 代入，也得到同一能量。

以焦点为原点的近、远距离不能误当成椭圆半轴；图中的不对称位置正好确定偏心率。
\end{solution}

\item \begin{minipage}[t]{\linewidth}
【提高】正电荷圆环轴线上 $z>0$ 的场强最大位置和最大值是什么？证明这是全局最大值，而不仅是驻点。

\end{minipage}
% 题设配图；与紧随其后的解答分别编号。
\ExpansionDiagram{fig:audit-m-q-058}{带电圆环与轴上观察点}{
\draw[->](-.4,0)--(3.5,0)node[right]{$ r $};\draw[->](0,-0.4)--(0,3.8)node[above]{$ z $};\draw[thick](0,.3)ellipse(1.65 and .45);\node[left]at(-1.6,.4){$Q$};\draw[->](0,.3)--(1.65,.3)node[midway,below]{$R$};\fill(0,2.5)circle(2pt)node[right]{$P(z)$};\draw[dashed](-1.65,.3)--(0,2.5);
}

\begin{solution}
图\ref{fig:audit-m-q-058}中各电荷元的横向场分量相消，轴向分量相加后再对$z$求极值。
图\ref{fig:ex-2-7}把轴上点到圆环电荷元的距离写成 $\sqrt{R^2+z^2}$，投影后得到轴向场。

正电环的轴向场为
\[
E_z(z)=\frac{Q}{4\pi\varepsilon_0}\frac{z}{(R^2+z^2)^{3/2}},\qquad z>0.
\]
常数因子为正，只需求后面函数的最大值。求导并通分，得到
\[
\frac{\dd}{\dd z}\frac{z}{(R^2+z^2)^{3/2}}
=\frac1{(R^2+z^2)^{3/2}}
-\frac{3z}{2}(R^2+z^2)^{-5/2}(2z)
=\frac{R^2-2z^2}{(R^2+z^2)^{5/2}}.
\]
分母恒正，分子在 $z=R/\sqrt2$ 由正变负，故这是唯一极大点。再看两端，$z\to0^+$ 和 $z\to\infty$ 时场都趋零，所以该点是整个正轴的最大值。

\begin{figure}[H]\centering
\begin{tikzpicture}[scale=1]
\draw[bookcurve](0,0)ellipse(.45 and 1.55);\draw[bookaxis](-1,0)--(4.4,0)node[right]{$z$};
\fill(3,0)circle(2pt)node[above]{$P$};\fill(0,0)circle(1.5pt)node[below left]{$O$};
\draw(0,0)--(0,1.55)node[midway,left]{$R$};\draw[bookguide](0,1.55)--(3,0);\node[above left]at(0,1.65){$Q>0$};
\draw[<->](0,-.65)--(3,-.65)node[midway,below]{$z$};
\end{tikzpicture}
\caption{题2.7：均匀带电圆环与轴上观察点。}\label{fig:ex-2-7}
\end{figure}

代回极值点，几何因子为
\[\frac{R/\sqrt2}{(R^2+R^2/2)^{3/2}}
=\frac{R}{\sqrt2}\frac{2^{3/2}}{3^{3/2}R^3}
=\frac2{3\sqrt3R^2}.\]
故$E_{\max}=[Q/(4\pi\varepsilon_0)]\,2/(3\sqrt3R^2)$，具有场强的量纲。

求得的最大位置 $z=R/\sqrt2$ 比环半径小；环面处为零、远处也趋零，结合导数符号便能判定全局最大。
\end{solution}

\item \begin{minipage}[t]{\linewidth}
【提高】圆盘轴线上取 $z>0$，从电势重新推导电场；求 $z\to0^+$ 与 $z\to\infty$ 的电势、场强极限，并说明无限大圆盘不能继续使用无穷远为零的同一个电势极限。

\end{minipage}
% 题设配图；与紧随其后的解答分别编号。
\ExpansionDiagram{fig:audit-m-q-059}{带电圆盘分成同心环带}{
\draw[->](-.4,0)--(3.5,0)node[right]{$ r $};\draw[->](0,-0.4)--(0,3.8)node[above]{$ z $};\draw[thick](0,.3)ellipse(1.65 and .45);\draw[dashed](0,.3)ellipse(.8 and .22);\node[left]at(-1.6,.4){$\sigma$};\draw[->](0,.3)--(1.65,.3)node[midway,below]{$R$};\fill(0,2.5)circle(2pt)node[right]{$P(z)$};\draw[dashed](-1.65,.3)--(0,2.5);
}

\begin{solution}
图\ref{fig:audit-m-q-059}与图\ref{fig:ex-2-8}中的圆盘按同心细环分割。半径$\rho$、宽$\dd\rho$的细环面积为$2\pi\rho\dd\rho$，电荷为$\dd q=2\pi\sigma\rho\dd\rho$；各电荷元到轴上观察点的距离同为$\sqrt{z^2+\rho^2}$。电势先按标量相加，再对观察点坐标求负导数。以无穷远为零，对$z>0$有
\[
\Phi(z)=\frac{\sigma}{2\varepsilon_0}\int_0^R\frac{\rho\dd\rho}{\sqrt{z^2+\rho^2}}.
\]
令$u=z^2+\rho^2$，积分时$z$固定，$\dd u=2\rho\dd\rho$，两端对应$u=z^2,z^2+R^2$，所以
\[\Phi(z)=\frac{\sigma}{4\varepsilon_0}
\int_{z^2}^{z^2+R^2}u^{-1/2}\dd u
=\frac{\sigma}{4\varepsilon_0}[2\sqrt u]_{z^2}^{z^2+R^2}
=\frac{\sigma}{2\varepsilon_0}(\sqrt{R^2+z^2}-z).\]
这里$\sqrt{z^2}=z$使用了$z>0$。再对观察点坐标求负导数，
\[E_z=-\frac{\sigma}{2\varepsilon_0}
\left[\frac12(R^2+z^2)^{-1/2}(2z)-1\right]
=\frac{\sigma}{2\varepsilon_0}\left(1-\frac z{\sqrt{R^2+z^2}}\right).\]
故$z\to0^+$时，电势趋于$\sigma R/(2\varepsilon_0)$，场趋于$\sigma/(2\varepsilon_0)$。

\begin{figure}[H]\centering
\begin{tikzpicture}[scale=1]
\fill[black!5](0,0)ellipse(.55 and 1.65);\draw(0,0)ellipse(.55 and 1.65);\draw[bookguide](0,0)ellipse(.35 and 1.05);
\draw[bookaxis](-1,0)--(4.4,0)node[right]{$z$};\fill(3,0)circle(2pt)node[above]{$P$};
\draw(0,0)--(0,1.65)node[midway,left]{$R$};\node[left]at(-.55,.7){$\sigma$};
\draw[bookguide](0,1.05)--(3,0);\draw[<->](0,-.65)--(3,-.65)node[midway,below]{$z$};\node[below left]at(0,0){$O$};
\end{tikzpicture}
\caption{题2.8：均匀带电圆盘，可分成同轴细圆环后积分。}\label{fig:ex-2-8}
\end{figure}

远处展开 $\sqrt{R^2+z^2}=z+R^2/(2z)+O(z^{-3})$，可得
\[
\Phi\sim\frac{\sigma R^2}{4\varepsilon_0z},\qquad
E_z\sim\frac{\sigma R^2}{4\varepsilon_0z^2}.
\]
它们均趋零，且等于总电荷 $\pi R^2\sigma$ 的点电荷首项。若固定 $z$ 让 $R\to\infty$，场有有限极限，电势却随 $R$ 发散；无限平面只能另选有限位置作零点。有限两点间电势差仍趋于 $-\sigma(z_2-z_1)/(2\varepsilon_0)$。

远处整盘可缩成一点电荷；靠近盘面时，各半径细环的贡献都要保留，所以不能用同一个点电荷近似。
\end{solution}

\item \begin{minipage}[t]{\linewidth}
【提高】对 $y=ax^2$，$a>0$，求任意 $x$ 处的单位切向、主法向和曲率；以恒定速率 $v$ 通过该点时，加速度的直角坐标分量是什么？

\end{minipage}
% 题设配图；与紧随其后的解答分别编号。
\ExpansionDiagram{fig:audit-m-q-060}{抛物线上的任意观察点与弧长方向}{
\draw[->](-.4,0)--(3,0)node[right]{$ x $};\draw[->](0,-0.3)--(0,2.6)node[above]{$ y $};\draw[thick,domain=-2:2,samples=70]plot(\x,{.4*\x*\x});\fill(1.2,.576)circle(2pt)node[below right]{$P(x,ax^2)$};\draw[->](1.2,.576)--(2.2,1.536)node[right]{$\vct T$};\draw[->](1.2,.576)--(.24,1.576)node[left]{$\vct N$};
}

\begin{solution}
图\ref{fig:audit-m-q-060}中先取沿$x$增大方向的切线，再求指向曲线凹侧的主法线。
图\ref{fig:ex-2-9}中的速度沿单位切向，恒速加速度沿指向弯曲侧的主法向。

取沿 $x$ 增大的运动方向，切向量为 $(1,2ax)$，归一化得
\[
\vct T=\frac{(1,2ax)}{\sqrt{1+4a^2x^2}}.
\]
记$D=1+4a^2x^2$，先用乘积法则对$\vct T=(1,2ax)D^{-1/2}$求导：
\[\frac{\dd\vct T}{\dd x}
=(0,2a)D^{-1/2}+(1,2ax)\left(-\frac12D^{-3/2}\right)(8a^2x)
=\frac{(-4a^2x,2a)}{D^{3/2}}.\]
再由$\dd s/\dd x=\sqrt D$换成弧长导数，
\[\frac{\dd\vct T}{\dd s}=\frac{(-4a^2x,2a)}{D^2},\qquad
\left|\frac{\dd\vct T}{\dd s}\right|
=\frac{2a\sqrt D}{D^2}=\frac{2a}{D^{3/2}}.\]
取模得到曲率，再将弧长导数归一化得到主法向：
\[
\kappa=\frac{2a}{(1+4a^2x^2)^{3/2}},\qquad
\vct N=\frac{(-2ax,1)}{\sqrt{1+4a^2x^2}}.
\]
因速率恒定，切向加速度为零，$\vct a=v^2\kappa\vct N$，所以直角坐标分量为
\[
a_x=-\frac{4a^2xv^2}{(1+4a^2x^2)^2},\qquad
 a_y=\frac{2av^2}{(1+4a^2x^2)^2}.
\]
在顶点恢复 $(0,2av^2)$。若逆向通过同一点，$\vct T$ 变号，主法向及同速率下的加速度不变。

\begin{figure}[H]\centering
\begin{tikzpicture}[scale=.95]
\draw[bookaxis](-2.3,0)--(2.8,0)node[right]{$x$};\draw[bookaxis](0,-.3)--(0,2.8)node[above]{$y$};
\draw[bookcurve,domain=-2:2,samples=70]plot(\x,{.5*\x*\x});\coordinate(P)at(1.2,.72);\fill(P)circle(2pt)node[right]{$P(x,ax^2)$};
\draw[bookguide](.4,-.24)--(2,1.68);\draw[bookaxis,thick](P)--(1.78,1.416)node[above]{$\vct v$};
\node[above left]at(-1.5,1.5){$y=ax^2$};\node[below left]at(0,0){$O$};
\draw[bookaxis](P)--++(-1.05,.7)node[above]{$\vct N$};
\end{tikzpicture}
\caption{题2.9解图：实箭头为切向速度，辅助箭头为指向弯曲一侧的主法向；恒速加速度沿主法向。}\label{fig:ex-2-9}
\end{figure}

所得加速度两分量的比例与图中的法向 $(-2ax,1)$ 相同；在顶点水平分量为零，竖直分量向上。
\end{solution}

\item \begin{minipage}[t]{\linewidth}
【提高】对场 $\vct F=C(-y,x)$，分别用参数积分和格林定理求逆时针椭圆 $x=a\cos t,y=b\sin t$ 上的环流，$a,b>0$。

\end{minipage}
% 题设配图；与紧随其后的解答分别编号。
\ExpansionDiagram{fig:audit-m-q-061}{逆时针椭圆边界及内部区域}{
\draw[->](-.4,0)--(3,0)node[right]{$ x $};\draw[->](0,-1.8)--(0,2)node[above]{$ y $};\draw[thick,->](2.4,0)arc(0:330:2.4 and 1.2);\draw(330:2.4 and 1.2)arc(330:360:2.4 and 1.2);\node[below]at(2.4,0){$a$};\node[left]at(0,1.2){$b$};
}

\begin{solution}
图\ref{fig:audit-m-q-061}中边界为逆时针，格林定理相应取正的面积方向。
沿图\ref{fig:ex-2-10}的逆时针方向参数化；从长轴正端出发时，速度先向上。

取 $x=a\cos t,y=b\sin t$，$0\le t\le2\pi$。沿图中的逆时针走向，有
\[
\dd\vct r=(-a\sin t,b\cos t)\dd t,\qquad
\vct F=C(-b\sin t,a\cos t).
\]
两向量点乘，得到
\[
\vct F\cdot\dd\vct r=Cab(\sin^2t+\cos^2t)\dd t=Cab\dd t.
\]
直接积分为$\int_0^{2\pi}Cab\dd t=Cab[t]_0^{2\pi}=2\pi Cab$。

\begin{figure}[H]\centering
\begin{tikzpicture}[scale=1]
\draw[bookaxis](-2.6,0)--(2.9,0)node[right]{$x$};\draw[bookaxis](0,-1.7)--(0,2)node[above]{$y$};\draw[bookcurve](0,0)ellipse(2 and 1.1);
\draw[bookaxis,domain=20:65,samples=30]plot({2*cos(\x)},{1.1*sin(\x)});\node[right]at(1.9,.8){$C$};
\node[below right]at(2,0){$a$};\node[above right]at(0,1.1){$b$};\node[below left]at(0,0){$O$};
\end{tikzpicture}
\caption{题2.10：椭圆边界取逆时针方向。}\label{fig:ex-2-10}
\end{figure}

另用格林定理，$\partial_xF_y-\partial_yF_x=C-(-C)=2C$。椭圆面积为 $\pi ab$，故环流为 $\iint2C\dd S=2\pi Cab$，与参数积分一致。两个计算都使用逆时针方向；反向绕行时应整体变号。

\end{solution}

\item \begin{minipage}[t]{\linewidth}
【提高】$f=x^2+y^2+2z^2$。在点 $(1,2,1)$ 求梯度、沿 $(1,1,1)$ 的单位方向导数、最大方向导数及该等值面的切平面。

\end{minipage}
% 题设配图；与紧随其后的解答分别编号。
\ExpansionDiagram{fig:audit-m-q-062}{等值面上一点的梯度与指定方向}{
\draw[rotate=-8.112,gray](0,0)ellipse(2.87 and 1.49);
\begin{scope}[x={(1cm,0cm)},y={(-.4cm,.3cm)},z={(0cm,.7cm)}]
\draw[dashed,domain=0:360,samples=100]plot({sqrt(7)*cos(\x)},{sqrt(7)*sin(\x)},0);
\fill(1,2,1)circle(2pt)node[below right]{$P(1,2,1)$};
\draw[->,thick](1,2,1)--(2,3,2)node[right]{$(1,1,1)$};
\draw[->,thick](1,2,1)--(1.6,3.2,2.2)node[left]{$\nabla f$};
\end{scope}
\node[below]at(0,-1.8){$x^2+y^2+2z^2=7$};

}

\begin{solution}
图\ref{fig:audit-m-q-062}中给定方向须归一化，等值面的法向则由该点梯度给出。
图\ref{fig:ex-2-11}先固定所经过的等值面和题给方向；梯度法向与题给方向一般不同。

逐项求偏导，$\nabla f=(2x,2y,4z)$，在 $P=(1,2,1)$ 处为 $(2,4,4)$。把题给方向先归一化为 $\vct n=(1,1,1)/\sqrt3$，方向导数为
\[
D_{\vct n}f=\nabla f(P)\cdot\vct n=\frac{10}{\sqrt3}.
\]
由柯西--施瓦茨不等式，单位方向上的导数最多为 $|\nabla f(P)|=6$，取等号的方向为 $(1,2,2)/3$。

\begin{figure}[H]\centering
\begin{tikzpicture}[x={(.7cm,-.25cm)},y={(.7cm,.25cm)},z={(0cm,.95cm)},scale=.95]
\begin{scope}[x=1cm,y=1cm]\draw[bookcurve](0,0)ellipse(2.61916 and 2.00842);\end{scope}
\draw[bookaxis](0,0,0)--(3.6,0,0)node[below right]{$x$};
\draw[bookaxis](0,0,0)--(0,3.6,0)node[below right,yshift=-4pt]{$y$};
\draw[bookaxis](0,0,0)--(0,0,2.6)node[above]{$z$};
\draw[black!60,domain=0:360,variable=\t,samples=90]plot({sqrt(7)*cos(\t)},{sqrt(7)*sin(\t)},0);
\draw[bookguide,domain=0:360,variable=\t,samples=90]plot({sqrt(7)*cos(\t)},0,{sqrt(3.5)*sin(\t)});
\draw[bookguide,domain=0:360,variable=\t,samples=90]plot(0,{sqrt(7)*cos(\t)},{sqrt(3.5)*sin(\t)});
\fill(1,2,1)circle(2pt)node[above left,fill=white,inner sep=1pt]{$P(1,2,1)$};
\draw[bookaxis,thick](1,2,1)--(1.65,2.65,1.65)node[above right]{$(1,1,1)$ 方向};
\node[left]at(-1.8,-1.8,1.8){$S:f=f(P)$};\node[below left]at(0,0,0){$O$};
\end{tikzpicture}
\caption{题2.11：经过点 $P$ 的椭球等值面，箭头表示方向 $(1,1,1)$。}\label{fig:ex-2-11}
\end{figure}

点 $P$ 所在等值面的函数值为 $1+4+2=7$。梯度垂直于切平面，故平面方程为
\[
2(x-1)+4(y-2)+4(z-1)=0,
\quad\text{即 }x+2y+2z=7.
\]
将 $P$ 代回等于7，法向也与所求梯度平行。

切平面必须垂直 $\nabla f(P)=(2,4,4)$，而方向导数取它沿单位 $(1,1,1)$ 的投影；两项不能共用未归一化的方向。
\end{solution}

\item \begin{minipage}[t]{\linewidth}
【提高】场 $\vct A=C(-y,x,0)$，求任意位置、半径 $R$ 的倾斜圆周环流。圆面法向与正 $z$ 轴夹角为 $\beta$，边界方向依右手规则。

\end{minipage}
% 题设配图；与紧随其后的解答分别编号。
\ExpansionDiagram{fig:audit-m-q-063}{倾斜圆面的法向与右手边界方向}{
\begin{scope}[x={(1cm,0cm)},y={(-.35cm,.35cm)},z={(0cm,1cm)}]\draw[->](0,0,0)--(0,0,2.3)node[above]{$z$};\draw[thick,domain=0:360,samples=100]plot({1.6*cos(\x)},{1.6*.8*sin(\x)},{1.6*.6*sin(\x)});\draw[->,thick,domain=15:85,samples=30]plot({1.6*cos(\x)},{1.6*.8*sin(\x)},{1.6*.6*sin(\x)});\draw[->,thick](0,0,0)--(0,-1.2,1.6)node[right]{$\vct n$};\draw[->](0,0,0)--(1.6,0,0)node[midway,below]{$R$};\end{scope}\node[right]at(1,2){$\beta=\angle(\vct n,\vct e_z)$};\node[below]at(0,-1.5){圆面水平有向投影为$\pi R^2\cos\beta$};
}

\begin{solution}
图\ref{fig:audit-m-q-063}中边界方向与圆面法向依右手规则联系，旋度通量只取法向投影。
图\ref{fig:ex-2-12}用右手规则把边界走向与圆面法向配对，通量中保留带符号的 $\cos\beta$。

任意位置直接沿倾斜圆周参数积分较繁，旋度却是常向量：
\[
\nabla\times\vct A=(0,0,\partial_x(Cx)-\partial_y(-Cy))=(0,0,2C).
\]
取圆周所围平圆盘为跨接面，单位法向与正 $z$ 夹角为 $\beta$。按题给右手定向，斯托克斯定理给出
\[
\oint_C\vct A\cdot\dd\vct r
=\iint_S(\nabla\times\vct A)\cdot\vct n\dd S
=2\pi CR^2\cos\beta.
\]
场的旋度均匀，故平移圆心不改变结果。$\beta=0$ 恢复水平逆时针圆的正环流；$\beta=\pi/2$ 时为零。这里保留带符号的 $\cos\beta$，不能取绝对值。

\begin{figure}[H]\centering
\begin{tikzpicture}[scale=1]
\draw[bookcurve,rotate=-32](0,0)ellipse(2 and .65);\draw[bookaxis](0,0)--(0,2.4)node[above]{$z$};
\draw[bookaxis,thick](0,0)--(1.2,1.92)node[right]{$\vct n$};\draw[->](0,.85)arc(90:58:.85);\node at(.35,1.05){$\beta$};
\draw[bookaxis,rotate=-32,domain=15:60,samples=30]plot({2*cos(\x)},{.65*sin(\x)});
\draw(0,0)--(-1.696,1.06)node[midway,below left]{$R$};\node[below]at(0,-1.3){圆心位置任意};
\end{tikzpicture}
\caption{题2.12：倾斜圆面、单位法向与边界方向。$\vct n$ 和边界方向按右手规则对应。}\label{fig:ex-2-12}
\end{figure}

把圆面投影到 $xy$ 平面后，有向面积为 $\pi R^2\cos\beta$，再乘常旋度 $2C$ 即为环流。
\end{solution}

\item \begin{minipage}[t]{\linewidth}
【提高】对 $\vct A=C\vct r/r^3$，求 $r>0$ 处散度；在 $R_1<r<R_2$ 球壳上比较体积分和外、内边界通量。解释为什么含原点球体不能按同样的无源推理处理。

\end{minipage}
% 题设配图；与紧随其后的解答分别编号。
\ExpansionDiagram{fig:audit-m-q-064}{球壳内、外边界的外法向相反}{
\fill[gray!10,even odd rule](0,0)circle(1.7)(0,0)circle(.8);\draw(0,0)circle(1.7);\draw(0,0)circle(.8);\draw[->](1.7,0)--(2.6,0)node[right]{$\vct n_2$};\draw[->](.8,0)--(.2,0)node[midway,above]{$\vct n_1$};\node[below]at(.5,-.6){$R_1$};\node[below]at(1.4,-1.1){$R_2$};
}

\begin{solution}
图\ref{fig:audit-m-q-064}中内球面的外法向朝球心，两个边界的有向通量必须采用相反径向。
图\ref{fig:ex-2-13}明确了积分区域是球壳：内边界的区域外法向指向球心。

场只有径向分量 $A_r=C/r^2$，在 $r>0$ 区域可用球坐标散度式，
\[
\nabla\cdot\vct A=\frac1{r^2}\frac{\dd}{\dd r}(r^2A_r)
=\frac1{r^2}\frac{\dd C}{\dd r}=0.
\]
球壳上的体积分为零。外边界的外法向为 $+\vct e_r$，内边界的外法向为 $-\vct e_r$，指向空腔。两边界通量分别为
\[
\Phi_{\rm out}=\frac C{R_2^2}4\pi R_2^2=4\pi C,\qquad
\Phi_{\rm in}=-\frac C{R_1^2}4\pi R_1^2=-4\pi C.
\]
相加为零，与体积分相符。

\begin{figure}[H]\centering
\begin{tikzpicture}[scale=1]
\draw(0,0)circle(1.7);\draw(0,0)circle(.85);\draw[bookaxis](.85,0)--(.2,0)node[above]{$\vct n_{\rm in}$};
\draw[bookaxis](1.7,0)--(2.45,0)node[right]{$\vct n_{\rm out}$};
\draw(0,0)--(120:1.7)node[midway,left]{$R_2$};\draw(0,0)--(230:.85)node[midway,below left]{$R_1$};
\fill(0,0)circle(1.6pt)node[below right]{$O$};\node at(0,1.25){积分区域};
\end{tikzpicture}
\caption{题2.13：球壳截面。区域外法向在内边界指向球心，在外边界背离球心。}\label{fig:ex-2-13}
\end{figure}

若改为包含原点的实心球，场在原点发散，不满足在整个体积内连续可微的条件。球面通量仍为 $4\pi C$；不能把仅在 $r>0$ 得到的零散度结论未经处理延伸到原点。

两球面大小不同但径向场按 $r^{-2}$ 变化，通量的模相同；图中相反的法向使两项恰好抵消。
\end{solution}

\item \begin{minipage}[t]{\linewidth}
【选做】令复势 $W=z^3$，求 $\Phi,\Psi$、二维电场形式、两族等值线的正交关系；在点 $(1,1)$ 判断电场方向。
\end{minipage}
% 题设配图；与紧随其后的解答分别编号。
\ExpansionDiagram{fig:audit-m-q-065}{复势把同一个点映射为两种实标量}{
\draw[->](-.4,0)--(3.5,0)node[right]{$ x $};\draw[->](0,-0.3)--(0,2.6)node[above]{$ y $};\fill(1.5,1.5)circle(2pt)node[above right]{$P(1,1)$};\draw[dashed](1.5,0)--(1.5,1.5)--(0,1.5);\node[right]at(3.8,1){$W=z^3=\Phi+\ii\Psi$};
}

\begin{solution}
图\ref{fig:audit-m-q-065}中的$\Phi,\Psi$是同一解析函数的实、虚部，先展开再求各自梯度。
图\ref{fig:hw-2-14}画出经过 $(1,1)$ 的两条等值线，用来检查下面算得的场方向与正交性。

按二项式展开，
\[
W=(x+\ii y)^3=(x^3-3xy^2)+\ii(3x^2y-y^3),
\]
所以 $\Phi=x^3-3xy^2$、$\Psi=3x^2y-y^3$。由 $\vct E=-\nabla\Phi$，得到 $\vct E=(-3x^2+3y^2,6xy)$。

\begin{figure}[H]\centering
\begin{tikzpicture}[scale=2]
\draw[bookaxis](.2,0)--(2.3,0)node[right]{$x$};\draw[bookaxis](0,.1)--(0,2)node[above]{$y$};
\draw[bookcurve]plot coordinates{(0.6669617,1.0713715) (0.6806375,1.0648457) (0.6944196,1.0586661) (0.7082867,1.0528351) (0.7222176,1.0473533) (0.7361917,1.0422191) (0.7501888,1.0374293) (0.7641898,1.0329790) (0.7781765,1.0288618) (0.7921319,1.0250699) (0.8060403,1.0215946) (0.8198870,1.0184261) (0.8336590,1.0155540) (0.8473443,1.0129673) (0.8609324,1.0106545) (0.8744139,1.0086040) (0.8877806,1.0068041) (0.9010257,1.0052429) (0.9141431,1.0039088) (0.9271280,1.0027904) (0.9399762,1.0018764) (0.9526847,1.0011558) (0.9652509,1.0006180) (0.9776730,1.0002530) (0.9899500,1.0000508) (1.0020812,1.0000022) (1.0140664,1.0000980) (1.0259058,1.0003299) (1.0376003,1.0006896) (1.0491506,1.0011695) (1.0605580,1.0017623) (1.0718239,1.0024611) (1.0829501,1.0032594) (1.0939382,1.0041510) (1.1047903,1.0051302) (1.1155083,1.0061914) (1.1260945,1.0073296) (1.1365511,1.0085399) (1.1468802,1.0098177) (1.1570842,1.0111588) (1.1671654,1.0125592) (1.1771262,1.0140150) (1.1869688,1.0155227) (1.1966957,1.0170790) (1.2063090,1.0186807) (1.2158111,1.0203249) (1.2252042,1.0220089) (1.2344907,1.0237299) (1.2436726,1.0254855) (1.2527522,1.0272736) (1.2617316,1.0290918) (1.2706128,1.0309382) (1.2793980,1.0328108) (1.2880891,1.0347079) (1.2966880,1.0366277) (1.3051968,1.0385687) (1.3136173,1.0405293) (1.3219513,1.0425082) (1.3302006,1.0445041) (1.3383670,1.0465156) (1.3464521,1.0485416) (1.3544577,1.0505810) (1.3623853,1.0526328) (1.3702365,1.0546959) (1.3780128,1.0567694) (1.3857158,1.0588525) (1.3933470,1.0609443) (1.4009076,1.0630441) (1.4083992,1.0651510) (1.4158231,1.0672646) (1.4231806,1.0693839) (1.4304729,1.0715086) (1.4377013,1.0736378) (1.4448671,1.0757713) (1.4519714,1.0779083) (1.4590154,1.0800484) (1.4660002,1.0821911) (1.4729269,1.0843361) (1.4797965,1.0864828) (1.4866102,1.0886310) (1.4933688,1.0907802)};
\draw[dashed]plot coordinates{(1.0907802,1.4933688) (1.0878160,1.4840322) (1.0848541,1.4745903) (1.0818954,1.4650403) (1.0789411,1.4553795) (1.0759922,1.4456049) (1.0730500,1.4357136) (1.0701160,1.4257024) (1.0671916,1.4155682) (1.0642784,1.4053075) (1.0613781,1.3949170) (1.0584927,1.3843929) (1.0556241,1.3737316) (1.0527747,1.3629292) (1.0499467,1.3519816) (1.0471429,1.3408848) (1.0443659,1.3296344) (1.0416190,1.3182260) (1.0389054,1.3066549) (1.0362287,1.2949164) (1.0335929,1.2830056) (1.0310023,1.2709173) (1.0284615,1.2586464) (1.0259757,1.2461874) (1.0235502,1.2335349) (1.0211911,1.2206831) (1.0189051,1.2076262) (1.0166991,1.1943583) (1.0145810,1.1808734) (1.0125592,1.1671654) (1.0106429,1.1532283) (1.0088423,1.1390559) (1.0071682,1.1246421) (1.0056326,1.1099811) (1.0042483,1.0950671) (1.0030296,1.0798946) (1.0019916,1.0644587) (1.0011510,1.0487547) (1.0005257,1.0327787) (1.0001351,1.0165278) (1.0000000,1.0000000) (1.0001428,0.9831945) (1.0005874,0.9661122) (1.0013593,0.9487558) (1.0024854,0.9311302) (1.0039938,0.9132427) (1.0059139,0.8951037) (1.0082761,0.8767269) (1.0111111,0.8581295) (1.0144497,0.8393330) (1.0183222,0.8203632) (1.0227578,0.8012505) (1.0277836,0.7820301) (1.0334238,0.7627417) (1.0396989,0.7434299) (1.0466245,0.7241430) (1.0542106,0.7049328) (1.0624608,0.6858538) (1.0713715,0.6669617) (1.0809320,0.6483123) (1.0911236,0.6299605) (1.1019209,0.6119584) (1.1132915,0.5943541) (1.1251972,0.5771911) (1.1375950,0.5605068) (1.1504385,0.5443319) (1.1636786,0.5286907) (1.1772656,0.5136004) (1.1911497,0.4990717) (1.2052822,0.4851090) (1.2196165,0.4717113) (1.2341086,0.4588726) (1.2487175,0.4465831) (1.2634058,0.4348293) (1.2781395,0.4235952) (1.2928883,0.4128628) (1.3076251,0.4026127) (1.3223263,0.3928247) (1.3369714,0.3834780) (1.3515428,0.3745518) (1.3660254,0.3660254)};
\fill(1,1)circle(1pt)node[below left]{$P(1,1)$};
\draw[bookaxis,thick](1,1)--(1,1.65)node[above]{$\vct E$};\draw[bookaxis](1,1)--(1.65,1)node[right]{$\nabla\Psi$};
\draw(1.12,1)--(1.12,1.12)--(1,1.12);
\node[left]at(.7,1.4){$\Phi=-2$};\node[below]at(1.4,.8){$\Psi=2$};
\end{tikzpicture}
\caption{题2.14解图：实线 $\Phi=-2$ 与虚线 $\Psi=2$ 在 $P$ 正交。$\vct E=(0,6)$ 向上，沿 $\Psi$ 等值线的切线。}\label{fig:hw-2-14}
\end{figure}

再计算两族等值线的法向：
\[
\nabla\Phi=(3x^2-3y^2,-6xy),\qquad
\nabla\Psi=(6xy,3x^2-3y^2).
\]
两者点积为零，因此在梯度非零处，两族曲线相互正交；同时 $\vct E\cdot\nabla\Psi=0$，说明 $\Psi$ 的等值线可作为场线。在 $(1,1)$，电场为 $(0,6)$，沿正 $y$。原点梯度为零，不能用正则等值线的切线关系判断该处方向。

在 $P$，$\nabla\Phi=(0,-6)$、$\nabla\Psi=(6,0)$，两条法向互相垂直；电场向上，正好沿 $\Psi=2$ 曲线的切线。
\end{solution}

\item \begin{minipage}[t]{\linewidth}
【综合】小珠沿光滑竖直轨道 $y=x^4/(4L^3)$ 运动，从 $x=a$ 处静止释放，$a\ll L$。求首阶周期的积分表达式及其对 $a$ 的依赖，解释为什么原点稳定却没有非零的小振动极限频率。

\end{minipage}
% 题设配图；与紧随其后的解答分别编号。
\ExpansionDiagram{fig:audit-m-q-066}{四次轨道在原点附近比抛物线更平坦}{
\draw[->](-.4,0)--(3,0)node[right]{$ x $};\draw[->](0,-0.4)--(0,2.4)node[above]{$ y $};\draw[thick,domain=-2:2,samples=80]plot(\x,{.09*(\x)^4});\fill(1.7,{.09*1.7^4})circle(2pt)node[above right]{$A$};\draw[dashed](1.7,0)--(1.7,{.09*1.7^4});\node[below]at(1.7,0){$a$};\node[left]at(-1.7,1.8){$ y=x^4/(4L^3) $};\draw[->](2.6,1.7)--(2.6,.7)node[right]{$g$};
}

\begin{solution}
图\ref{fig:audit-m-q-066}中高度的首项为$x^4$，所以小振幅周期应从能量积分求尺度。
图\ref{fig:ex-2-15}把两侧转折点与轨道最低点对应到能量积分的上下限。

约束给出 $\dot y=(x^3/L^3)\dot x$，因而
\[
T=\frac m2(1+x^6/L^6)\dot x^2,\qquad U=\frac{mgx^4}{4L^3}.
\]
由释放初态能量 $mga^4/(4L^3)$，可解出
\[
\dot x^2=\frac{g(a^4-x^4)}{2L^3(1+x^6/L^6)}.
\]
利用往返对称，取从$0$到$a$的运动支计算四分之一周期，完整积分为
\[T=4\sqrt{\frac{2L^3}{g}}\int_0^a
\frac{\sqrt{1+x^6/L^6}}{\sqrt{a^4-x^4}}\dd x.\]
当$a\ll L$时，分子根号相对$1$的修正一致为$O((a/L)^6)$。令$x=au$，则$\dd x=a\dd u$、$\sqrt{a^4-x^4}=a^2\sqrt{1-u^4}$，两端变为$0,1$，所以首阶周期为
\[
T\simeq\frac4a\sqrt{\frac{2L^3}{g}}\int_0^1\frac{\dd u}{\sqrt{1-u^4}}.
\]
在$u=1$附近，$1-u^4=(1-u)(1+u+u^2+u^3)\sim4(1-u)$，故端点核与可积的$(1-u)^{-1/2}$同阶。积分为有限无量纲常数，因此$T\propto a^{-1}$。原点是四次势能的严格极小点，因而稳定；势能两导数为$U'=mgx^3/L^3$、$U''=3mgx^2/L^3$，故$U''(0)=0$，恢复力没有线性首项，非零简谐极限频率不存在。

\begin{figure}[H]\centering
\begin{tikzpicture}[scale=.95]
\draw[bookaxis](-2.5,0)--(2.8,0)node[right]{$x$};\draw[bookaxis](0,-.3)--(0,2.6)node[above]{$y$};
\draw[bookcurve,domain=-2.1:2.1,samples=90]plot(\x,{.11*(\x)^4});\coordinate(P)at(1.8,1.155);\fill(P)circle(2.3pt)node[above left]{静止释放};
\draw[bookguide](1.8,0)--(P);\node[below]at(1.8,0){$a$};\draw[bookaxis,thick](P)--(1.8,.35)node[right]{$mg$};
\node[left]at(-1.4,1.8){$y=x^4/(4L^3)$};\node[below left]at(0,0){$O$};
\end{tikzpicture}
\caption{题2.15：四次曲轨与静止释放的位置。}\label{fig:ex-2-15}
\end{figure}

从最低点到右侧转折点只走一个四分之一周期；横坐标换元 $x=au$ 后，图中行程尺度与四次势能共同产生 $T\propto a^{-1}$。
\end{solution}

\item \begin{minipage}[t]{\linewidth}
【综合】水平光滑直杆绕固定端以恒定角速度 $\Omega$ 转动，小珠可沿杆无摩擦滑动。初始距轴 $r_0$、相对杆静止。求 $r(t)$、杆的横向作用力及维持匀速转动的装置对小珠所做的功。检查能量变化与角动量变化的关系。

\end{minipage}
% 题设配图；与紧随其后的解答分别编号。
\ExpansionDiagram{fig:audit-m-q-067}{绕固定端转动的光滑杆与可滑动小珠}{
\fill(0,0)circle(2pt)node[below]{$O$};\draw[very thick](0,0)--(4,2.3);\fill(2,1.15)circle(3pt)node[above]{$m$};\draw[->](0,-.6)--(2.7,-.6)node[right]{$x$};\draw[->](.8,0)arc(0:29.9:.8);\node[right]at(.9,.25){$\Omega t$};\draw[<->](.1,.5)--(1.8,1.48)node[midway,above]{$r$};
}

\begin{solution}
图\ref{fig:audit-m-q-067}中杆的转角由外部装置规定，小珠的独立自由度是沿杆距离$r$。
图\ref{fig:ex-2-16}按径向、横向分解速度和加速度；光滑杆的支持力只有横向分量。

取极坐标且$\dot\theta=\Omega$。在水平面内，无摩擦的杆只提供横向约束力；竖直支持力与重力相抵。径向实合力为零，所以$\ddot r-r\Omega^2=0$。若$\Omega\ne0$，特征根为$\pm\Omega$，通解$r=A\ee^{\Omega t}+B\ee^{-\Omega t}$。初值给出
\[A+B=r_0,\qquad\Omega(A-B)=0,\qquad A=B=\frac{r_0}{2}.\]
由双曲函数定义得到
\[
r=r_0\cosh\Omega t,\qquad\dot r=\Omega r_0\sinh\Omega t.
\]
$\Omega=0$时径向方程为$\ddot r=0$，两初值给出$r=r_0$，也与上述公式的极限一致。横向加速度为$2\Omega\dot r$，故杆作用力为$N=2m\Omega\dot r$。在题给正转方向下，小珠启动后向外滑，$N$沿转动切向。

\begin{figure}[H]\centering
\begin{tikzpicture}[scale=.95]
\draw[bookaxis](-.5,0)--(4,0)node[right]{$x$};\draw[bookaxis](0,-.4)--(0,3.1)node[above]{$y$};
\draw[thick](0,0)--(3.2,2.1);\coordinate(P)at(2,1.3125);\fill(P)circle(3pt)node[above left]{$m$};
\draw[<->](.15,-.25)--(2.15,1.0625)node[midway,below right]{$r$};\draw[bookaxis](P)--(2.8,1.8375)node[right]{$\vct e_r$};
\draw[bookaxis](P)--(1.475,2.1125)node[above]{$\vct e_\theta$};\draw[bookaxis](1,0)arc(0:65:1)node[above]{$\Omega$};
\node[below left]at(0,0){$O$};\node[below]at(2.7,-.55){水平面俯视图};
\end{tikzpicture}
\caption{题2.16：匀速转杆与沿杆滑动的小珠，固定端为 $O$。}\label{fig:ex-2-16}
\end{figure}

惯性系速度平方为 $\dot r^2+\Omega^2r^2$，所以动能相对初态的增量是
\[
W=\frac{m\Omega^2r_0^2}{2}(\sinh^2\Omega t+\cosh^2\Omega t-1)
=m\Omega^2r_0^2\sinh^2\Omega t.
\]
支持力功率为$N\Omega r=2m\Omega^2r\dot r$，从初态到当前状态的功也可直接积分：
\[W=2m\Omega^2\int_{r_0}^{r}s\dd s
=m\Omega^2[s^2]_{r_0}^{r}
=m\Omega^2(r^2-r_0^2),\]
代入$r=r_0\cosh\Omega t$，用$\cosh^2-1=\sinh^2$，与动能差相同。对动能求导并用$\ddot r=\Omega^2r$，有
\[\dot E=m\dot r\ddot r+m\Omega^2r\dot r
=2m\Omega^2r\dot r.\]
角动量$L_z=m\Omega r^2$给出$\dot L_z=2m\Omega r\dot r$，所以$\dot E=\Omega\dot L_z$。光滑条件使杆的作用力垂直于杆，惯性系速度却含横向分量，支持力仍可做功。

杆的横向速度为 $\Omega r$，与横向支持力同向，功率为 $N\Omega r$。这条图上的投影解释了约束力为何做功。
\end{solution}

\item \begin{minipage}[t]{\linewidth}
【提高】两个半径 $R$ 的同轴均匀带电圆环分别位于 $z=d,-d$，电荷为 $Q,-Q$，$Q,d>0$。写轴上电势与电场，求原点电场，并在 $z\gg R,d$ 时保留第一个非零远场项。若先把两环都近似为位于原点的点电荷，会丢失什么信息？

\end{minipage}
% 题设配图；与紧随其后的解答分别编号。
\ExpansionDiagram{fig:audit-m-q-068}{位置分开的正、负带电圆环}{
\draw[->](0,-1.8)--(0,2.4)node[above]{$z$};\draw[thick](0,1)ellipse(1.5 and .35);\draw[thick](0,-1)ellipse(1.5 and .35);\node[left]at(-1.5,1){$+Q$};\node[left]at(-1.5,-1){$-Q$};\node[right]at(0,1.5){$d$};\node[right]at(0,-.65){$-d$};\fill(0,0)circle(1.4pt)node[right]{$O$};
}

\begin{solution}
图\ref{fig:audit-m-q-068}中两圆环位于不同的$z$坐标，先叠加完整电势再作远场展开。
图\ref{fig:ex-2-17}中观察点到两环中心的轴向距离分别为 $z-d,z+d$，两电势按电荷符号叠加。

将两环到观察点的距离分别写为 $\sqrt{R^2+(z-d)^2}$ 与 $\sqrt{R^2+(z+d)^2}$，叠加电势得
\[
\Phi=k_eQ\left(\frac1{\sqrt{R^2+(z-d)^2}}-\frac1{\sqrt{R^2+(z+d)^2}}\right).
\]
对每个核的链式求导为
\[\frac{\dd}{\dd z}[R^2+(z\mp d)^2]^{-1/2}
=-\frac12[R^2+(z\mp d)^2]^{-3/2}\,2(z\mp d).\]
因此对叠加电势取负导数，
\[
E_z=k_eQ\left(\frac{z-d}{[R^2+(z-d)^2]^{3/2}}-
\frac{z+d}{[R^2+(z+d)^2]^{3/2}}\right).
\]
令 $z=0$，得 $E_z(0)=-2k_eQd/(R^2+d^2)^{3/2}$，方向从正环指向负环。

\begin{figure}[H]\centering
\begin{tikzpicture}[scale=.95]
\draw[bookaxis](-3,0)--(4.1,0)node[right]{$z$};
\foreach\x/\q/\z in{-1.5/-Q/-d,1.5/+Q/d}{\draw[bookcurve](\x,0)ellipse(.35 and 1.3);\node[above]at(\x,1.4){$\q$};\node[below]at(\x,-.15){$\z$};}
\draw(1.5,0)--(1.5,1.3)node[midway,right]{$R$};\fill(3.4,0)circle(2pt)node[above]{$P$};\node[below]at(0,0){$O$};
\end{tikzpicture}
\caption{题2.17：同轴等半径圆环，正电环位于 $z=d$，负电环位于 $z=-d$。}\label{fig:ex-2-17}
\end{figure}

对正的大$z$，先提取距离中的$z$，再使用
$(1+w)^{-1/2}=1-w/2+3w^2/8+O(w^3)$，
其中$w=\mp2d/z+(R^2+d^2)/z^2$。保留括号内至$z^{-2}$阶，得到
\begin{align*}
\frac1{\sqrt{R^2+(z\mp d)^2}}
&=\frac1z\left[1\mp\frac{2d}{z}+\frac{R^2+d^2}{z^2}\right]^{-1/2}\\
&=\frac1z\left[1\pm\frac d z+\frac{d^2-R^2/2}{z^2}+O(z^{-3})\right].
\end{align*}
两式相减时，$1/z$项和$1/z^3$项都消去，故
\[
\Phi=\frac{2k_eQd}{z^2}+O(z^{-4}),\qquad
E_z=\frac{4k_eQd}{z^3}+O(z^{-5}).
\]
这对应偶极矩$p_z=2Qd$。若先把两环都移到原点，总电荷为零的替代模型会把偶极矩也消去，因而丢失首个非零远场项。

在原点两环对正试探电荷的场都指向负 $z$；远处电荷总量抵消，图中电荷间距 $2d$ 则保留在偶极首项。
\end{solution}

\end{enumerate}

\chapter{线性代数与物理系统的结构}
一个物块位移以后，往往会同时改变另一个物块的受力。直接追踪每个坐标容易把方程混在一起；正常模则寻找能保持自身形状的运动。本章从线性组合和换基走到正常模，讨论质量不同、耦合较弱或对称性被破坏后的变化。

\section{线性组合、独立性与基}
\subsection{线性结构为什么反复出现在物理中}
在平衡点附近的小振动、线性电路、坐标变换以及微小扰动分析中，多个量之间常由线性关系连接。线性意味着叠加原理成立：若 $\vct x_1,\vct x_2$ 分别产生响应 $A\vct x_1,A\vct x_2$，则 $a\vct x_1+b\vct x_2$ 产生 $aA\vct x_1+bA\vct x_2$。因此可先分别求各分量的响应，再按原系数叠加。

向量 $\vct v_1,\ldots,\vct v_k$ 的线性组合为 $\sum c_j\vct v_j$。若存在不全为零的系数使这个组合为零，则这些向量线性相关；否则线性无关。相关表示某些方向可以由其他方向表示，并非所有向量都提供新信息。

能生成整个空间且线性无关的一组向量称为基。空间的每个向量在给定基下都有唯一分量，基向量的个数称为维数。唯一性可以直接证明：若同一向量有两种表示，相减得到基向量的一个零线性组合；线性无关迫使所有系数差为零。

\begin{example}
判断 $\vct v_1=(1,1),\vct v_2=(1,-1)$ 是否为平面的基，并表示 $(a,b)$。
\end{example}
% 题设配图；与紧随其后的解答分别编号。
\ExpansionDiagram{fig:audit-m-q-069}{给定的两个基候选向量与待表示向量}{
\draw[->](-.4,0)--(3.8,0)node[right]{$ x $};\draw[->](0,-2.3)--(0,2.6)node[above]{$ y $};\draw[->,thick](0,0)--(1.5,1.5)node[above]{$\vct v_1$};\draw[->,thick](0,0)--(1.5,-1.5)node[below]{$\vct v_2$};\draw[->](0,0)--(2.5,.8)node[right]{$(a,b)$};
}

\begin{solution}
图\ref{fig:audit-m-q-069}中的两基候选方向不平行，待表示向量须写成这两个方向的线性组合。
沿两根基矢分解目标向量，如图\ref{fig:eg-3-1}。比较直角坐标分量，即可求出分解系数。

检验是否成基，需要同时看能否生成平面以及表示是否唯一。设
\[
c_1(1,1)+c_2(1,-1)=(a,b),
\]
比较分量得 $c_1+c_2=a$、$c_1-c_2=b$。相加、相减分别得到
\[
c_1=\frac{a+b}{2},\qquad c_2=\frac{a-b}{2}.
\]
对任意 $a,b$ 都有且只有这一组系数，故两向量是平面的一组基。也可组成列矩阵 $P=\begin{pmatrix}1&1\\1&-1\end{pmatrix}$，其行列式为 $-2\ne0$，给出相同判断。

\begin{figure}[H]\centering
\begin{tikzpicture}[scale=.95]
\draw[bookaxis](-.4,0)--(4.1,0)node[right]{$x$};\draw[bookaxis](0,-1.5)--(0,2.8)node[above]{$y$};
\draw[bookaxis](0,0)--(1,1)node[above left]{$\vct v_1$};\draw[bookaxis](0,0)--(1,-1)node[below right]{$\vct v_2$};
\draw[bookaxis,thick](0,0)--(2,2)node[midway,above left]{$c_1\vct v_1$};
\draw[bookaxis,thick](2,2)--(3,1)node[midway,above right]{$c_2\vct v_2$};
\draw[bookaxis,thick](0,0)--(3,1)node[right]{$(a,b)$};\draw[bookguide](0,0)--(1,-1)--(3,1);
\end{tikzpicture}
\caption{例3.1：沿两个独立基方向分解目标向量。图取 $(a,b)=(3,1)$，对应 $c_1=2,c_2=1$；分量方程给出一般系数。}\label{fig:eg-3-1}
\end{figure}

再代回得 $((a+b+a-b)/2,(a+b-a+b)/2)=(a,b)$。这组基的两个方向相互正交，但模都为 $\sqrt2$；若需要正交单位基，还应各除以 $\sqrt2$。

两基方向独立，目标向量的分解系数唯一；代数计算对任意目标都给出了这一组系数。
\end{solution}

\section{矩阵、线性变换与方程组}
\subsection{矩阵的列记录基向量的像}
设 $A$ 为 $m\times n$ 矩阵，$\vct y=A\vct x$，则 $y_i=\sum_jA_{ij}x_j$。矩阵的第 $j$ 列，就是第 $j$ 个标准基向量经过线性变换后的像。按输入坐标对这些列作线性组合，就得到任意向量的像，也得到矩阵乘法的规则。

先做变换 $B$ 再做变换 $A$，合成矩阵为 $AB$，因为 $A(B\vct x)=(AB)\vct x$。一般 $AB\ne BA$，变换顺序有物理和几何意义。例如伸缩与旋转通常不能交换，先伸长某一轴再转向，与先转向再沿固定轴伸长不同。

单位矩阵 $I$ 不改变向量。转置 $A^{\mathsf T}$ 交换行与列，满足 $(AB)^{\mathsf T}=B^{\mathsf T}A^{\mathsf T}$。右侧顺序反转，因为原来内外层的映射关系发生对应变化；直接从分量求和也可验证。

\subsection{消元与秩}
线性方程组写成 $A\vct x=\vct b$。对增广矩阵 $[A\mid\vct b]$ 做交换行、非零倍乘一行、将一行的倍数加到另一行，均保持解集。这些操作将方程整理成主元与自由变量明确的阶梯形。

矩阵的秩是独立列的数目，也等于独立行的数目与消元后的主元数。对 $n$ 个未知数的系统，存在解当且仅当 $\rank A=\rank[A\mid\vct b]$；如果共同秩为 $n$，解唯一；如果小于 $n$，存在自由变量，解构成一个仿射集合。

齐次系统 $A\vct x=0$ 总有零解。若 $\rank A=r<n$，其解空间维数为 $n-r$。物理上，这些自由方向可能对应尚未施加的约束、对称性或自由运动，而不一定是错误。但增广矩阵出现 $0=c\ne0$ 时，模型内的方程确实互相矛盾，应检查建模假设和符号。

\begin{example}
解 $x+y+z=2,\ 2x+y+3z=5,\ x-y+z=0$。
\end{example}
% 题设配图；与紧随其后的解答分别编号。
\ExpansionDiagram{fig:audit-m-q-070}{三个平面共同约束同一个未知点}{
\begin{scope}[x={(1cm,0cm)},y={(-.45cm,.3cm)},z={(0cm,.8cm)}]\filldraw[fill=gray!10](0,0,2)--(2,0,0)--(2,1,-1)--(0,1,1)--cycle;\filldraw[fill=gray!25,opacity=.7](0,0,1.6667)--(2,0,.3333)--(2,1,0)--(0,1,1.3333)--cycle;\draw[thick](0,0,0)--(2,0,-2)--(2,1,-1)--(0,1,1)--cycle;\end{scope}\node[right,align=left]at(2.6,1){$x+y+z=2$\\$2x+y+3z=5$\\$x-y+z=0$};
}

\begin{solution}
图\ref{fig:audit-m-q-070}中三个平面的公共点满足全部方程。消去 $x$ 后，只须先在图\ref{fig:eg-3-2}的 $yz$ 平面内求两条直线的交点，再回代求 $x$。

将三式按顺序记为(1)(2)(3)。先用(1)消去另两式中的 $x$：
\[
\text{(2)}-2\text{(1)}:\quad-y+z=1,
\qquad
\text{(3)}-\text{(1)}:\quad-2y=-2.
\]
由后一式得到 $y=1$；代入前一式，得到 $z=2$。最后代回(1)，$x=2-y-z=-1$。故解为 $(-1,1,2)$。

\begin{figure}[H]\centering
\begin{tikzpicture}[scale=1.2]
\draw[bookaxis](-.3,0)--(2.8,0)node[right]{$y$};\draw[bookaxis](0,-.2)--(0,3.8)node[above]{$z$};
\draw[bookcurve](0,1)--(2.5,3.5)node[right]{$-y+z=1$};\draw[bookcurve](1,-.1)--(1,3.5)node[above]{$y=1$};
\fill(1,2)circle(1.8pt)node[right]{$(1,2)$};\draw[bookguide](0,2)--(1,2);\node[left]at(0,2){$2$};\node[below]at(1,0){$1$};
\end{tikzpicture}
\caption{例3.2：消元把三元问题化为 $yz$ 平面中两直线求交；唯一交点确定 $y,z$，再回代求 $x$。}\label{fig:eg-3-2}
\end{figure}

这些消元操作相当于把增广矩阵化为
\[
\left(\begin{array}{ccc|c}1&1&1&2\\0&-1&1&1\\0&-2&0&-2\end{array}\right).
\]
右下二阶块的行列式为 $(-1)\cdot0-1\cdot(-2)=2$，原系数矩阵可逆，所以解唯一。逐条核对：$-1+1+2=2$、$-2+1+6=5$、$-1-1+2=0$，三式均成立。

图中唯一交点为 $(y,z)=(1,2)$，与代数消元相同。若两约束平行或重合，唯一性判断便会改变。
\end{solution}

\section{行列式、逆与几何缩放}
\subsection{行列式的意义与运算性质}
二维矩阵的两列张成有向平行四边形，因此 $\det\begin{pmatrix}a&b\\c&d\end{pmatrix}=ad-bc$。三维行列式就是列向量的混合积。更高维行列式可由对各列的多线性、交换两列变号及 $\det I=1$ 定义。

行列式为零说明列向量相关，变换把空间压到较低维。非零行列式的绝对值为体积缩放因子，负号表示定向翻转。由线性映射的合成或多线性性质可得 $\det(AB)=\det A\det B$。交换两行变号，一行乘 $c$ 使行列式乘 $c$，一行加上另一行的倍数不改变行列式。

因此行列式可以用消元计算，但必须记录行交换与倍乘所造成的因子。只把矩阵消到对角线后相乘，却忘记此前交换行造成的符号，是常见错误。

\subsection{可逆性与求解}
方阵 $A$ 可逆，当且仅当它的秩满、行列式非零，也等价于 $A\vct x=\vct b$ 对任意右侧有唯一解。二维逆矩阵为
\begin{equation}
A^{-1}=\frac1{ad-bc}\begin{pmatrix}d&-b\\-c&a\end{pmatrix}.
\end{equation}
公式只在分母非零时成立。精确代数中的“可逆”不自动意味着数值计算容易，但行列式的绝对值也不能单独判断敏感程度。例如$A=cI$的行列式为$c^n$，取很小的$c\ne0$时它很小，各方向却只是同等缩放，相对误差没有因此被放大。需要比较矩阵在不同方向上的相对缩放，并同时考虑数据和参数的单位；第四章的灵敏度分析将具体讨论这种退化。

对具体方程组，消元求解通常比先求整个逆矩阵更直接。写 $\vct x=A^{-1}\vct b$ 有概念上的价值，但计算机实际求解时不必显式构造 $A^{-1}$。这一区分在第四章的数值线性代数中还会出现。

\section{正交、投影与换基}
\subsection{正交基如何简化分量}
若单位基两两正交，向量分量直接由点积 $x_j=\vct e_j\cdot\vct x$ 得到。若基不正交，就不能把这些点积直接当成坐标分量，因为各基方向之间会互相混入。

从独立向量构造正交基，可用 Gram--Schmidt 方法。先令 $\vct e_1=\vct a_1/|\vct a_1|$；对第二向量减去已有方向的投影，$\vct b_2=\vct a_2-(\vct a_2\cdot\vct e_1)\vct e_1$，再归一化。后续每一步都减去全部已构造方向的投影。独立性保证剩余向量非零。

\begin{example}
将 $(1,1,0)$ 与 $(1,0,1)$ 正交化。
\end{example}
% 题设配图；与紧随其后的解答分别编号。
\ExpansionDiagram{fig:audit-m-q-071}{把第二向量分成投影与垂直剩余量}{
\draw[->,thick](0,0)--(3,0)node[right]{$\vct a=(1,1,0)$};
\draw[->,thick](0,0)--(1.5,2.59808)node[above]{$\vct b=(1,0,1)$};
\draw[dashed](1.5,2.59808)--(1.5,0);
\fill(1.5,0)circle(1.5pt)node[below]{$P=\vct a/2$};
\draw(1.3,0)--(1.3,.2)--(1.5,.2);
\draw(.6,0)arc(0:60:.6);
\node at(.9,.35){$60^\circ$};
\node[below]at(1.5,-.8){在$\vct a,\vct b$张成的平面内作图};

}

\begin{solution}
图\ref{fig:audit-m-q-071}中$P$沿第一向量方向，第二向量减去这段投影后才与第一方向正交。
图\ref{fig:eg-3-3}把第二输入向量拆成已有方向上的投影与垂直余量。下面逐项算出图中的两段长度。

令 $\vct a_1=(1,1,0)$、$\vct a_2=(1,0,1)$。先取第一单位方向
\[
\vct e_1=\frac{\vct a_1}{|\vct a_1|}=\frac{(1,1,0)}{\sqrt2}.
\]
第二向量在此方向上的标量投影为 $\vct a_2\cdot\vct e_1=1/\sqrt2$，投影向量为 $(1/2,1/2,0)$。减去它，得到
\[
\vct b_2=\vct a_2-(\vct a_2\cdot\vct e_1)\vct e_1
=(1/2,-1/2,1),\quad |\vct b_2|=\sqrt{3/2}.
\]
再归一化，得到 $\vct e_2=(1,-1,2)/\sqrt6$。点积 $\vct e_1\cdot\vct e_2=(1-1)/\sqrt{12}=0$，两个模均为1。

\begin{figure}[H]\centering
\begin{tikzpicture}[scale=1.75]
\draw[bookaxis](-.25,0)--(2.3,0)node[right]{$\vct e_1$};\draw[bookaxis](0,-.25)--(0,1.8)node[above]{$\vct e_2$};
\draw[bookaxis,thick](0,0)--(1.414214,0)node[below]{$\vct a_1$};
\draw[bookaxis,thick](0,0)--(.707107,1.224745)node[above right]{$\vct a_2$};
\draw[bookguide](.707107,1.224745)--(.707107,0);\draw[bookaxis](0,0)--(0,1.224745)node[left]{$\vct b_2$};
\draw(.587107,0)--(.587107,.12)--(.707107,.12);\node[right]at(1.05,.75){$|\vct b_2|=\sqrt{3/2}$};
\node[below]at(.6,-.28){投影长度 $1/\sqrt2$};
\end{tikzpicture}
\caption{例3.3：在两输入向量张成的平面内，以构造出的 $\vct e_1,\vct e_2$ 为直角轴。去掉水平投影后得到竖直余量 $\vct b_2$。}\label{fig:eg-3-3}
\end{figure}

若题意只要求一组正交而不要求单位化的向量，可保留 $\vct a_1,\vct b_2$。归一化只改长度，不改已经得到的方向。

归一化只改变图中余量的长度，不改变其垂直方向。用这两根单位轴表示原向量，就得到 QR 分解的上三角系数。
\end{solution}

\subsection{正交矩阵}
将正交单位基作为矩阵 $Q$ 的列，有 $Q^{\mathsf T}Q=I$，所以 $Q^{-1}=Q^{\mathsf T}$。于是
\[
|Q\vct x|^2=\vct x^{\mathsf T}Q^{\mathsf T}Q\vct x=|\vct x|^2,
\]
点积也保持不变。正交变换保持几何长度与角度。行列式为 $+1$ 的情形保持定向，为 $-1$ 的情形反转定向；旋转与镜像都属于正交变换，但不是同一种操作。

\subsection{同一对象在不同基下的矩阵}
新基在旧坐标中的列组成矩阵 $P$。同一向量的旧、新分量满足 $\vct x=P\vct x'$。若线性变换旧矩阵为 $A$，则 $P\vct y'=AP\vct x'$，所以
\begin{equation}
A'=P^{-1}AP.
\end{equation}
这是相似变换。正交换基时变为 $Q^{\mathsf T}AQ$。矩阵中的数字改变了，物理映射并未改变，正如旋转坐标不会改变一个力本身。

\section{非正交坐标与度量矩阵}
\subsection{坐标分量不一定就是投影}
取平面基 $\vct b_1=(1,0),\vct b_2=(1,1)$。向量 $\vct r=c_1\vct b_1+c_2\vct b_2$ 的直角分量为 $(c_1+c_2,c_2)$，所以 $(x,y)$ 在这个基下的坐标是 $(x-y,y)$。但 $\vct r\cdot\vct b_1=x$，它不是一般的 $c_1$；非正交基下，点积投影与坐标系数必须分开。
\begin{figure}[htbp]\centering
\begin{tikzpicture}[scale=1.2]
\draw[bookaxis](0,0)--(2.8,0)node[right]{$x$};\draw[bookaxis](0,0)--(0,1.8)node[above]{$y$};
\draw[bookaxis,thick](0,0)--(1,0)node[below]{$\vct b_1$};\draw[bookaxis,thick](0,0)--(1,1)node[above left]{$\vct b_2$};
\draw[bookguide](1,0)--(2,1)--(1,1);\draw[bookaxis,thick](0,0)--(2,1)node[right]{$(2,1)$};
\end{tikzpicture}
\caption{非正交基下，向量 $(2,1)$ 的坐标为 $(1,1)$，但长度不是 $\sqrt2$。}
\label{fig:orig-059}
\end{figure}

将基向量作为 $P$ 的列，$\vct r=P\vct c$，则物理长度满足
\[
|\vct r|^2=\vct c^{\mathsf T}G\vct c,\qquad G=P^{\mathsf T}P.
\]
$G$ 称为这个基的度量矩阵。两向量内积为 $\vct c^{\mathsf T}G\vct d$。本例 $G=\begin{pmatrix}1&1\\1&2\end{pmatrix}$，坐标 $(1,1)$ 的长度平方为5。只有正交单位基才有 $G=I$。

\begin{example}
同一物理旋转在非正交基下得到矩阵 $A'$，它是否仍满足 $A'^{\mathsf T}A'=I$？
\end{example}
% 题设配图；与紧随其后的解答分别编号。
\ExpansionDiagram{fig:audit-m-q-072}{非正交基的平行四边形与物理旋转}{
\draw[->](-.4,0)--(4,0)node[right]{$ x $};\draw[->](0,-0.3)--(0,2.6)node[above]{$ y $};\draw[->,thick](0,0)--(2,0)node[above]{$\vct e_1$};\draw[->,thick](0,0)--(2,2)node[above]{$\vct e_2$};\draw[dashed](2,0)--(4,2)--(2,2);\draw[->](1,0)arc(0:90:1);\node[above left]at(.7,.7){物理旋转};
}

\begin{solution}
图\ref{fig:audit-m-q-072}中基矢的长度和夹角进入度量矩阵，物理旋转保持的是这一度量定义的长度。
图\ref{fig:eg-3-4}在原来的欧氏平面中画旋转，同时用非正交基分解旋转后的向量。

设旧基为正交单位基，新基的列组成 $P$，旋转的旧矩阵为 $R$。新坐标下 $A'=P^{-1}RP$，而物理长度的平方是 $\vct x'^{\mathsf T}G\vct x'$，其中 $G=P^{\mathsf T}P$。因此应检验
\begin{align*}
A'^{\mathsf T}GA'&=(P^{-1}RP)^{\mathsf T}P^{\mathsf T}P(P^{-1}RP)\\
&=P^{\mathsf T}R^{\mathsf T}RP=P^{\mathsf T}P=G.
\end{align*}
中间使用了 $(P^{-1})^{\mathsf T}P^{\mathsf T}=I$ 以及 $R^{\mathsf T}R=I$。

\begin{figure}[H]\centering
\begin{tikzpicture}[scale=1.5]
\draw[bookaxis](-1.2,0)--(2.7,0)node[right]{$x$};\draw[bookaxis](0,-.4)--(0,1.8)node[above]{$y$};
\draw[bookguide](0,0)circle(1);\draw[bookaxis](0,0)--(2,1)node[above right]{$\vct b_2=(2,1)$};
\draw[bookaxis,thick](0,0)--(1,0)node[below]{$\vct v$};\draw[bookaxis,thick](0,0)--(0,1)node[above left]{$R\vct v$};
\draw[bookaxis,dashed](2,1)--(0,1)node[midway,above]{$-2\vct b_1$};
\draw[->](.5,0)arc(0:90:.5);\node[right]at(.6,.45){$90^\circ$};
\end{tikzpicture}
\caption{例3.4：同一个单位物理向量旋转前后都在单位圆上；旋转后新坐标为 $(-2,1)$，因为 $R\vct v=-2\vct b_1+\vct b_2$。}\label{fig:eg-3-4}
\end{figure}

举例取 $P=\begin{pmatrix}1&2\\0&1\end{pmatrix}$，旋转为逆时针 $90^\circ$，则
\[
G=\begin{pmatrix}1&2\\2&5\end{pmatrix},\qquad
A'=\begin{pmatrix}-2&-5\\1&2\end{pmatrix}.
\]
向量新坐标 $(1,0)$ 变成 $(-2,1)$，普通坐标平方和由1变成5，但按 $G$ 计算仍为 $4-8+5=1$。所以一般不能要求 $A'^{\mathsf T}A'=I$；保持的是由 $G$ 表示的物理内积。

旋转前后两向量的模均为1。新坐标的两基矢不正交，长度平方中须保留它们的交叉内积。
\end{solution}

\subsection{Gram--Schmidt 的分解意义}
对满列秩矩阵 $A=(\vct a_1,\ldots,\vct a_n)$，逐列正交化得到单位向量组成的 $Q$。每个原向量由已经构造的前几列组合而成，因此 $A=QR$，$R$ 为上三角矩阵，且 $R_{ij}=\vct q_i\cdot\vct a_j$。这就是 QR 分解的一种构造。

对最小二乘$A\vct x\approx\vct b$，先把数据分解为列空间内、外两部分：
\[
\vct b=QQ^{\mathsf T}\vct b+\vct b_\perp,
\qquad Q^{\mathsf T}\vct b_\perp=0.
\]
代入$A=QR$，残差为$Q(R\vct x-Q^{\mathsf T}\vct b)-\vct b_\perp$。两项正交，且$Q^{\mathsf T}Q=I$，所以
\[
|A\vct x-\vct b|^2
=|R\vct x-Q^{\mathsf T}\vct b|^2+|\vct b_\perp|^2.
\]
第二项不随参数改变，满列秩又保证上三角矩阵$R$可逆。因此唯一的最小二乘解由$R\vct x=Q^{\mathsf T}\vct b$确定，最小残差的模为$|\vct b_\perp|$。实际求解时，对$R$作回代即可。

\section{特征值、特征向量与对角化}
特征方向在线性变换后只改变尺度，在线性动力学中对应保持形状的运动。这些方向能否组成一组基，决定能否用它们解耦方程。

\subsection{寻找不会被混合的方向}
若 $A\vct u=\lambda\vct u$ 且 $\vct u\ne0$，则 $\vct u$ 为特征向量，$\lambda$ 为特征值。变换在该方向上只伸缩，不混入其他方向。存在非零解要求 $\det(A-\lambda I)=0$。

先解特征多项式，再对每个特征值解齐次系统。这两步缺一不可。一个特征值的代数重数是它在多项式中的重复次数，几何重数是独立特征向量的数目，二者不能混为一谈。

\begin{example}
求 $A=\begin{pmatrix}2&-1\\-1&2\end{pmatrix}$ 的特征结构。
\end{example}
% 题设配图；与紧随其后的解答分别编号。
\ExpansionDiagram{fig:audit-m-q-073}{矩阵的两列给出两坐标基矢的像}{
\draw[->](-.4,0)--(3,0)node[right]{$ x $};\draw[->](0,-1.8)--(0,3)node[above]{$ y $};\draw[->,thick](0,0)--(2,-1)node[right]{$A\vct e_1$};\draw[->,thick](0,0)--(-1,2)node[left]{$A\vct e_2$};
}

\begin{solution}
图\ref{fig:audit-m-q-073}中的两箭头是给定矩阵作用于坐标基的结果，特征方向还需由平行条件求出。
图\ref{fig:orig-062}把特征方向与单位圆的像画在一起；先由特征方程求出两条轴及各自伸缩率。

先求特征多项式：
\[
\det(A-\lambda I)=(2-\lambda)^2-1
=(\lambda-1)(\lambda-3).
\]
对 $\lambda_1=1$，方程 $(A-I)\vct u=0$ 给出 $u_x=u_y$；归一化可取 $\vct u_1=(1,1)^{\mathsf T}/\sqrt2$。对 $\lambda_2=3$，方程给出 $u_y=-u_x$，可取 $\vct u_2=(1,-1)^{\mathsf T}/\sqrt2$。

\begin{figure}[H]\centering
\begin{tikzpicture}[scale=.85]
\draw[bookaxis](-2.6,0)--(2.8,0)node[right]{$x$};\draw[bookaxis](0,-2.6)--(0,2.8)node[above]{$y$};
\draw[dashed](0,0)circle(1);\draw[thick,rotate=-45](0,0)ellipse(3 and 1);
\draw[bookaxis](0,0)--(45:1.35)node[above right]{$\vct u_1$};\draw[bookaxis](0,0)--(-45:2.9)node[below right]{$\vct u_2$};
\end{tikzpicture}
\caption{对称变换把单位圆变成椭圆；两条主轴是相互正交的特征方向。}
\label{fig:orig-062}
\end{figure}

将它们作为列组成
\[
Q=\frac1{\sqrt2}\begin{pmatrix}1&1\\1&-1\end{pmatrix},
\quad Q^{\mathsf T}Q=I,
\quad Q^{\mathsf T}AQ=\begin{pmatrix}1&0\\0&3\end{pmatrix}.
\]
因此沿同相方向的向量保持长度，沿反相方向伸长3倍。迹为4、行列式为3，分别等于两特征值之和与积。这里 $A$ 为实对称矩阵，才可进一步选到完整的正交特征基。

像椭圆在同相轴的半轴长为1，在反相轴的半轴长为3。这是两个特征值的几何核对。
\end{solution}

如果有 $n$ 个独立特征向量，将它们作为 $P$ 的列，则 $AP=P\Lambda$，因此 $P^{-1}AP=\Lambda$ 为对角矩阵。反过来，可对角化必须存在完整特征基。矩阵 $\begin{pmatrix}1&1\\0&1\end{pmatrix}$ 虽有二重特征值1，却只有一个独立特征方向，不能对角化；单位矩阵虽然同样有重根，却可以用任何基对角化。

\subsection{实对称矩阵的特殊性质}
若 $A=A^{\mathsf T}$，不同特征值的特征向量正交。证明为
\[
\lambda\vct u^{\mathsf T}\vct v=(A\vct u)^{\mathsf T}\vct v
=\vct u^{\mathsf T}A\vct v=\mu\vct u^{\mathsf T}\vct v.
\]
当 $\lambda\ne\mu$ 时点积必为零。实对称矩阵还能保证所有特征值为实数，并存在完整正交特征基。

可以从二次型在单位球面的极值理解后一个结论。实连续函数$\vct x^{\mathsf T}A\vct x$在紧的单位球面上取到最大值。引入约束$\sum_i x_i^2=1$，对每个实分量求导：
\[
\partial_{x_i}(\vct x^{\mathsf T}A\vct x)
=\sum_jA_{ij}x_j+\sum_jx_jA_{ji}
=2(A\vct x)_i=2\lambda x_i.
\]
于是极值点$\vct u$满足$A\vct u=\lambda\vct u$，给出一个实单位特征向量。它的正交补空间在$A$下保持不变，因为对$\vct v\perp\vct u$，有$\vct u^{\mathsf T}A\vct v=(A\vct u)^{\mathsf T}\vct v=0$。在这个低一维空间重复同一过程，递归到一维，就得到$n$个两两正交的特征向量。

\section{线性变换的几何与特征方向}
\subsection{面积缩放与剪切}
矩阵 $A=\begin{pmatrix}1&s\\0&1\end{pmatrix}$ 将 $(x,y)$ 变为 $(x+sy,y)$。水平位置随高度改变，称为剪切。行列式为1，面积保持不变，但角度和长度一般改变。面积保持并不等于正交，也不等于每个方向都没有变形。
\begin{figure}[htbp]\centering
\begin{tikzpicture}[scale=1.6]
\draw[thick](0,0)rectangle(1,1);\draw[bookaxis](-.1,0)--(1.35,0);\draw[bookaxis](0,-.1)--(0,1.3);
\draw[bookaxis](1.65,.5)--(2.2,.5)node[midway,above]{$A$};
\begin{scope}[shift={(2.6,0)}]\draw[thick](0,0)--(1,0)--(1.7,1)--(.7,1)--cycle;\draw[bookaxis](-.1,0)--(2,0);\draw[bookaxis](0,-.1)--(0,1.3);\end{scope}
\end{tikzpicture}
\caption{单位正方形经剪切成为同面积的平行四边形。图取 $s=0.7$。}
\label{fig:orig-061}
\end{figure}

实对称矩阵可以选择正交特征轴。二维非奇异情形下，它把单位圆变成主轴沿特征方向的椭圆，半轴长为特征值的绝对值；负特征值还反转相应轴向，零特征值则使图形退化。对$\begin{pmatrix}2&-1\\-1&2\end{pmatrix}$，同相方向伸长$1$倍，反相方向伸长$3$倍。轴长与轴方向分别由特征值和特征向量确定。

\subsection{等势线与稳定性图像}
正定二次型的等值线为闭合椭圆，势能在所有方向离开原点时升高。不定二次型的等值线为双曲线，至少一个方向上势能降低。观察等值线可帮助理解矩阵特征值符号，但边际零方向仍需高阶项分析。
\begin{figure}[htbp]\centering
\begin{tikzpicture}[scale=.9]
\begin{scope}\foreach\r in{.5,.9,1.3}{\draw(0,0)circle(\r);}\draw[bookaxis](-1.6,0)--(1.6,0);\draw[bookaxis](0,-1.6)--(0,1.6);\node at(0,-2){$U=x^2+y^2$};\end{scope}
\begin{scope}[shift={(4.6,0)}]
\foreach\c in{.4,.8}{\draw[domain=-1.15:1.15,samples=60]plot({sqrt(\x*\x+\c)},\x);\draw[domain=-1.15:1.15,samples=60]plot({-sqrt(\x*\x+\c)},\x);\draw[dashed,domain=-1.15:1.15,samples=60]plot(\x,{sqrt(\x*\x+\c)});\draw[dashed,domain=-1.15:1.15,samples=60]plot(\x,{-sqrt(\x*\x+\c)});}
\draw[bookaxis](-1.7,0)--(1.7,0);\draw[bookaxis](0,-1.7)--(0,1.7);\node at(0,-2){$U=x^2-y^2$};\end{scope}
\end{tikzpicture}
\caption{正定势能与鞍形势能的等值线。右图实线为正等值，虚线为负等值。}
\label{fig:orig-063}
\end{figure}

\section{二次型与平衡点稳定性}
二次型描述平衡点附近不同方向的势能变化。正负性只在允许的独立位移中判断，出现零方向时再检查更高阶势能。

\subsection{从泰勒展开得到刚度矩阵}
设平衡点取为 $\vct x=0$，势能足够光滑且一阶导数为零。二阶展开为
\begin{equation}
U(\vct x)=U_0+\tfrac12\vct x^{\mathsf T}K\vct x+O(|\vct x|^3),
\end{equation}
其中 $K_{ij}=\partial^2U/\partial x_i\partial x_j$ 是 Hessian 矩阵；在连续二阶偏导条件下它是对称的。线性恢复力为 $-K\vct x$。

用正交矩阵对角化 $K$，二次型变为 $\tfrac12\sum\kappa_jq_j^2$。所有 $\kappa_j>0$ 时称正定，每个方向都有上升势能；存在负特征值时沿相应方向势能降低，平衡点在保守正质量小振动模型中线性不稳定。

零特征值不能单独决定完整稳定性。例如 $U=x^4+y^2$ 在原点仍是局部极小，而 $U=-x^4+y^2$ 沿 $x$ 方向下降；二者的二阶矩阵相同。二阶近似看不见的方向，必须进一步检查高阶项、约束或物理对称性。

\subsection{二维正定条件与配方}
对 $Q=ax^2+2bxy+cy^2$，若 $a>0$，可配方为
\begin{equation}
Q=a\left(x+\frac ba y\right)^2+\left(c-\frac{b^2}{a}\right)y^2.
\end{equation}
因此二维对称矩阵正定当且仅当 $a>0,ac-b^2>0$。这比只检查两个对角元素为正严格得多，因为交叉项可能导致某个斜方向的二次型为负。

\begin{example}
对 $U=k(x^2+2\eta xy+y^2)/2$，$k>0$，判定二阶稳定性。
\end{example}
% 题设配图；与紧随其后的解答分别编号。
\ExpansionDiagram{fig:audit-m-q-074}{二次势能沿两个独立对角方向的变化}{
\draw[->](-.4,0)--(3.3,0)node[right]{$ x $};\draw[->](0,-2.2)--(0,2.4)node[above]{$ y $};\draw[dashed](-1.5,-1.5)--(1.5,1.5)node[above]{$x=y$};\draw[dotted](-1.5,1.5)--(1.5,-1.5)node[below]{$x=-y$};\node[right]at(3.5,.9){$U=\frac{k}{2}(x^2+2\eta xy+y^2)$};
}

\begin{solution}
图\ref{fig:audit-m-q-074}中的两个对角方向可分别检验势能曲率；正定要求所有非零方向都具有正的二次能量。
将坐标转到同相、反相方向后，可用图\ref{fig:eg-3-6}中的闭合等势线或鞍形等势线理解系数的正负。

先用正交坐标 $q_+=(x+y)/\sqrt2$、$q_-=(x-y)/\sqrt2$ 消去交叉项。逆变换代入给出
\[
x^2+y^2=q_+^2+q_-^2,\qquad 2xy=q_+^2-q_-^2,
\]
故
\[
U=\frac k2[(1+\eta)q_+^2+(1-\eta)q_-^2].
\]
当 $|\eta|<1$，两系数为正，原点是严格势能极小点；配合正定动能时，原点稳定。当 $|\eta|>1$，至少有一系数为负，沿该方向有任意小的位移使势能下降，原点不稳定。

\begin{figure}[H]\centering
\begin{tikzpicture}[scale=.9]
\begin{scope}\foreach\r in{.5,.9,1.3}{\draw(0,0)circle(\r);}\draw[bookaxis](-1.6,0)--(1.6,0);\draw[bookaxis](0,-1.6)--(0,1.6);\node at(0,-2){$U=x^2+y^2$};\end{scope}
\begin{scope}[shift={(4.6,0)}]
\foreach\c in{.4,.8}{\draw[domain=-1.15:1.15,samples=60]plot({sqrt(\x*\x+\c)},\x);\draw[domain=-1.15:1.15,samples=60]plot({-sqrt(\x*\x+\c)},\x);\draw[dashed,domain=-1.15:1.15,samples=60]plot(\x,{sqrt(\x*\x+\c)});\draw[dashed,domain=-1.15:1.15,samples=60]plot(\x,{-sqrt(\x*\x+\c)});}
\draw[bookaxis](-1.7,0)--(1.7,0);\draw[bookaxis](0,-1.7)--(0,1.7);\node at(0,-2){$U=x^2-y^2$};\end{scope}
\end{tikzpicture}
\caption{例3.6解图：化到特征方向并缩放坐标后，正定情形为闭合等势线，不定情形为双曲线；零方向需检查高阶势能。}
\label{fig:eg-3-6}
\end{figure}

当$\eta=1$，$q_-$方向二次项为零；当$\eta=-1$，$q_+$方向为零。若题式就是完整势能，沿零方向有自由运动$q_0(t)=q_0(0)+\dot q_0(0)t$。因此只扰动位置而保持静止时会停在另一平衡位置；若允许任意小的初速度扰动，则会逐渐离开原点，原点作为一个静止状态不具有李雅普诺夫稳定性。

若题式只是势能的局部二阶展开，零方向的恢复作用还取决于高阶项，例如正四次项与负四次项给出不同结论。此时二阶判据不足，须保留高阶势能再判断。

两个系数都正时等势线闭合；一正一负时出现降低势能的方向。恰有零系数时，图像沿该方向不再闭合，需另作分析。
\end{solution}

\section{两个耦合振子的正常模}
\subsection{从受力得到矩阵}
两质量 $m$ 的物体由三根刚度 $k$ 的轻弹簧接在固定端之间。位移 $x_1,x_2$ 相对平衡位置测量，方向统一向右。中间弹簧形变量为 $x_2-x_1$，因此
\begin{align}
m\ddot x_1&=-kx_1+k(x_2-x_1),\\
m\ddot x_2&=-kx_2-k(x_2-x_1).
\end{align}
写成矩阵即 $m\ddot{\vct x}+K\vct x=0$，$K=k\begin{pmatrix}2&-1\\-1&2\end{pmatrix}$。负的非对角元素表示一个物体的位移会改变另一个物体的受力，它不是负刚度弹簧。

\begin{figure}[htbp]
\centering
\begin{tikzpicture}[scale=.83]
\draw[thick](0,-.5)--(0,.5) (8,-.5)--(8,.5);
\draw[decorate,decoration={coil,segment length=4pt,amplitude=3pt}](0,0)--(2,0);
\draw(2,-.3)rectangle(3,.3);\node at(2.5,0){$m$};
\draw[decorate,decoration={coil,segment length=4pt,amplitude=3pt}](3,0)--(5,0);
\draw(5,-.3)rectangle(6,.3);\node at(5.5,0){$m$};
\draw[decorate,decoration={coil,segment length=4pt,amplitude=3pt}](6,0)--(8,0);
\foreach\x in{1,4,7}{\node at(\x,.55){$k$};}
\draw[-{Stealth}](2.2,-.7)--(3.2,-.7)node[right]{$x_1$};
\draw[-{Stealth}](5.2,-.7)--(6.2,-.7)node[right]{$x_2$};
\end{tikzpicture}
\caption{固定端双振子的弹簧连接}\label{fig:two-masses}
\end{figure}

设正常模为 $\vct x=\vct u\cos(\omega t+\varphi)$，代入得到 $K\vct u=m\omega^2\vct u$。正常模因此是刚度矩阵的特征方向，其特征值除以质量才是角频率平方。由此前计算，
\begin{equation}
\omega_1=\sqrt{k/m},\qquad\omega_2=\sqrt{3k/m}.
\end{equation}
低频模的两物体同相，中间弹簧不变形；高频模两物体反相，中间弹簧参与恢复作用，回复力较大，因而频率较高。

\textbf{算例：对称双振子的同相运动。} 两个质量均为$m$的物块由外侧弹簧接固定端，中间另有弹簧，各弹簧刚度均为 $k$。若二者始终同向等位移，中间弹簧是否伸长？求该运动的角频率。
% 题设配图；与紧随其后的解答分别编号。
\ExpansionDiagram{fig:audit-m-q-075}{固定端双振子的连接及给定初态}{
\draw[thick](0,-.5)--(0,.8);\draw[exspring](0,.15)--(1.4,.15);\node[above]at(.7,.4){$ k $};\draw[thick](6.7,-.5)--(6.7,.8);\draw[exspring](5.3,.15)--(6.7,.15);\node[above]at(6,.4){$ k $};\draw(1.4,-.3)rectangle(2.4,.6);\node at(1.9,.15){$ m $};\draw(4.3,-.3)rectangle(5.3,.6);\node at(4.8,.15){$ m $};\draw[exspring](2.4,.15)--(4.3,.15);\node[above]at(3.35,.4){$ k $};\draw[->](1.4,-.8)--(2.7,-.8)node[right]{$x_1$};\draw[->](4.3,-.8)--(5.6,-.8)node[right]{$x_2$};\node[above]at(3.35,1.2){$x_1(0)=x_2(0)=a$};
}

\begin{solution}
按题设图\ref{fig:audit-m-q-075}的位移正方向，令两物块均从平衡位置向右移动$x$，其初态与解图\ref{fig:warm-3}一致。
中间弹簧的伸长为$x_2-x_1=x-x=0$，两端弹簧对各自物块的力均为$-kx$。
把这组位移代入完整方程，
\[
 m\ddot x_1=-kx_1+k(x_2-x_1),\qquad
 m\ddot x_2=-kx_2-k(x_2-x_1),
\]
两式都化为$m\ddot x+kx=0$，所以等位移的条件能够保持。令$x\propto\ee^{\lambda t}$，特征方程为$m\lambda^2+k=0$，两根$\lambda=\pm\ii\sqrt{k/m}$给出
角频率$\omega=\sqrt{k/m}$。
也可把两块视为这一运动方向上的总质量$2m$：
两端弹簧的总势能为$kx^2=\tfrac12(2k)x^2$，仍得到$\omega^2=2k/(2m)=k/m$。

\begin{figure}[H]\centering
\begin{tikzpicture}[scale=.9]
\draw[thick](0,-.5)--(0,.7) (7.5,-.5)--(7.5,.7);\draw[exspring](0,0)--(2,0);\draw[exspring](3,0)--(4.5,0);\draw[exspring](5.5,0)--(7.5,0);
\draw(2,-.4)rectangle(3,.4);\draw(4.5,-.4)rectangle(5.5,.4);\node at(2.5,0){$m$};\node at(5,0){$m$};
\draw[bookaxis](2,.9)--(3.1,.9)node[midway,above]{$a$};\draw[bookaxis](4.5,.9)--(5.6,.9)node[midway,above]{$a$};
\node[below]at(3.75,-.75){中间弹簧伸长 $x_2-x_1=0$};
\end{tikzpicture}
\caption{对称双振子的同相运动：两物块同向等位移，中间弹簧不变形，各物块只受外侧弹簧回复力。}\label{fig:warm-3}
\end{figure}
\end{solution}

\subsection{正常坐标与任意初值}
取$q_1=(x_1+x_2)/\sqrt2,q_2=(x_1-x_2)/\sqrt2$。分别把两条受力方程相加、相减，再除以$\sqrt2$，中间弹簧项在和式中抵消，在差式中相加：
\[
m\ddot q_1=-kq_1,\qquad m\ddot q_2=-3kq_2.
\]
于是方程变为$\ddot q_j+\omega_j^2q_j=0$。其特征根为$\pm\ii\omega_j$，通解为$q_j=A_j\cos\omega_jt+B_j\sin\omega_jt$。将任意初值变换为正常坐标后，$A_j=q_j(0)$、$\omega_jB_j=\dot q_j(0)$，因此
\begin{equation}
q_j(t)=q_j(0)\cos\omega_jt+\frac{\dot q_j(0)}{\omega_j}\sin\omega_jt,
\end{equation}
再用 $x_1=(q_1+q_2)/\sqrt2,x_2=(q_1-q_2)/\sqrt2$ 还原。方程分离以后，初始位移和速度也要变换到正常坐标。

\begin{example}
只将第一个物体向右拉开 $a$，第二个在平衡点，二者由静止释放。求运动与总能量。
\end{example}
% 题设配图；与紧随其后的解答分别编号。
\ExpansionDiagram{fig:audit-m-q-076}{固定端双振子的连接及给定初态}{
\draw[thick](0,-.5)--(0,.8);\draw[exspring](0,.15)--(1.4,.15);\node[above]at(.7,.4){$ k $};\draw[thick](6.7,-.5)--(6.7,.8);\draw[exspring](5.3,.15)--(6.7,.15);\node[above]at(6,.4){$ k $};\draw(1.4,-.3)rectangle(2.4,.6);\node at(1.9,.15){$ m $};\draw(4.3,-.3)rectangle(5.3,.6);\node at(4.8,.15){$ m $};\draw[exspring](2.4,.15)--(4.3,.15);\node[above]at(3.35,.4){$ k $};\draw[->](1.4,-.8)--(2.7,-.8)node[right]{$x_1$};\draw[->](4.3,-.8)--(5.6,-.8)node[right]{$x_2$};\node[above]at(3.35,1.2){$x_1(0)=a,\quad x_2(0)=0$};
}

\begin{solution}
图\ref{fig:audit-m-q-076}所给初态中，只有第一物块偏离平衡位置，需将其分解到两个正常模上。连接和位移正方向沿用图\ref{fig:two-masses}，所得运动曲线见图\ref{fig:eg-3-7}。

三弹簧的连接已给出 $\omega_1=\sqrt{k/m}$、$\omega_2=\sqrt{3k/m}$。取正交正常坐标
\[
q_1=\frac{x_1+x_2}{\sqrt2},\qquad
q_2=\frac{x_1-x_2}{\sqrt2}.
\]
每个正常坐标满足 $\ddot q_j+\omega_j^2q_j=0$。其特征根为 $\lambda=\pm\ii\omega_j$，实通解可写成
\[
q_j=C_j\cos\omega_jt+D_j\sin\omega_jt,
\qquad
\dot q_j=-\omega_j C_j\sin\omega_jt+\omega_j D_j\cos\omega_jt.
\]
初值 $x_1=a,x_2=0$、两速度为零，给出 $q_j(0)=a/\sqrt2$、$\dot q_j(0)=0$。于是 $C_j=a/\sqrt2$、$\omega_jD_j=0$；两频率均非零，故 $D_j=0$，从而
\[
q_1=\frac a{\sqrt2}\cos\omega_1t,\qquad
q_2=\frac a{\sqrt2}\cos\omega_2t.
\]
逆变换得到
\[
x_1=\frac a2(\cos\omega_1t+\cos\omega_2t),\qquad
x_2=\frac a2(\cos\omega_1t-\cos\omega_2t).
\]
初态第一根与中间弹簧的形变量分别为 $a,-a$，第三根为零，所以总能量为 $ka^2$。也可在正常坐标中相加：$E_1=ka^2/4$、$E_2=3ka^2/4$，和仍为 $ka^2$。模态坐标振幅相同，不表示两模能量相同。

\begin{figure}[H]\centering
\begin{tikzpicture}[x=.7cm,y=1.6cm]
\draw[bookaxis](0,0)--(10.8,0)node[right]{$\tau=t\sqrt{k/m}$};\draw[bookaxis](0,-1.2)--(0,1.3)node[above]{$x_j/a$};
\draw[bookcurve,domain=0:10,samples=200]plot(\x,{(cos(\x r)+cos(sqrt(3)*\x r))/2});
\draw[dashed,domain=0:10,samples=200]plot(\x,{(cos(\x r)-cos(sqrt(3)*\x r))/2});
\fill(0,1)circle(1.5pt)node[left]{$1$};\foreach\t in{2,4,6,8,10}{\node[below]at(\t,0){$\t$};}
\draw[thick](1.2,1.2)--(1.8,1.2)node[right]{$x_1/a$};\draw[dashed](5.3,1.2)--(5.9,1.2)node[right]{$x_2/a$};
\end{tikzpicture}
\caption{例3.7：两个正常模叠加后得到两物块的完整运动。曲线在原点时刻分别从1和0出发，且初始切线均水平。}\label{fig:eg-3-7}
\end{figure}

两条曲线共同包含两种频率；它们的和只保留同相模，差只保留反相模，可以从图文两种表示核对逆变换。
\end{solution}

\section{质量不同、零频模与质量归一化}
\subsection{一般质量矩阵}
若两质量不同，或左右弹簧刚度不同，写 $M\ddot{\vct x}+K\vct x=0$，其中
\[
M=\diag(m_1,m_2),\qquad
K=\begin{pmatrix}k_1+k_c&-k_c\\-k_c&k_2+k_c\end{pmatrix}.
\]
正常模满足广义特征方程 $K\vct u=\omega^2M\vct u$。频率由
\begin{equation}
(k_1+k_c-m_1\omega^2)(k_2+k_c-m_2\omega^2)-k_c^2=0
\end{equation}
求得。这里不能直接对 $K$ 求特征值再开根号，质量已经进入每个方向的惯性响应。

先看动能，便能知道为什么需要质量归一化。对角质量矩阵下，
\[
T=\tfrac12m_1\dot x_1^2+\tfrac12m_2\dot x_2^2.
\]
同样大小的坐标速度具有不同的动能权重。若令 $z_j=\sqrt{m_j}x_j$，恰有 $T=(\dot z_1^2+\dot z_2^2)/2$，各方向的动能系数变为相同。改用 $m_jx_j$ 则仍留下 $1/m_j$，达不到这个目的。

对一般实对称正定质量矩阵，可先将它正交对角化，再对正特征值取平方根，由此定义 $M^{1/2}$。写 $\vct z=M^{1/2}\vct x$，将 $\vct x=M^{-1/2}\vct z$ 代回方程，再左乘 $M^{-1/2}$，得
\[
\ddot{\vct z}+B\vct z=0,\qquad B=M^{-1/2}KM^{-1/2}.
\]
由于 $B$ 实对称，存在正交矩阵 $U$，满足 $U^{\mathsf T}BU=\Lambda$。令 $\vct z=U\vct q$，则各正常坐标满足 $\ddot q_j+\omega_j^2q_j=0$。把两次变换合并，$\vct x=S\vct q$，其中 $S=M^{-1/2}U$。

逆变换必须按原顺序倒回来：先乘 $M^{1/2}$，再乘 $U^{\mathsf T}$。因此
\[
\vct q=U^{\mathsf T}M^{1/2}\vct x=S^{\mathsf T}M\vct x,
\quad S^{\mathsf T}MS=I,\quad S^{\mathsf T}KS=\Lambda.
\]
速度也按同样方式变换。此时 $T=\sum_j\dot q_j^2/2$、$U_{\rm pot}=\sum_j\omega_j^2q_j^2/2$，方程与能量同时分离。检查这两个二次型，比只核对频率更容易发现漏掉质量权重的错误。

当 $M=\diag(m,2m)$ 时，
\[
M^{1/2}=\diag(\sqrt m,\sqrt{2m}),\qquad
M^{-1/2}=\diag(1/\sqrt m,1/\sqrt{2m}).
\]
若两个模态列向量为 $\vct u_\pm=(1,r_\pm)^{\mathsf T}$，质量加权正交条件给出
\[
\vct u_-^{\mathsf T}M\vct u_+=0\quad\Longrightarrow\quad1+2r_-r_+=0.
\]
普通点积中的系数2缺失，正是它不适合直接判断这两个模态正交性的原因。

\subsection{自由双物体：整体平移与相对振动}
两物体仅由一根弹簧 $k$ 相连，不接固定端，刚度矩阵为 $k\begin{pmatrix}1&-1\\-1&1\end{pmatrix}$。整体同向平移不会改变弹簧长度，因此有零频模。

取质心坐标 $X=(m_1x_1+m_2x_2)/(m_1+m_2)$、相对位移 $\xi=x_2-x_1$，直接相加与相减牛顿方程得
\begin{equation}
\ddot X=0,\qquad \ddot\xi+k\left(\frac1{m_1}+\frac1{m_2}\right)\xi=0.
\end{equation}
令约化质量$\mu_r=m_1m_2/(m_1+m_2)$，相对振动角频率为$\sqrt{k/\mu_r}$。动能的分离可从逆变换核对：
\[
x_1=X-\frac{m_2}{m_1+m_2}\xi,\qquad
x_2=X+\frac{m_1}{m_1+m_2}\xi.
\]
对两式求导，再代入$m_1\dot x_1^2/2+m_2\dot x_2^2/2$。交叉项的系数为
$-m_1m_2/(m_1+m_2)+m_2m_1/(m_1+m_2)=0$，
$\dot\xi^2$的系数则为
\[
\frac{m_1m_2^2+m_2m_1^2}{2(m_1+m_2)^2}
=\frac{\mu_r}{2}.
\]
所以$T=(m_1+m_2)\dot X^2/2+\mu_r\dot\xi^2/2$。零频方向的解为常量加上线性时间项，表示整体匀速平移，单独一个常量只能表示其中静止的特例。

\section{不同质量双振子的完整求解}
仍取三根相同弹簧 $k$，但质量为 $m,2m$。令 $\lambda=m\omega^2/k$，广义特征方程变为
\[
\det\begin{pmatrix}2-\lambda&-1\\-1&2-2\lambda\end{pmatrix}=0,
\quad 2\lambda^2-6\lambda+3=0.
\]
两个根为 $\lambda_\pm=(3\pm\sqrt3)/2$，频率为 $\omega_\pm=\sqrt{(k/m)\lambda_\pm}$。以第一分量为1，特征形状为 $\vct u_\pm=(1,r_\pm)^{\mathsf T}$，$r_\pm=2-\lambda_\pm=(1\mp\sqrt3)/2$。

低频模 $r_->0$，两物体同向但幅值不同；高频模 $r_+<0$，二者反向。质量较大的物体并非在所有模中都简单地幅值减半。运动形状由惯性和全部弹簧的联立方程决定。

普通点积 $\vct u_-\cdot\vct u_+$ 不为零，但质量内积为
\[
\vct u_-^{\mathsf T}M\vct u_+=m(1+2r_-r_+)=0,
\]
因为 $r_-r_+=-1/2$。将各列除以 $\sqrt{m(1+2r_\pm^2)}$，得到质量归一化矩阵 $S$。令 $\vct x=S\vct q$，正常坐标由 $\vct q=S^{\mathsf T}M\vct x$ 求出，不是简单用 $S^{\mathsf T}\vct x$。在这种归一化下，动能为 $\dot{\vct q}^{\mathsf T}\dot{\vct q}/2$。

\begin{example}
设初始 $x_1=a,x_2=0$，速度均为零。用未归一化模态表示运动。
\end{example}
% 题设配图；与紧随其后的解答分别编号。
\ExpansionDiagram{fig:audit-m-q-077}{固定端双振子的连接及给定初态}{
\draw[thick](0,-.5)--(0,.8);\draw[exspring](0,.15)--(1.4,.15);\node[above]at(.7,.4){$ k $};\draw[thick](6.7,-.5)--(6.7,.8);\draw[exspring](5.3,.15)--(6.7,.15);\node[above]at(6,.4){$ k $};\draw(1.4,-.3)rectangle(2.4,.6);\node at(1.9,.15){$ m $};\draw(4.3,-.3)rectangle(5.3,.6);\node at(4.8,.15){$ m $};\draw[exspring](2.4,.15)--(4.3,.15);\node[above]at(3.35,.4){$ k $};\draw[->](1.4,-.8)--(2.7,-.8)node[right]{$x_1$};\draw[->](4.3,-.8)--(5.6,-.8)node[right]{$x_2$};\node[above]at(3.35,1.2){$x_1(0)=a,\quad x_2(0)=0$};
}

\begin{solution}
图\ref{fig:audit-m-q-077}使用未归一化振型，系数须由给定初位移确定。在图\ref{fig:orig-065}中，每个振型均把第一物体的位移取为单位，第二物体的带符号位移为 $r_\pm$。

这里的两质量为 $m,2m$，不能沿用等质量系统的等幅分解。正文已求出
\[
\vct u_\pm=(1,r_\pm)^{\mathsf T},\qquad
r_-=(1+\sqrt3)/2,\quad r_+=(1-\sqrt3)/2.
\]
先保留两模的正弦与余弦项，写成
\[
\vct x=\sum_{j=\pm}\vct u_j(c_j\cos\omega_jt+d_j\sin\omega_jt).
\]
在 $t=0$ 求导，初速度条件成为
\[
\omega_-d_-\vct u_-+\omega_+d_+\vct u_+=0.
\]
两振型线性无关且两频率均非零，因此 $d_-=d_+=0$。再代入初始位置，得到
\[
c_-+c_+=a,\qquad r_-c_-+r_+c_+=0.
\]
第二式减去第一式的 $r_+$ 倍，得到 $(r_--r_+)c_-=-r_+a$。使用 $r_--r_+=\sqrt3$，可得
\[
c_-=\frac{\sqrt3-1}{2\sqrt3}a,\qquad
c_+=\frac{\sqrt3+1}{2\sqrt3}a.
\]
所以
\begin{align*}
x_1&=\frac a{2\sqrt3}[(\sqrt3-1)\cos\omega_-t+(\sqrt3+1)\cos\omega_+t],\\
x_2&=\frac a{2\sqrt3}(\cos\omega_-t-\cos\omega_+t).
\end{align*}
直接检查 $x_1(0)=a,x_2(0)=0$，且两速度均为零。初始能量仍为 $ka^2$，因为初态弹簧形变与质量无关；改变质量改变的是演化频率与模态权重。

\begin{figure}[H]\centering
\begin{tikzpicture}[scale=.85]
\foreach\y/\name in{0/低频模,-1.7/高频模}{\draw(0,\y)circle(.23);\draw(3,\y)circle(.32);\node[above]at(0,\y+.3){$m$};\node[above]at(3,\y+.4){$2m$};\node[left]at(-.6,\y){\name};}
\draw[bookaxis,thick](.3,0)--(1.1,0);\draw[bookaxis,thick](3.4,0)--(4.493,0);
\draw[bookaxis,thick](.3,-1.7)--(1.1,-1.7);\draw[bookaxis,thick](2.6,-1.7)--(2.307,-1.7);
\end{tikzpicture}
\caption{质量 $m,2m$ 的模态箭头。每行以第一物体位移为尺度，第二位移比例分别约为1.366和 $-0.366$。}
\label{fig:orig-065}
\end{figure}

\end{solution}

\section{弱耦合、失谐与能量交换}
\subsection{先从对称性找运动形状}
两物块质量均为 $m>0$，左右外侧弹簧刚度为 $k>0$，中间弹簧刚度为 $\kappa>0$；$\kappa=0$ 的独立振子情形可另取极限。位移从平衡位置向右量。若两物块同位移，中间弹簧不伸长，频率应为 $\omega_-=\sqrt{k/m}$；若两者反向等幅位移，中间弹簧的形变量是单个物块位移的两倍，频率应为 $\omega_+=\sqrt{(k+2\kappa)/m}$。这两个猜测可代入受力方程逐一验证。

矩阵写法为
\[
m\ddot{\vct x}+\begin{pmatrix}k+\kappa&-\kappa\\-\kappa&k+\kappa\end{pmatrix}\vct x=0.
\]
系统具有交换两个物块的对称性，所以先尝试 $(1,1)$、$(1,-1)$。矩阵用于确认结果，并为破坏对称性后的计算作准备。

\subsection{完整解与缓慢包络}
只将第一物块拉开$a$，两物块由静止释放。令$q_-=(x_1+x_2)/2$、$q_+=(x_1-x_2)/2$，将受力方程相加、相减，得到
\[\ddot q_-+\omega_-^2q_-=0,\qquad
\ddot q_++\omega_+^2q_+=0.\]
每个方程取指数解，特征根为$\pm\ii\omega_\pm$，所以
$q_\pm=C_\pm\cos\omega_\pm t+D_\pm\sin\omega_\pm t$。
初位移给出$C_-=C_+=a/2$，初速度给出$\omega_-D_-=\omega_+D_+=0$，故两个$D$均为零。由$x_1=q_-+q_+$、$x_2=q_--q_+$还原，得到
\begin{align}
x_1&=\frac a2(\cos\omega_-t+\cos\omega_+t)
=a\cos\bar\omega t\cos\frac{\Delta\omega t}{2},\\
x_2&=\frac a2(\cos\omega_-t-\cos\omega_+t)
=a\sin\bar\omega t\sin\frac{\Delta\omega t}{2},
\end{align}
其中 $\bar\omega=(\omega_++\omega_-)/2$，$\Delta\omega=\omega_+-\omega_-$。三角恒等变形是精确的；把前后两个因子分别看作快振动与慢包络，需要 $\Delta\omega\ll\bar\omega$。

令$\epsilon=\kappa/k\ll1$，则$\omega_+=\omega_0\sqrt{1+2\epsilon}$、$\omega_-=\omega_0$。把
$\sqrt{1+u}=1+u/2-u^2/8+O(u^3)$中的$u$取为$2\epsilon$，再相减，得到
\[
\Delta\omega=\omega_0\left(\epsilon-\frac{\epsilon^2}{2}+O(\epsilon^3)\right),\qquad\omega_0=\sqrt{k/m}.
\]
第一次接近完全的能量转移出现在 $t\simeq\pi/\Delta\omega$，所需时间约为振动周期的 $1/(2\epsilon)$ 倍。只模拟几个快周期，可能根本看不到转移。

定义各外侧振子的局部能量$E_j=m\dot x_j^2/2+kx_j^2/2$，耦合弹簧能量$E_c=\kappa(x_2-x_1)^2/2$。精确守恒的是$E_1+E_2+E_c$，本题初值给出其总量$(k+\kappa)a^2/2$。

令$c=\cos(\Delta\omega t/2)$、$s=\sin(\Delta\omega t/2)$，对两位移作乘积求导：
\[\dot x_1=-a\bar\omega\sin\bar\omega t\,c
-\frac{a\Delta\omega}{2}\cos\bar\omega t\,s,\]
\[\dot x_2=a\bar\omega\cos\bar\omega t\,s
+\frac{a\Delta\omega}{2}\sin\bar\omega t\,c.\]
弱耦合下$\Delta\omega/\bar\omega=O(\epsilon)$，所以每式第二项比典型第一项小一阶。一个快周期内，慢因子$c,s$的变化也为$O(\epsilon)$。取首阶平均时可固定它们，快正弦和余弦平方的平均均为$1/2$，故
\[\langle E_1\rangle_{\rm fast}
=\frac{a^2}{4}(m\bar\omega^2+k)c^2+O(\epsilon ka^2),\qquad
\langle E_2\rangle_{\rm fast}
=\frac{a^2}{4}(m\bar\omega^2+k)s^2+O(\epsilon ka^2).\]
再用$m\bar\omega^2=k[1+O(\epsilon)]$，得到弱耦合首阶结果
\[
E_1\simeq E_*\cos^2(\Delta\omega t/2),\quad
E_2\simeq E_*\sin^2(\Delta\omega t/2),\quad E_*=ka^2/2.
\]
忽略的相对修正为 $O(\epsilon)$。包络的零点不保证某物块的瞬时位移与速度同时为零，因此上述交换图像也有精度范围。

\begin{figure}[htbp]\centering
\begin{tikzpicture}[x=1.5cm,y=2cm]
\draw[bookaxis](0,0)--(6.7,0)node[right]{$\Delta\omega t$};\draw[bookaxis](0,0)--(0,1.35)node[above]{$E_j/E_*$};
\draw[bookcurve,domain=0:6.283,samples=100]plot(\x,{cos(deg(\x/2))^2});
\draw[dashed,domain=0:6.283,samples=100]plot(\x,{sin(deg(\x/2))^2});
\node[below]at(3.1416,0){$\pi$};\node[below]at(6.283,0){$2\pi$};\node[left]at(0,1){$1$};
\node at(1,1.15){实线：$E_1$};\node at(4,1.15){虚线：$E_2$};
\end{tikzpicture}
\caption{弱耦合首阶的平均局部能量。精确总能量还包括中间弹簧的储能。}
\label{fig:orig-066}
\end{figure}

\subsection{轻微不对称为什么能明显改变结果}
将左右外侧弹簧分别改成 $k(1+\delta)$、$k(1-\delta)$，$|\delta|<1$。刚度矩阵为
\[
K=(k+\kappa)I+\begin{pmatrix}k\delta&-\kappa\\-\kappa&-k\delta\end{pmatrix}.
\]
后一个矩阵的平方为 $d^2I$，其中 $d=\sqrt{(k\delta)^2+\kappa^2}$，所以
\begin{equation}\label{eq:detuned-frequencies}
\omega_\pm^2=\frac{k+\kappa\pm d}{m}.
\end{equation}
为写出初值的投影，记上面的无迹矩阵为$H$，它满足$H^2=d^2I$。两个模态的投影矩阵可取
\[P_\pm=\frac12\left(I\pm\frac Hd\right),\qquad
HP_\pm=\pm dP_\pm,\quad P_++P_-=I.\]
将初位移$\vct x_0=(a,0)^{\mathsf T}$代入，
\[P_\pm\vct x_0
=\frac a2\begin{pmatrix}1\pm k\delta/d\\\mp\kappa/d\end{pmatrix}.\]
每个投影各自按$\omega_\pm$振动；初速度为零，使两模态的正弦系数都为零。因此
$\vct x(t)=P_-\vct x_0\cos\omega_-t+P_+\vct x_0\cos\omega_+t$，逐分量得到
\begin{align}
x_1&=\frac a2\left[\left(1-\frac{k\delta}{d}\right)\cos\omega_-t+
\left(1+\frac{k\delta}{d}\right)\cos\omega_+t\right],\\
x_2&=\frac{a\kappa}{2d}(\cos\omega_-t-\cos\omega_+t).
\end{align}
第二物块位移可化为$(a\kappa/d)\sin\bar\omega t\sin(\Delta\omega t/2)$，其慢包络最大值为$a\kappa/d$。当$|\delta|\ll1$且$\kappa/k\ll1$时，按上面对快振动平均的方法，
\[\langle E_2\rangle_{\rm fast}
\simeq\frac{ka^2}{2}\frac{\kappa^2}{d^2}
\sin^2\frac{\Delta\omega t}{2}.\]
总能量的首阶值仍为$ka^2/2$，故最大平均能量转移比例为
\begin{equation}\label{eq:transfer-fraction}
\eta\simeq\frac{\kappa^2}{\kappa^2+(k\delta)^2}.
\end{equation}
所谓“轻微不对称”应与耦合强度比较。若 $|\delta|=\kappa/k$，转移比例已约降为一半；若 $|\delta|\gg\kappa/k$，两个振子虽仍接近同频，能量转移也会受到明显抑制。

\subsection{频率误差与相位积累}
若以 $\Delta\omega\simeq\epsilon\omega_0$ 代替精确频差，频差误差为 $O(\epsilon^2\omega_0)$。经过时间 $t$，包络相位误差增长到 $O(\epsilon^2\omega_0t)$。在一次交换时间 $t=O(1/(\epsilon\omega_0))$ 内，它为 $O(\epsilon)$；若把同一近似延长到 $t=O(1/(\epsilon^2\omega_0))$，误差就可能达到1的量级。

近似式的适用范围还受观察时间限制。第四章的计算实验将比较精确运动与近似包络随时间产生的差别。

\subsection{双振子的初值与失谐}
\textbf{同相初值。} 固定端三弹簧双振子均取 $m,k$，初值为 $(a,a)$、速度为零，求运动及能量。

% 题设配图；与紧随其后的解答分别编号。
\ExpansionDiagram{fig:audit-m-q-078}{固定端双振子的连接及给定初态}{
\draw[thick](0,-.5)--(0,.8);\draw[exspring](0,.15)--(1.4,.15);\node[above]at(.7,.4){$ k $};\draw[thick](6.7,-.5)--(6.7,.8);\draw[exspring](5.3,.15)--(6.7,.15);\node[above]at(6,.4){$ k $};\draw(1.4,-.3)rectangle(2.4,.6);\node at(1.9,.15){$ m $};\draw(4.3,-.3)rectangle(5.3,.6);\node at(4.8,.15){$ m $};\draw[exspring](2.4,.15)--(4.3,.15);\node[above]at(3.35,.4){$ k $};\draw[->](1.4,-.8)--(2.7,-.8)node[right]{$x_1$};\draw[->](4.3,-.8)--(5.6,-.8)node[right]{$x_2$};\node[above]at(3.35,1.2){$x_1(0)=x_2(0)=a$};
}

\begin{solution}
题设图\ref{fig:audit-m-q-078}给出$x_1(0)=x_2(0)=a$，两初速度都为零；解图\ref{fig:class-3-1}把两物块的等位移及不变的中间距离画出。
同相方向是正常模，因此两物块始终同位移$x(t)$，中间弹簧伸长$x_2-x_1=0$。
每条运动方程均化为$m\ddot x+kx=0$。记$\omega_1=\sqrt{k/m}$，通解为$x=A\cos\omega_1t+B\sin\omega_1t$；初态给出$A=a$、$\omega_1B=0$，故$B=0$。因此
\[
 x_1(t)=x_2(t)=a\cos(\omega_1t),\qquad \omega_1=\sqrt{k/m}.
\]
为写出任意时刻的能量，令$\chi=\omega_1t$。两物块总动能和两端弹簧总势能分别为
\[
 T=2\cdot\frac12m a^2\omega_1^2\sin^2\chi
   =ka^2\sin^2\chi,\qquad
 U=2\cdot\frac12ka^2\cos^2\chi=ka^2\cos^2\chi.
\]
所以$E=T+U=ka^2$。图中不变形的中间弹簧始终不储能。

\begin{figure}[H]
\centering
\begin{tikzpicture}[scale=.83]
\draw[thick](0,-.5)--(0,.5) (8,-.5)--(8,.5);
\draw[decorate,decoration={coil,segment length=4pt,amplitude=3pt}](0,0)--(2,0);
\draw(2,-.3)rectangle(3,.3);\node at(2.5,0){$m$};
\draw[decorate,decoration={coil,segment length=4pt,amplitude=3pt}](3,0)--(5,0);
\draw(5,-.3)rectangle(6,.3);\node at(5.5,0){$m$};
\draw[decorate,decoration={coil,segment length=4pt,amplitude=3pt}](6,0)--(8,0);
\foreach\x in{1,4,7}{\node at(\x,.55){$k$};}
\draw[-{Stealth}](2.2,-.7)--(3.2,-.7)node[right]{$x_1$};
\draw[-{Stealth}](5.2,-.7)--(6.2,-.7)node[right]{$x_2$};
\draw[bookaxis](2.2,1.05)--(3.2,1.05)node[midway,above]{$a$};\draw[bookaxis](5.2,1.05)--(6.2,1.05)node[midway,above]{$a$};
\end{tikzpicture}
\caption{第三章同相初值解图：两位移始终相等，中间弹簧不形变，对应正常坐标中的同相轴。}\label{fig:class-3-1}
\end{figure}
\end{solution}

\textbf{反相初值。} 改为初值 $(a,-a)$、速度为零，运动与能量怎样变化？

% 题设配图；与紧随其后的解答分别编号。
\ExpansionDiagram{fig:audit-m-q-079}{固定端双振子的连接及给定初态}{
\draw[thick](0,-.5)--(0,.8);\draw[exspring](0,.15)--(1.4,.15);\node[above]at(.7,.4){$ k $};\draw[thick](6.7,-.5)--(6.7,.8);\draw[exspring](5.3,.15)--(6.7,.15);\node[above]at(6,.4){$ k $};\draw(1.4,-.3)rectangle(2.4,.6);\node at(1.9,.15){$ m $};\draw(4.3,-.3)rectangle(5.3,.6);\node at(4.8,.15){$ m $};\draw[exspring](2.4,.15)--(4.3,.15);\node[above]at(3.35,.4){$ k $};\draw[->](1.4,-.8)--(2.7,-.8)node[right]{$x_1$};\draw[->](4.3,-.8)--(5.6,-.8)node[right]{$x_2$};\node[above]at(3.35,1.2){$x_1(0)=a,\quad x_2(0)=-a$};
}

\begin{solution}
图\ref{fig:audit-m-q-079}的初位移为$(a,-a)$。按解图\ref{fig:class-3-2}中两箭头的方向，中间弹簧的有符号伸长为$x_2-x_1=-2a$。
令$x_1=x,x_2=-x$代入第一物块的方程，
\[
 m\ddot x=-kx+k(-x-x)=-3kx.
\]
第二条方程在同一代入下也成立。特征方程$m\lambda^2+3k=0$给出$\omega_2=\sqrt{3k/m}$，于是$x=A\cos\omega_2t+B\sin\omega_2t$。初位移$x(0)=a$给出$A=a$，初速度$\dot x(0)=\omega_2B=0$给出$B=0$，因此
\[
 x_1=a\cos(\omega_2t),\qquad x_2=-a\cos(\omega_2t),
 \qquad\omega_2=\sqrt{3k/m}.
\]
在释放时刻，两端弹簧总储能为$ka^2$，中间弹簧储能为$\tfrac12k(2a)^2=2ka^2$，
故总能量$E=3ka^2$。任意时刻令$x=x_1$，则
$T=m\dot x^2$、$U=kx^2+2kx^2=3kx^2$；代入上式仍得$T+U=3ka^2$。

\begin{figure}[H]\centering
\begin{tikzpicture}[scale=.9]
\draw[thick](0,-.6)--(0,.6) (8,-.6)--(8,.6);\draw[exspring](0,0)--(2,0);\draw[exspring](3,0)--(5,0);\draw[exspring](6,0)--(8,0);
\draw(2,-.4)rectangle(3,.4);\draw(5,-.4)rectangle(6,.4);\node at(2.5,0){$m$};\node at(5.5,0){$m$};
\draw[bookaxis](2,.95)--(3.1,.95)node[midway,above]{$a$};\draw[bookaxis](6,.95)--(4.9,.95)node[midway,above]{$-a$};
\node[below]at(4,-.75){中间弹簧形变 $-2a$，储能 $2ka^2$};
\end{tikzpicture}
\caption{第三章反相初值解图：两端弹簧各形变大小 $a$，中间弹簧形变大小 $2a$，总势能为 $3ka^2$。}\label{fig:class-3-2}
\end{figure}
\end{solution}

\textbf{加入失谐。} 对正文的失谐模型，$\delta\ne0$ 时同向等幅初值还能否只激发一个频率？

% 题设配图；与紧随其后的解答分别编号。
\ExpansionDiagram{fig:audit-m-q-080}{失谐写入两端弹簧的不同刚度}{
\draw[thick](0,-.5)--(0,.8);\draw[exspring](0,.15)--(1.4,.15);\node[above]at(.7,.4){$ k(1+\delta) $};\draw[thick](6.7,-.5)--(6.7,.8);\draw[exspring](5.3,.15)--(6.7,.15);\node[above]at(6,.4){$ k(1-\delta) $};\draw(1.4,-.3)rectangle(2.4,.6);\node at(1.9,.15){$ m $};\draw(4.3,-.3)rectangle(5.3,.6);\node at(4.8,.15){$ m $};\draw[exspring](2.4,.15)--(4.3,.15);\node[above]at(3.35,.4){$ \kappa $};\draw[->](1.4,-.8)--(2.7,-.8)node[right]{$x_1$};\draw[->](4.3,-.8)--(5.6,-.8)node[right]{$x_2$};
}

\begin{solution}
图\ref{fig:audit-m-q-080}的两端刚度相差$2k\delta$。把同相向量代入刚度矩阵，便可判断图\ref{fig:class-3-3}中的两个方向是否平行。令$\epsilon=\kappa/k$，则
\[
 \frac Kk\binom11=
 \begin{pmatrix}1+\epsilon+\delta&-\epsilon\\-\epsilon&1+\epsilon-\delta\end{pmatrix}
 \binom11=\binom{1+\delta}{1-\delta}.
\]
若$(1,1)$是振型，两分量必须等于同一个比例因子，要求
$1+\delta=1-\delta$，即$\delta=0$。
因此在题设$\delta\ne0$时，初始加速度已不沿同相方向，运动不能始终保持等幅同向。
对正定、等质量的刚度矩阵，可用两个正交特征向量展开初始位移；
若其中某个投影为零，初始位移就必须沿另一个特征向量，与上述计算矛盾。
故同相零速度初态含有两个正常模；两个特征值不同，便出现两个频率。

\begin{figure}[H]\centering
\begin{tikzpicture}[scale=1.4]
\draw[bookaxis](-.2,0)--(2.6,0)node[right]{$x_1$};\draw[bookaxis](0,-.2)--(0,2.4)node[above]{$x_2$};
\draw[bookaxis,thick](0,0)--(1.4,1.4)node[above]{$(1,1)$};\draw[bookaxis,dashed](0,0)--(1.8,1)node[right]{$(1+\delta,1-\delta)$};
\draw[bookguide](0,0)--(2.1,2.1);\node[below]at(1.6,-.45){示意 $\delta>0$};
\end{tikzpicture}
\caption{第三章失谐初值解图：刚度矩阵作用后的方向不再与同相位移平行，因此同向等幅初态不再是单一正常模。}\label{fig:class-3-3}
\end{figure}
\end{solution}

\section{多振子链与边界条件}
局部耦合给出链内递推，端部约束选择允许振型。先求一般空间形式，再用固定、自由或周期边界确定离散频率。

考虑 $n$ 个相同质量 $m$，由相同弹簧 $k$ 连接成链，两端固定。规定 $x_0=x_{n+1}=0$，方程为
\begin{equation}
m\ddot x_j=k(x_{j-1}-2x_j+x_{j+1}),\qquad j=1,\ldots,n.
\end{equation}
试探空间形状 $u_j=\sin(j\theta)$，利用 $\sin[(j-1)\theta]+\sin[(j+1)\theta]=2\cos\theta\sin j\theta$，得到 $m\omega^2=2k(1-\cos\theta)$。两端条件要求 $\theta_\ell=\ell\pi/(n+1)$，所以
\begin{equation}
\omega_\ell=2\sqrt{k/m}\sin\frac{\ell\pi}{2(n+1)},\qquad \ell=1,\ldots,n.
\end{equation}
归一化常数可由有限等比和确定。对$\theta_\ell=\ell\pi/(n+1)$，有
$\sum_{j=0}^{n}\ee^{2\ii j\theta_\ell}=0$，取实部得$\sum_{j=1}^n\cos(2j\theta_\ell)=-1$。于是
\[\sum_{j=1}^n\sin^2(j\theta_\ell)
=\frac n2-\frac12\sum_{j=1}^n\cos(2j\theta_\ell)=\frac{n+1}{2}.\]
对应单位特征向量可取$u_j^{(\ell)}=\sqrt{2/(n+1)}\sin[j\ell\pi/(n+1)]$。这些不同频率的振型因刚度矩阵实对称而正交，完整的$n$个方向构成正常模基。$n=2$时恢复前面的同相与反相模。

边界条件参与决定允许的 $\theta$。若两端自由，整体平移的零频模会重新出现，端点方程也不再具有与内部相同的系数2。只记住频率公式而忽略端点条件，会把不同物理系统混为一谈。

\begin{figure}[htbp]
\centering
\begin{tikzpicture}[x=1.1cm,y=.6cm]
\draw[dashed](0,0)--(4,0);\draw[dashed](0,-3)--(4,-3);\draw[dashed](0,-6)--(4,-6);
\draw[thick](0,0)--(1,.707)--(2,1)--(3,.707)--(4,0);
\draw[thick](0,-3)--(1,-2)--(2,-3)--(3,-4)--(4,-3);
\draw[thick](0,-6)--(1,-5.293)--(2,-7)--(3,-5.293)--(4,-6);
\foreach\x in{1,2,3}{\fill(\x,{sin(\x*45)})circle(1.5pt);\fill(\x,{-3+sin(\x*90)})circle(1.5pt);\fill(\x,{-6+sin(\x*135)})circle(1.5pt);}
\node[right]at(4.3,0){$\ell=1$};\node[right]at(4.3,-3){$\ell=2$};\node[right]at(4.3,-6){$\ell=3$};
\end{tikzpicture}
\caption{固定端三振子链的三个离散正常模}\label{fig:chain-modes}
\end{figure}
图中折线仅连接各物体的位移值，描述的是离散模态形状，不能把折线斜率直接解释成相邻弹簧的空间倾角。三个图的垂直尺度相同，整体乘以任意常数仍代表同一个正常模方向。

\section{转动惯量张量与主轴}
同一刚体绕不同方向转动，动能和角动量的关系也会改变。先由质量微元构造张量，再用主轴把方向耦合分离。

\subsection{从质点动能推导惯量矩阵}
刚体绕参考点纯转动时，位置 $\vct r$ 的质元速度为 $\vct\omega\times\vct r$。用叉积模的恒等式，
\[
T=\frac12\int[\omega^2r^2-(\vct\omega\cdot\vct r)^2]\dd m
=\frac12\vct\omega^{\mathsf T}I\vct\omega,
\]
其中
\begin{equation}
I_{ij}=\int(r^2\delta_{ij}-r_ir_j)\dd m.
\end{equation}
例如$I_{xx}=\int(y^2+z^2)\dd m$、$I_{xy}=-\int xy\dd m$。它是实对称矩阵。角动量也由质元积分得到。对内层速度用三重积恒等式展开，
\[
\begin{aligned}
\vct L&=\int\vct r\times(\vct\omega\times\vct r)\dd m\\
&=\int[r^2\vct\omega-(\vct r\cdot\vct\omega)\vct r]\dd m
=I\vct\omega.
\end{aligned}
\]
因此$\vct L$与$\vct\omega$一般不平行。对称矩阵的特征方向称为惯量主轴，特征值称为主惯量；沿主轴转动时二者平行。

\begin{example}
均匀细杆长 $\ell$、质量 $M$，沿 $x$ 轴放置，原点在质心。求惯量张量及任意过质心轴的惯量。
\end{example}
% 题设配图；与紧随其后的解答分别编号。
\ExpansionDiagram{fig:audit-m-q-081}{细杆沿横轴，取过质心的任意转轴}{
\draw[very thick](-2,0)--(2,0);\fill(0,0)circle(2pt)node[below]{$C$};\draw[->](0,-1.6)--(0,2)node[above]{$z$};\draw[<->](-2,-.6)--(2,-.6)node[midway,below]{$\ell$};\draw[->](0,0)--(1.5,1.5)node[right]{$\vct n$};
}

\begin{solution}
图\ref{fig:audit-m-q-081}中的轴向单位矢量$\vct n$决定到轴距离，惯量张量须按质点位置积分。
图\ref{fig:orig-068}中的垂线给出质元到转轴的真实距离，不能把沿杆的位置参数直接当成力臂。

杆上质元写成 $\vct r=(s,0,0)$，$\dd m=(M/\ell)\dd s$，$-\ell/2\le s\le\ell/2$。由惯量张量定义，
\[
I_{xx}=\int(y^2+z^2)\dd m=0,\qquad
I_{yy}=I_{zz}=\frac M\ell\int_{-\ell/2}^{\ell/2}s^2\dd s
=\frac M\ell\left[\frac{s^3}{3}\right]_{-\ell/2}^{\ell/2}
=\frac{M\ell^2}{12}.
\]
每个交叉项都含 $y$ 或 $z$，故均为零，因而 $I=(M\ell^2/12)\diag(0,1,1)$。任意单位轴向 $\vct n=(n_x,n_y,n_z)$ 的轴惯量为
\[
I_n=\vct n^{\mathsf T}I\vct n
=\frac{M\ell^2}{12}(n_y^2+n_z^2)
=\frac{M\ell^2}{12}(1-n_x^2).
\]
若轴与杆夹角为 $\alpha$，$n_x=\cos\alpha$，故 $I_n=M\ell^2\sin^2\alpha/12$。从距离法看，质元到轴距离为 $|s|\sin\alpha$，把它平方后积分也给出同式。

\begin{figure}[H]\centering
\begin{tikzpicture}[scale=1]
\draw[line width=1.4pt](-2,0)--(2,0);\fill(0,0)circle(1.4pt)node[below]{$O$};
\draw[bookaxis](-2.3,0)--(2.8,0)node[right]{杆的主轴};
\draw[bookaxis,thick](-1.7,-1.7)--(1.9,1.9)node[above]{$\vct n$};
\coordinate(A)at(1.6,0);\coordinate(H)at(.8,.8);\fill(A)circle(1.4pt)node[below]{$\dd m$};
\draw[bookguide](A)--(H);\draw(.66,.66)--(.80,.52)--(.94,.66);
\node[right]at(1.15,.45){$d_\perp$};\draw(.55,0)arc(0:45:.55);\node at(23:.83){$\alpha$};
\end{tikzpicture}
\caption{细杆相对倾斜轴的距离。杆上位置 $s$ 的质元到轴距离为 $|s|\sin\alpha$，轴惯量由这一距离的平方积分得到。}
\label{fig:orig-068}
\end{figure}

当轴垂直杆时恢复 $M\ell^2/12$；当轴沿杆时为零，这是理想零厚度模型的结果。真实杆有截面，沿杆轴的惯量虽小但不为零。若角速度同时含沿杆和垂直杆的分量，角动量只含后者的贡献，一般不与角速度平行。

图中 $d_\perp=|s|\sin\alpha$，平方后积分得到 $\sin^2\alpha$ 因子，与张量投影的结果相同。
\end{solution}

\subsection{平行轴定理的矩阵形式}
以质心为原点的质元位置为$\vct r$，换到相距$\vct a$的参考点时新位置为$\vct r-\vct a$。先展开惯量积分中的矩阵核：
\[
\begin{aligned}
&|\vct r-\vct a|^2I_3-(\vct r-\vct a)(\vct r-\vct a)^{\mathsf T}\\
&\quad=(r^2I_3-\vct r\vct r^{\mathsf T})+(a^2I_3-\vct a\vct a^{\mathsf T})\\
&\qquad-2(\vct r\cdot\vct a)I_3
+\vct r\vct a^{\mathsf T}+\vct a\vct r^{\mathsf T}.
\end{aligned}
\]
最后三个交叉项积分后均含$\int\vct r\dd m$，由质心条件为零；不含$\vct r$的两项积分后乘总质量$M$。于是
\begin{equation}
I'=I_{\rm cm}+M(a^2I_3-\vct a\vct a^{\mathsf T}),
\end{equation}
其中 $I_3$ 为三维单位矩阵。对单位轴向 $\vct n$ 取二次型，增加量为 $M[a^2-(\vct a\cdot\vct n)^2]=Md_\perp^2$，恢复熟悉的平行轴定理。平行轴项使用的是两轴之间的垂直距离。

\section{矩阵指数与线性微分方程（选读）}
\subsection{常系数系统的演化矩阵}
对 $\dot{\vct y}=A\vct y$，定义
\[
\ee^{At}=I+At+\frac{A^2t^2}{2!}+\cdots.
\]
级数及其导数在有限时间闭区间上一致收敛，可以逐项求导：
\[\frac{\dd\ee^{At}}{\dd t}
=\sum_{n=1}^{\infty}\frac{A^nnt^{n-1}}{n!}
=A\sum_{j=0}^{\infty}\frac{A^jt^j}{j!}=A\ee^{At}.\]
又有$\ee^{A0}=I$，所以$\vct y(t)=\ee^{At}\vct y(0)$满足方程和初值。当$A=P\Lambda P^{-1}$可对角化时，$A^n=P\Lambda^nP^{-1}$，代回级数得到$\ee^{At}=P\diag(\ee^{\lambda_jt})P^{-1}$。特征方向各自独立增长或衰减，换回原变量后表现为耦合。

这个公式要求 $A$ 不随时间变化。对一般 $A(t)$，不能未经检查就写成 $\ee^{\int A\dd t}$，因为不同时间的矩阵可能不交换，求导的次序会发生作用。

\begin{example}
$A=\begin{pmatrix}0&-\omega\\\omega&0\end{pmatrix}$。求演化矩阵。
\end{example}
% 题设配图；与紧随其后的解答分别编号。
\ExpansionDiagram{fig:audit-m-q-082}{线性演化矩阵把位置转为垂直速度}{
\draw[->](-.4,0)--(3.6,0)node[right]{$ x $};\draw[->](0,-0.3)--(0,2.8)node[above]{$ y $};\draw[->,thick](0,0)--(2,0)node[below]{$\vct y_0$};\draw[->,thick](2,0)--(2,1.6)node[right]{$A\vct y_0$};\draw[dashed](2,0)arc(0:70:2);
}

\begin{solution}
图\ref{fig:audit-m-q-082}中初始速度垂直初始状态，矩阵指数仍应由方程与单位初值确定。
图\ref{fig:eg-3-10}显示该矩阵场处处沿圆周切向。先由幂级数求演化矩阵，再与图中的转角核对。

直接计算 $A^2=-\omega^2I$，所以任意偶次幂与奇次幂分别为
\[
A^{2n}=(-1)^n\omega^{2n}I,\qquad
A^{2n+1}=(-1)^n\omega^{2n}A.
\]
在指数级数中按奇偶分组，得
\begin{align*}
\ee^{At}&=I\sum_{n=0}^\infty\frac{(-1)^n(\omega t)^{2n}}{(2n)!}\\
&\quad+\frac A\omega\sum_{n=0}^\infty\frac{(-1)^n(\omega t)^{2n+1}}{(2n+1)!}\\
&=\begin{pmatrix}\cos\omega t&-\sin\omega t\\\sin\omega t&\cos\omega t\end{pmatrix}.
\end{align*}
这就是旋转矩阵 $R(\omega t)$，其转置乘自身为单位阵，故状态长度保持不变。还可检查 $R(0)=I$、$\dot R=AR$，两条件恰好确定演化矩阵。上面的分组写法暂设 $\omega\ne0$；若 $\omega=0$，则 $A=0$，直接有 $\ee^{At}=I$。

\begin{figure}[H]\centering
\begin{tikzpicture}[scale=1.2]
\draw[bookaxis](-1.9,0)--(2.5,0)node[right]{$y_1$};\draw[bookaxis](0,-1.8)--(0,2.1)node[above]{$y_2$};
\draw[bookguide](0,0)circle(1.4);\draw[bookaxis,thick](0,0)--(1.4,0)node[below right]{$\vct y(0)$};
\draw[bookaxis,thick](0,0)--(55:1.4)node[above right]{$\vct y(t)$};
\draw[->](.6,0)arc(0:55:.6);\node at(28:.95){$\omega t$};
\draw[bookaxis](55:1.4)--({1.4*cos(55)-.75*sin(55)},{1.4*sin(55)+.75*cos(55)})node[above left]{$A\vct y$};
\end{tikzpicture}
\caption{例3.10：$\omega>0$ 时速度 $A\vct y$ 沿圆周切向，演化矩阵把状态逆时针转过 $\omega t$，保持长度。}\label{fig:eg-3-10}
\end{figure}

矩阵在 $t=0$ 为单位阵，转角随时间增加为 $\omega t$；这同时检验正弦项的符号和初始条件。
\end{solution}

\subsection{复特征值的实解与相图}
若实矩阵有特征值 $\lambda=\alpha+\ii\beta$，特征向量为 $\vct u+\ii\vct v$，则两条实解为
\begin{align*}
\vct y_1&=\ee^{\alpha t}(\vct u\cos\beta t-\vct v\sin\beta t),\\
\vct y_2&=\ee^{\alpha t}(\vct u\sin\beta t+\vct v\cos\beta t).
\end{align*}
实部 $\alpha$ 控制增减，虚部 $\beta$ 控制振动。对二维线性系统，实负根给衰减方向，一正一负根形成鞍点，纯虚共轭根通常给中心，具有负实部的复根给稳定焦点。若存在重根与非平凡 Jordan 块，还要检查多项式时间因子。
\begin{figure}[htbp]\centering
\begin{tikzpicture}[scale=.75]
\begin{scope}\draw[bookaxis](-1.7,0)--(1.7,0);\draw[bookaxis](0,-1.7)--(0,1.7);\draw(0,0)circle(1.2);\draw[bookaxis](30:1.2)arc(30:70:1.2);\node at(0,-2){中心};\end{scope}
\begin{scope}[shift={(4.2,0)}]\draw[bookaxis](-1.7,0)--(1.7,0);\draw[bookaxis](0,-1.7)--(0,1.7);\draw[domain=0:17,samples=300]plot({1.4*exp(-.13*\x)*cos(deg(\x))},{1.4*exp(-.13*\x)*sin(deg(\x))});\node at(0,-2){稳定焦点};\end{scope}
\begin{scope}[shift={(8.4,0)}]\draw[bookaxis](-1.7,0)--(1.7,0);\draw[bookaxis](0,-1.7)--(0,1.7);\foreach\c in{.2,.5}{\draw[domain=.35:1.5,samples=50]plot(\x,{\c/\x});\draw[domain=.35:1.5,samples=50]plot(-\x,{-\c/\x});}\node at(0,-2){鞍点};\end{scope}
\end{tikzpicture}
\caption{二维线性系统的典型相图。中心保持轨道，焦点逐渐趋近原点，鞍点含增长方向。}
\label{fig:orig-069}
\end{figure}

\subsection{重根并不总给两个纯指数}
对$A=\lambda I+N$，若$N^2=0$且$N\ne0$，矩阵$\lambda I$与$N$交换，因此可分解指数。再按幂级数计算，二次及更高次的$N$均为零，故
\[
\ee^{At}=\ee^{\lambda t}\ee^{Nt}
=\ee^{\lambda t}\left(I+tN+\frac{t^2N^2}{2!}+\cdots\right)
=\ee^{\lambda t}(I+tN).
\]
例如$\begin{pmatrix}\lambda&1\\0&\lambda\end{pmatrix}$的解包含$t\ee^{\lambda t}$。仅把重特征值重复写两次，不能恢复全部初值解。

临界阻尼简谐振子的状态矩阵为 $A=\begin{pmatrix}0&1\\-1&-2\end{pmatrix}$，有 $(A+I)^2=0$。因此 $\ee^{At}=\ee^{-t}[I+t(A+I)]$。初值 $(1,0)$ 给出 $q=(1+t)\ee^{-t},p=-t\ee^{-t}$，直接恢复第一章临界阻尼解。Jordan 结构解释了它为何多出一个时间因子。

\section{强迫耦合系统与反共振}
耦合系统的受迫运动由外力的施加位置和各振型共同决定。先建立频域方程，再区分固有频率附近的共振与某个坐标响应为零的反共振。

\subsection{从正常模到频率响应}
机械系统 $M\ddot{\vct x}+B\dot{\vct x}+K\vct x=\Rea[\vct F\ee^{\ii\omega t}]$ 的稳态相量满足
\[
(K-\omega^2M+\ii\omega B)\vct X=\vct F.
\]
计算这个线性方程组即可得到每个物体的幅值与相位。不能对每个物体分别使用一个孤立单振子的分母，因为耦合项同时连接所有坐标。

若阻尼为 $B=\alpha M+\beta K$，质量归一化正常坐标仍可将方程拆开，各模阻尼系数为 $\alpha+\beta\omega_j^2$。对一般阻尼矩阵，刚度正常模下的阻尼项未必对角，模态之间仍可能耦合。所谓“先求正常模后逐模响应”，需要检查这个条件。

\begin{example}
固定端双振子质量与弹簧均相同，忽略阻尼，只给第一物体驱动力 $F_0\cos\omega t$。求远离共振时的稳态，指出第一物体可能不动的频率。
\end{example}
% 题设配图；与紧随其后的解答分别编号。
\ExpansionDiagram{fig:audit-m-q-083}{外力只加在固定端双振子的第一物体上}{
\draw[thick](0,-.5)--(0,.8);\draw[exspring](0,.15)--(1.4,.15);\node[above]at(.7,.4){$ k $};\draw[thick](6.7,-.5)--(6.7,.8);\draw[exspring](5.3,.15)--(6.7,.15);\node[above]at(6,.4){$ k $};\draw(1.4,-.3)rectangle(2.4,.6);\node at(1.9,.15){$ m $};\draw(4.3,-.3)rectangle(5.3,.6);\node at(4.8,.15){$ m $};\draw[exspring](2.4,.15)--(4.3,.15);\node[above]at(3.35,.4){$ k $};\draw[->](1.4,-.8)--(2.7,-.8)node[right]{$x_1$};\draw[->](4.3,-.8)--(5.6,-.8)node[right]{$x_2$};\draw[->,thick](1.9,1.25)--(3.2,1.25)node[right]{$F_0\cos\omega t$};
}

\begin{solution}
图\ref{fig:audit-m-q-083}中只有第一分量含外力，联立两条相量方程后才能判断某物体的反共振。
图\ref{fig:orig-070}用带符号系数区分同相与反相；零点与两侧的极点来自不同的代数条件。

位移方向都取向右。第一物块受外侧弹簧力 $-kx_1$ 和中间弹簧力 $k(x_2-x_1)$，第二物块类似。因此
\[
m\ddot x_1+2kx_1-kx_2=F_0\cos\omega t,
\qquad m\ddot x_2+2kx_2-kx_1=0.
\]
取纯谐特解后，记$d=2k-m\omega^2$，相量方程为$dX_1-kX_2=F_0$、$-kX_1+dX_2=0$。将第一式乘$d$、第二式乘$k$相加，得$(d^2-k^2)X_1=dF_0$；将第一式乘$k$、第二式乘$d$相加，得$(d^2-k^2)X_2=kF_0$。两步均未除以$d$，故也保留$d=0$的反共振情形。于是
\[
D=d^2-k^2=(k-m\omega^2)(3k-m\omega^2),
\quad X_1=\frac{dF_0}{D},\quad X_2=\frac{kF_0}{D}.
\]
要求 $D\ne0$，故不含两个固有频率 $\sqrt{k/m},\sqrt{3k/m}$。令第一位移为零，需要 $d=0$，即 $\omega^2=2k/m$。此时 $D=-k^2\ne0$，
\[
X_1=0,\qquad X_2=-F_0/k.
\]
第一物体上的中间弹簧力与外力逐时抵消，第二物体仍在运动。负相量表示反相，物理幅值取其绝对值。

\begin{figure}[H]\centering
\begin{tikzpicture}[x=1.4cm,y=.58cm]
\begin{scope}\clip(0,-3)rectangle(4,3);
\foreach\a/\b in{0/.91,1.09/2.91,3.09/4}{\draw[thick,domain=\a:\b,samples=100]plot(\x,{(2-\x)/((2-\x)^2-1)});\draw[dashed,domain=\a:\b,samples=100]plot(\x,{1/((2-\x)^2-1)});}
\end{scope}
\draw[bookaxis](0,0)--(4.4,0)node[right]{$s=m\omega^2/k$};\draw[bookaxis](0,-3.2)--(0,3.3)node[above]{$kX_j/F_0$};
\draw[dotted](1,-3)--(1,3) (3,-3)--(3,3);\foreach\t in{1,2,3,4}{\node[below]at(\t,0){$\t$};}\fill(2,0)circle(1.5pt);
\end{tikzpicture}
\caption{无阻尼双振子的带符号稳态系数，实线为第一物体，虚线为第二物体；共振附近已截去，$s=2$ 处第一物体反共振。}
\label{fig:orig-070}
\end{figure}

这些式子描述选定的纯谐响应。无阻尼系统从任意初态启动还可含不衰减的自由振动，不能靠等待自动消去。共振处有限纯谐特解失效；改变驱动位置或观察位置，也会改变反共振出现在哪个输出上。

图中实线在 $s=2$ 穿过零点，而虚线仍为负；它表示第二物体运动所产生的弹簧力抵消了第一物体上的外力。
\end{solution}

\subsection{非谐波驱动的变参数公式}
对常系数一阶系统$\dot{\vct y}=A\vct y+\vct f(t)$，令$\vct y=\ee^{At}\vct c(t)$，乘积求导为
\[\dot{\vct y}=A\ee^{At}\vct c+\ee^{At}\dot{\vct c}.\]
代回原方程，两边的$A\ee^{At}\vct c$相消，故$\ee^{At}\dot{\vct c}=\vct f$。左乘$\ee^{-At}$，再用$\vct c(0)=\vct y(0)$积分，得到
\[\vct c(t)=\vct y(0)+\int_0^t\ee^{-As}\vct f(s)\dd s.\]
乘回$\ee^{At}$，利用同一矩阵指数相乘可合并为$\ee^{A(t-s)}$，于是
\[
\vct y(t)=\ee^{At}\vct y(0)+\int_0^t\ee^{A(t-s)}\vct f(s)\dd s.
\]
第一项是初值的自由演化，积分项是每个时刻输入经过剩余时间传播后的叠加。这个式子将初值响应和强迫响应明确分开，也提供数值积分中的解析检查基准。

\clearpage
\section*{习题与解答}\addcontentsline{toc}{section}{习题与解答}
\begin{enumerate}[label=3.\arabic*]
\item \begin{minipage}[t]{\linewidth}
【基础】对正文三元方程组写出系数矩阵，计算行列式并解释唯一解；当右侧任意改变时，唯一性是否保持？
\end{minipage}
% 题设配图；与紧随其后的解答分别编号。
\ExpansionDiagram{fig:audit-m-q-084}{同一系数矩阵对应三个平面法向}{
\begin{scope}[x={(1cm,0cm)},y={(-.45cm,.3cm)},z={(0cm,.8cm)}]\filldraw[fill=gray!10](0,0,2)--(2,0,0)--(2,1,-1)--(0,1,1)--cycle;\filldraw[fill=gray!25,opacity=.7](0,0,1.6667)--(2,0,.3333)--(2,1,0)--(0,1,1.3333)--cycle;\draw[thick](0,0,0)--(2,0,-2)--(2,1,-1)--(0,1,1)--cycle;\end{scope}\node[right,align=left]at(2.6,1){$x+y+z=2$\\$2x+y+3z=5$\\$x-y+z=0$};
}

\begin{solution}
图\ref{fig:audit-m-q-084}中改变右端相当于平移三个平面，法向与系数矩阵保持不变。
例3.2的图\ref{fig:hw-3-1}给出本方程组消元后的交点。这里用行列式检验该交点的唯一性是否依赖右侧。

将未知量按 $(x,y,z)$ 排列，系数矩阵为
\[
A=\begin{pmatrix}1&1&1\\2&1&3\\1&-1&1\end{pmatrix}.
\]
用第二行减第一行的两倍、第三行减第一行，这些行变换不改变行列式，得
\[
\det A=\det\begin{pmatrix}1&1&1\\0&-1&1\\0&-2&0\end{pmatrix}
=(-1)\times0-1\times(-2)=2.
\]
因此 $A$ 可逆，$\rank A=3$。对任意右侧列 $\vct b$，均有且仅有一个解 $\vct x=A^{-1}\vct b$。右侧改变只会改变这个解的位置，系数矩阵、秩和可逆性都未改变。也可从齐次方程看唯一性：若有两个解，其差满足 $A\Delta\vct x=0$；可逆性迫使差为零。

\begin{figure}[H]\centering
\begin{tikzpicture}[scale=1.2]
\draw[bookaxis](-.3,0)--(2.8,0)node[right]{$y$};\draw[bookaxis](0,-.2)--(0,3.8)node[above]{$z$};
\draw[bookcurve](0,1)--(2.5,3.5)node[right]{$-y+z=1$};\draw[bookcurve](1,-.1)--(1,3.5)node[above]{$y=1$};
\fill(1,2)circle(1.8pt)node[right]{$(1,2)$};\draw[bookguide](0,2)--(1,2);\node[left]at(0,2){$2$};\node[below]at(1,0){$1$};
\end{tikzpicture}
\caption{题3.1解图：消元把三元问题化为 $yz$ 平面中两直线求交；唯一交点确定 $y,z$，再回代求 $x$。}\label{fig:hw-3-1}
\end{figure}

改变右侧会平移消元后的约束线，但系数所决定的独立方向不变；行列式非零保证它们仍有唯一交点。
\end{solution}

\item \begin{minipage}[t]{\linewidth}
【基础】分别讨论 $x+y=2,2x+2y=c$ 的解，给出系数矩阵和增广矩阵的秩；写出有解时的解空间方向与一个特解。
\end{minipage}
% 题设配图；与紧随其后的解答分别编号。
\ExpansionDiagram{fig:audit-m-q-085}{比例相同的两条直线可能重合或平行}{
\draw[->](-.4,0)--(4,0)node[right]{$ x $};\draw[->](0,-0.3)--(0,3)node[above]{$ y $};\draw[thick](-.3,2.3)--(2.3,-.3)node[below left]{$x+y=2$};\draw[dashed](.7,2.8)--(3.8,-.3);\node[right]at(3,1.8){$2x+2y=c$};
}

\begin{solution}
图\ref{fig:audit-m-q-085}中两直线法向平行，右端条件决定它们重合还是分离。
图\ref{fig:hw-3-2}区分两约束重合与平行不交。下面的行变换正好识别这两种情形。

系数矩阵与增广矩阵分别为
\[
A=\begin{pmatrix}1&1\\2&2\end{pmatrix},\qquad
(A\mid\vct b)=\begin{pmatrix}1&1&2\\2&2&c\end{pmatrix}.
\]
第二行减第一行的两倍后，增广矩阵变为
\[
\begin{pmatrix}1&1&2\\0&0&c-4\end{pmatrix}.
\]
系数矩阵始终秩为1。若 $c=4$，增广矩阵也秩为1，只剩独立条件 $x+y=2$。以 $y=s$ 为自由参数，解为 $(x,y)=(2,0)+s(-1,1)$。其中 $(2,0)$ 是一个特解，$(-1,1)$ 给出齐次解空间的方向。

\begin{figure}[H]\centering
\begin{tikzpicture}[scale=.9]
\foreach\a in{0,5.8}{\begin{scope}[shift={(\a,0)}]\draw[bookaxis](-.3,0)--(3.7,0)node[right]{$x$};\draw[bookaxis](0,-.3)--(0,3.8)node[above]{$y$};\draw[bookcurve](-.2,2.2)--(2.2,-.2);\end{scope}}
\draw[bookaxis](.5,1.5)--(1,1);\fill(2,0)circle(1.6pt)node[below]{$(2,0)$};\node at(1.5,3.45){$c=4$：两线重合};
\draw[dashed](5.8,3)--(8.8,0);\node at(7.3,3.45){$c=6$：两线平行};
\node[above right]at(6.3,1.5){$x+y=2$};\node[above right]at(8,1.1){$x+y=3$};
\end{tikzpicture}
\caption{题3.2解图：相关的系数只确定同一法向。右侧一致时有一条解直线，右侧不一致时无交点；左图箭头与齐次方向平行。}\label{fig:hw-3-2}
\end{figure}

若 $c\ne4$，增广矩阵秩为2，第二行要求 $0=c-4$，矛盾，因而无解。有解时整条仿射直线不过原点；经过原点的齐次解直线只提供方向，不能代替原方程的解集。

有解时，整条直线组成仿射解集。特解固定直线上的一个点，齐次解给出沿直线移动的方向。
\end{solution}

\item \begin{minipage}[t]{\linewidth}
【基础】将 $(1,1,0),(1,0,1)$ 正交归一化，并用叉积补出右手正交单位基。

\end{minipage}
% 题设配图；与紧随其后的解答分别编号。
\ExpansionDiagram{fig:audit-m-q-086}{Gram--Schmidt投影和第三方向的右手约定}{
\draw[->,thick](0,0)--(3,0)node[right]{$\vct a=(1,1,0)$};
\draw[->,thick](0,0)--(1.5,2.59808)node[above]{$\vct b=(1,0,1)$};
\draw[dashed](1.5,2.59808)--(1.5,0);
\fill(1.5,0)circle(1.5pt)node[below]{$P=\vct a/2$};
\draw(1.3,0)--(1.3,.2)--(1.5,.2);
\draw(.6,0)arc(0:60:.6);
\node at(.9,.35){$60^\circ$};
\node[below]at(1.5,-.8){在$\vct a,\vct b$张成的平面内作图};

}

\begin{solution}
图\ref{fig:audit-m-q-086}中先消去第二向量在第一方向的投影，再归一化；第三方向由有序叉积确定。
投影分解见图\ref{fig:eg-3-3}，空间定向见图\ref{fig:ex-3-3}。先算两单位方向，再用叉积固定第三方向。

令 $\vct a=(1,1,0)$、$\vct b=(1,0,1)$。先归一化第一向量：$\vct e_1=\vct a/|\vct a|=(1,1,0)/\sqrt2$。从第二向量中扣除沿 $\vct e_1$ 的分量，得
\[
\vct b_\perp=\vct b-(\vct b\cdot\vct e_1)\vct e_1
=(1,0,1)-\tfrac12(1,1,0)=(1/2,-1/2,1).
\]
其模为 $\sqrt6/2$，故 $\vct e_2=(1,-1,2)/\sqrt6$。为固定空间定向，取
\[
\vct e_3=\vct e_1\times\vct e_2
=\frac{(2,-2,-2)}{\sqrt{12}}
=\frac{(1,-1,-1)}{\sqrt3}.
\]
直接点乘可得三个模均为1、两两点积均为零；且 $(\vct e_1\times\vct e_2)\cdot\vct e_3=1$，因此是右手正交单位基。单独把第三向量变号仍保持正交，但会改变定向。

\begin{figure}[H]\centering
\begin{tikzpicture}[scale=.95]
\draw[bookaxis](0,0)--(3,0)node[right]{$x$};\draw[bookaxis](0,0)--(-1.4,-1)node[below left]{$y$};\draw[bookaxis](0,0)--(0,2.5)node[above]{$z$};
\draw[bookguide](1.5,0)--(.8,-.5)--(-.7,-.5);
\draw[bookguide](1.5,0)--(1.5,1.4)--(0,1.4);
\draw[bookaxis,thick](0,0)--(.8,-.5)node[below left]{$\vct a$};
\draw[bookaxis,thick](0,0)--(1.5,1.4)node[above right]{$\vct b=(1,0,1)$};
\node[below right]at(0,0){$O$};
\draw[bookaxis,dashed](0,0)--(2.2,-.9)node[below right]{$\vct e_3\parallel(1,-1,-1)$};
\end{tikzpicture}
\caption{题3.3解图：实箭头是输入向量，虚箭头给出由叉积确定的第三正交方向。三维投影不保持图上夹角。}\label{fig:ex-3-3}
\end{figure}

第三方向垂直前两输入向量张成的平面；右手规则固定了符号，单独变号会改成左手基。
\end{solution}

\item \begin{minipage}[t]{\linewidth}
【基础】固定端双振子取 $m=1\,\mathrm{kg},k=4\,\mathrm{N/m}$，只将第一个拉开 $0.10\,\mathrm m$ 由静止释放。求完整位移与初始总能量。

\end{minipage}
% 题设配图；与紧随其后的解答分别编号。
\ExpansionDiagram{fig:audit-m-q-087}{数值初态下的固定端双振子}{
\draw[thick](0,-.5)--(0,.8);\draw[exspring](0,.15)--(1.4,.15);\node[above]at(.7,.4){$ 4 $};\draw[thick](6.7,-.5)--(6.7,.8);\draw[exspring](5.3,.15)--(6.7,.15);\node[above]at(6,.4){$ 4 $};\draw(1.4,-.3)rectangle(2.4,.6);\node at(1.9,.15){$ 1 $};\draw(4.3,-.3)rectangle(5.3,.6);\node at(4.8,.15){$ 1 $};\draw[exspring](2.4,.15)--(4.3,.15);\node[above]at(3.35,.4){$ 4 $};\draw[->](1.4,-.8)--(2.7,-.8)node[right]{$x_1$};\draw[->](4.3,-.8)--(5.6,-.8)node[right]{$x_2$};\node[above]at(3.35,1.2){$x_1(0)=0.10,\quad x_2(0)=0$};
}

\begin{solution}
图\ref{fig:audit-m-q-087}中三弹簧均参与势能，初态应投影到同相、反相两个振型。
图\ref{fig:ex-3-4}把三个弹簧的伸长依次标为 $x_1,x_2-x_1,-x_2$；刚度矩阵的每一项都可由此核对。

两位移都以平衡位置为零、向右为正。中间弹簧的伸长为 $x_2-x_1$，所以
\[
\begin{pmatrix}\ddot x_1\\\ddot x_2\end{pmatrix}
+4\begin{pmatrix}2&-1\\-1&2\end{pmatrix}
\begin{pmatrix}x_1\\x_2\end{pmatrix}=0.
\]
沿 $(1,1)$、$(1,-1)$ 的刚度特征值分别为4、12，角频率为 $2$、$2\sqrt3\,\mathrm{rad/s}$。初始位移满足
\[
\begin{pmatrix}0.10\\0\end{pmatrix}
=0.05\begin{pmatrix}1\\1\end{pmatrix}
+0.05\begin{pmatrix}1\\-1\end{pmatrix}.
\]
每个模态的通解为$q_j=C_j\cos\omega_jt+D_j\sin\omega_jt$。上面的初位移分解给出$C_1=C_2=0.05\,\mathrm m$；初速度为零给出$\omega_jD_j=0$，故两个$D_j$均为零。还原物理坐标，得
\begin{align*}
x_1(t)&=0.05(\cos2t+\cos2\sqrt3t)\,\mathrm m,\\
x_2(t)&=0.05(\cos2t-\cos2\sqrt3t)\,\mathrm m.
\end{align*}
初始时外侧第一根和中间弹簧各形变0.10 m，另一根不变形，故 $E=\tfrac12k(0.10)^2+\tfrac12k(0.10)^2=0.040\,\mathrm J$。代入 $t=0$ 可核对两位移和两初速。中间弹簧储能属于整个系统，不能要求每个物块的某个局部能量单独守恒。

\begin{figure}[H]\centering
\begin{tikzpicture}[scale=1.05]
\draw[thick](0,-.55)--(0,.55) (5.5,-.55)--(5.5,.55);
\draw[exspring](0,0)--(0.95,0);\node[above]at(0.475,.28){$k$};
\draw[thick](0.95,-.4)rectangle(1.95,.4);\node at(1.45,0){$m$};
\draw[bookaxis](0.95,-.85)--(2.3,-.85)node[right]{$x_1$};
\draw[exspring](1.95,0)--(3.4000000000000004,0);\node[above]at(2.6750000000000003,.28){$k$};
\draw[thick](3.4000000000000004,-.4)rectangle(4.4,.4);\node at(3.9000000000000004,0){$m$};
\draw[bookaxis](3.4000000000000004,-.85)--(4.75,-.85)node[right]{$x_2$};
\draw[exspring](4.4,0)--(5.5,0);\node[above]at(4.95,.28){$k$};
\node[below]at(.48,-1.3){$x_1$};\node[below]at(2.67,-1.3){$x_2-x_1$};\node[below]at(4.95,-1.3){$-x_2$};
\end{tikzpicture}
\caption{题3.4解图：上方标刚度，下方依次标三根弹簧的有向伸长；弹性势能对各伸长平方后相加。}\label{fig:ex-3-4}
\end{figure}

初态在图中只有第一根与中间弹簧形变，因而总能量包含两项 $ka^2/2$；完整位移的形状还可与图\ref{fig:eg-3-7}比较。
\end{solution}

\item \begin{minipage}[t]{\linewidth}
【提高】自由双物体取 $m_1=1\,\mathrm{kg},m_2=2\,\mathrm{kg},k=6\,\mathrm{N/m}$，初始位移 $x_1=0.10\,\mathrm m,x_2=0$，速度为零。求两频率和两物体的完整运动。

\end{minipage}
% 题设配图；与紧随其后的解答分别编号。
\ExpansionDiagram{fig:audit-m-q-088}{没有外侧固定弹簧的不等质量双物体}{
\draw(1.4,-.3)rectangle(2.4,.6);\node at(1.9,.15){$ 1 $};\draw(4.3,-.3)rectangle(5.3,.6);\node at(4.8,.15){$ 2 $};\draw[exspring](2.4,.15)--(4.3,.15);\node[above]at(3.35,.4){$ 6 $};\draw[->](1.4,-.8)--(2.7,-.8)node[right]{$x_1$};\draw[->](4.3,-.8)--(5.6,-.8)node[right]{$x_2$};\node[above]at(3.3,1.2){$x_1(0)=0.10,\quad x_2(0)=0$};
}

\begin{solution}
图\ref{fig:audit-m-q-088}中外侧没有固定端，整体平移是零频运动，另一自由度才改变弹簧长度。
图\ref{fig:ex-3-5}只有一根内弹簧，没有固定端，因此先找质心运动和相对运动。

记相对伸长 $\xi=x_2-x_1$，则 $m_1\ddot x_1=k\xi$、$m_2\ddot x_2=-k\xi$。两式相加给出质心匀速，作差给出
\[
\ddot\xi+k\left(\frac1{m_1}+\frac1{m_2}\right)\xi=0.
\]
因此整体平移模频率为零，相对模角频率为 $\sqrt{6(1+1/2)}=3\,\mathrm{rad/s}$。本题初速为零，质心位置
\[
X=\frac{x_1+2x_2}{3}=\frac1{30}\,\mathrm m
\]
保持不变。这也可由$\ddot X=0$两次积分得到：$\dot X=V_C$、$X=X_0+V_Ct$，初态确定$V_C=0,X_0=1/30\,\mathrm m$。相对方程的特征根为$\pm3\ii$，故$\xi=A\cos3t+B\sin3t$；由$\xi(0)=-0.10\,\mathrm m,\dot\xi(0)=3B=0$，得$A=-0.10\,\mathrm m,B=0$。再由$x_1=X-2\xi/3,x_2=X+\xi/3$还原：
\begin{align*}
x_1(t)&=\left(\frac1{30}+\frac1{15}\cos3t\right)\mathrm m,\\
x_2(t)&=\left(\frac1{30}-\frac1{30}\cos3t\right)\mathrm m.
\end{align*}
两物体的振幅比为2，方向相反，使总动量始终为零。总能量为初始弹性势能 $k\xi(0)^2/2=0.030\,\mathrm J$。零频模不代表本题质心位移为零，而代表质心没有回复力；是否匀速移动还由初值决定。

\begin{figure}[H]\centering
\begin{tikzpicture}[scale=1.05]
\draw[thick](0,-.4)rectangle(1,.4) (3.2,-.4)rectangle(4.6,.4);\node at(.5,0){$m_1$};\node at(3.9,0){$m_2$};
\draw[exspring](1,0)--(3.2,0);\node[above]at(2.1,.3){$k$};\draw[bookguide](-.5,-.45)--(5.1,-.45);
\draw[bookaxis](0,-.95)--(1.35,-.95)node[right]{$x_1$};\draw[bookaxis](3.2,-.95)--(4.6,-.95)node[right]{$x_2$};
\node[above]at(2.3,1){两端均不接固定物};
\end{tikzpicture}
\caption{题3.5：仅由一根弹簧相连的自由双物体，水平面光滑。}\label{fig:ex-3-5}
\end{figure}

相对模中两物体位移方向相反、振幅比为2，正好使质量加权位移之和为零；图中的整体平移模则不改变弹簧伸长。
\end{solution}

\item \begin{minipage}[t]{\linewidth}
【提高】固定端三振子链取相同 $m,k$，求三个角频率和归一化模态；解释每个模中的节点和符号。

\end{minipage}
% 题设配图；与紧随其后的解答分别编号。
\ExpansionDiagram{fig:audit-m-q-089}{三物块四弹簧链的固定端边界}{
\draw[thick](0,-.4)--(0,.7);\draw[thick](7,-.4)--(7,.7);\foreach \x/\j in {1.4/1,3.5/2,5.6/3}{\draw(\x-.35,-.2)rectangle(\x+.35,.5);\node at(\x,.15){$m$};\node[below]at(\x,-.5){$x_\j$};}\draw[exspring](0,.15)--(1.05,.15);\draw[exspring](1.75,.15)--(3.15,.15);\draw[exspring](3.85,.15)--(5.25,.15);\draw[exspring](5.95,.15)--(7,.15);\node[above]at(3.5,1){四条弹簧均为$k$};
}

\begin{solution}
图\ref{fig:audit-m-q-089}中两端固定，因此振型必须同时满足两个端部位移为零。
连接图\ref{fig:ex-3-6}决定三对角矩阵；图\ref{fig:chain-modes}则用于检查求出的三个模态符号和节点。

设三个位移为 $x_1,x_2,x_3$，固定端位移为零。每个物块相邻的两根弹簧给出
\[
M=mI_3,\qquad K=k\begin{pmatrix}2&-1&0\\-1&2&-1\\0&-1&2\end{pmatrix}.
\]
令 $\lambda=m\omega^2/k$，沿首行展开特征行列式，得
\[
\det(K/k-\lambda I_3)
=(2-\lambda)[(2-\lambda)^2-2]=0.
\]
三个根为 $2-\sqrt2,2,2+\sqrt2$，所以角频率依次为
\[
\sqrt{\frac km}\sqrt{2-\sqrt2},\qquad
\sqrt{\frac{2k}m},\qquad
\sqrt{\frac km}\sqrt{2+\sqrt2}.
\]
将各根代回首行 $(2-\lambda)u_1-u_2=0$，再用第三行及归一化条件，可取
\[
\vct u_1=\begin{pmatrix}1/2\\1/\sqrt2\\1/2\end{pmatrix},\quad
\vct u_2=\begin{pmatrix}1/\sqrt2\\0\\-1/\sqrt2\end{pmatrix},\quad
\vct u_3=\begin{pmatrix}1/2\\-1/\sqrt2\\1/2\end{pmatrix}.
\]
这里采用欧氏单位长度；如要质量归一化，再各除以 $\sqrt m$。第一模全部同向，第二模中间物块为节点、两端反相，第三模中间物块与两端反相。外侧固定端始终是节点。模态整体变号只改变时间相位，不产生新的模态。

\begin{figure}[H]\centering
\begin{tikzpicture}[scale=.98]
\draw[thick](0,-.55)--(0,.55) (7.95,-.55)--(7.95,.55);
\draw[exspring](0,0)--(0.95,0);\node[above]at(0.475,.28){$k$};
\draw[thick](0.95,-.4)rectangle(1.95,.4);\node at(1.45,0){$m$};
\draw[bookaxis](0.95,-.85)--(2.3,-.85)node[right]{$x_1$};
\draw[exspring](1.95,0)--(3.4000000000000004,0);\node[above]at(2.6750000000000003,.28){$k$};
\draw[thick](3.4000000000000004,-.4)rectangle(4.4,.4);\node at(3.9000000000000004,0){$m$};
\draw[bookaxis](3.4000000000000004,-.85)--(4.75,-.85)node[right]{$x_2$};
\draw[exspring](4.4,0)--(5.8500000000000005,0);\node[above]at(5.125,.28){$k$};
\draw[thick](5.8500000000000005,-.4)rectangle(6.8500000000000005,.4);\node at(6.3500000000000005,0){$m$};
\draw[bookaxis](5.8500000000000005,-.85)--(7.2,-.85)node[right]{$x_3$};
\draw[exspring](6.8500000000000005,0)--(7.95,0);\node[above]at(7.4,.28){$k$};
\end{tikzpicture}
\caption{题3.6：三物体、四弹簧的固定端振子链。}\label{fig:ex-3-6}
\end{figure}

第二模的中间物体是节点，最低模全同向，最高模相邻反相。把这些图形特征逐行代回矩阵，可检查特征向量。
\end{solution}

\item \begin{minipage}[t]{\linewidth}
【提高】对 $U=k(x^2+2\eta xy+y^2)/2$ 讨论稳定性；在 $\eta=1$ 的零方向上分别补加正、负四次项，说明二阶结论的局限。
\end{minipage}
% 题设配图；与紧随其后的解答分别编号。
\ExpansionDiagram{fig:audit-m-q-090}{零二次曲率方向需要继续考察高阶项}{
\draw[->](-.4,0)--(3.3,0)node[right]{$ x $};\draw[->](0,-2.2)--(0,2.4)node[above]{$ y $};\draw[dashed](-1.5,-1.5)--(1.5,1.5)node[above]{$x=y$};\draw[dotted](-1.5,1.5)--(1.5,-1.5)node[below]{$x=-y$};\node[right]at(3.5,.9){$U=\frac{k}{2}(x^2+2\eta xy+y^2)$};
}

\begin{solution}
图\ref{fig:audit-m-q-090}中二次型在临界参数下可能出现零方向，此时须沿该方向保留四次项。
先用同相、反相方向化二次型；图\ref{fig:hw-3-7}专门检查 $\eta=1$ 时的零方向。

假定 $k>0$。用正交坐标 $q_+=(x+y)/\sqrt2,q_-=(x-y)/\sqrt2$，二次势能化为
\[
U=\frac k2\big[(1+\eta)q_+^2+(1-\eta)q_-^2\big].
\]
当 $|\eta|<1$，两系数均正，原点为严格极小点；当 $|\eta|>1$，有一个负系数，沿相应方向离开原点会降低势能，原点不稳定。当 $|\eta|=1$，势能半正定，有一个零方向，线性模型沿该方向没有回复力。

\begin{figure}[H]\centering
\begin{tikzpicture}[x=1.4cm,y=1.5cm]
\draw[bookaxis](-1.5,0)--(1.7,0)node[right]{$s=x-y$};\draw[bookaxis](0,-1.4)--(0,1.5)node[above]{$\Delta U$};
\draw[bookcurve,domain=-1.05:1.05,samples=90]plot(\x,{(\x)^4});\draw[dashed,domain=-1.05:1.05,samples=90]plot(\x,{-(\x)^4});
\node[right]at(1.05,1.05){$+\alpha s^4$};\node[right]at(1.05,-1.05){$-\alpha s^4$};\fill(0,0)circle(1.3pt);
\end{tikzpicture}
\caption{题3.7解图：沿二次势能的零方向检查四阶项。两条曲线在原点二阶导数同为零，局部稳定性却相反；纵轴按 $\alpha$ 缩放。}\label{fig:hw-3-7}
\end{figure}

具体取 $\eta=1$，则 $U=k(x+y)^2/2$，零方向为 $x=-y$。分别补上 $+\alpha(x-y)^4$ 和 $-\alpha(x-y)^4$，$\alpha>0$。前者在所有非零的小位移处都比原点势能高，原点成为严格极小点；后者沿 $x=-y$ 降低，原点不稳定。两种势能在原点的二阶导数完全相同。故出现零特征值时，二阶展开只能找出需要继续检查的方向，必须结合高阶项判断完整系统。

两条曲线在原点的二阶曲率同为零。沿零方向保留四次项后，一条势能升高，另一条降低，稳定性才得以区分。
\end{solution}

\item \begin{minipage}[t]{\linewidth}
【提高】均匀细杆长 $\ell$、质量 $M$，求过一端、与杆夹角 $45^\circ$ 的轴的惯量。分别用直接积分和平行轴定理求解。

\end{minipage}
% 题设配图；与紧随其后的解答分别编号。
\ExpansionDiagram{fig:audit-m-q-091}{过杆端且倾斜四十五度的转轴}{
\draw[very thick](0,0)--(4,0);\fill(0,0)circle(2pt)node[below]{$O$};\draw[->](0,0)--(2.3,2.3)node[above]{$\vct n$};\draw(.7,0)arc(0:45:.7);\node[right]at(.8,.3){$45^\circ$};\draw[dashed](2.5,0)--(1.25,1.25);\node[below]at(2.5,0){$s$};\node[below]at(4,0){$\ell$};
}

\begin{solution}
图\ref{fig:audit-m-q-091}中杆上位置$s$到转轴的垂距为$s\sin45^\circ$，平行轴路线也应使用同一条轴。
图\ref{fig:ex-3-8}先画质元到端点轴的垂距，再画两平行轴的最短距离。两次计算分别使用这两条垂线。

取沿杆的长度坐标 $s\in[0,\ell]$，线密度为 $M/\ell$。质元到所给轴的垂距为 $s\sin45^\circ$，直接积分得
\[
I=\frac M\ell\int_0^\ell s^2\sin^245^\circ\dd s
=\frac{M}{2\ell}\left[\frac{s^3}{3}\right]_0^\ell
=\frac{M\ell^2}{6}.
\]
改用平行轴定理时，先过质心作与原轴平行的轴。用中心长度坐标 $s\in[-\ell/2,\ell/2]$，该轴的惯量为
\[
I_C=\frac M\ell\int_{-\ell/2}^{\ell/2}s^2\sin^245^\circ\dd s
=\frac{M}{2\ell}\left[\frac{s^3}{3}\right]_{-\ell/2}^{\ell/2}
=\frac{M\ell^2}{24}.
\]
两轴间的最短距离为 $d=(\ell/2)\sin45^\circ$，不是 $\ell/2$。因此 $I=I_C+Md^2=M\ell^2/24+M\ell^2/8=M\ell^2/6$。轴倾斜后，两次计算都必须使用到轴的垂距。

\begin{figure}[H]\centering
\begin{tikzpicture}[scale=1.15]
\coordinate(O)at(0,0);\coordinate(C)at(2,0);\coordinate(H)at(1,1);
\draw[line width=1.4pt](O)--(4,0);\fill(C)circle(1.7pt)node[below]{$C$};\fill(O)circle(1.7pt)node[below]{$O$};
\draw[bookaxis](-.3,-.3)--(2.8,2.8)node[above]{$\mathcal A$};\draw[bookguide](1.6,-.4)--(4,2)node[right]{$\mathcal A_C$};
\draw[bookguide](C)--(H);\draw(.86,.86)--(1,.72)--(1.14,.86);\node[right]at(1.5,.7){$d$};
\draw[<->](0,-.55)--(2,-.55)node[midway,below]{$\ell/2$};\draw(.6,0)arc(0:45:.6);\node at(22.5:1){$45^\circ$};
\end{tikzpicture}
\caption{题3.8解图：从质心向端点轴作垂线，得到两平行轴间距 $d=(\ell/2)\sin45^\circ$。平行轴项必须使用这条最短距离。}\label{fig:ex-3-8}
\end{figure}

图中 $OC=\ell/2$ 沿杆，不垂直转轴；只有它在转轴法向的投影 $d$ 才能进入平行轴项。
\end{solution}

\item \begin{minipage}[t]{\linewidth}
【提高】对基矩阵 $P=\begin{pmatrix}1&1\\0&1\end{pmatrix}$，求度量 $G$ 和逆时针 $90^\circ$ 物理旋转在该基下的矩阵 $A'$。验证 $A'^{\mathsf T}GA'=G$。

\end{minipage}
% 题设配图；与紧随其后的解答分别编号。
\ExpansionDiagram{fig:audit-m-q-092}{给定基矩阵的两列是非正交基矢}{
\draw[->](-.4,0)--(4,0)node[right]{$ x $};\draw[->](0,-0.3)--(0,2.6)node[above]{$ y $};\draw[->,thick](0,0)--(2,0)node[above]{$\vct e_1$};\draw[->,thick](0,0)--(2,2)node[above]{$\vct e_2$};\draw[dashed](2,0)--(4,2)--(2,2);\draw[->](1,0)arc(0:90:1);\node[above left]at(.7,.7){物理旋转};
}

\begin{solution}
图\ref{fig:audit-m-q-092}中$\vct e_1=(1,0)$、$\vct e_2=(1,1)$由$P$的两列给出，换基后须同步变换度量。
图\ref{fig:ex-3-9}把真实旋转和新基分解同时画出，可直接检查 $A'$ 的第一列。

物理向量与新坐标满足 $\vct r=P\vct r'$，所以 $|\vct r|^2=\vct r'^{\mathsf T}G\vct r'$，其中
\[
G=P^{\mathsf T}P=\begin{pmatrix}1&1\\1&2\end{pmatrix},\qquad
P^{-1}=\begin{pmatrix}1&-1\\0&1\end{pmatrix}.
\]
在原正交基下，逆时针旋转90度为 $R=\begin{pmatrix}0&-1\\1&0\end{pmatrix}$。由于先用 $P$ 还原物理向量，再旋转，最后用 $P^{-1}$ 写回新坐标，故
\[
A'=P^{-1}RP=\begin{pmatrix}-1&-2\\1&1\end{pmatrix}.
\]
直接验证：
\[
GA'=\begin{pmatrix}0&-1\\1&0\end{pmatrix},\qquad
A'^{\mathsf T}GA'=\begin{pmatrix}1&1\\1&2\end{pmatrix}=G.
\]
而 $A'^{\mathsf T}A'=\begin{pmatrix}2&3\\3&5\end{pmatrix}\ne I$。旋转仍保持真实长度；变化的是非正交坐标下计算长度的式子，不能再用坐标平方的简单和。

\begin{figure}[H]\centering
\begin{tikzpicture}[scale=1.3]
\draw[bookaxis](-.5,0)--(2.8,0)node[right]{$x$};\draw[bookaxis](0,-.4)--(0,2.3)node[above]{$y$};
\draw[bookaxis](0,0)--(1,1)node[above right]{$\vct b_2$};\draw[bookaxis,thick](0,0)--(1,0)node[below]{$\vct b_1$};
\draw[bookaxis,thick](0,0)--(0,1)node[above left]{$R\vct b_1$};\draw[bookaxis,dashed](1,1)--(0,1)node[midway,above]{$-\vct b_1$};
\draw[->](.5,0)arc(0:90:.5);\node[right]at(.65,.4){$90^\circ$};\draw[bookguide](0,0)circle(1);
\end{tikzpicture}
\caption{题3.9解图：旋转后的第一基矢满足 $R\vct b_1=-\vct b_1+\vct b_2$，直接读出 $A'$ 的第一列 $(-1,1)^{\mathsf T}$。}\label{fig:ex-3-9}
\end{figure}

把 $(-1,1)$ 代回度量得到 $1-2+2=1$，与图中单位圆半径相同；按普通平方和却得到2。
\end{solution}

\item \begin{minipage}[t]{\linewidth}
【综合】两个摆球分别由长 $\ell$ 的轻质刚杆悬挂，摆球质量均为 $m$，悬点高度可不同。轻弹簧刚度为 $\kappa$，连接在距各自悬点 $h_1,h_2$ 的位置；悬点位置使静止时弹簧水平且不伸长。小摆角从竖直向下方向同向量起。由动能、势能建立质量与刚度矩阵，求两个频率与模态形状。为什么最低频模不一定等幅？

\end{minipage}
% 题设配图；与紧随其后的解答分别编号。
\ExpansionDiagram{fig:audit-m-q-093}{不同连接高度使弹簧伸长依赖两种摆角权重}{
\fill(0,2.5)circle(2pt)node[above]{$O_1$};\fill(3.7,3)circle(2pt)node[above]{$O_2$};\draw[thick](0,2.5)--(0,0);\draw[thick](3.7,3)--(3.7,.5);\fill(0,0)circle(3pt)node[below]{$m$};\fill(3.7,.5)circle(3pt)node[below]{$m$};\fill(0,1.25)circle(1.6pt);\fill(3.7,1.25)circle(1.6pt);\draw[exspring](0,1.25)--(3.7,1.25);\node[above]at(1.85,1.4){$\kappa$};\draw[<->](-.5,2.5)--(-.5,1.25)node[midway,left]{$h_1$};\draw[<->](4.3,3)--(4.3,1.25)node[midway,right]{$h_2$};\node[right]at(.1,.5){$\ell$};\node[right]at(3.8,.8){$\ell$};\node[below]at(1.85,-.6){静止位置，弹簧水平且为原长};
}

\begin{solution}
图\ref{fig:audit-m-q-093}中弹簧连接点的位移分别含$h_1\theta_1,h_2\theta_2$，不能直接以两摆球位移之差替代。
图\ref{fig:ex-3-10}画出平衡杆位及一阶水平位移，弹簧两端位移之差就是所需的几何约束。

小角下两弹簧连接点的水平位移分别为 $h_1\theta_1,h_2\theta_2$，竖直位移为二阶小量。因弹簧静止时水平且不伸长，求二次势能只需其伸长的一阶项 $\Delta\ell=h_2\theta_2-h_1\theta_1$。于是
\begin{align*}
T&=\tfrac12m\ell^2(\dot\theta_1^2+\dot\theta_2^2),\\
U&=\tfrac12mg\ell(\theta_1^2+\theta_2^2)
+\tfrac12\kappa(h_2\theta_2-h_1\theta_1)^2.
\end{align*}
按 $T=\dot{\vct\theta}^{\mathsf T}M\dot{\vct\theta}/2$、$U=\vct\theta^{\mathsf T}K\vct\theta/2$ 读出
\[
M=m\ell^2I,\qquad
K=mg\ell I+\kappa\begin{pmatrix}h_1^2&-h_1h_2\\-h_1h_2&h_2^2\end{pmatrix}.
\]
设 $h_1^2+h_2^2>0$。弹簧贡献的矩阵乘 $(h_2,h_1)^{\mathsf T}$ 为零，乘 $(h_1,-h_2)^{\mathsf T}$ 为原向量的 $h_1^2+h_2^2$ 倍。因此
\begin{align*}
\omega_-^2&=g/\ell,&\vct u_-&\propto(h_2,h_1),\\
\omega_+^2&=g/\ell+\frac{\kappa(h_1^2+h_2^2)}{m\ell^2},
&\vct u_+&\propto(h_1,-h_2).
\end{align*}
低频模满足 $h_1\theta_1=h_2\theta_2$，连接点水平位移相同。摆长虽相同，连接臂长却可不同，所以摆球的振幅一般不等；$h_1=h_2$ 时才成为通常的等幅同相模。若 $h_1=h_2=0$，弹簧不起耦合作用，两频率相同，任意两个独立方向均可取作模态。

\begin{figure}[H]\centering
\begin{tikzpicture}[scale=1.05]
\draw[thick](-.3,2.8)--(.5,2.8) (4.2,2.3)--(5,2.3);
\draw[bookguide](0,2.8)--(0,-.3) (4.5,2.3)--(4.5,-.8);
\draw[thick](0,2.8)--(.65,0);\draw[thick](4.5,2.3)--(4.9,-.5);\fill(.65,0)circle(2.6pt)node[right]{$m$};\fill(4.9,-.5)circle(2.6pt)node[right]{$m$};
\coordinate(A)at(.325,1.4);\coordinate(B)at(4.628571,1.4);\draw[exspring](A)--(B)node[midway,above]{$\kappa$};
\draw[bookaxis](0,1.05)--(.325,1.05)node[below]{$h_1\theta_1$};\draw[bookaxis](4.5,1.05)--(4.628571,1.05);\node[right]at(4.9,1.1){$h_2\theta_2$};
\draw[<->](-.5,1.4)--(-.5,2.8)node[midway,left]{$h_1$};\draw[<->](5.25,1.4)--(5.25,2.3)node[midway,right]{$h_2$};
\draw(0,2.2)arc(-90:-76.93:.6);\node[right]at(.2,2.15){$\theta_1$};\draw(4.5,1.75)arc(-90:-81.87:.55);\node[right]at(4.72,1.95){$\theta_2$};
\node[below]at(2.4,-1.05){$\Delta\ell=h_2\theta_2-h_1\theta_1$（一阶）};
\end{tikzpicture}
\caption{题3.10解图：虚线为平衡杆位，实线为小角位移的示意。两连接点位移之差给出弹簧伸长；几何图只保留一阶水平位移。}\label{fig:ex-3-10}
\end{figure}

最低模要求连接点水平位移相同，而非摆角相同；图中的不同臂长使 $h_1\theta_1=h_2\theta_2$ 一般对应不等幅摆球运动。
\end{solution}

\item \begin{minipage}[t]{\linewidth}
【选做】求系统 $\dot{\vct y}=\begin{pmatrix}-1&4\\0&-1\end{pmatrix}\vct y$、初值 $(0,1)$ 的解。两个特征值都为负，是否保证每个分量或欧氏长度始终单调减小？
\end{minipage}
% 题设配图；与紧随其后的解答分别编号。
\ExpansionDiagram{fig:audit-m-q-094}{初始状态及矩阵给出的瞬时速度}{
\draw[->](-.4,0)--(4.5,0)node[right]{$ y_1 $};\draw[->](0,-0.3)--(0,2.6)node[above]{$ y_2 $};\fill(0,1)circle(2pt)node[left]{$(0,1)$};\draw[->,thick](0,1)--(3,.25)node[right]{$A\vct y(0)=(4,-1)$};\draw[dashed](0,0)circle(1);
}

\begin{solution}
图\ref{fig:audit-m-q-094}中初速度带有明显的水平分量，特征值的符号并不能单独判断欧氏长度的逐时变化。
图\ref{fig:hw-3-11}比较两个分量及欧氏长度；先精确求解，再检查是否逐时单调。

将方程按分量写开，先解$\dot y_2=-y_2$。乘积分因子$\ee^t$后，$(\ee^ty_2)'=0$，积分得$\ee^ty_2=C$，初值$y_2(0)=1$确定$C=1$。所以$y_2=\ee^{-t}$。代入另一式，
\[
\dot y_1+y_1=4\ee^{-t},\qquad
\frac{\dd}{\dd t}(\ee^t y_1)=4.
\]
从$0$到$t$积分，$\ee^ty_1(t)-y_1(0)=\int_0^t4\dd s=[4s]_0^t=4t$。再用$y_1(0)=0$，得到$y_1=4t\ee^{-t}$。两个特征值均为$-1$，解的两个分量长期趋于零。乘积求导却给出$\dot y_1=4\ee^{-t}-4t\ee^{-t}=4(1-t)\ee^{-t}$，在$0<t<1$为正，第一分量先增长，在$t=1$达到$4/\ee$。

\begin{figure}[H]\centering
\begin{tikzpicture}[x=1.3cm,y=1.7cm]
\draw[bookaxis](0,0)--(5.6,0)node[right]{$t$};\draw[bookaxis](0,0)--(0,1.9);
\draw[bookcurve,domain=0:5,samples=130]plot(\x,{4*\x*exp(-\x)});\draw[dashed,domain=0:5,samples=100]plot(\x,{exp(-\x)});
\draw[dotted,domain=0:5,samples=130]plot(\x,{exp(-\x)*sqrt(16*(\x)^2+1)});
\draw[bookguide](1,0)--(1,{4/exp(1)});\fill(1,{4/exp(1)})circle(1.7pt)node[above right]{$4/\ee$};\node[below]at(1,0){$1$};
\node[align=left]at(3.5,1.6){实线：$y_1$\\虚线：$y_2$\\点线：$|\vct y|$};
\end{tikzpicture}
\caption{题3.11解图：实线$y_1$在$t=1$达到最大值；点线$|\vct y|$先略减小，随后增大，再趋于零。两者的极大时刻不同。}\label{fig:hw-3-11}
\end{figure}

欧氏长度的平方为$L^2=(1+16t^2)\ee^{-2t}$。求导得
\[
\frac{\dd L^2}{\dd t}
=32t\ee^{-2t}-2(1+16t^2)\ee^{-2t}
=-2\ee^{-2t}(16t^2-16t+1).
\]
令括号为零，二次公式给出$t_\pm=(16\pm\sqrt{16^2-4\cdot16})/32=(2\pm\sqrt3)/4$。
因此在$0<t<t_-$时长度减小，在$t_-<t<t_+$时增大，在$t>t_+$时又减小。初始长度为1，$t=1$时为$\sqrt{17}/\ee>1$，说明中间的增长已经超过初始值。图中的初始下降很短，不能仅凭曲线整体形状把它概括为从初始时刻就增长。负实部特征值决定长期衰减；耦合项仍可在有限时间内使某分量或总长度增大。

在图中的 $t=1$ 处，第一分量已增至 $4/\ee$，总长度也超过初值1；长期衰减与短时增长可以同时出现。
\end{solution}

\item \begin{minipage}[t]{\linewidth}
【提高】固定端同质量双振子的驱动力改为只作用于第二物体。求两位移相量，并判断同样的 $\omega^2=2k/m$ 处，哪个物体静止。

\end{minipage}
% 题设配图；与紧随其后的解答分别编号。
\ExpansionDiagram{fig:audit-m-q-095}{外力改到第二物体，连接关系保持不变}{
\draw[thick](0,-.5)--(0,.8);\draw[exspring](0,.15)--(1.4,.15);\node[above]at(.7,.4){$ k $};\draw[thick](6.7,-.5)--(6.7,.8);\draw[exspring](5.3,.15)--(6.7,.15);\node[above]at(6,.4){$ k $};\draw(1.4,-.3)rectangle(2.4,.6);\node at(1.9,.15){$ m $};\draw(4.3,-.3)rectangle(5.3,.6);\node at(4.8,.15){$ m $};\draw[exspring](2.4,.15)--(4.3,.15);\node[above]at(3.35,.4){$ k $};\draw[->](1.4,-.8)--(2.7,-.8)node[right]{$x_1$};\draw[->](4.3,-.8)--(5.6,-.8)node[right]{$x_2$};\draw[->,thick](4.8,1.25)--(6.1,1.25)node[right]{$F_0\cos\omega t$};
}

\begin{solution}
按图\ref{fig:audit-m-q-095}和图\ref{fig:ex-3-12}的连接，回复力矩阵不变，驱动力改为进入第二条方程。求出两个位移相量后，再由各自分子的零点判断哪个物体发生反共振。

取 $x_j=\Rea(X_j\ee^{\ii\omega t})$，令 $d=2k-m\omega^2$。驱动力只在第二物块上，相量方程为
\[
\begin{pmatrix}d&-k\\-k&d\end{pmatrix}
\begin{pmatrix}X_1\\X_2\end{pmatrix}=\begin{pmatrix}0\\F_0\end{pmatrix}.
\]
在 $D=d^2-k^2\ne0$ 时，逆矩阵给出
\[
\begin{pmatrix}X_1\\X_2\end{pmatrix}
=\frac1D\begin{pmatrix}d&k\\k&d\end{pmatrix}
\begin{pmatrix}0\\F_0\end{pmatrix}
=\frac1D\begin{pmatrix}kF_0\\dF_0\end{pmatrix},
\]
即
\[
X_1=\frac{kF_0}{D},\qquad X_2=\frac{dF_0}{D}.
\]
代入 $\omega^2=2k/m$，有 $d=0,D=-k^2$，因此 $X_1=-F_0/k,X_2=0$。这次静止的是受驱动的第二物体；第一物体通过中间弹簧对第二物体施力，逐时抵消作用于第二物体的驱动力。固有频率仍由同一个 $D=0$ 决定，驱动位置改变了响应的分子及反共振所在的输出。以上静止指选定纯谐特解，无阻尼任意初值下还可叠加自由振动。

\begin{figure}[H]\centering
\begin{tikzpicture}[scale=1]
\draw[thick](0,-.55)--(0,.55) (5.5,-.55)--(5.5,.55);
\draw[exspring](0,0)--(0.95,0);\node[above]at(0.475,.28){$k$};
\draw[thick](0.95,-.4)rectangle(1.95,.4);\node at(1.45,0){$m$};
\draw[bookaxis](0.95,-.85)--(2.3,-.85)node[right]{$x_1$};
\draw[exspring](1.95,0)--(3.4000000000000004,0);\node[above]at(2.6750000000000003,.28){$k$};
\draw[thick](3.4000000000000004,-.4)rectangle(4.4,.4);\node at(3.9000000000000004,0){$m$};
\draw[bookaxis](3.4000000000000004,-.85)--(4.75,-.85)node[right]{$x_2$};
\draw[bookaxis,thick](3.9,.8)--(5.35,.8)node[above]{$F\cos\omega t$};
\draw[exspring](4.4,0)--(5.5,0);\node[above]at(4.95,.28){$k$};
\end{tikzpicture}
\caption{题3.12：驱动力改作用于第二物体，连接方式与位移正方向保持不变。}\label{fig:ex-3-12}
\end{figure}

\end{solution}

\item \begin{minipage}[t]{\linewidth}
【提高】对 $M,K$ 的质量归一化模态 $S$，设 $B=\alpha M+\beta K$。推导模态阻尼系数和阻尼比，并给出各模欠阻尼的条件。
\end{minipage}
% 题设配图；与紧随其后的解答分别编号。
\ExpansionDiagram{fig:audit-m-q-096}{各模态在质量内积下作正交投影}{
\draw[->](0,0)--(3,0)node[right]{$\vct s_1$};\draw[->](0,0)--(1.1,2.1)node[above]{$\vct s_2$};\node[right,align=left]at(3.6,1){$S^{\mathsf T}MS=I$\\$B=\alpha M+\beta K$};
}

\begin{solution}
图\ref{fig:audit-m-q-096}中正交性按质量内积定义，阻尼矩阵也须作同一模态变换。
图\ref{fig:hw-3-13}把质量归一化后的每个正常坐标画成独立单位质量振子，用来核对模态系数。

从 $M\ddot{\vct x}+B\dot{\vct x}+K\vct x=\vct f$ 出发，令 $\vct x=S\vct q$，左乘 $S^{\mathsf T}$。质量归一化给出
\[
S^{\mathsf T}MS=I,\qquad
S^{\mathsf T}KS=\diag(\omega_j^2).
\]
故比例阻尼变为
\[
S^{\mathsf T}BS=\alpha I+\beta\diag(\omega_j^2)
=\diag(d_j),\qquad d_j=\alpha+\beta\omega_j^2.
\]
各模独立满足$\ddot q_j+d_j\dot q_j+\omega_j^2q_j=F_j$，$F_j$是$S^{\mathsf T}\vct f$的第$j$分量。对$\omega_j>0$，阻尼比为$\zeta_j=d_j/(2\omega_j)$。齐次特征方程为
\[\lambda^2+d_j\lambda+\omega_j^2=0,
\qquad\left(\lambda+\frac{d_j}{2}\right)^2
=\frac{d_j^2}{4}-\omega_j^2.\]
两根为$-d_j/2\pm\sqrt{d_j^2/4-\omega_j^2}$。要求复根且实部为负，分别给出$d_j^2<4\omega_j^2$及$d_j>0$，合并得到$0<d_j<2\omega_j$。$d_j=0$是无阻尼振动，$d_j=2\omega_j$为临界阻尼，$d_j<0$则存在增长。

\begin{figure}[H]\centering
\begin{tikzpicture}[scale=.95]
\draw[thick](0,-.7)--(0,.8);\draw[exspring](0,.35)--(3,.35)node[midway,above]{$\omega_j^2$};
\draw(0,-.35)--(1,-.35);\draw(1,-.6)--(1,-.1)--(2,-.1);\draw(1,-.6)--(2,-.6);\draw(1.6,-.55)--(1.6,-.15);\draw(1.6,-.35)--(3,-.35);\node[below]at(1.5,-.7){$d_j=\alpha+\beta\omega_j^2$};
\draw[thick](3,-.6)rectangle(4,.65);\node at(3.5,0){$1$};\draw[bookaxis](4,0)--(5.5,0)node[right]{$F_j$};\draw[bookaxis](3,-1.3)--(4.8,-1.3)node[right]{$q_j$};
\end{tikzpicture}
\caption{题3.13解图：质量归一化后，第 $j$ 个正常坐标相当于单位质量、刚度 $\omega_j^2$、阻尼 $d_j$ 的独立振子。图中系数属于模态坐标。}\label{fig:hw-3-13}
\end{figure}

若 $\omega_j=0$，该模变为 $\ddot q_j+\alpha\dot q_j=F_j$，应直接求解，不能套含 $1/\omega_j$ 的阻尼比。也不能仅分别看 $\alpha,\beta$ 的符号，而不检查实际的各个 $d_j$。

图上刚度是 $\omega_j^2$，阻尼是 $d_j$，因此临界条件直接为 $d_j=2\omega_j$；零刚度模应另行处理。
\end{solution}

\item \begin{minipage}[t]{\linewidth}
【提高】两质点各质量 $m$，位置为 $\pm(a,b,0)$。求相对原点的惯量张量、主轴与主惯量；说明零主惯量方向的几何原因。

\end{minipage}
% 题设配图；与紧随其后的解答分别编号。
\ExpansionDiagram{fig:audit-m-q-097}{两点质量关于原点对称，连线也是过原点的轴}{
\draw[->](-.4,0)--(3,0)node[right]{$ x $};\draw[->](0,-2.1)--(0,2.6)node[above]{$ y $};\fill(1.6,1.2)circle(3pt)node[above right]{$m:(a,b,0)$};\fill(-1.6,-1.2)circle(3pt)node[below left]{$m:(-a,-b,0)$};\draw[dashed](-2.2,-1.65)--(2.2,1.65);
}

\begin{solution}
图\ref{fig:audit-m-q-097}中两点的外积相同；沿连线转动时每个质点到轴的距离均为零。
图\ref{fig:ex-3-14}先标出穿过两质点的轴，再在其垂直平面内选另两轴。

点质量位于 $\vct r$ 时的惯量张量为 $m(r^2I_3-\vct r\vct r^{\mathsf T})$。两个位置互为相反数，外积相同，故
\[
I=2m\begin{pmatrix}b^2&-ab&0\\-ab&a^2&0\\0&0&a^2+b^2\end{pmatrix}.
\]
假定 $a^2+b^2>0$，令 $r_0=\sqrt{a^2+b^2}$，取单位轴
\[
\vct n_1=(a,b,0)/r_0,\quad
\vct n_2=(-b,a,0)/r_0,\quad
\vct n_3=(0,0,1).
\]
矩阵分别作用在这三个向量上，得到特征值 $0,2mr_0^2,2mr_0^2$。第一轴沿两质点的连线，各点到轴的距离均为零，所以主惯量为零。另两个主惯量相等，在垂直第一轴的平面内任取一对正交轴也都是主轴。若 $a=b=0$，两个理想质点均在原点，惯量张量整体为零，所有轴等价。

\begin{figure}[H]\centering
\begin{tikzpicture}[scale=.95]
\draw[bookaxis](-2.8,0)--(3,0)node[right]{$x$};\draw[bookaxis](0,-2.1)--(0,2.2)node[above]{$y$};
\draw[bookguide](-2,-1.3)--(2,1.3);\fill(2,1.3)circle(3pt)node[above]{$m,\ (a,b,0)$};\fill(-2,-1.3)circle(3pt)node[below]{$m,\ (-a,-b,0)$};
\draw[bookguide](2,0)--(2,1.3)--(0,1.3);\node[below]at(2,0){$a$};\node[left]at(0,1.3){$b$};
\draw(0,0)circle(.13);\fill(0,0)circle(1pt);\node[above left]at(-.13,.13){$z$};\node[below right]at(0,0){$O$};
\draw[bookaxis,thick](0,0)--(2.55,1.6575)node[right]{$\vct n_1$};\draw[bookaxis,dashed](0,0)--(-1.15,1.76923)node[left]{$\vct n_2$};
\end{tikzpicture}
\caption{题3.14解图：连线方向为零主惯量轴 $\vct n_1$，平面内垂直方向为 $\vct n_2$，第三轴为出纸面的 $z$ 轴。图示取 $a,b>0$。}\label{fig:ex-3-14}
\end{figure}

连线轴上两质点的垂距都为零；另外任意一根垂直轴到各点的距离均为 $\sqrt{a^2+b^2}$，解释了两重主惯量。
\end{solution}

\item \begin{minipage}[t]{\linewidth}
【综合】弱耦合双振子取 $\kappa/k=0.02$。求 $\delta=0,0.02,0.06$ 时精确频率与第二物块包络上界；若要求首阶最大能量转移比例至少为90\%，求 $|\delta|$ 的允许范围。估计用首阶频差预报多次交换时的累积误差。

\end{minipage}
% 题设配图；与紧随其后的解答分别编号。
\ExpansionDiagram{fig:audit-m-q-098}{耦合与失谐共同决定两个振型}{
\draw[thick](0,-.5)--(0,.8);\draw[exspring](0,.15)--(1.4,.15);\node[above]at(.7,.4){$ k(1+\delta) $};\draw[thick](6.7,-.5)--(6.7,.8);\draw[exspring](5.3,.15)--(6.7,.15);\node[above]at(6,.4){$ k(1-\delta) $};\draw(1.4,-.3)rectangle(2.4,.6);\node at(1.9,.15){$ m $};\draw(4.3,-.3)rectangle(5.3,.6);\node at(4.8,.15){$ m $};\draw[exspring](2.4,.15)--(4.3,.15);\node[above]at(3.35,.4){$ \kappa $};\draw[->](1.4,-.8)--(2.7,-.8)node[right]{$x_1$};\draw[->](4.3,-.8)--(5.6,-.8)node[right]{$x_2$};\node[above]at(3.4,1.2){$\kappa/k=0.02$};
}

\begin{solution}
图\ref{fig:audit-m-q-098}中失谐在两端刚度中取相反符号，频率和包络应由同一刚度矩阵求出。
图\ref{fig:ex-3-15}把失谐写成两外侧刚度之差，先读出对角元，再算频差与混合程度。

令 $\omega_0=\sqrt{k/m}$、$\epsilon=\kappa/k=0.02$、$d=\sqrt{\epsilon^2+\delta^2}$。无量纲刚度矩阵的特征值为 $1+\epsilon\pm d$，故
\[
\omega_\pm/\omega_0=\sqrt{1+\epsilon\pm d}.
\]
在 $x_1(0)=a,x_2(0)=0$、初速为零时，谱分解给出
\[
x_2=\frac{a\epsilon}{2d}(\cos\omega_-t-\cos\omega_+t).
\]
用和差化积，两个正弦因子之一描述慢包络，其上界为 $a\epsilon/d$。代入三组参数：
\[
\begin{array}{c|ccc}
\delta&\omega_-/\omega_0&\omega_+/\omega_0&\text{包络上界}\\\hline
0&1&1.0198039&a\\
0.02&0.9958493&1.0238575&a/\sqrt2\\
0.06&0.9781383&1.0407908&a/\sqrt{10}
\end{array}
\]
弱耦合首阶的最大平均能量转移比例为 $\eta\simeq\epsilon^2/(\epsilon^2+\delta^2)$，要求其不小于0.9，给出 $\delta^2\le\epsilon^2/9$，即 $|\delta|\lesssim0.006667$。这是平均能量近似的判据；包络上界来自线性精确解，也不意味着快速振动每次都恰好触及上界。

\begin{figure}[H]\centering
\begin{tikzpicture}[scale=1.05]
\draw[thick](0,-.55)--(0,.55) (5.5,-.55)--(5.5,.55);
\draw[exspring](0,0)--(0.95,0);\node[above]at(0.475,.28){$k(1+\delta)$};
\draw[thick](0.95,-.4)rectangle(1.95,.4);\node at(1.45,0){$m$};
\draw[bookaxis](0.95,-.85)--(2.3,-.85)node[right]{$x_1$};
\draw[exspring](1.95,0)--(3.4000000000000004,0);\node[above]at(2.6750000000000003,.28){$\kappa$};
\draw[thick](3.4000000000000004,-.4)rectangle(4.4,.4);\node at(3.9000000000000004,0){$m$};
\draw[bookaxis](3.4000000000000004,-.85)--(4.75,-.85)node[right]{$x_2$};
\draw[exspring](4.4,0)--(5.5,0);\node[above]at(4.95,.28){$k(1-\delta)$};
\end{tikzpicture}
\caption{题3.15：失谐来自两外侧弹簧刚度不同，中间弹簧满足 $\kappa/k=0.02$。}\label{fig:ex-3-15}
\end{figure}

以 $\delta=0$ 检查长时间误差，精确频差为
\[
\Delta\omega=\omega_0(\sqrt{1+2\epsilon}-1)
=\omega_0(\epsilon-\epsilon^2/2+O(\epsilon^3)).
\]
若仅用$\epsilon\omega_0$，频差误差约为$\epsilon^2\omega_0/2$。经历$n$次单程转移，时间量级为$t\sim n\pi/(\epsilon\omega_0)$，所以两个模之间的相位差$\Delta\omega t$的误差约为$n\pi\epsilon/2$。包络中的慢正弦使用相位$\Delta\omega t/2$，其相位误差为前者的一半，即$n\pi\epsilon/4$。即使耦合很弱，交换次数足够多后仍须使用精确频率。

失谐既增大频差又降低转移包络；可把表中的三组幅值与图\ref{fig:hw-4-7}的包络上界曲线交叉核对。
\end{solution}

\item \begin{minipage}[t]{\linewidth}
【综合】一维势能 $U(q)=\alpha q^2/2+\beta q^4/4$，$\beta>0$，动能为 $m\dot q^2/2$。分别讨论 $\alpha>0,\alpha=0,\alpha<0$ 的平衡位置与稳定性，求稳定平衡附近存在时的非零线性振动频率。
\end{minipage}
% 题设配图；与紧随其后的解答分别编号。
\ExpansionDiagram{fig:audit-m-q-099}{二次项符号可变，而四次项系数为正}{
\draw[->](-.4,0)--(3.2,0)node[right]{$ q $};\draw[->](0,-1.8)--(0,2.6)node[above]{$ U $};\draw[thick,domain=-1.7:1.7,samples=60]plot(\x,{.23*(\x)^4});\draw[dashed,domain=-1.7:1.7]plot(\x,{.45*\x*\x});\draw[dotted,domain=-1.7:1.7]plot(\x,{-.45*\x*\x});\node[right]at(2,2){$\beta q^4/4$};\node[right]at(2,.9){$\alpha q^2/2$};
}

\begin{solution}
图\ref{fig:audit-m-q-099}中虚线表示二次项的两种符号，平衡和稳定性需由两项相加后的总势能判断。
图\ref{fig:hw-3-16}画出三种势能形状；每个平衡点和频率都需在相应曲线的极值处求导。

平衡要求 $U'(q)=q(\alpha+\beta q^2)=0$，曲率为 $U''(q)=\alpha+3\beta q^2$。在平衡位置 $q_e$ 附近令 $q=q_e+\xi$，若曲率为正，线性方程为 $m\ddot\xi+U''(q_e)\xi=0$。

\begin{figure}[H]\centering
\begin{tikzpicture}[x=.9cm,y=1.3cm]
\foreach\s/\a/\title in{0/1/\alpha>0,4/0/\alpha=0,8/-1/\alpha<0}{\begin{scope}[shift={(\s,0)}]
\draw[bookaxis](-1.6,0)--(1.7,0)node[right]{$q$};\draw[bookaxis](0,-.5)--(0,2.1)node[above]{$U$};
\draw[bookcurve,domain=-1.4:1.4,samples=90]plot(\x,{\a*(\x)^2/2+(\x)^4/4});\node at(0,2.5){$\title$};\end{scope}}
\fill(0,0)circle(1.5pt);\fill(4,0)circle(1.5pt);\fill(8,0)circle(1.5pt);\fill(7,-.25)circle(1.5pt);\fill(9,-.25)circle(1.5pt);
\end{tikzpicture}
\caption{题3.16解图：示意取 $\beta=1$、$\alpha=1,0,-1$。正曲率单谷经过平坦四次谷变为双谷；频率由各平衡点处的曲率计算。}\label{fig:hw-3-16}
\end{figure}

当 $\alpha>0$，只有 $q_e=0$，曲率 $\alpha>0$，为稳定极小点，角频率 $\sqrt{\alpha/m}$。当 $\alpha=0$，势能为正四次函数，原点仍为严格极小点，但曲率为零，没有非零线性振动频率；有限振幅的周期必须保留四次势能计算。

当 $\alpha<0$，原点曲率为负，不稳定；另有 $q_\pm=\pm\sqrt{-\alpha/\beta}$。在两点均有 $U''(q_\pm)=\alpha-3\alpha=-2\alpha>0$，故均稳定，小振动角频率同为 $\sqrt{-2\alpha/m}$。频率应在新的平衡点展开得到，不能继续使用原点处的负曲率。

中图虽在原点形成严格极小值，切线及二阶曲率都为零；右图必须移到两个新谷底计算正曲率。
\end{solution}

\item \begin{minipage}[t]{\linewidth}
【综合】主物体 $m$ 由弹簧 $k$ 接固定端，并受 $F\cos\omega t$ 驱动；辅助物体 $\mu$ 仅通过弹簧 $\kappa$ 接主物体。忽略阻尼，求两个位移相量。选择 $\kappa$ 使主物体在指定非零频率 $\omega_d$ 下的纯谐稳态位移为零，求辅助物体位移。说明这一频率与固有频率的关系。

\end{minipage}
% 题设配图；与紧随其后的解答分别编号。
\ExpansionDiagram{fig:audit-m-q-100}{辅助物体只连接主物体，不另接固定端}{
\draw[thick](0,-.5)--(0,.8);\draw[exspring](0,.15)--(1.4,.15);\node[above]at(.7,.4){$ k $};\draw(1.4,-.3)rectangle(2.4,.6);\node at(1.9,.15){$ m $};\draw(4.3,-.3)rectangle(5.3,.6);\node at(4.8,.15){$ \mu $};\draw[exspring](2.4,.15)--(4.3,.15);\node[above]at(3.35,.4){$ \kappa $};\draw[->](1.4,-.8)--(2.7,-.8)node[right]{$x_1$};\draw[->](4.3,-.8)--(5.6,-.8)node[right]{$x_2$};\draw[->,thick](1.9,1.25)--(3.2,1.25)node[right]{$F_0\cos\omega t$};
}

\begin{solution}
图\ref{fig:audit-m-q-100}中辅助物体没有独立的接地弹簧，它受的弹簧力只由两物体的相对位移决定。
图\ref{fig:ex-3-17}的辅助物体不接固定端，只感受两物体之间的相对弹簧力。

以平衡位置为零，位移 $x,y$ 均向右为正。两物体的方程为
\begin{align*}
m\ddot x+(k+\kappa)x-\kappa y&=F\cos\omega t,\\
\mu\ddot y+\kappa y-\kappa x&=0.
\end{align*}
取纯谐响应，记$A=k+\kappa-m\omega^2$、$B=\kappa-\mu\omega^2$，相量方程为$AX-\kappa Y=F$、$-\kappa X+BY=0$。第一式乘$B$、第二式乘$\kappa$后相加，得$(AB-\kappa^2)X=BF$；第一式乘$\kappa$、第二式乘$A$后相加，得$(AB-\kappa^2)Y=\kappa F$。记$D=AB-\kappa^2$，在$D\ne0$时得到
\begin{align*}
D&=(k+\kappa-m\omega^2)(\kappa-\mu\omega^2)-\kappa^2,\\
X&=\frac{F(\kappa-\mu\omega^2)}D,\qquad Y=\frac{F\kappa}D.
\end{align*}
指定非零频率 $\omega_d$，选 $\kappa=\mu\omega_d^2$，则 $D(\omega_d)=-\kappa^2\ne0$，于是 $X=0,Y=-F/\kappa$。主物体不动时，辅助弹簧对它的力为 $\kappa y=-F\cos\omega_dt$，逐时抵消驱动力。

\begin{figure}[H]\centering
\begin{tikzpicture}[scale=1.05]
\draw[thick](0,-.55)--(0,.55);\draw[exspring](0,0)--(1.6,0);\node[above]at(.8,.3){$k$};
\draw[thick](1.6,-.5)rectangle(2.8,.5);\node at(2.2,0){$m$};\draw[exspring](2.8,0)--(4.6,0);\node[above]at(3.7,.3){$\kappa$};
\draw[thick](4.6,-.4)rectangle(5.5,.4);\node at(5.05,0){$\mu$};\draw[bookguide](.2,-.55)--(5.9,-.55);
\draw[bookaxis,thick](2.2,.8)--(3.8,.8)node[above]{$F\cos\omega t$};
\draw[bookaxis](1.6,-1)--(3,-1)node[right]{$x$};\draw[bookaxis](4.6,-1)--(5.9,-1)node[right]{$y$};
\end{tikzpicture}
\caption{题3.17：辅助物体只通过弹簧 $\kappa$ 接主物体，自身不接固定端。}\label{fig:ex-3-17}
\end{figure}

两固有频率由 $D=0$ 决定，吸振频率满足 $D=-\kappa^2\ne0$，故不在其中。任意初态还可叠加不衰减的自由振动，零位移只描述上述纯谐特解。

当主物体静止时，图中辅助弹簧的回复力由 $y$ 决定；调谐后它与外力大小相等、方向相反，形成位移零点。
\end{solution}

\end{enumerate}

\chapter{AI 辅助物理建模与计算实验}
一个振子的周期怎样随振幅变化？两个几乎相同的振子为何有时能交换大部分能量，有时却几乎各振各的？前三章给出了处理这些问题的数学工具。本章把模型变成可运行的实验，用参数变化检查推导，并研究近似失效以后会发生什么。

\section{从一个问题得到可运行的实验}
\subsection{把问题交代完整}
“模拟一个单摆”还没有确定计算任务。摆长、重力加速度、初始角度、初速度、阻尼，以及希望观察的物理量，都会改变程序。先在纸上写出模型，再把这些条件交给能够生成或运行代码的 AI 工具。若工具只能生成代码，可要求输出一个在浏览器中运行的完整 HTML 文件；若使用能执行代码的环境，则同时保存代码和数据。

本章配套的离线实验页包含单摆、共振、耦合、抛体和弹簧拟合五个实验。打开 HTML 文件即可改变参数，无须连接模型接口。它给出可复算的参考实现；学生仍需按正文设计一次参数比较，记录自己的预测与观察。纸面图保留主要关系，便于离开电脑后复习。

一次计算委托应包括模型方程、变量的单位、初值、终止条件和所需输出。还应说明怎样检查结果。例如要求同时给出无阻力极限，通常比单独要求“提高准确率”更具体。

\begin{figure}[htbp]\centering
\begin{tikzpicture}[>=Stealth]
\foreach\x/\t in{0/预测,2.7/模型,5.4/计算,8.1/检验}{\node[draw,minimum width=1.8cm,minimum height=.7cm]at(\x,0){\t};}
\draw[->](.9,0)--(1.8,0);\draw[->](3.6,0)--(4.5,0);\draw[->](6.3,0)--(7.2,0);
\draw[->](8.1,-.35)--(8.1,-1.15)--(2.7,-1.15)--(2.7,-.35);
\node[fill=white]at(5.4,-1.15){检查假设或实现};
\end{tikzpicture}
\caption{计算结果与预测不一致时，回到具体方程和假设检查。不能仅凭结果看起来熟悉就结束实验。}
\label{fig:orig-082}
\end{figure}

\subsection{先预测，再看结果}
每次运行前写两三句可以被结果否定的预测。例如：“无阻尼单摆的机械能应保持不变；初始角度从 $5^\circ$ 增大到 $60^\circ$ 后，周期应增加。”这种预测可以核查。只写“曲线应合理”没有给出判断依据。

运行后保存参数、图和数据。若修改了方程、单位或初值，给新结果另起文件名；否则以后无法知道两幅图的差别来自哪里。AI 生成的解释也需要与输出数据对照。声称“收紧容差后结果不变”，应附上前后两个数值及其差，而不只是一句结论。

\section{简谐振子的数值误差}
\subsection{解析基准与相位误差}
取无量纲方程 $q''+q=0$，$q(0)=1,q'(0)=0$。准确解为 $q=\cos\tau,p=q'=-\sin\tau$，能量为 $\mathcal E=(q^2+p^2)/2=1/2$。请 AI 生成位移曲线、相图和能量曲线，并计算 $0\le\tau\le100$ 上的最大位移误差。

检验至少涉及两个不同的量。能量接近常数，只能约束相图上的半径；沿圆周走得过快或过慢，仍会有相位误差。若程序用 $\widetilde q=\cos[(1+\eta)\tau]$、$\widetilde p=-\sin[(1+\eta)\tau]$，两变量组成的“能量”恒为 $1/2$，却不满足原来的 $q'=p$。能量检查无法单独发现这个错误。

\begin{figure}[H]\centering
\begin{tikzpicture}[scale=1.5]
\draw[bookaxis](-1.4,0)--(1.7,0)node[right]{$q$};\draw[bookaxis](0,-1.4)--(0,1.6)node[above]{$p$};\draw(0,0)circle(1);
\fill(.862319,.506366)circle(1.4pt)node[above right]{准确状态};\fill(.892005,-.452026)circle(1.4pt)node[below right]{错误状态};
\draw[bookguide](0,0)--(.862319,.506366) (0,0)--(.892005,-.452026);\draw[->](30.42205:1.25)arc(30.42205:-26.87373:1.25);\node[right]at(1.25,.03){$1\,\mathrm{rad}$};
\end{tikzpicture}
\caption{题4.1所用的相图检验：在 $\tau=100$，准确与错误状态都在单位圆上，错误状态却沿运动方向超前1 rad。}\label{fig:phase-error}
\end{figure}

\subsection{给 AI 的第一份计算任务}
\begin{quote}\small
求解 $q'=p,p'=-q$，初值为 $(1,0)$，计算到 $\tau=100$。分别绘制 $q(\tau)$、$(q,p)$ 和能量误差。把数值位移与 $\cos\tau$ 比较，报告最大绝对误差。用两组逐渐收紧的精度设置重复计算，列出结果。输出图像注明无量纲变量；保存可以重新运行的完整文件。请说明实际采用的方程和精度设置，不用动画是否平滑判断准确性。
\end{quote}

若使用固定步长，就比较 $h,h/2,h/4$；若使用自适应求解器，就收紧容差，并限制最大内部步长。显示点数只决定保存或绘制多少点，增加显示点数未必改善积分精度。初学时不必先证明算法的阶数，但要能读出关心的物理量是否已经稳定到所需位数。


\begin{example}[用同方向过零事件测量周期]\label{ex:ai-event-period}
对无量纲振子$q''+q=0$，初值$q(0)=1,q'(0)=0$，先在$0\le\tau\le10$上对精确解$q=\cos\tau$按间隔$h$采样。让AI寻找前两次由正到负的过零事件，在线性插值后计算周期。推导事件时刻公式，区分事件定位误差与数值积分误差，并比较$h=0.4,0.2,0.1,0.05$的结果。 见图~\ref{fig:ai-event-period-q}。
\end{example}

\begin{figure}[H]\centering
\begin{tikzpicture}[>=Stealth,line cap=round,line join=round]
\draw[->](-.1,0)--(8.8,0)node[right]{$\tau$};\draw[->](0,-1.35)--(0,1.4)node[above]{$q$};
\draw[thick,domain=0:8.4,samples=160]plot(\x,{cos(57.29578*\x)});
\foreach\x in{1.570796,7.853982}{\fill(\x,0)circle(1.7pt);\draw[dashed](\x,0)--(\x,-1.6);}
\node[above right]at(1.570796,0){$\tau_1$};\node[above right]at(7.853982,0){$\tau_2$};
\draw[<->](1.570796,-1.55)--(7.853982,-1.55)node[midway,below]{$T=\tau_2-\tau_1$};
\end{tikzpicture}
\caption{相邻两次同方向穿越相隔一个完整周期；反方向过零不能混入同一事件序列。}\label{fig:ai-event-period-q}
\end{figure}
\begin{solution}
图~\ref{fig:ai-event-period-q}选取两次同方向过零事件。若只按$q$的正负找全部零点，相邻零点相差半个周期；这里要求$q$由正变负，并同时核对$p=q'<0$，才会将两次向下穿越配成一个完整周期。精确时刻为
\[\tau_1=\frac\pi2,\qquad \tau_2=\frac{5\pi}{2},\qquad T=2\pi.
\]
设一个事件落在$\tau_k,\tau_{k+1}=\tau_k+h$之间，样本满足$q_k>0,q_{k+1}\le0$。通过图~\ref{fig:ai-event-period-s}中$A,B$的直线为
\[L(\tau)=q_k+\frac{q_{k+1}-q_k}{h}(\tau-\tau_k).
\]
令$L(\tau_*)=0$，逐项移项得到
\[\boxed{\tau_*=\tau_k-h\frac{q_k}{q_{k+1}-q_k}}.
\]
分母是相邻样本的差，不能只用初始振幅代替。因为这里穿越斜率不为零，公式在该局部区间有唯一的插值根。

\begin{figure}[H]\centering
\begin{tikzpicture}[>=Stealth,line cap=round,line join=round]
\begin{scope}[x=5cm,y=3cm]
\draw[->](-.46,0)--(.32,0)node[right]{$\tau-\pi/2$};\draw[->](0,-.3)--(0,.46)node[above]{$q$};
\draw[thick,domain=-.42:.27,samples=100]plot(\x,{-sin(57.29578*\x)});
\draw[dashed](-.370796,.362358)--(.029204,-.029200);
\fill(-.370796,.362358)circle(.8pt)node[above left]{$A$};\fill(.029204,-.029200)circle(.8pt)node[below right]{$B$};
\fill(-.0006254,0)circle(.7pt);\draw[dotted](-.0006254,0)--(-.12,-.18)node[left]{$\tau_*$};
\draw[dotted](0,0)--(.12,-.18)node[right]{$\tau_0=\pi/2$};\end{scope}
\end{tikzpicture}
\caption{实线为精确余弦信号，虚线为两样本的割线；两根很接近，图中用引线分开标注。}\label{fig:ai-event-period-s}
\end{figure}

误差可由插值余项估计。若区间内$|q''|\le M$，线性插值满足$|q-L|\le Mh^2/8$；再设真根附近$|q'|\ge c>0$，由中值定理有
\[|\tau_* -\tau_0|\le\frac{Mh^2}{8c}.
\]
若样本本身还带有大小不超过$\delta_q$的积分误差，则事件时间误差另有$O(\delta_q/c)$。所以收紧ODE容差与缩小输出网格，控制的是两类误差；只改变一个参数未必能使周期继续收敛。

下表直接对精确函数$q=\cos\tau$采样，再对前两次向下过零线性插值，不含ODE积分误差。可要求AI先复现这张表，再换用它生成的数值积分样本。
\begin{center}\small
\begin{tabular}{rrrr}\toprule
$h$ & $\tau_{1,*}$ & $\tau_{2,*}$ & $T_*$\\\midrule
0.40 & 1.570170880 & 7.853311746 & 6.283140866 \\
0.20 & 1.570678304 & 7.854102728 & 6.283424424 \\
0.10 & 1.570781990 & 7.853978336 & 6.283196347 \\
0.05 & 1.570797178 & 7.853982918 & 6.283185740 \\
\bottomrule\end{tabular}\end{center}
这里$h=0.4$时的周期误差小于$h=0.2$，却不表示较粗网格的算法更好。每个事件在网格单元中的位置不同，两次事件误差相减时可能抵消。本题真根处$q''=0$，上述$O(h^2)$界也不是最紧的阶数。判断事件处理是否可靠，应分别比较$\tau_{1,*},\tau_{2,*}$与精确值，再查看一组递减步长下的误差变化。

若只把最接近零的网格点当作事件时刻，单次误差一般为$O(h)$；两次取整误差还可能产生周期的锯齿变化。对非线性单摆，可沿用同方向过零事件，但应以周期积分作独立基准，不能把$2\pi$继续当作任意振幅的精确周期。接近分界轨线时周期很长，观察窗口也须覆盖足够多次完整事件。

委托AI编写程序时，应写明过零方向和样本间的插值方法，并区分输出间隔与内部积分步长。程序须保留方程、初值、原始样本及事件时刻表，以便单独复算事件位置。曲线平滑或能量近似守恒，只能说明部分数值性质，周期误差仍须按事件时刻比较。

\end{solution}


\section{实验一：单摆的小角近似能用到哪里}
单摆周期的振幅依赖由解析展开和数值积分分别求得。无阻尼模型给出独立的周期表达式，参数变化再用于区分近似误差与数值误差。

\subsection{模型与周期的独立推导}
取摆长 $\ell$、质量 $m$，摆角 $\theta$ 从竖直向下方向量起。忽略阻力，方程为
\[
\ddot\theta+\frac g\ell\sin\theta=0,\qquad\theta(0)=\theta_0,\quad\dot\theta(0)=0.
\]
以下取$0<\theta_0<\pi$，讨论绕最低点的往复摆动；负初角由$\theta\mapsto-\theta$还原。若$\theta_0=0$，所给零初速度使摆保持静止，小角周期应理解为振幅趋零的极限。令$\tau=t\sqrt{g/\ell}$，方程变成$\theta''+\sin\theta=0$，唯一的非平凡初值参数是$\theta_0$。小角解为$\theta_{\rm lin}=\theta_0\cos\tau$。

\begin{figure}[htbp]\centering
\begin{tikzpicture}[scale=1.1]
\draw[thick](-.7,0)--(.8,0);\draw[bookguide](0,0)--(0,-2.8);\draw[thick](0,0)--(1.3,-2.25);\fill(1.3,-2.25)circle(3pt);\node[right]at(.75,-1.2){$\ell$};\draw[->](0,-.65)arc(-90:-60:.65);\node at(.3,-.87){$\theta$};\draw[bookaxis](1.3,-2.25)--(1.3,-3.1)node[below]{$mg$};\draw[bookaxis](1.3,-2.25)--(.8,-1.38)node[left]{张力};
\end{tikzpicture}
\caption{摆角与受力方向。计算使用弧度，界面可显示角度，但必须在进入三角函数前转换。}
\label{fig:orig-084}
\end{figure}

将方程乘$\theta'$，有
\[\frac{\dd}{\dd\tau}\left(\frac{\theta'^2}{2}-\cos\theta\right)
=\theta'\theta''+\sin\theta\,\theta'=0.\]
用初态积分得到$\theta'^2=2(\cos\theta-\cos\theta_0)$，故无量纲四分之一周期为
\[\frac{\mathcal T}{4}=\int_0^{\theta_0}
\frac{\dd\theta}{\sqrt{2(\cos\theta-\cos\theta_0)}}.\]
令$\sin(\theta/2)=s\sin u$、$s=\sin(\theta_0/2)$，则
\[\dd\theta=\frac{2s\cos u}{\sqrt{1-s^2\sin^2u}}\dd u,\qquad
\sqrt{2(\cos\theta-\cos\theta_0)}=2s\cos u.\]
两端变为$u=0,\pi/2$，约去相同因子，得到
\begin{equation}\label{eq:pendulum-period}
\mathcal T=4\int_0^{\pi/2}\frac{\dd u}{\sqrt{1-s^2\sin^2u}},\qquad
\frac T{T_0}=\frac{\mathcal T}{2\pi},\quad T_0=2\pi\sqrt{\ell/g}.
\end{equation}
被积函数展开为$1+s^2\sin^2u/2+3s^4\sin^4u/8+O(s^6)$。两个所需的角积分为
\[\int_0^{\pi/2}\sin^2u\dd u
=\left[\frac u2-\frac{\sin2u}{4}\right]_0^{\pi/2}=\frac\pi4,\]
\[\int_0^{\pi/2}\sin^4u\dd u
=\left[\frac{3u}{8}-\frac{\sin2u}{4}+\frac{\sin4u}{32}\right]_0^{\pi/2}
=\frac{3\pi}{16}.\]
因此$T/T_0=1+s^2/4+9s^4/64+O(s^6)$。再用
$s=\theta_0/2-\theta_0^3/48+O(\theta_0^5)$，有
\[s^2=\frac{\theta_0^2}{4}-\frac{\theta_0^4}{48}+O(\theta_0^6),\qquad
s^4=\frac{\theta_0^4}{16}+O(\theta_0^6).\]
四阶系数为$-1/192+9/1024=11/3072$，所以
\[
\frac T{T_0}=1+\frac{\theta_0^2}{16}+\frac{11\theta_0^4}{3072}+O(\theta_0^6).
\]
这里 $\theta_0$ 必须用弧度。周期积分和时间积分是两条不同的计算路线，适合互相检查。

\subsection{把实验交给 AI 实现}
\begin{quote}\small
制作无阻尼单摆实验。使用 $\theta'=p,p'=-\sin\theta$ 和给定初始角度、零初速度。提供 $5^\circ$ 至 $120^\circ$ 的角度控制，显示准确模型与小角模型的角度曲线、相图和能量误差。周期由首次向下穿越 $\theta=0$ 的时刻乘4得到，并用独立的周期积分核对。输出角度、周期比和两种精度下的差值表。保留线性模型曲线供比较，说明界面角度如何转换成弧度。
\end{quote}

首次向下过零时刻可用相邻积分步之间的连续插值定位，避免直接从稀疏采样中找峰。动画帧间隔、保存数据的间隔和积分内部步长应分别设置。

\begin{table}[htbp]\centering\caption{单摆周期比与二阶近似；参考值由式\eqref{eq:pendulum-period}计算}
\begin{tabular}{rrr}\toprule
初始角度 & $T/T_0$ & $1+\theta_0^2/16$\\\midrule
$5^\circ$ &1.000476&1.000476\\
$15^\circ$&1.004301&1.004284\\
$30^\circ$&1.017409&1.017135\\
$60^\circ$&1.073182&1.068539\\
$90^\circ$&1.180341&1.154213\\
$120^\circ$&1.372881&1.274156\\\bottomrule
\end{tabular}\end{table}

\begin{figure}[htbp]\centering
\begin{tikzpicture}[x=.065cm,y=5cm]
\draw[bookaxis](0,0)--(130,0)node[right]{$\theta_0/(^\circ)$};\draw[bookaxis](0,0)--(0,.44)node[above]{$T/T_0-1$};
\draw[bookcurve]plot coordinates{(0,0)(5,.000476)(15,.004301)(30,.017409)(60,.073182)(90,.180341)(120,.372881)};
\draw[dashed,domain=0:120,samples=100]plot(\x,{(\x*pi/180)^2/16});
\foreach\x in{30,60,90,120}{\draw(\x,.005)--(\x,-.005)node[below]{$\x$};}
\node at(80,.37){实线：参考值连线};\node at(85,.1){虚线：二阶近似};
\end{tikzpicture}
\caption{周期的振幅依赖。比较曲线时应检查具体误差，图上接近不等于达到某个精度要求。}
\label{fig:orig-085}
\end{figure}

\subsection{周期误差与振幅范围}
若允许周期相对误差约1\%，首阶估计 $\theta_0^2/16\lesssim0.01$ 给出 $\theta_0\lesssim0.4\,\mathrm{rad}$，约 $23^\circ$。若要求几十个周期之后摆角仍与线性解接近，允许角度还要更小，因为相位误差会累积。

周期误差、最大角度误差与能量漂移是不同的误差指标，同为1\%也不能互换。相位误差的累积见习题4.2。

\section{实验二：共振峰为什么不在同一处}
同一受迫振子的位移、速度和输入功率具有不同频率响应。比较峰位之前，应先定义观察量，并让暂态衰减到所要求的精度。

\subsection{同时观察位移、速度与功率}
取 $q''+2\zeta q'+q=\cos\Omega\tau$，位移幅值为 $A$，速度幅值为 $V=\Omega A$。平均输入功率与平均耗散功率相等，为 $\overline P=\zeta\Omega^2A^2$。为在同一图上比较，可分别除以各自的最大值；但确定峰值前，须保留各量原始定义。

\begin{quote}\small
生成受迫振动实验，阻尼比可调。分别画出稳态位移幅值、速度幅值、平均耗散功率和位移滞后相位随频率比的曲线。在选定频率处再积分零初值的运动，显示瞬态到稳态的过程。把有限时间估计的振幅与解析稳态幅值比较。相位采用 atan2 计算；请明确测量窗口，并通过延长等待时间检查瞬态是否足够小。
\end{quote}

\begin{figure}[htbp]\centering
\begin{tikzpicture}[x=4cm,y=2cm]
\draw[bookaxis](0,0)--(2.1,0)node[right]{$\Omega$};\draw[bookaxis](0,0)--(0,1.25)node[above]{归一化响应};
\draw[bookcurve,domain=0:2,samples=150]plot(\x,{sqrt(.5904)/sqrt((1-\x*\x)^2+.72*\x*\x)});
\draw[dashed,domain=0:2,samples=150]plot(\x,{sqrt(.72)*\x/sqrt((1-\x*\x)^2+.72*\x*\x)});
\draw[dotted,thick,domain=0:2,samples=150]plot(\x,{.72*\x*\x/((1-\x*\x)^2+.72*\x*\x)});
\draw[bookguide](.8,0)--(.8,1.08);\draw[bookguide](1,0)--(1,1.08);\node[below]at(.8,0){$0.8$};\node[below]at(1.1,0){$1$};
\node[anchor=west]at(1.3,1.55){实线：位移};\node[anchor=west]at(1.3,1.3){虚线：速度};\node[anchor=west]at(1.3,1.05){点线：功率};
\end{tikzpicture}
\caption{$\zeta=\sqrt{0.18}$ 时分别归一化的三种响应。位移峰在0.8，速度与功率峰在1。}
\label{fig:orig-086}
\end{figure}

\subsection{等待时间也是实验条件}
欠阻尼齐次项含因子 $\ee^{-\zeta\tau}$。若希望瞬态包络相对于其初值降到 $10^{-3}$，需等待 $\tau\gtrsim\ln1000/\zeta$。它相对于稳态的占比还取决于初始瞬态系数，不能只凭这个时间估计就宣称误差一定小于千分之一。

\begin{figure}[htbp]\centering
\begin{tikzpicture}[x=.13cm,y=1.2cm]
\draw[bookaxis](0,0)--(65,0)node[right]{$\tau$};\draw[bookaxis](0,-1.2)--(0,1.4);
\draw[bookcurve,domain=0:60,samples=450]plot(\x,{exp(-.08*\x)*cos(deg(\x))});
\draw[dashed,domain=0:60,samples=100]plot(\x,{exp(-.08*\x)});\draw[dashed,domain=0:60,samples=100]plot(\x,{-exp(-.08*\x)});
\node at(43,1){瞬态及其包络};
\end{tikzpicture}
\caption{等待时间不足时，测得的峰值还带有初始条件的影响。图示衰减振动，包络按 $\ee^{-0.08\tau}$ 衰减。}
\label{fig:orig-087}
\end{figure}

以 $\zeta=0.1$ 为例，位移峰频率比为0.989949，功率峰为1，半功率点约为0.904988和1.104988。若程序报告四者完全相同，应检查测量对象和频率扫描分辨率。先用粗扫描找峰，再在附近细化；一个稀疏网格的最大点只给出峰的位置区间。

\section{实验三：弱耦合与失谐}
耦合与失谐共同决定能量交换的幅度和时间尺度。模态频差给出所需观察时间，过短的时间窗口可能掩盖已经存在的能量交换。

\subsection{用时间尺度选择观察窗口}
沿用第三章，取 $\tau=t\sqrt{k/m}$、$q_j=x_j/a$，则
\[
\vct q''+\begin{pmatrix}1+\epsilon+\delta&-\epsilon\\-\epsilon&1+\epsilon-\delta\end{pmatrix}\vct q=0,
\quad\vct q(0)=(1,0)^{\mathsf T},\quad\vct q'(0)=0.
\]
参数 $\epsilon=\kappa/k$，$\delta$ 为左右刚度的不对称程度。观察窗口至少应覆盖一次能量交换，并同时显示两种时间尺度。

\begin{quote}\small
建立双振子实验，控制参数为 epsilon 和 delta。显示两个位移、第二振子的包络、各外侧振子能量、中间弹簧能量及总能量漂移。先用矩阵正常模给出解析运动，再独立积分方程并比较。观察时间按精确频差设置，至少覆盖两个包络峰。保持 epsilon=0.02，分别比较 delta=0、0.02、0.06。请把精确关系与弱耦合近似标明。
\end{quote}

\begin{table}[htbp]\centering\caption{$\epsilon=0.02$ 时失谐的影响；时间为无量纲时间}
\begin{tabular}{rrrr}\toprule
$\delta$ & 包络最大值 & 首次包络峰时刻 & 首阶转移比例\\\midrule
0 &1.000000&158.635&1.00\\
0.02&0.707107&112.167&0.50\\
0.06&0.316228&50.143&0.10\\\bottomrule
\end{tabular}\end{table}

\begin{figure}[htbp]\centering
\begin{tikzpicture}[x=2cm,y=2.5cm]
\draw[bookaxis](0,0)--(4.2,0)node[right]{$|\delta|/\epsilon$};\draw[bookaxis](0,0)--(0,1.25)node[above]{$\eta$};
\draw[bookcurve,domain=0:4,samples=100]plot(\x,{1/(1+\x*\x)});\foreach\x/\y in{0/1,1/.5,3/.1}{\fill(\x,\y)circle(1.5pt);}
\node[below]at(1,0){$1$};\node[below]at(3,0){$3$};\node[left]at(0,.5){$0.5$};
\end{tikzpicture}
\caption{首阶最大平均能量转移比例。决定抑制程度的是失谐与耦合之比。}
\label{fig:orig-088}
\end{figure}

\subsection{一个容易误判的结果}
失谐增大后，表中的首次包络峰反而提前。这与转移变弱并不矛盾：频差增大使包络变化更快，而模态混合减弱使包络幅度变小。时间和幅度反映不同的参数关系。

检查能量时应使用完整势能
\[
\mathcal E=\tfrac12(p_1^2+p_2^2)+\tfrac12(1+\delta)q_1^2+
\tfrac12(1-\delta)q_2^2+\tfrac12\epsilon(q_2-q_1)^2.
\]
计算总能量时须保留最后一项所表示的中间弹簧势能。还可将 $\epsilon$ 减半，比较解析频差给出的交换时间与数值结果。能量能否大部分转移取决于失谐与耦合之比，仅凭两个振子的固有频率接近不能作出判断。

\section{实验四：有阻力抛体与最佳发射角}
二次阻力下，最佳发射角随无量纲阻力参数改变。射程由落地事件确定，轨迹积分与角度寻优分别控制误差。

\subsection{先消去无关的量纲参数}
考虑重力和二次阻力 $\vct F_d=-c|\vct v|\vct v$。取 $v_0/g$ 为时间尺度，$v_0^2/g$ 为长度尺度，写无量纲位置 $(X,Y)$、速度 $(U,V)$，得
\begin{align}
X'&=U,&Y'&=V,\\
U'&=-\beta\sqrt{U^2+V^2}\,U,&V'&=-1-\beta\sqrt{U^2+V^2}\,V,
\end{align}
其中 $\beta=cv_0^2/(mg)$。从地面以角度 $\alpha$ 发射，初值为 $(0,0,\cos\alpha,\sin\alpha)$。只记录向下穿越 $Y=0$ 的事件。

\begin{quote}\small
生成二次阻力抛体实验，使用正文无量纲方程，角度范围5至85度，beta范围0至3。轨迹在首次向下穿越地面时终止；利用事件定位给出落地时间和射程。画射程随角度的曲线，先粗扫描，再在最大值附近细化。beta=0时与 sin(2 alpha) 比较。报告两组精度下的射程与最佳角度，并输出完整数据。
\end{quote}

\subsection{调试一份有问题的结果}
若所有射程都是零，先检查程序是否把 $t=0,Y=0$ 当作落地。若运动在最高点停止，检查终止条件是否误写成 $V=0$。若二次阻力被写为 $(-\beta U|U|,-\beta V|V|)$，它描述的是另一种分量模型，不能拿来代替沿总速度反向的阻力。定位错误时应逐项核对方程。

无量纲能量 $\mathcal E=(U^2+V^2)/2+Y$ 满足
\[
\mathcal E'=-\beta(U^2+V^2)^{3/2}.
\]
应比较能量变化与右侧的时间积分。只观察能量下降，还不足以确认阻力的方向和大小写对了。

\begin{table}[htbp]\centering\caption{同高发射与落地、二次阻力模型的参考结果}
\begin{tabular}{rrrr}\toprule
$\beta$ & 最佳角度 & 最大无量纲射程 & $45^\circ$ 射程\\\midrule
0&$45.000^\circ$&1.000000&1.000000\\
0.2&$43.912^\circ$&0.868607&0.868026\\
1&$40.978^\circ$&0.593607&0.589012\\
3&$37.299^\circ$&0.357710&0.349059\\\bottomrule
\end{tabular}\end{table}

\begin{figure}[htbp]\centering
\begin{tikzpicture}[x=.11cm,y=1.7cm]
\draw[bookaxis](0,0)--(90,0)node[right]{$\alpha/(^\circ)$};\draw[bookaxis](0,0)--(0,1.3)node[above]{$\mathcal R$};
\draw[bookcurve,domain=0:90,samples=100]plot(\x,{sin(2*\x)});\draw[bookguide](45,0)--(45,1);\node[below]at(45,0){$45$};\node at(68,1.1){$\beta=0$};
\end{tikzpicture}
\caption{无阻力射程曲线是程序的基准。加入阻力后，应重新计算曲线及其最大值，不能只缩小纵轴。}
\label{fig:orig-089}
\end{figure}

这组数值表明，在所取模型与参数下，最佳角度随阻力增强而下降。它们不构成对所有阻力规律、风场和起落高度的证明。无阻力高处发射的图解推导与极限检验见习题4.10。

\section{实验五：从数据识别模型遗漏}
拟合的参数依赖所假设的模型。先写出测量量、单位和误差，再比较残差与参数不确定度，判断数据是否支持增加新的物理项。

\subsection{一组带单位的模拟测量}
测量弹簧伸长量 $x$ 与拉力 $F$。配套 CSV 中的七组数值是为本实验构造的模拟数据，不是真实实验记录。设各拉力测量的独立标准差均为 $0.02\,\mathrm N$，伸长量误差暂忽略：
\begin{center}\small
\begin{tabular}{rrrrrrrr}\toprule
$x/\mathrm m$&0&0.01&0.02&0.03&0.04&0.05&0.06\\
$F/\mathrm N$&0.20&1.01&1.83&2.70&3.65&4.71&5.87\\\bottomrule
\end{tabular}\end{center}
先尝试$F=b+kx$，允许仪器存在零点偏置$b$。残差为$r_i=F_i-b-kx_i$，目标为$\chi^2=\sum_i r_i^2/\sigma_i^2$。本题等权$\sigma_i=\sigma$，由$\partial_b r_i=-1$、$\partial_k r_i=-x_i$，
\[\frac{\partial\chi^2}{\partial b}
=-\frac2{\sigma^2}\sum_i r_i,\qquad
\frac{\partial\chi^2}{\partial k}
=-\frac2{\sigma^2}\sum_i x_ir_i.\]
令两偏导为零，得
\[
\sum_i r_i=0,\qquad\sum_i x_ir_i=0.
\]
第一式给出$b=\bar F-k\bar x$。代入第二式，
\[\sum_i x_i(F_i-\bar F)-k\sum_i x_i(x_i-\bar x)=0.\]
用$\sum_i(F_i-\bar F)=\sum_i(x_i-\bar x)=0$将$x_i$改为$x_i-\bar x$，得到
\[
k=\frac{\sum_i(x_i-\bar x)(F_i-\bar F)}{\sum_i(x_i-\bar x)^2},\qquad b=\bar F-k\bar x.
\]
这些关系可用来检查 AI 给出的参数与残差是否彼此一致。

\subsection{拟合之后要看什么}
\begin{quote}\small
读取配套弹簧数据，核对列名和单位。先拟合 $F=b+kx$，再拟合 $F=b+kx+cx^3$；给出参数单位、预测值、残差和按0.02 N归一化的残差平方和。画两种模型的残差图。为改善尺度，允许先使用 $u=x/0.06$，但最终参数换回 SI 单位。区分模拟数据和真实测量，不把残差较小直接解释成物理模型已被证明。提出一个能区分两模型的后续测量方案。
\end{quote}

线性拟合给出 $b=0.04250\,\mathrm N$、$k=93.6786\,\mathrm{N/m}$。两端残差为正，中段多为负，且明显超过设定测量误差。加入三次项后，得到
\[
b=0.20300\,\mathrm N,\quad k=79.6863\,\mathrm{N/m},\quad
c=4115.38\,\mathrm{N/m^3}.
\]
本组数据的缩放设计矩阵、分数消元及逐点残差见图\ref{fig:ex-4-9}所配习题解答，可据此复算上述参数。归一化残差平方和由297.634降为0.400；自由度分别为5和4。后一数值偏小，说明这组构造数据的离散程度小于设定误差所对应的尺度，其解释依赖于测量不确定度的估计。拟合优度只描述给定数据和误差假设下的一致性。

\begin{figure}[htbp]\centering
\begin{tikzpicture}[x=1.4cm,y=8cm]
\draw[bookaxis](0,0)--(6.7,0)node[right]{$x/\mathrm{cm}$};\draw[bookaxis](0,-.2)--(0,.27)node[above]{$r/\mathrm N$};
\draw[dotted](0,.02)--(6,.02);\draw[dotted](0,-.02)--(6,-.02);
\draw[bookcurve]plot coordinates{(0,.1575)(1,.030714)(2,-.086071)(3,-.152857)(4,-.139643)(5,-.016429)(6,.206786)};
\draw[dashed]plot coordinates{(0,-.003)(1,.006022)(2,.000352)(3,-.004703)(4,-.003835)(5,.008264)(6,-.003099)};
\node at(3,.22){实线：线性；虚线：含三次项};
\end{tikzpicture}
\caption{残差图把原始曲线上不易发现的系统偏差显露出来。点线表示 $\pm0.02\,\mathrm N$，不是拟合参数的置信区间。}
\label{fig:orig-090}
\end{figure}

\subsection{参数可信度与后续测量}
在线性参数模型$\vct F=A\vct\vartheta+\vct\varepsilon$中，列向量可取$1,x,x^3$。设$A$满列秩，测量误差均值为零且相互独立，已知协方差为$\Sigma=\diag(\sigma_i^2)$，权矩阵$W=\Sigma^{-1}$。令$\widetilde A=W^{1/2}A$、$\widetilde{\vct F}=W^{1/2}\vct F$，加权目标就成为普通的长度平方
\[
\chi^2=|\widetilde{\vct F}-\widetilde A\vct\vartheta|^2.
\]
按第三章QR分解的正交残差证明，极小值处的残差垂直于$\widetilde A$的每一列，所以
\[
\begin{aligned}
\widetilde A^{\mathsf T}(\widetilde{\vct F}-\widetilde A\widehat{\vct\vartheta})&=0,\\
A^{\mathsf T}WA\widehat{\vct\vartheta}&=A^{\mathsf T}W\vct F.
\end{aligned}
\]
$W$正定且$A$满列秩，使$A^{\mathsf T}WA$可逆；解出参数，再代入数据模型，得到
\[\widehat{\vct\vartheta}=(A^{\mathsf T}WA)^{-1}A^{\mathsf T}W\vct F,\qquad
\widehat{\vct\vartheta}-\vct\vartheta=(A^{\mathsf T}WA)^{-1}A^{\mathsf T}W\vct\varepsilon.\]
对误差作协方差变换，
\[\begin{aligned}
\operatorname{Cov}(\widehat{\vct\vartheta})
&=(A^{\mathsf T}WA)^{-1}A^{\mathsf T}W\Sigma WA(A^{\mathsf T}WA)^{-1}\\
&=(A^{\mathsf T}WA)^{-1},
\end{aligned}\]
最后一步使用$W\Sigma W=W$。这是已给绝对测量方差时的结果。若只知道相对权重$W_0$，并假设$\Sigma=\sigma^2W_0^{-1}$，则协方差还应乘未知的$\sigma^2$；以残差估计时可取$\widehat\sigma^2=\vct r^{\mathsf T}W_0\vct r/(N-p)$，$p$为拟合参数个数。两种方差约定须分清。实际计算宜用正交分解，公式中的逆矩阵用于说明误差传播。

后续可在安全的更大伸长量处增加测量，或测量正、负伸长量来区分偏置与奇次非线性。新测量应独立于拟合所用数据。把多项式阶数不断升高，总能减小训练残差，却未必改善新位置的预测。

\section{选读：采样让两个信号看起来相同}
对采样时刻 $t_n=n/f_s$，频率相差整数倍 $f_s$ 的复指数具有相同样本。取 $f_s=100\,\mathrm{Hz}$，则
\[
\cos(2\pi70t_n)=\cos(2\pi30t_n),\qquad
\sin(2\pi70t_n)=-\sin(2\pi30t_n).
\]
请 AI 画出连续曲线和采样点，再分别对样本求频谱。仅凭这些点，程序无法判断连续信号来自30 Hz还是70 Hz；需要额外的频带信息或改变采样条件。

\begin{figure}[htbp]\centering
\begin{tikzpicture}[x=7cm,y=1cm]
\draw[bookaxis](0,0)--(1.12,0)node[right]{$t/\mathrm s$};\draw[bookaxis](0,-1.3)--(0,1.4);
\draw[bookcurve,domain=0:1,samples=120]plot(\x,{cos(360*\x)});\draw[dashed,domain=0:1,samples=300]plot(\x,{cos(1800*\x)});
\foreach\n in{0,...,4}{\fill({\n/4},{cos(90*\n)})circle(2pt);}
\end{tikzpicture}
\caption{用较低频率示意混叠：1 Hz实线与5 Hz虚线在4 Hz采样时取值相同。增加画图曲线的细腻程度不能补回缺失的信息。}
\label{fig:orig-091}
\end{figure}

记录时间为 $T_{\rm rec}$ 时，常用离散频率格点间隔为 $1/T_{\rm rec}$。零填充可以使频谱显示更密，却没有延长真实观测。若有两个很接近的频率，应先考虑增加观测时间。报告频谱时还要注明幅值归一化、采样率、记录长度和是否使用窗函数。

\section{事件定位与计算步长}
\subsection{积分到某个时刻与积分到某个事件不同}
抛体落地、单摆通过平衡位置和粒子越过界面，停止条件都是一个状态函数$g(t,\vct y)=0$。数值程序通常先在相邻时间节点发现符号变化，再定位中间的事件。即使节点上的状态很准确，直接把第一个负高度节点当作落地时刻，仍会引入量级为步长的时间误差。

事件条件还应包括方向。例如抛体从$y=0$向上发射，初始时刻已经满足$y=0$，但落地事件要求从正高度向负高度穿越，且$v_y<0$。只写“遇到$y=0$便结束”，程序可能在发射瞬间停止。对擦边接触，$g$可在零点两侧同号，单凭符号变化又可能漏掉事件。

记真事件时刻为$t_0$，数值状态使事件函数增加误差$\delta g$，求出的根为$t_0+\delta t$。在$t_0$处展开，并用$g(t_0,\vct y(t_0))=0$，一阶关系为
\[
0\simeq g(t_0,\vct y(t_0))+\dot g(t_0)\delta t+\delta g
=\dot g(t_0)\delta t+\delta g.
\]
这里$\dot g$是沿轨迹的总时间导数。若它不为零，便可解出
\[
 \delta t\simeq-\frac{\delta g}{\dot g}.
\]
因此接近切向穿越时，时间误差会放大。把积分相对误差设得很小，不一定能保证事件时间有同样精度；还要检查$\dot g$的大小与事件定位算法。

\begin{example}线性插值事件的独立误差
竖直抛体满足$y(t)=v_0t-gt^2/2$，$v_0,g>0$。真正落地时刻为$T=2v_0/g$。程序已精确得到相邻节点$t_-=T-a$、$t_+=T+b$的高度，其中$0<a<T$、$b>0$、$a+b=h$，并用两点直线与$y=0$的交点估计落地时刻$\widehat T$。求事件时间误差。
\ExpansionDiagram{fig:dep-ai-event-q}{节点高度准确，弦线的零点仍偏离曲线零点}{
\draw[->](0,0)--(6,0)node[right]{$t$};\draw[->](0,-1.3)--(0,2.4)node[above]{$y$};\draw[thick,domain=.3:5.2,samples=80]plot(\x,{1.8-.11*(\x+.1)^2});\fill(2.5,1.0564)circle(1.5pt)node[above]{$t_-$};\fill(4.8,-.8411)circle(1.5pt)node[below]{$t_+$};\draw[dashed](2.5,1.0564)--(4.8,-.8411);\node[below]at(3.945,0){$T$};}
\end{example}
\begin{solution}
图\ref{fig:dep-ai-event-q}中曲线下凹，弦线在真正零点处位于横轴下方，所以线性估计会提前。直接代入两节点，
\[
 y(t_-)=v_0a-\frac{ga^2}{2},\qquad
 y(t_+)=-v_0b-\frac{gb^2}{2}.
\]
记两节点高度为$y_-,y_+$，弦线斜率及其在$T$处的值分别为
\[
s=\frac{y_+-y_-}{a+b}=-v_0-\frac{g(b-a)}2,\qquad
L(T)=\frac{b y_-+a y_+}{a+b}=-\frac{gab}{2}.
\]
由$0=L(T)+s(\widehat T-T)$解出
\[
 \widehat T-T=-\frac{gab}{2v_0+g(b-a)}.
\]
只要$h\ll v_0/g$，分母保持远离零，首项为$-gab/(2v_0)$。由$ab\leq h^2/4$可见误差为$O(h^2)$，与节点积分器是几阶无关。
\ExpansionDiagram{fig:dep-ai-event-a}{相同网格宽度内，零点相对网格的位置改变插值误差}{
\draw[->](0,0)--(5.5,0)node[right]{$a/h$};\draw[->](0,0)--(0,2.7)node[above]{首项误差的绝对值};\draw[thick,domain=0:5,samples=70]plot(\x,{.32*\x*(5-\x)});\node[below]at(0,0){0};\node[below]at(2.5,0){$1/2$};\node[below]at(5,0){1};}
图\ref{fig:dep-ai-event-a}的形状来自$ab=h^2(a/h)(1-a/h)$。步长减半时零点在新网格中的位置也改变，实际误差比未必恰为4。要区分积分误差和事件插值误差，可以先用解析解提供精确节点值，再单独研究插值方法；然后固定事件算法，比较积分步长。这两组计算分别改变一个因素。
\end{solution}

线性阻力竖直运动提供另一个基准：$\dot v=-g-\kappa v$、$y(0)=0,v(0)=v_0$。令$w=v+g/\kappa$，方程成为$\dot w+\kappa w=0$；乘$\ee^{\kappa t}$并积分，得
$\ee^{\kappa t}w=w(0)=v_0+g/\kappa$。所以
\[v=\left(v_0+\frac g\kappa\right)\ee^{-\kappa t}-\frac g\kappa.\]
再积分高度，
\[y=\int_0^tv(s)\dd s
=\left(v_0+\frac g\kappa\right)
\left[-\frac{\ee^{-\kappa s}}\kappa\right]_0^t
-\frac g\kappa[s]_0^t,\]
即
\[
 y(t)=\frac{v_0+g/\kappa}{\kappa}(1-\ee^{-\kappa t})-\frac{gt}{\kappa}.
\]
取$v_0=10\,\mathrm{m/s}$、$g=10\,\mathrm{m/s^2}$、$\kappa=0.4\,\mathrm{s^{-1}}$，排除$t=0$的根后，落地时间为$1.7885391259\,\mathrm s$。用精确节点高度作线性事件插值可得：
\begin{center}\begin{tabular}{ccc}\toprule
步长／s&插值落地时间／s&时间误差／s\\\midrule
0.200&1.7874963599&$-1.0428\times10^{-3}$\\
0.100&1.7880784777&$-4.6065\times10^{-4}$\\
0.050&1.7883443970&$-1.9473\times10^{-4}$\\
0.025&1.7884716887&$-6.7437\times10^{-5}$\\\bottomrule
\end{tabular}\end{center}
\ExpansionDiagram{fig:dep-ai-drag}{线性阻力模型的高度与落地根；横轴为秒，高度以米计}{
\draw[->](0,0)--(5.6,0)node[right]{$t/\mathrm s$};\draw[->](0,-.7)--(0,3.3)node[above]{$y/\mathrm m$};\draw(2.5,-.06)--(2.5,.06)node[below=3pt]{1};\draw(5,-.06)--(5,.06)node[below=3pt]{2};\draw[thick,domain=0:1.87,samples=100]plot({2.5*\x},{.6*(87.5*(1-exp(-.4*\x))-25*\x)});\fill(4.47135,0)circle(1.6pt)node[below]{$T$};\draw[dashed](0,1.8)--(5,1.8);\node[left]at(0,1.8){3};}
图\ref{fig:dep-ai-drag}中的第二个零点才是落地事件。给AI实现时，应要求分别输出事件括区间、两端状态、插值方法和最终根，而不只输出一个飞行时间。此表的节点由解析式生成，没有积分误差；若程序给出不同结果，可以先从单位、初始零点和穿越方向查起。

\section{参数反演中的可识别性}
\subsection{残差很小并不保证参数唯一}
考虑恒定驱动与线性阻力模型$\dot v=a-\kappa v$，$v(0)=0$，未知参数为$a>0,\kappa>0$。沿用上一节的积分因子法，把常数驱动$-g$换成$a$、初速取零，得到$v(t)=a(1-\ee^{-\kappa t})/\kappa$；下例还将按这组初值逐步积分。只有一次速度测量时，每选一个$\kappa$都能调整$a$满足数据，所以不能同时唯一求出两参数。

很早的测量主要给$v\simeq at-a\kappa t^2/2$，对阻力的敏感度很弱；很晚的测量主要给终端速度$a/\kappa$，又只能识别参数比。把所有数据都放在同一个极端区间，即使点很多，也可能无法稳定区分$a$和$\kappa$。

有给定测量标准差$\sigma_i$时，可用灵敏度矩阵$J_{ij}=\partial v(t_i)/\partial\theta_j$讨论局部信息量。对$a$求偏导，$v$中的其余因子不变；对$\kappa$求偏导则同时对指数和分母使用链式法则与商法则。两列为
\[
\begin{aligned}
\frac{\partial v}{\partial a}&=\frac{1-\ee^{-\kappa t}}\kappa,\\
\frac{\partial v}{\partial\kappa}
&=a\frac{\kappa t\ee^{-\kappa t}-(1-\ee^{-\kappa t})}{\kappa^2}
=\frac{a[(1+\kappa t)\ee^{-\kappa t}-1]}{\kappa^2}.
\end{aligned}
\]
若两列在测量时刻近似成比例，$J^{\mathsf T}WJ$接近奇异，参数误差会很大。应先缩放参数与单位，再用奇异值或相关系数判断退化；原始矩阵某一列数值很小也可能只是单位不同。

\begin{example}两次速度测量的解析反演
粒子从静止开始满足$\dot v=a-\kappa v$。在$t_1=1\,\mathrm s$、$2t_1$测得速度$6.0,9.0\,\mathrm{m/s}$。求$a,\kappa$，预测$3t_1$时速度；若两次测量误差独立、标准差均为$0.03\,\mathrm{m/s}$，估计$\kappa$的一阶标准差。
\ExpansionDiagram{fig:dep-ai-inverse-q}{两次速度数据位于同一条向终端速度弯曲的曲线上}{
\draw[->](0,0)--(5.5,0)node[right]{$t/t_1$};\draw[->](0,0)--(0,3.2)node[above]{$v$};\draw[thick,domain=0:5,samples=90]plot(\x,{3*(1-exp(-ln(2)*\x))});\fill(1,1.5)circle(1.8pt)node[above left]{$v_1$};\fill(2,2.25)circle(1.8pt)node[above]{$v_2$};\draw[dashed](0,3)--(5,3)node[right]{$a/\kappa$};}
\end{example}
\begin{solution}
先求速度初值解。将方程写成$\dot v+\kappa v=a$，乘积分因子$\ee^{\kappa t}$，得到
\[\frac{\dd}{\dd t}(\ee^{\kappa t}v)=a\ee^{\kappa t}.\]
用$v(0)=0$从$0$到$t$积分，
\[\ee^{\kappa t}v(t)=a\int_0^t\ee^{\kappa s}\dd s
=a\left[\frac{\ee^{\kappa s}}{\kappa}\right]_0^t
=\frac a\kappa(\ee^{\kappa t}-1),\]
故$v(t)=a(1-\ee^{-\kappa t})/\kappa$。令$q=\ee^{-\kappa t_1}$，则$v_1=a(1-q)/\kappa$、$v_2=a(1-q^2)/\kappa$。图\ref{fig:dep-ai-inverse-q}的两点之比给
\[
 \frac{v_2}{v_1}=1+q,\qquad
 q=\frac12,\quad \kappa=\frac{\ln2}{t_1}=0.693147\,\mathrm{s^{-1}}.
\]
再由第一点求$a=\kappa v_1/(1-q)=12\ln2\,\mathrm{m/s^2}\simeq8.31777\,\mathrm{m/s^2}$。第三点应为$v_3=v_1(1+q+q^2)=10.5\,\mathrm{m/s}$。这个第三次预测没有参与确定参数，可以用于下一轮独立测量。
\ExpansionDiagram{fig:dep-ai-inverse-a}{反演先由比值确定衰减因子，再由尺度确定驱动强度}{
\node[draw]at(0,0){$v_2/v_1$};\node[draw]at(2.5,0){$q$};\node[draw]at(5,0){$\kappa$};\draw[->](.8,0)--(2.1,0);\draw[->](2.9,0)--(4.4,0);\node[draw]at(5,-1.4){$a$};\draw[->](5,-.4)--(5,-1);\draw[->](0,-1.4)--(4.5,-1.4)node[midway,below]{$v_1$提供尺度};}
图\ref{fig:dep-ai-inverse-a}中正阻力模型要求$0<q<1$，即$1<v_2/v_1<2$。如果数据比值在这个区间外，不应让优化器通过任意负阻力或极端参数掩盖模型不相容。

误差传播从$\kappa=-t_1^{-1}\ln q$求偏导，其中
$\partial q/\partial v_2=1/v_1$、$\partial q/\partial v_1=-v_2/v_1^2$。分别乘外层对数导数$-1/(t_1q)$，得到
\[
 \frac{\partial\kappa}{\partial v_2}=-\frac1{t_1qv_1},\qquad
 \frac{\partial\kappa}{\partial v_1}=\frac{v_2}{t_1qv_1^2}.
\]
所以$\sigma_\kappa=\sigma_v\sqrt{1+(v_2/v_1)^2}/(t_1qv_1)\simeq0.01803\,\mathrm{s^{-1}}$。两个数据恰好解两个参数时，拟合残差为零，但不代表参数没有误差；本题误差来自已给的测量标准差，不能再用零残差估计零方差。
\end{solution}

\section{静摩擦切换与分段运动}
\subsection{程序需要区分运动状态和静止状态}
库仑摩擦不是把$-\mu mg\operatorname{sgn}v$在$v=0$处随便补一个值即可。静止时，若外力$|F|\leq\mu_smg$，静摩擦调节为$-F$，速度保持零；外力超过阈值才开始滑动。滑动时摩擦大小为$\mu_kmg$、方向反对速度。速度到零后，须重新检查静止是否能维持。

交给AI生成数值模型时，状态量应包含“粘着”与“滑动”两类；滑动段定位$v=0$事件，再依据当时外力决定继续静止或立即反向。简单把每步负速度夹成零，可能丢失真实反向运动；把符号函数任意平滑，则改变了低速摩擦模型。两者都不能仅凭曲线平滑就当作原题解。

\begin{example}正弦驱动力下的停滑区间
水平滑块初始静止，外力$F(t)=0.6mg\sin\Omega t$，静摩擦系数$\mu_s=0.5$、动摩擦系数$\mu_k=0.3$。从$t=0$开始，求第一次启动、第一次停止及下一次反向启动的时刻，写出第一滑动段的速度。用$\theta=\Omega t$表示结果。
\ExpansionDiagram{fig:dep-ai-friction-q}{驱动力的正弦曲线与静摩擦上下限}{
\draw[->](0,0)--(6.8,0)node[right]{$\theta$};\draw[->](0,-1.7)--(0,1.8)node[above]{$F/(mg)$};\draw[thick,domain=0:6.283185,samples=150]plot(\x,{2*.6*sin(\x r)});\draw[dashed](0,1)--(6.3,1);\draw[dashed](0,-1)--(6.3,-1);\node[right]at(6.3,1){$0.5$};\node[right]at(6.3,-1){$-0.5$};}
\end{example}
\begin{solution}
图\ref{fig:dep-ai-friction-q}中第一次超过静摩擦上限发生在$\theta_0=\arcsin(5/6)\simeq0.985111$之后。启动临界时速度仍为零，开始向右滑动后满足
\[
 \frac{\dd v}{\dd t}=g(0.6\sin\theta-0.3).
\]
用$\dd\theta=\Omega\dd t$并带入启动初值$v(\theta_0)=0$，积分为
\[\frac{\Omega v}{g}=\int_{\theta_0}^{\theta}(0.6\sin u-0.3)\dd u
=[-0.6\cos u-0.3u]_{\theta_0}^{\theta}.\]
代入两端，得
\[
 \frac{\Omega v}{g}=0.6(\cos\theta_0-\cos\theta)-0.3(\theta-\theta_0).
\]
记上式右端为$H(\theta)$。启动根$H(\theta_0)=0$不能用作停止时刻。由$H'=0.6\sin\theta-0.3$，速度先增至$\theta=5\pi/6$，随后在所考察区间内严格减小。计算得
\[
H(3.74)\simeq9.3602\times10^{-4}>0,\qquad
H(3.75)\simeq-5.4687\times10^{-3}<0.
\]
因此首次停止根位于$[3.74,3.75]$。二分该区间至宽度小于$10^{-8}$，得$\theta_s\simeq3.741466$。此时$F/(mg)=0.6\sin\theta_s\simeq-0.338723$，仍在$[-0.5,0.5]$内，所以滑块可以停住，由静摩擦抵消负驱动力。
\ExpansionDiagram{fig:dep-ai-friction-a}{第一正向滑动段结束后存在有限粘着区间}{
\draw[->](0,0)--(6.3,0)node[right]{$\theta$};\draw[->](0,-.7)--(0,1.6)node[above]{$\Omega v/g$};\draw[very thick](0,0)--(.985111,0);\draw[thick,domain=.985111:3.741466,samples=100]plot(\x,{.6*(cos(.985111 r)-cos(\x r))-.3*(\x-.985111)});\draw[very thick](3.741466,0)--(4.126703,0);\draw[thick,domain=4.126703:5.2,samples=60]plot(\x,{.6*(cos(4.126703 r)-cos(\x r))+.3*(\x-4.126703)});\draw[dashed](3.741466,-.35)--(3.741466,.8);\draw[dashed](4.126703,-.35)--(4.126703,.8);\node[above]at(3.8,.9){停滑};}
下一次超过负静摩擦阈值在$\theta_r=\pi+\arcsin(5/6)\simeq4.126703$。因此图\ref{fig:dep-ai-friction-a}中的静止区间为$[\theta_s,\theta_r]$。反向滑动开始后摩擦改为向右，用$v(\theta_r)=0$积分，
\[\frac{\Omega v}{g}=\int_{\theta_r}^{\theta}(0.6\sin u+0.3)\dd u
=[-0.6\cos u+0.3u]_{\theta_r}^{\theta},\]
故局部速度为
\[
 \frac{\Omega v}{g}=0.6(\cos\theta_r-\cos\theta)+0.3(\theta-\theta_r).
\]
三个实际时刻分别为$\theta_0/\Omega,\theta_s/\Omega,\theta_r/\Omega$。编写数值程序时，可逐段用上述积分式核对速度，再按静摩擦阈值接续状态。停下后的有限粘着区间应被保留，平滑插值会把这段运动规律抹去。
\end{solution}

\section{事件终点的灵敏度与参数优化}
\subsection{参数改变同时移动轨迹和终止时刻}
对微分方程$\dot{\vct y}=\vct f(t,\vct y,\lambda)$，固定时刻的参数灵敏度为$\vct S=\partial\vct y/\partial\lambda$。对原方程求参数导数，得到与原轨迹一并积分的切线方程
\[
 \dot{\vct S}=\frac{\partial\vct f}{\partial\vct y}\vct S
 +\frac{\partial\vct f}{\partial\lambda},\qquad
 \vct S(0)=\frac{\partial\vct y_0}{\partial\lambda}.
\]
原方程中没有显含参数，并不意味着结果对它不敏感；参数可能只出现在初始位置、速度或方向中。

若终点由$g(T,\vct y(T),\lambda)=0$确定，链式法则给出
\[0=g_tT_\lambda+g_{\vct y}(\vct S+\vct fT_\lambda)+g_\lambda.\]
合并所有$T_\lambda$项，只要事件穿越速度$g_t+g_{\vct y}\vct f\ne0$，就可解得
\[
 \frac{\dd T}{\dd\lambda}
 =-\frac{g_{\vct y}\vct S+g_\lambda}
 {g_t+g_{\vct y}\vct f}\bigg|_{t=T}.
\]
分母是沿轨迹的事件穿越速度。最终观测量若为$x(T,\lambda)$，其总导数是
\[
 \frac{\dd x(T,\lambda)}{\dd\lambda}
 =S_x(T)+\dot x(T)\frac{\dd T}{\dd\lambda}.
\]
遗漏后一项，相当于把两条不同参数轨迹强行比较在同一时刻，得到的最优参数可能错误。

让AI生成切线方程时，应明确各偏导中保持不变的变量，再用独立的有限参数差分核对。差分间隔过大会留下截断误差，过小则放大舍入和事件定位误差。直接积分切线方程可减少两条接近轨迹相减产生的噪声；原轨迹自身的积分误差仍需控制。

\begin{example}线性阻力下的最佳发射角
抛体满足$\dot v_x=-\kappa v_x$、$\dot v_y=-g-\kappa v_y$，初速$v_0$固定，发射角为$\theta\in(0,\pi/2)$，初始高度为零。以第一次向下返回$y=0$为终点，推导射程对发射角的导数。取$v_0=10\,\mathrm{m/s}$、$g=10\,\mathrm{m/s^2}$、$\kappa=0.4\,\mathrm{s^{-1}}$，求最优角度，并给出无阻力极限。
\ExpansionDiagram{fig:dep-ai-sensitivity-q}{同一初速下，改变角度同时改变落地点与飞行时间}{
\draw[->](0,0)--(6.5,0)node[right]{$x$};\draw[->](0,0)--(0,3.2)node[above]{$y$};\draw[thick](0,0)..controls(2,2.7)and(4,2)..(5.5,0);\draw[dashed](0,0)..controls(1.7,3.5)and(3.8,2.6)..(5.1,0);\node[above]at(3,2.5){$\theta+\dd\theta$};\node[below]at(5.5,0){$R$};}
\end{example}
\begin{solution}
令$a=v_0\sin\theta$、$b=v_0\cos\theta$。横向速度方程乘积分因子后为$(\ee^{\kappa t}v_x)'=0$，初值给出$v_x=b\ee^{-\kappa t}$。竖向将常数项移入未知量，
\[\frac{\dd}{\dd t}\left[\ee^{\kappa t}\left(v_y+\frac g\kappa\right)\right]=0,\]
由$v_y(0)=a$得到$v_y=(a+g/\kappa)\ee^{-\kappa t}-g/\kappa$。再用$x(0)=y(0)=0$对速度积分：
\[x=b\int_0^t\ee^{-\kappa s}\dd s
=b\left[-\frac{\ee^{-\kappa s}}\kappa\right]_0^t,\]
\[y=\left(a+\frac g\kappa\right)
\left[-\frac{\ee^{-\kappa s}}\kappa\right]_0^t
-\frac g\kappa[s]_0^t.\]
因此图\ref{fig:dep-ai-sensitivity-q}的轨迹解析式为
\[
 x(t)=\frac b\kappa(1-\ee^{-\kappa t}),\quad
 y(t)=\frac{a+g/\kappa}{\kappa}(1-\ee^{-\kappa t})-\frac{gt}{\kappa}.
\]
非零落地根满足$y(T(\theta),\theta)=0$。固定$t$时，$\partial_\theta a=b$，所以$S_y=\partial_\theta y=b(1-\ee^{-\kappa t})/\kappa$。对事件方程作总导数，
\[0=\frac{\dd}{\dd\theta}y(T(\theta),\theta)
=S_y(T)+v_y(T)T_\theta.\]
落地速度$v_y(T)=(a+g/\kappa)\ee^{-\kappa T}-g/\kappa<0$，因此可除以它，得到
\[
 T_\theta=-\frac{b(1-\ee^{-\kappa T})}
 {\kappa[(a+g/\kappa)\ee^{-\kappa T}-g/\kappa]}.
\]
射程为$R=x(T(\theta),\theta)$。由$\partial_\theta b=-a$，固定时刻的$S_x=-a(1-\ee^{-\kappa t})/\kappa$；再加上终点时间项$v_x(T)T_\theta$，总导数为
\[
 R_\theta=-\frac{a}{\kappa}(1-\ee^{-\kappa T})
 +b\ee^{-\kappa T}T_\theta.
\]
\ExpansionDiagram{fig:dep-ai-sensitivity-a}{先求落地事件，再把终点时间变化加到射程导数}{
\node[draw](event)at(0,0){$y(T)=0$};\node[draw](time)at(3.2,0){$T_\theta=-S_y/v_y$};\node[draw](range)at(7,0){$R_\theta=S_x+v_xT_\theta$};\draw[->](event.east)--(time.west);\draw[->](time.east)--(range.west);\node[below]at(3.5,-.7){以非零、向下穿越根为终点};}
图\ref{fig:dep-ai-sensitivity-a}把事件导数单独保留。求最优角度时，还须确认驻点的性质。令$u=\kappa T$、$\beta=\kappa v_0/g$，落地条件给出
\[
h(u)\equiv\frac{u}{1-\ee^{-u}}-1=\beta\sin\theta.
\]
商法则给出
\[h'(u)=\frac{(1-\ee^{-u})-u\ee^{-u}}{(1-\ee^{-u})^2}
=\frac{1-(1+u)\ee^{-u}}{(1-\ee^{-u})^2}>0,\]
最后用到$\ee^u>1+u$。因此发射角增大时，非零飞行时间严格增大，可改以$u$为参数求最大射程。由$\beta\cos\theta=\sqrt{\beta^2-h^2}$，有
\[
\left(\frac{\kappa^2R}{g}\right)^2
=\beta^2(1-\ee^{-u})^2-(u-1+\ee^{-u})^2.
\]
记$s=1-\ee^{-u}$，则$s'=\ee^{-u}$、$(u-s)'=s$。对右侧两平方项逐项求导，得到
\[2\beta^2s\ee^{-u}-2(u-s)s
=2s[\beta^2\ee^{-u}-u+s].\]
将$s$代回并合并，得到
\[
\frac{\dd}{\dd u}\left(\frac{\kappa^2R}{g}\right)^2
=2(1-\ee^{-u})[1-u+(\beta^2-1)\ee^{-u}].
\]
记方括号为$G(u)$。对于$u>0$，$G'(u)=-1+(1-\beta^2)\ee^{-u}<0$；在$u=0$处$G(0)=\beta^2>0$，在竖直发射对应的终点$u_v$处，
$G(u_v)=\beta[(1+\beta)\ee^{-u_v}-1]<0$，最后的不等式正是落地竖直速度为负。因而有效区间只有一个驻点，射程先增后减，该点给出全局最大值。

本题$\beta=0.4$，只需在$[0.48,0.49]$内求$G(u)=0$。用二分法将区间缩至$10^{-10}$以内，得到$u\simeq0.480462068$。再依次计算$T=u/\kappa$、$\theta=\arcsin[h(u)/\beta]$与射程，结果为
\[
 \theta_{\rm opt}\simeq0.705593884\,\mathrm{rad}
 =40.4275516^\circ,\quad
 T\simeq1.20115517\,\mathrm s,\quad R\simeq7.26024518\,\mathrm m.
\]
在$45^\circ$时射程约$7.17528806\,\mathrm m$，导数约$-2.10619\,\mathrm{m/rad}$，因此它已经在射程下降一侧。数值导数中的角度单位必须为弧度；若使用度，需另乘$\pi/180$。

当$\kappa\to0$，解析式的差商应用指数展开或稳定的$\operatorname{expm1}$实现，避免两个接近数相减。无阻力时$T=2v_0\sin\theta/g$、$T_\theta=2v_0\cos\theta/g$，两项合并给$R_\theta=2v_0^2\cos2\theta/g$，最优角恢复$\pi/4$。只保留固定时刻导数$S_x$则对整个角域均为负，显然无法恢复这一极限。
\end{solution}

同一方法还能用于拟合中的参数Jacobian和约束事件到达时间。若终点恰好擦过事件面，分母趋零，灵敏度可能很大甚至失去普通可微性；此时应改变事件参数化或分段讨论，不能要求优化程序在所有参数处都返回平滑导数。

\section{随机抽样的误差与物理测度}
\subsection{随机点代表面积、体积还是粒子数}
用AI编写复杂边界积分的抽样程序，首先要给出积分测度。平面圆盘的面积元为$r\dd r\dd\theta$，均匀面积分布对应的半径概率密度是$2r/R^2$；直接均匀抽取$r$会使圆心附近的点过密。若$U,V$在$[0,1]$独立均匀，则
\[
 r=R\sqrt U,\qquad\theta=2\pi V
\]
给出均匀面积点。这一平方根来自半径的累积概率
\[P(r\le s)=\int_0^s\frac{2r}{R^2}\dd r
=\left[\frac{r^2}{R^2}\right]_0^s=\frac{s^2}{R^2}.\]
球内均匀体积的半径累积概率则为
\[P(r\le s)=\int_0^s\frac{3r^2}{R^3}\dd r
=\left[\frac{r^3}{R^3}\right]_0^s=\frac{s^3}{R^3},\]
故反演需$r=RU^{1/3}$。角向概率元为$\sin\theta\dd\theta\dd\phi/(4\pi)$；令$c=\cos\theta$，$\dd c=-\sin\theta\dd\theta$，可见$c$在$[-1,1]$均匀、$\phi$在$[0,2\pi)$均匀。因此球内采样可另取独立均匀数$V,W$，令$\cos\theta=2V-1$、$\phi=2\pi W$。

积分$I=\int f(x)p(x)\dd x$用按$p$抽取的独立同分布点估计为
\[
\widehat I=\frac1N\sum_{i=1}^N f(X_i).
\]
每项均值都是$I$，故$\langle\widehat I\rangle=I$。若单点方差$\sigma_f^2$有限，均值的方差为
\[
\begin{aligned}
\operatorname{Var}(\widehat I)
&=\frac1{N^2}\sum_{i,j}\operatorname{Cov}(f(X_i),f(X_j))\\
&=\frac1{N^2}\sum_{i=1}^N\sigma_f^2
=\frac{\sigma_f^2}{N}.
\end{aligned}
\]
第二步用独立性消去了$i\ne j$的协方差。标准误差是方差的平方根，因此按$N^{-1/2}$下降；样本数增加四倍只使它减半，这与确定性二阶积分的步长规律不同。

样本方差$s_f^2=\sum(f_i-\bar f)^2/(N-1)$可估计$\operatorname{Var}_p(f)$，但前提是样本近乎独立、二阶矩有限。连续轨迹上相邻时间点往往高度相关，不能仅因文件里有很多行便把它们当作大量独立样本。固定随机种子便于复现同一次计算，多次使用同一条序列也不会产生新的独立证据。

\begin{example}非均匀圆盘的惯量及比值误差
半径$R$的薄圆盘面密度为$\rho(r)=\rho_0(1+2r^2/R^2)$。用均匀面积抽样求总质量$M$和绕中心垂直轴的转动惯量$I$，写出同一批点估计$I/M$时的一阶标准误差。给出可用来核对程序的解析值。
\ExpansionDiagram{fig:dep-ai-disk-q}{同心环带面积随半径增加，均匀半径不能代替均匀面积}{
\draw[thick](0,0)circle(1.8);\draw[dashed](0,0)circle(.6);\draw[dashed](0,0)circle(1.2);\draw[->](0,0)--(1.8,0)node[right]{$R$};\draw[->](.6,.2)--(1.2,.2)node[midway,above]{环带};\node[right]at(2.5,.8){$r=R\sqrt U$};\node[right]at(2.5,.1){$\rho=\rho_0(1+2U)$};}
\end{example}
\begin{solution}
图\ref{fig:dep-ai-disk-q}中$Z=r^2/R^2$在$[0,1]$上均匀。若各点为$Z_i$，面积为$\pi R^2$，所以
\[
 \widehat M=\pi R^2\rho_0\overline{(1+2Z)},\qquad
 \widehat I=\pi R^4\rho_0\overline{(Z+2Z^2)}.
\]
均匀$Z$的矩由幂积分给出：
\[\langle Z^n\rangle=\int_0^1z^n\dd z
=\left[\frac{z^{n+1}}{n+1}\right]_0^1=\frac1{n+1}\quad(n\ge0).\]
取$n=1,2$分别得$1/2,1/3$，故
\[M=\pi R^2\rho_0(1+2/2)=2\pi\rho_0R^2,\]
\[I=\pi R^4\rho_0(1/2+2/3)=\frac{7\pi\rho_0R^4}{6},\qquad
K=\frac IM=\frac{7R^2}{12}.\]
\ExpansionDiagram{fig:dep-ai-disk-a}{质量与惯量来自同一批点，误差具有相关性}{
\node[draw]at(0,0){$Z_i$};\node[draw]at(3,1){$b_i=\rho_i$};\node[draw]at(3,-1){$a_i=\rho_ir_i^2$};\draw[->](.4,0)--(1.8,1);\draw[->](.4,0)--(1.7,-1);\node[draw]at(6,0){$\widehat K=\bar a/\bar b$};\draw[->](4,1)--(5.1,.3);\draw[->](4,-1)--(5.1,-.3);}
按图\ref{fig:dep-ai-disk-a}，记单点贡献为$a=\rho r^2=\rho_0R^2(Z+2Z^2)$、$b=\rho=\rho_0(1+2Z)$，则$\widehat I=\pi R^2\bar a$、$\widehat M=\pi R^2\bar b$，比值为$\widehat K=\bar a/\bar b$。在两样本均值各自的期望附近作一阶展开，
\[
 \widehat K-K\simeq\frac{(\bar a-\langle a\rangle)
 -K(\bar b-\langle b\rangle)}{\langle b\rangle},
 \quad\operatorname{Var}(\widehat K)\simeq
 \frac{\operatorname{Var}(a-Kb)}{N\langle b\rangle^2}.
\]
把质量、惯量误差当作独立量相加会遗漏协方差。这里
\[
 a-Kb=\rho_0R^2\left(2Z^2-\frac Z6-\frac7{12}\right),\quad
 \langle b\rangle=2\rho_0,
\]
括号中多项式的均值为零，其平方的均值为
\begin{align*}
\int_0^1\left(2z^2-\frac z6-\frac7{12}\right)^2\dd z
&=\int_0^1\left(4z^4-\frac23z^3-\frac{83}{36}z^2
+\frac7{36}z+\frac{49}{144}\right)\dd z\\
&=\left[\frac45z^5-\frac16z^4-\frac{83}{108}z^3
+\frac7{72}z^2+\frac{49}{144}z\right]_0^1\\
&=\frac45-\frac16-\frac{83}{108}+\frac7{72}+\frac{49}{144}
=\frac{653}{2160}.
\end{align*}
乘回$\rho_0^2R^4$，再除以$N\langle b\rangle^2=4N\rho_0^2$，便得$\operatorname{Var}(\widehat K)\simeq653R^4/(8640N)$，所以标准误差约为$0.2749R^2/\sqrt N$。有限$N$下比值本身还可能有$O(1/N)$偏差；上述式子描述其主导随机误差。

若按质量分布生成粒子点，$Z$的概率密度为$(1+2z)/2$，累积函数为
\[P(Z\le z)=\int_0^z\frac{1+2u}{2}\dd u
=\left[\frac{u+u^2}{2}\right]_0^z=\frac{z+z^2}{2}.\]
令其等于均匀随机数$U$，得到$Z^2+Z-2U=0$；取$[0,1]$内的根，$Z=(-1+\sqrt{1+8U})/2$。此时点已经按质量加权，求$I/M$直接平均$r^2$；若再乘一次$\rho$，会重复加权。给AI的任务中应写清抽样测度和每点权重。
\end{solution}

\subsection{成对抽样可以减少误差，但要按相同成本比较}
若一次随机取$U$后同时计算$f(U)$和$f(1-U)$，两值负相关时可降低均值方差，这称为对偶抽样。记单点方差为$\sigma_f^2$、同对协方差为$c_f$，一对平均$G=[f(U)+f(1-U)]/2$的方差为
\[\operatorname{Var}(G)=\frac14(\sigma_f^2+\sigma_f^2+2c_f)
=\frac{\sigma_f^2+c_f}{2}.\]
若总共作$N$次函数求值，即有$N/2$个独立配对，最终均值方差为
\[\frac{\operatorname{Var}(G)}{N/2}=\frac{\sigma_f^2+c_f}{N}.\]
普通$N$点均值的方差为$\sigma_f^2/N$，故$c_f<0$时配对才有益。两法必须按同样的函数求值次数比较。

\begin{example}杆的惯量积分中的对偶抽样
均匀细杆一端为转轴，长度、总质量均已归一化，惯量积分为$I=\int_0^1x^2\dd x$。比较$N$次普通均匀抽样与$N/2$组成对抽样的方差，$N$为偶数。
\ExpansionDiagram{fig:dep-ai-pair-q}{同一对点关于杆中点对称，求值次数仍计为两次}{
\draw[very thick](0,0)--(6,0);\fill(0,0)circle(2pt)node[below]{转轴};\fill(1.5,0)circle(2pt)node[above]{$U$};\fill(4.5,0)circle(2pt)node[above]{$1-U$};\draw[dashed](3,-.5)--(3,.7);\node[below]at(6,0){1};}
\end{example}
\begin{solution}
图\ref{fig:dep-ai-pair-q}中$U$在$[0,1]$上均匀。由
\[\langle U^n\rangle=\int_0^1u^n\dd u
=\left[\frac{u^{n+1}}{n+1}\right]_0^1=\frac1{n+1},\]
单点函数$f(U)=U^2$的均值为$1/3$，二阶矩为$\langle U^4\rangle=1/5$，方差为$1/5-1/9=4/45$。普通$N$点平均的方差为$4/(45N)$。

一对的平均值为$G=[U^2+(1-U)^2]/2=U^2-U+1/2$，其均值仍为$1/3$。展开平方后，
\[
\begin{aligned}
\langle G^2\rangle&=\int_0^1\left(u^4-2u^3+2u^2-u+\frac14\right)\dd u\\
&=\left[\frac{u^5}{5}-\frac{u^4}{2}+\frac{2u^3}{3}-\frac{u^2}{2}+\frac u4\right]_0^1\\
&=\frac15-\frac12+\frac23-\frac12+\frac14=\frac7{60}.
\end{aligned}
\]
因此$\operatorname{Var}(G)=7/60-(1/3)^2=1/180$。独立的成对样本共有$N/2$组，所以
\[
 \operatorname{Var}(\widehat I_{\rm pair})=\frac{1/180}{N/2}
 =\frac1{90N}.
\]
\ExpansionDiagram{fig:dep-ai-pair-a}{单点函数变化范围较大，成对平均后波动缩小}{
\draw[->](0,0)--(5,0)node[right]{$U$};\draw[->](0,0)--(0,3)node[above]{函数值};\draw[dashed,domain=0:4.5,samples=80]plot(\x,{2.5*(\x/4.5)^2});\draw[thick,domain=0:4.5,samples=80]plot(\x,{2.5*((\x/4.5)^2-\x/4.5+.5)});\node[right]at(3.3,2.2){$U^2$};\node[right]at(3.4,.35){$G$};\node[below]at(4.5,0){1};}
图\ref{fig:dep-ai-pair-a}中成对平均的波动较小。在相同$N$次函数求值下，方差降为普通法的$1/8$，标准误差降为$1/\sqrt8$。这一收益来自本题函数的相关结构；若$f(U)$与$f(1-U)$正相关，配对未必有利。程序应输出成对平均值的样本方差，不能把同组的两个点误当成独立点计算误差条。
\end{solution}

\section{稀有事件的加权抽样}
\subsection{改变抽样分布以后必须保留密度比}
若目标概率密度为$p(x)$，希望求$I=\int f(x)p(x)\dd x$，可改按便于抽样的$q(x)$生成点，但每点需乘权重$p/q$：
\[
 I=\int f(x)\frac{p(x)}{q(x)}q(x)\dd x.
\]
要求在$f(x)p(x)$非零的区域，$q(x)$也非零。只把粒子“抽得更热”而不重新加权，得到的是另一个物理系统，不是原来低概率事件的准确估计。

令单点加权贡献为$Y=f(X)p(X)/q(X)$，其中$X$按$q$抽取。其均值为$I$，二阶矩为
\[
\langle Y^2\rangle_q
=\int\left(\frac{fp}{q}\right)^2q\dd x
=\int\frac{f^2p^2}{q}\dd x.
\]
若此积分有限，$N$个独立加权样本的均值方差为
\[
\operatorname{Var}(\widehat I_q)
=\frac1N\left(\int\frac{f^2p^2}{q}\dd x-I^2\right).
\]
使稀有区域出现更多样本可能降低方差，也可能因少数极大权重而恶化。输出时应同时保留权重、加权样本值及其二阶矩；用未加权直方图代替加权结果，会改变归一化和尾部概率。

\begin{example}二维热粒子的高能尾概率
二维经典平动粒子以$\kb T$为能量单位，$X=E/(\kb T)$的概率密度为$p(x)=\ee^{-x}$，$x\geq0$。求$P(X>8)$。比较直接抽样与按$q_\lambda(x)=\lambda\ee^{-\lambda x}$抽样时的相对方差，并取$\lambda=1/8$估算达到10\%相对标准误差所需的样本数。
\ExpansionDiagram{fig:dep-ai-tail-q}{较宽的抽样分布增加高能区样本，但每个样本需重新加权}{
\draw[->](0,0)--(7,0)node[right]{$x$};\draw[->](0,0)--(0,2.6)node[above]{概率密度};\draw[thick,domain=0:16,samples=120]plot({.4*\x},{2.2*exp(-\x)});\draw[dashed,domain=0:16,samples=120]plot({.4*\x},{2.2*.125*exp(-.125*\x)});\draw[dotted](3.2,0)--(3.2,1.4)node[above]{$x=8$};\node[right]at(.6,1.2){$p$};\node[right]at(4,.5){$q_{1/8}$};}
\end{example}
\begin{solution}
目标概率为
\[P=\int_8^\infty\ee^{-x}\dd x
=[-\ee^{-x}]_8^\infty=\ee^{-8}\simeq3.3546\times10^{-4}.\]
普通抽样把每点记为0或1，单点方差为$P(1-P)$，所以平均估计的相对方差为$(1-P)/(NP)$。要求相对标准误差不超过0.1，需要$N\geq(\ee^8-1)/0.01\simeq2.98\times10^5$。

按图\ref{fig:dep-ai-tail-q}的较宽分布抽样，单点加权值是
\[
 Y=\frac1\lambda\ee^{-(1-\lambda)X}\,\boldsymbol1_{X>8}.
\]
其均值为
\[\langle Y\rangle_q=\int_8^\infty
\frac{\ee^{-(1-\lambda)x}}\lambda\,
\lambda\ee^{-\lambda x}\dd x
=\int_8^\infty\ee^{-x}\dd x=P.\]
二阶矩在$0<\lambda<2$时有限：
\[
 \langle Y^2\rangle_q=\frac1\lambda\int_8^\infty
 \ee^{-(2-\lambda)x}\dd x
 =\frac1\lambda\left[-\frac{\ee^{-(2-\lambda)x}}{2-\lambda}\right]_8^\infty
 =\frac{\ee^{-8(2-\lambda)}}{\lambda(2-\lambda)}.
\]
因此$N$点平均的相对方差为
\[
 \frac1N\left[\frac{\ee^{8\lambda}}{\lambda(2-\lambda)}-1\right].
\]
取$\lambda=1/8$，括号约为10.5980，故所需样本数约为1060。该样本数是标准误差条件，不是对每次随机结果都保证10\%误差的确定性上界。
\ExpansionDiagram{fig:dep-ai-tail-a}{改变抽样概率后，按$p/q$把每点贡献还原到原来的热分布}{
\node[draw]at(0,0){$q_\lambda$样本};\node[draw]at(3,0){乘$p/q_\lambda$};\node[draw]at(6,0){原分布积分};\draw[->](1,0)--(1.8,0);\draw[->](4.1,0)--(4.9,0);\node[below]at(3,-.6){较多尾部样本不等于改变物理温度};}
图\ref{fig:dep-ai-tail-a}区分抽样分布与目标热分布。对一般阈值$A>0$，相对二阶矩为$Q(\lambda)=\ee^{A\lambda}/[\lambda(2-\lambda)]$。取对数再求导，
\[\frac{\dd}{\dd\lambda}\ln Q
=A-\frac1\lambda+\frac1{2-\lambda}=0.\]
乘$\lambda(2-\lambda)$并移项，得到
\[A\lambda^2-2(A+1)\lambda+2=0,\qquad
\lambda=\frac{A+1\pm\sqrt{A^2+1}}{A}.\]
加号根大于$2$，只有减号根属于允许区间。又有
$(\ln Q)''=1/\lambda^2+1/(2-\lambda)^2>0$，所以它是唯一极小点：
\[\lambda_* =\frac{A+1-\sqrt{A^2+1}}{A}
=\frac2{A+1+\sqrt{A^2+1}}.\]
当$A$很大时$\lambda_*\sim1/A$，因而本题取$1/8$。
\end{solution}

\subsection{让AI生成可复算的抽样实验}
提交给AI的计算任务应包含目标密度、抽样密度、变量变换、单位及误差定义。例如本节尾概率实验应要求：分别运行直接法和加权法，保存每次独立运行的样本数、随机种子、估计均值、样本方差和标准误差，并与$\ee^{-8}$比较。重跑时改变独立随机序列，才可观察估计量的运行间分布。

比较两套物理参数时，若使用同一批基础均匀随机数，两输出差值的方差为$\operatorname{Var}(A-B)=\operatorname{Var}(A)+\operatorname{Var}(B)-2\operatorname{Cov}(A,B)$。正相关能使差值更稳定，但也意味着两次结果并不独立。对同一轨迹长时间取样，相关时间又会减少有效样本数；仅增加输出频率不会创造新的独立轨迹。

报告中宜把单次估计、运行间散布和理论标准误差分开。解析基准用于核对模型与抽样过程；误差条说明有限样本造成的波动。它们回答的问题不同。若平均值长期偏离基准，优先检查密度比、Jacobian和边界筛选，再增加样本量；系统性的测度错误不会随$N$增大而自动消失。

\clearpage
\section*{习题与解答}\addcontentsline{toc}{section}{习题与解答}
\begin{enumerate}[label=4.\arabic*,itemsep=3pt]
\item \begin{minipage}[t]{\linewidth}
【基础】有人画出 $q=\cos(1.01\tau),p=-\sin(1.01\tau)$，能量恒为 $1/2$。检验它是否满足 $q'=p,p'=-q$，并求 $\tau=100$ 时与准确解的相位差。
\end{minipage}
% 题设配图；与紧随其后的解答分别编号。
\ExpansionDiagram{fig:audit-m-q-109}{候选解在单位圆上的状态与转动方向}{
\draw[->](-.4,0)--(2.5,0)node[right]{$ q $};\draw[->](0,-1.9)--(0,2)node[above]{$ p $};\draw[thick](0,0)circle(1.4);\fill(1.4,0)circle(2pt)node[above right]{$(1,0)$};\draw[->,thick](1.4,0)--(1.4,-1.1)node[right]{$(q^\prime,p^\prime)$};\node[left]at(-1.4,0){$q^2+p^2=1$};
}

\begin{solution}
图\ref{fig:audit-m-q-109}中沿圆周运动只保证$q^2+p^2$不变，角速度还必须满足给定微分方程。
图\ref{fig:hw-4-1}已把本题在 $\tau=100$ 的两个状态画在同一单位圆上。先检查方程，再解释圆上转角。

逐项求导有
\[
q'=-1.01\sin(1.01\tau)=1.01p,\qquad
p'=-1.01\cos(1.01\tau)=-1.01q.
\]
原方程的两个残差分别为 $q'-p=0.01p$、$p'+q=-0.01q$，所以不满足方程。初值 $(q,p)=(1,0)$ 虽然正确，却不能保证随后运动正确。

\begin{figure}[H]\centering
\begin{tikzpicture}[scale=1.5]
\draw[bookaxis](-1.4,0)--(1.7,0)node[right]{$q$};\draw[bookaxis](0,-1.4)--(0,1.6)node[above]{$p$};\draw(0,0)circle(1);
\fill(.862319,.506366)circle(1.4pt)node[above right]{准确状态};\fill(.892005,-.452026)circle(1.4pt)node[below right]{错误状态};
\draw[bookguide](0,0)--(.862319,.506366) (0,0)--(.892005,-.452026);\draw[->](30.42205:1.25)arc(30.42205:-26.87373:1.25);\node[right]at(1.25,.03){$1\,\mathrm{rad}$};
\end{tikzpicture}
\caption{题4.1解图：在 $\tau=100$，错误状态比准确状态沿顺时针运动方向超前1 rad，两状态的半径仍同为1。}\label{fig:hw-4-1}
\end{figure}

准确解的相位是 $\tau$，给定曲线的相位是 $1.01\tau$，在 $\tau=100$ 超前1 rad。两者都沿单位圆运动，因而 $(q^2+p^2)/2=1/2$ 无法区分转动快慢。还需检查相位事件或方程残差；本题在初始点已有 $p'(0)=-1.01\ne-1$。

两条半径一样长只说明能量相同；它们沿运动方向相差的1 rad才是累计相位误差。
\end{solution}

\item \begin{minipage}[t]{\linewidth}
【基础】对 $\ell=1\,\mathrm m,g=9.8\,\mathrm{m/s^2},\theta_0=0.2\,\mathrm{rad}$，估计单摆周期；经过100个小角周期，估计真实运动落后的相位。说明两种误差为何不同。

\end{minipage}
% 题设配图；与紧随其后的解答分别编号。
\ExpansionDiagram{fig:audit-m-q-110}{单摆从有限摆角处由静止释放}{
\fill(0,2.5)circle(2pt)node[above]{$O$};\draw[dashed](0,2.5)--(0,-.4);\draw[thick](0,2.5)--(1.2,.1);\fill(1.2,.1)circle(3pt)node[right]{$m$};\draw(0,1.7)arc(-90:-63.435:.8);\node[right]at(.25,1.5){$\theta_0$};\node[right]at(.8,1.2){$\ell$};\draw[->](2,1.9)--(2,.8)node[right]{$g$};
}

\begin{solution}
图\ref{fig:audit-m-q-110}中$\theta_0$决定周期修正；累计相位误差还取决于观察时间。
图\ref{fig:ex-4-2}以竖直向下为摆角零点，释放点的 $\theta_0$ 决定周期修正。

小角周期为 $T_0=2\pi\sqrt{\ell/g}=2.00709\,\mathrm s$。用单摆周期的二阶展开，
\[
T/T_0\simeq1+\theta_0^2/16=1.0025,\qquad
T\simeq2.01211\,\mathrm s.
\]
在 $t=100T_0$ 时，小角解的相位已走过 $200\pi$，实际相位约为 $2\pi t/T$，所以落后量为
\[
\Delta\varphi\simeq200\pi\left(1-\frac1{1.0025}\right)
\simeq1.567\,\mathrm{rad}.
\]
若同样只保留振幅二阶，$\Delta\varphi\simeq200\pi\theta_0^2/16=\pi/2$。周期的相对误差约0.25\%，每一周期只差很少；相位记录了累计走过的周期数，误差会随观察时间增长。这里的数值仍是二阶近似估计，不能作为长时间精确单摆轨迹。

\begin{figure}[H]\centering
\begin{tikzpicture}[scale=1]
\draw[thick](-.5,0)--(.5,0);\fill(0,0)circle(1.7pt);\draw[bookguide](0,0)--(0,-2.8);
\draw[thick](0,0)--(1.15,-2.36);\fill(1.15,-2.36)circle(3.4pt)node[right]{$m$};\node[right]at(.65,-1.15){$\ell$};
\draw[->](0,-.8)arc(-90:-64:.8);\node at(.4,-1){$\theta_0$};\draw[bookaxis](2,-.4)--(2,-1.2)node[below]{$\vct g$};
\node[below]at(1,-2.9){$\dot\theta(0)=0$};
\end{tikzpicture}
\caption{题4.2、4.5：从角度 $\theta_0$ 静止释放的无阻尼单摆。}\label{fig:ex-4-2}
\end{figure}

图中完整往返一次才是一个周期；微小的单次周期差重复100次后，才累积成题目所问的相位落后。
\end{solution}

\item \begin{minipage}[t]{\linewidth}
【基础】$\zeta=0.02$ 时，要使欠阻尼瞬态包络降为初值的 $10^{-3}$，至少等待多少无量纲时间？折合多少个 $2\pi$ 的时间单位？
\end{minipage}
% 题设配图；与紧随其后的解答分别编号。
\ExpansionDiagram{fig:audit-m-q-111}{在对数纵轴上比较瞬态包络与阈值}{
\draw[->](-.4,0)--(5,0)node[right]{$ \tau $};\draw[->](0,-3)--(0,1)node[above]{$ \ln(A/A_0) $};\draw[thick](0,0)--(4.6,-2.76);\draw[dashed](0,-2.4)--(4.8,-2.4);\node[left]at(0,-2.4){$-\ln1000$};\draw[dotted](4,0)--(4,-2.4);\node[above]at(4,0){$\tau_*$};\node[below]at(2,-.6){斜率$-\zeta$};
}

\begin{solution}
图\ref{fig:audit-m-q-111}中比较的是自由瞬态的包络，不是某一次经过零点的瞬时位移。
图\ref{fig:hw-4-3}把指数衰减改用对数纵轴，等待时间由阈值交点读出。

在本章无量纲时间下，欠阻尼自由分量的包络为 $A\ee^{-\zeta\tau}$。要求相对于初始包络衰减至 $10^{-3}$，有
\[
\ee^{-0.02\tau}\le10^{-3}
\quad\Longrightarrow\quad
\tau\ge\frac{\ln1000}{0.02}=345.388.
\]
除以 $2\pi$ 得54.970个时间单位；若按完整单位等待，至少55个。这里 $2\pi$ 是无阻尼固有周期对应的无量纲时间，并非任意驱动力的周期。若要保证瞬态相对稳态的误差，还应乘上二者初始幅度系数之比，再设阈值。

\begin{figure}[H]\centering
\begin{tikzpicture}[x=.017cm,y=.38cm]
\draw[bookaxis](0,0)--(440,0)node[right]{$\tau$};\draw[bookaxis](0,-8.7)--(0,1)node[above]{$\ln(A/A_0)$};
\draw[bookcurve](0,0)--(410,-8.2);\draw[bookguide](0,{-ln(1000)})--({ln(1000)/.02},{-ln(1000)})--({ln(1000)/.02},0);
\fill({ln(1000)/.02},{-ln(1000)})circle(1.8pt);\node[left]at(0,{-ln(1000)}){$-\ln1000$};\node[above]at({ln(1000)/.02},0){$345.388$};
\node[right]at(110,-1){斜率 $-0.02$};
\end{tikzpicture}
\caption{题4.3解图：在对数纵轴上，指数包络变为直线。与阈值 $-\ln1000$ 的交点给出最短等待时间。}\label{fig:hw-4-3}
\end{figure}

交点横坐标除以 $2\pi$ 就得到等待的周期单位数；直线斜率对应衰减率，不能用显示点数代替。
\end{solution}

\item \begin{minipage}[t]{\linewidth}
【基础】对二次阻力写出 $\beta$ 的无量纲推导。保持物体和介质不变，把初速加倍，$\beta$ 怎样改变？为什么不能认为有量纲射程必定变成四倍？

\end{minipage}
% 题设配图；与紧随其后的解答分别编号。
\ExpansionDiagram{fig:audit-m-q-112}{二次阻力与重力使用不同的速度尺度}{
\draw[->](-.4,0)--(5.8,0)node[right]{$ x $};\draw[->](0,-0.3)--(0,3.1)node[above]{$ y $};\fill(0,0)circle(2pt);\draw[->,thick](0,0)--(1.3,1)node[above]{$v_0$};\draw[->](.6,0)arc(0:37.6:.6);\node[right]at(.65,0.2){$\alpha$};\draw[dashed](0,0)..controls(2,2.5)and(4,2)..(5,0);\draw[->](6.2,2.8)--(6.2,1.9)node[right]{$g$};\fill(1.484375,1.335938)circle(1.5pt);\draw[->,thick](1.484375,1.335938)--(.734375,.91255)node[left]{$\vct F_d$};\node[below]at(3,-.65){$\vct F_d=-c|\vct v|\vct v$};
}

\begin{solution}
图\ref{fig:audit-m-q-112}中二次阻力随速率平方变化，改变初速也会改变无量纲参数。
图\ref{fig:ex-4-4}中的重力与反向阻力给出两种力尺度，二者之比就是无量纲参数。

二次阻力为 $\vct F_d=-c|\vct v|\vct v$，故 $[c]=[F]/[v]^2=\mathrm{kg/m}$。选择时间尺度 $t_0=v_0/g$、长度尺度 $\ell_0=v_0^2/g$，速度尺度为 $v_0$。牛顿方程除以重力尺度 $mg$ 后，阻力项系数为
\[
\beta=\frac{cv_0^2}{mg},
\]
分子、分母均为力，因而无量纲。保持 $c,m,g$ 不变而将 $v_0$ 加倍，$\beta$ 变为四倍。

\begin{figure}[H]\centering
\begin{tikzpicture}[scale=1]
\fill(0,0)circle(3pt);\draw[bookaxis,thick](0,0)--(1.8,1.25)node[right]{$\vct v$};
\draw[bookaxis,thick](0,0)--(-1.35,-.9375)node[left]{$\vct F_d=-c|\vct v|\vct v$};
\draw[bookaxis,thick](0,0)--(0,-1.8)node[below]{$m\vct g$};\node[above left]at(0,0){$m$};
\draw[bookaxis](2.6,-1.3)--(3.8,-1.3)node[right]{$x$};\draw[bookaxis](2.6,-1.3)--(2.6,0)node[above]{$y$};
\end{tikzpicture}
\caption{题4.4：二次阻力与速度反向；重力始终竖直向下。}\label{fig:ex-4-4}
\end{figure}

有量纲射程为 $R=(v_0^2/g)\mathcal R(\alpha,\beta)$。同角度加倍初速后，射程比为
\[
\frac{R(2v_0)}{R(v_0)}
=4\frac{\mathcal R(\alpha,4\beta)}{\mathcal R(\alpha,\beta)}.
\]
长度尺度增加四倍，同时相对阻力增强，无量纲轨迹也改变，所以不能只保留前面的因子4。

阻力箭头的大小随速率平方改变。加倍初速会改变整条轨迹上的相对阻力，因而无量纲射程也会变化。
\end{solution}

\item \begin{minipage}[t]{\linewidth}
【实验】复算单摆 $30^\circ,60^\circ,90^\circ$ 的周期。分别以事件积分、能量积分求值，列出两路线之差及二阶近似相对误差。 装置见图\ref{fig:ex-4-2}。
\end{minipage}
% 题设配图；与紧随其后的解答分别编号。
\ExpansionDiagram{fig:audit-m-q-113}{不同释放角对应不同的完整摆动周期}{
\fill(0,2.5)circle(2pt)node[above]{$O$};\draw[dashed](0,2.5)--(0,-.4);\draw[thick](0,2.5)--(1.2,.1);\fill(1.2,.1)circle(3pt)node[right]{$m$};\draw(0,1.7)arc(-90:-63.435:.8);\node[right]at(.25,1.5){$\theta_0$};\node[right]at(.8,1.2){$\ell$};\draw[->](2,1.9)--(2,.8)node[right]{$g$};\draw[dashed](0,2.5)--(2.5,2.5);\draw[dashed](0,2.5)--(2.165,1.25);
}

\begin{solution}
图\ref{fig:audit-m-q-113}中从释放点到最低点只是一周期的四分之一，事件法须采用正确的周期倍数。
装置仍见图\ref{fig:hw-4-5}。由释放点到最低点是四分之一周期，这一几何对称性把首次过零事件转为完整周期。

用无量纲时间 $\tau=t\sqrt{g/\ell}$，从 $\theta(0)=\theta_0,\theta'(0)=0$ 积分 $\theta''+\sin\theta=0$。记录第一次向下穿过 $\theta=0$ 的时刻 $\tau_q$，由运动对称性，周期比为 $T/T_0=2\tau_q/\pi$。必须记录下降事件，以免把初始点或反向过零当成所求事件。

\begin{figure}[H]\centering
\begin{tikzpicture}[scale=1]
\draw[thick](-.5,0)--(.5,0);\fill(0,0)circle(1.7pt);\draw[bookguide](0,0)--(0,-2.8);
\draw[thick](0,0)--(1.15,-2.36);\fill(1.15,-2.36)circle(3.4pt)node[right]{$m$};\node[right]at(.65,-1.15){$\ell$};
\draw[->](0,-.8)arc(-90:-64:.8);\node at(.4,-1){$\theta_0$};\draw[bookaxis](2,-.4)--(2,-1.2)node[below]{$\vct g$};
\node[below]at(1,-2.9){$\dot\theta(0)=0$};
\end{tikzpicture}
\caption{题4.5解图：三次运算仅改变释放角 $\theta_0=30^\circ,60^\circ,90^\circ$；由释放点到最低点的首次下降过零为四分之一周期。}\label{fig:hw-4-5}
\end{figure}

另由能量守恒，$\theta'^2=2(\cos\theta-\cos\theta_0)$。首次下行时$\theta'<0$，从释放点积到最低点并反转积分界，得
\[\tau_q=\int_{\theta_0}^0\frac{-\dd\theta}{\sqrt{2(\cos\theta-\cos\theta_0)}}
=\frac1{\sqrt2}\int_0^{\theta_0}\frac{\dd\theta}{\sqrt{\cos\theta-\cos\theta_0}}.\]
令$k=\sin(\theta_0/2)$、$\sin(\theta/2)=k\sin u$，则
\[\cos\theta-\cos\theta_0=2k^2\cos^2u,\qquad
\dd\theta=\frac{2k\cos u}{\sqrt{1-k^2\sin^2u}}\dd u.\]
两端变为$u=0,\pi/2$，分子与转折根式相约，得到光滑形式
\[
\frac T{T_0}=\frac2\pi\int_0^{\pi/2}
\frac{\dd u}{\sqrt{1-\sin^2(\theta_0/2)\sin^2u}}.
\]
运动方程采用五阶自适应 Runge--Kutta 积分，相对、绝对容差为 $10^{-11},10^{-13}$，最大无量纲步长0.02；能量积分采用自适应求积，绝对、相对容差均为 $10^{-13}$。复算结果为：
\begin{center}\small
\begin{tabular}{@{}c r r r@{}}\toprule
$\theta_0$&能量积分的 $T/T_0$&事件值减积分值&二阶式相对误差\\\midrule
$30^\circ$&1.017408798&$-1.1\times10^{-14}$&$-0.0269\%$\\
$60^\circ$&1.073182007&$+2.6\times10^{-14}$&$-0.4326\%$\\
$90^\circ$&1.180340599&$+3.2\times10^{-14}$&$-2.2136\%$\\\bottomrule
\end{tabular}
\end{center}
二阶式为 $1+\theta_0^2/16$，相对误差以能量积分值为分母。两条路线在所列精度下吻合，二阶近似的误差却随振幅明显增大。要改善这一误差，需要保留更高阶的振幅项或使用完整积分；只收紧程序容差没有作用。

周期积分的两个端点正对应图中的最低点和释放点；换元消去转折点根式后，才适合直接数值求积。
\end{solution}

\item \begin{minipage}[t]{\linewidth}
【实验】取 $\zeta=0.1$，求位移峰、速度峰与功率半高点。若频率网格步长为0.1，能否区分位移峰与速度峰？说明需要怎样细化。

\end{minipage}
% 题设配图；与紧随其后的解答分别编号。
\ExpansionDiagram{fig:audit-m-q-114}{位移响应、速度响应和功率峰分别测量}{
\draw[very thick](0,-.8)--(0,.9);\draw[exspring](0,.35)--(2.4,.35);\node[above]at(1.2,.45){$ k $};\draw[thick](2.4,-.65)rectangle(3.5,.65);\node at(2.95,0){$ m $};\draw(0,-.35)--(.8,-.35);\draw(.8,-.6)--(.8,-.1)--(1.7,-.1);\draw(.8,-.6)--(1.7,-.6);\draw(1.35,-.55)--(1.35,-.15);\draw(1.35,-.35)--(2.4,-.35);\node[below]at(1.2,-.7){$ b $};\draw[->,thick](3.5,0)--(5,0)node[right]{$ F_0\cos\omega t $};\draw[->](2.4,-1.2)--(4,-1.2)node[right]{$x$};\node[below]at(2.5,-1.8){$\zeta=0.1$};
}

\begin{solution}
图\ref{fig:audit-m-q-114}中同一振子有多种响应曲线，峰的位置必须按被测物理量分别求导。
图\ref{fig:ex-4-6}把峰顶局部放大，并另画功率半高交线；分别从相应响应式求这两类特征点。

取驱动幅值归一化，位移与速度幅值分别为
\[
A_x=\frac1{\sqrt{(1-\Omega^2)^2+4\zeta^2\Omega^2}},
\qquad A_v=\Omega A_x.
\]
记分母平方为$D=(1-\Omega^2)^2+4\zeta^2\Omega^2$。链式求导，
\[D'=2(1-\Omega^2)(-2\Omega)+8\zeta^2\Omega
=4\Omega(\Omega^2-1+2\zeta^2).\]
在$\zeta=0.1$时，$D'$于正根两侧由负变正，故位移峰为$\Omega_x=\sqrt{1-2\zeta^2}=\sqrt{0.98}=0.9899495$。速度幅值平方为$A_v^2=\Omega^2/D$，商法则给出
\[\frac{\dd A_v^2}{\dd\Omega}
=\frac{2\Omega D-\Omega^2D'}{D^2}
=\frac{2\Omega(1-\Omega^4)}{D^2}.\]
它在正频率$\Omega_v=1$两侧由正变负；平均功率正比$A_v^2$，功率峰也在$1$。

\begin{figure}[H]\centering
\begin{tikzpicture}
\begin{scope}[shift={(-34.78,-23.25)},x=37cm,y=25cm]
\draw[bookaxis](.94,.93)--(1.055,.93)node[right]{$\Omega$};\draw[bookaxis](.94,.93)--(.94,1.008);
\begin{scope}\clip(.94,.93)rectangle(1.055,1.008);
\draw[bookcurve,domain=.945:1.045,samples=120]plot(\x,{sqrt(.0396)/sqrt((1-(\x)^2)^2+.04*(\x)^2)});
\draw[dashed,domain=.945:1.045,samples=120]plot(\x,{.2*\x/sqrt((1-(\x)^2)^2+.04*(\x)^2)});
\end{scope}\node[left]at(.94,.93){$.93$};\node[left]at(.94,1){$1$};
\draw[bookguide](.9899495,.93)--(.9899495,1);\draw[bookguide](1,.93)--(1,1);
\node[below left]at(.9899495,.93){$.98995$};\node[below right]at(1,.93){$1$};\node[above]at(.994,1.016){各自按峰值归一化};
\end{scope}
\begin{scope}[shift={(.7,0)},x=7cm,y=1.9cm]
\draw[bookaxis](.7,0)--(1.34,0)node[right]{$\Omega$};\draw[bookaxis](.7,0)--(.7,1.15)node[above]{$P/P_{\max}$};
\draw[bookcurve,domain=.72:1.3,samples=130]plot(\x,{.04*(\x)^2/((1-(\x)^2)^2+.04*(\x)^2)});
\draw[bookguide](.7,.5)--(1.3,.5);\foreach\u in{.9049876,1.1049876}{\draw[bookguide](\u,0)--(\u,.5);\fill(\u,.5)circle(1.5pt);}
\node[below]at(.9049876,0){$.90499$};\node[below]at(1.1049876,0){$1.10499$};
\end{scope}
\end{tikzpicture}
\caption{题4.6解图：左图放大峰顶，实线为位移、虚线为速度，各自按最大值归一化；右图的水平半高线确定两半功率点。}\label{fig:ex-4-6}
\end{figure}

峰值处$A_v^2(1)=1/(4\zeta^2)$。半功率条件为$\Omega^2/D=1/(8\zeta^2)$，移项得
$(1-\Omega^2)^2=4\zeta^2\Omega^2$，即$|1-\Omega^2|=2\zeta\Omega$。低频一侧满足$\Omega^2+2\zeta\Omega-1=0$，高频一侧满足$\Omega^2-2\zeta\Omega-1=0$，分别取正根得到
\[
\Omega_\pm=\sqrt{1+\zeta^2}\pm\zeta
=0.9049876,\ 1.1049876.
\]
两种幅值峰只相差约0.01005，步长0.1不足以可靠区分，网格最大值往往都落在1。可先在0.95至1.05用0.001网格定位，再把局部步长减半，或对幅频式作一维极值求解。应比较细化前后峰位置，不能把显示点的横坐标直接当作准确极值。

\end{solution}

\item \begin{minipage}[t]{\linewidth}
【实验】取 $\epsilon=0.02$，改变 $\delta$，求第二物块包络最大值不超过初始位移一半时的条件。比较精确包络与首阶能量转移比例。

\end{minipage}
% 题设配图；与紧随其后的解答分别编号。
\ExpansionDiagram{fig:audit-m-q-115}{弱耦合双振子的给定初态与失谐}{
\draw[thick](0,-.5)--(0,.8);\draw[exspring](0,.15)--(1.4,.15);\node[above]at(.7,.4){$ k(1+\delta) $};\draw[thick](6.7,-.5)--(6.7,.8);\draw[exspring](5.3,.15)--(6.7,.15);\node[above]at(6,.4){$ k(1-\delta) $};\draw(1.4,-.3)rectangle(2.4,.6);\node at(1.9,.15){$ m $};\draw(4.3,-.3)rectangle(5.3,.6);\node at(4.8,.15){$ m $};\draw[exspring](2.4,.15)--(4.3,.15);\node[above]at(3.35,.4){$ \kappa $};\draw[->](1.4,-.8)--(2.7,-.8)node[right]{$x_1$};\draw[->](4.3,-.8)--(5.6,-.8)node[right]{$x_2$};\node[above]at(3.5,1.2){$x_1(0)=a,\quad x_2(0)=0$};
}

\begin{solution}
图\ref{fig:audit-m-q-115}中初始能量集中在第一物块，第二物块的幅值界由两个模态的叠加系数决定。
装置见图\ref{fig:ex-4-7}，由精确解得到的幅值界线见图\ref{fig:hw-4-7}。

记 $d=\sqrt{\epsilon^2+\delta^2}$。由正文的精确线性解，若初始位移为 $a$，
\[
\frac{x_2}{a}=\frac{\epsilon}{2d}(\cos\omega_-\tau-\cos\omega_+\tau)
=\frac\epsilon d\sin\frac{(\omega_++\omega_-)\tau}2
\sin\frac{(\omega_+-\omega_-)\tau}2.
\]
快振动的包络上界为 $\epsilon/d$，要求其不超过 $1/2$，得
\[
\frac{\epsilon}{\sqrt{\epsilon^2+\delta^2}}\le\frac12
\quad\Longleftrightarrow\quad
|\delta|\ge\sqrt3\epsilon=0.0346410.
\]
仍需处在系统稳定及所研究的弱失谐参数范围内。弱耦合首阶的最大平均能量转移比例为 $\eta\simeq\epsilon^2/d^2$，所以边界处为 $1/4$，超过边界时更小。前一包络公式在线性模型内精确，后一能量关系还用了周期平均与弱耦合近似；瞬时局部能量也含速度，不能只用任一时刻的位移平方替代。

\begin{figure}[H]\centering
\begin{tikzpicture}[scale=1.05]
\draw[thick](0,-.55)--(0,.55) (5.5,-.55)--(5.5,.55);
\draw[exspring](0,0)--(0.95,0);\node[above]at(0.475,.28){$k(1+\delta)$};
\draw[thick](0.95,-.4)rectangle(1.95,.4);\node at(1.45,0){$m$};
\draw[bookaxis](0.95,-.85)--(2.3,-.85)node[right]{$x_1$};
\draw[exspring](1.95,0)--(3.4000000000000004,0);\node[above]at(2.6750000000000003,.28){$\epsilon k$};
\draw[thick](3.4000000000000004,-.4)rectangle(4.4,.4);\node at(3.9000000000000004,0){$m$};
\draw[bookaxis](3.4000000000000004,-.85)--(4.75,-.85)node[right]{$x_2$};
\draw[exspring](4.4,0)--(5.5,0);\node[above]at(4.95,.28){$k(1-\delta)$};
\end{tikzpicture}
\caption{题4.7：固定 $\epsilon=0.02$，改变失谐量 $\delta$；两物块的质量相同。}\label{fig:ex-4-7}
\end{figure}

\begin{figure}[H]\centering
\begin{tikzpicture}[x=1.15cm,y=2.4cm]
\draw[bookaxis](0,0)--(5.2,0)node[right]{$|\delta|/\epsilon$};\draw[bookaxis](0,0)--(0,1.2)node[above]{包络上界 $/a$};
\draw[bookcurve,domain=0:4.8,samples=100]plot(\x,{1/sqrt(1+(\x)^2)});\draw[bookguide](0,.5)--(4.8,.5);
\draw[bookguide]({sqrt(3)},0)--({sqrt(3)},.5);\fill({sqrt(3)},.5)circle(1.7pt);\node[below]at({sqrt(3)},0){$\sqrt3$};\node[left]at(0,.5){$1/2$};
\draw[bookaxis]({sqrt(3)+.15},.25)--(4.7,.25)node[midway,below]{满足幅值要求};
\end{tikzpicture}
\caption{题4.7解图：精确包络上界与 $1/2$ 相交于 $|\delta|/\epsilon=\sqrt3$。首阶最大平均能量比例为纵坐标的平方。}\label{fig:hw-4-7}
\end{figure}

界线右侧的包络上界不超过初始位移的一半。在弱耦合首阶平均近似下，纵坐标的平方给出最大能量转移比例。
\end{solution}

\item \begin{minipage}[t]{\linewidth}
【实验】复算正文 $\beta=1$ 的最佳抛射角，检查无阻力极限及能量耗散。报告收紧精度前后的角度与射程，避免只提交动画截图。

\end{minipage}
% 题设配图；与紧随其后的解答分别编号。
\ExpansionDiagram{fig:audit-m-q-116}{含阻力抛体的发射方向与落地事件}{
\draw[->](-.4,0)--(5.8,0)node[right]{$ x $};\draw[->](0,-0.3)--(0,3.1)node[above]{$ y $};\fill(0,0)circle(2pt);\draw[->,thick](0,0)--(1.3,1)node[above]{$v_0$};\draw[->](.6,0)arc(0:37.6:.6);\node[right]at(.65,0.2){$\alpha$};\draw[dashed](0,0)..controls(2,2.5)and(4,2)..(5,0);\draw[->](6.2,2.8)--(6.2,1.9)node[right]{$g$};\node[above]at(3.5,3){$\beta=1$};
}

\begin{solution}
图\ref{fig:audit-m-q-116}中发射角改变轨迹和飞行时间，应先定位落地事件，再比较射程。
图\ref{fig:ex-4-8}由本题最佳角度的积分结果绘制，左图核对落地事件，右图核对能量收支。

设无量纲状态为 $(X,Y,U,V)$，二次阻力下满足
\begin{align*}
X'&=U,&Y'&=V,\\
U'&=-\beta\sqrt{U^2+V^2}\,U,&
V'&=-1-\beta\sqrt{U^2+V^2}\,V.
\end{align*}
初值为 $(0,0,\cos\alpha,\sin\alpha)$。记录 $Y$ 从正值向下穿越零的首次事件，其 $X$ 为射程；初始点的 $Y=0$ 不作为落地。先用角度扫描包围极大点，再在 $30^\circ$ 至 $50^\circ$ 内作有界一维极值计算。两次独立设置采用同样的角度终止容差 $10^{-7}$ 度：
\begin{center}\small
\begin{tabular}{@{}c c c r r@{}}\toprule
相对容差&绝对容差&最大步长&最佳角度（度）&最大射程\\\midrule
$10^{-6}$&$10^{-8}$&0.1&40.9775086&0.5936070541\\
$10^{-10}$&$10^{-12}$&0.01&40.9775090&0.5936070465\\\bottomrule
\end{tabular}
\end{center}
两射程相差约 $7.6\times10^{-9}$，支持报告 $\alpha_*\simeq40.9775^\circ$、$\mathcal R_{\max}\simeq0.593607$。角度极值附近曲线较平，末位角度比射程对求解设置更敏感。

\begin{figure}[H]\centering
\begin{tikzpicture}
\begin{scope}[x=7.4cm,y=8cm]
\draw[bookaxis](0,0)--(.7,0)node[right]{$X$};\draw[bookaxis](0,-.01)--(0,.25)node[above]{$Y$};
\draw[bookcurve]plot coordinates{(0.0000000,0.0000000) (0.0105423,0.0090586) (0.0209407,0.0177972) (0.0312004,0.0262228) (0.0413263,0.0343420) (0.0513230,0.0421613) (0.0611951,0.0496865) (0.0709467,0.0569234) (0.0805819,0.0638773) (0.0901044,0.0705534) (0.0995180,0.0769564) (0.1088262,0.0830911) (0.1180322,0.0889619) (0.1271392,0.0945729) (0.1361504,0.0999281) (0.1450684,0.1050313) (0.1538963,0.1098861) (0.1626364,0.1144959) (0.1712915,0.1188641) (0.1798638,0.1229937) (0.1883556,0.1268878) (0.1967693,0.1305492) (0.2051067,0.1339806) (0.2133699,0.1371846) (0.2215608,0.1401637) (0.2296812,0.1429202) (0.2377329,0.1454566) (0.2457173,0.1477748) (0.2536361,0.1498771) (0.2614908,0.1517655) (0.2692827,0.1534418) (0.2770131,0.1549079) (0.2846833,0.1561656) (0.2922944,0.1572167) (0.2998476,0.1580629) (0.3073440,0.1587056) (0.3147844,0.1591466) (0.3221697,0.1593873) (0.3295010,0.1594293) (0.3367788,0.1592739) (0.3440041,0.1589227) (0.3511774,0.1583769) (0.3582994,0.1576381) (0.3653707,0.1567075) (0.3723918,0.1555864) (0.3793632,0.1542763) (0.3862854,0.1527784) (0.3931588,0.1510940) (0.3999837,0.1492245) (0.4067605,0.1471711) (0.4134894,0.1449352) (0.4201707,0.1425181) (0.4268047,0.1399211) (0.4333915,0.1371456) (0.4399313,0.1341928) (0.4464242,0.1310642) (0.4528704,0.1277611) (0.4592699,0.1242849) (0.4656229,0.1206369) (0.4719293,0.1168187) (0.4781893,0.1128316) (0.4844028,0.1086772) (0.4905699,0.1043568) (0.4966905,0.0998719) (0.5027647,0.0952241) (0.5087924,0.0904149) (0.5147737,0.0854458) (0.5207084,0.0803185) (0.5265965,0.0750344) (0.5324380,0.0695952) (0.5382329,0.0640025) (0.5439811,0.0582579) (0.5496825,0.0523631) (0.5553372,0.0463198) (0.5609451,0.0401296) (0.5665061,0.0337942) (0.5720202,0.0273154) (0.5774873,0.0206948) (0.5829076,0.0139342) (0.5882808,0.0070354) (0.5936070,0.0000000)};\fill(.593607,0)circle(1.6pt);\node[below]at(.47,0){$\mathcal R=0.593607$};
\draw[bookaxis](0,0)--(.13,.11291);\node[above]at(.3,.25){$\alpha_*=40.9775^\circ$};
\end{scope}
\begin{scope}[shift={(6.4,0)},x=3.1cm,y=3.7cm]
\draw[bookaxis](0,0)--(1.3,0)node[right]{$\tau$};\draw[bookaxis](0,0)--(0,.61);
\draw[bookcurve]plot coordinates{(0.0000000,0.5000000) (0.0140615,0.4864168) (0.0281231,0.4737416) (0.0421846,0.4619057) (0.0562461,0.4508464) (0.0703077,0.4405058) (0.0843692,0.4308310) (0.0984307,0.4217730) (0.1124922,0.4132867) (0.1265538,0.4053307) (0.1406153,0.3978666) (0.1546768,0.3908589) (0.1687384,0.3842747) (0.1827999,0.3780838) (0.1968614,0.3722579) (0.2109230,0.3667710) (0.2249845,0.3615986) (0.2390460,0.3567184) (0.2531076,0.3521092) (0.2671691,0.3477517) (0.2812306,0.3436276) (0.2952921,0.3397199) (0.3093537,0.3360130) (0.3234152,0.3324918) (0.3374767,0.3291428) (0.3515383,0.3259529) (0.3655998,0.3229102) (0.3796613,0.3200031) (0.3937229,0.3172213) (0.4077844,0.3145546) (0.4218459,0.3119937) (0.4359075,0.3095299) (0.4499690,0.3071547) (0.4640305,0.3048605) (0.4780920,0.3026397) (0.4921536,0.3004855) (0.5062151,0.2983911) (0.5202766,0.2963503) (0.5343382,0.2943572) (0.5483997,0.2924061) (0.5624612,0.2904915) (0.5765228,0.2886083) (0.5905843,0.2867515) (0.6046458,0.2849164) (0.6187074,0.2830984) (0.6327689,0.2812931) (0.6468304,0.2794962) (0.6608919,0.2777037) (0.6749535,0.2759114) (0.6890150,0.2741156) (0.7030765,0.2723124) (0.7171381,0.2704982) (0.7311996,0.2686694) (0.7452611,0.2668224) (0.7593227,0.2649538) (0.7733842,0.2630601) (0.7874457,0.2611381) (0.8015073,0.2591844) (0.8155688,0.2571958) (0.8296303,0.2551692) (0.8436919,0.2531015) (0.8577534,0.2509894) (0.8718149,0.2488300) (0.8858764,0.2466203) (0.8999380,0.2443574) (0.9139995,0.2420383) (0.9280610,0.2396601) (0.9421226,0.2372202) (0.9561841,0.2347156) (0.9702456,0.2321436) (0.9843072,0.2295018) (0.9983687,0.2267873) (1.0124302,0.2239976) (1.0264918,0.2211304) (1.0405533,0.2181831) (1.0546148,0.2151533) (1.0686763,0.2120388) (1.0827379,0.2088373) (1.0967994,0.2055467) (1.1108609,0.2021648) (1.1249225,0.1986897)};\draw[dashed]plot coordinates{(0.0000000,0.0000000) (0.0140615,0.0135832) (0.0281231,0.0262584) (0.0421846,0.0380943) (0.0562461,0.0491536) (0.0703077,0.0594942) (0.0843692,0.0691690) (0.0984307,0.0782270) (0.1124922,0.0867133) (0.1265538,0.0946693) (0.1406153,0.1021334) (0.1546768,0.1091411) (0.1687384,0.1157253) (0.1827999,0.1219162) (0.1968614,0.1277421) (0.2109230,0.1332290) (0.2249845,0.1384014) (0.2390460,0.1432816) (0.2531076,0.1478908) (0.2671691,0.1522483) (0.2812306,0.1563724) (0.2952921,0.1602801) (0.3093537,0.1639870) (0.3234152,0.1675082) (0.3374767,0.1708572) (0.3515383,0.1740471) (0.3655998,0.1770898) (0.3796613,0.1799969) (0.3937229,0.1827787) (0.4077844,0.1854454) (0.4218459,0.1880063) (0.4359075,0.1904701) (0.4499690,0.1928453) (0.4640305,0.1951395) (0.4780920,0.1973603) (0.4921536,0.1995145) (0.5062151,0.2016089) (0.5202766,0.2036497) (0.5343382,0.2056428) (0.5483997,0.2075939) (0.5624612,0.2095085) (0.5765228,0.2113917) (0.5905843,0.2132485) (0.6046458,0.2150836) (0.6187074,0.2169016) (0.6327689,0.2187069) (0.6468304,0.2205038) (0.6608919,0.2222963) (0.6749535,0.2240886) (0.6890150,0.2258844) (0.7030765,0.2276876) (0.7171381,0.2295018) (0.7311996,0.2313306) (0.7452611,0.2331776) (0.7593227,0.2350462) (0.7733842,0.2369399) (0.7874457,0.2388619) (0.8015073,0.2408156) (0.8155688,0.2428042) (0.8296303,0.2448308) (0.8436919,0.2468985) (0.8577534,0.2490106) (0.8718149,0.2511700) (0.8858764,0.2533797) (0.8999380,0.2556426) (0.9139995,0.2579617) (0.9280610,0.2603399) (0.9421226,0.2627798) (0.9561841,0.2652844) (0.9702456,0.2678564) (0.9843072,0.2704982) (0.9983687,0.2732127) (1.0124302,0.2760024) (1.0264918,0.2788696) (1.0405533,0.2818169) (1.0546148,0.2848467) (1.0686763,0.2879612) (1.0827379,0.2911627) (1.0967994,0.2944533) (1.1108609,0.2978352) (1.1249225,0.3013103)};\draw[dotted](0,.5)--(1.15,.5);
\node[left]at(0,.5){$1/2$};\node[above]at(.62,.6){能量与累计耗散};\node at(.35,.43){$\mathcal E$};\node at(.85,.15){$D$};
\end{scope}
\end{tikzpicture}
\caption{题4.8解图：左图由 $\beta=1$、最佳角度的运动方程积分得到，落地点是下降过零事件；右图检验 $\mathcal E+D=1/2$，点线表示二者之和。}\label{fig:ex-4-8}
\end{figure}

无阻力时$U'=0,V'=-1$，积分并带初值，得$U=\cos\alpha,V=\sin\alpha-\tau$。再从零时刻积分位置，
\[X=\int_0^\tau\cos\alpha\dd s=\tau\cos\alpha,\qquad
Y=\int_0^\tau(\sin\alpha-s)\dd s
=[s\sin\alpha-s^2/2]_0^\tau.\]
非零落地根为$\tau_f=2\sin\alpha$，所以$\mathcal R=2\sin\alpha\cos\alpha=\sin2\alpha$。其导数$2\cos2\alpha$在$\pi/4$两侧由正变负，最大射程为$1$；相同事件程序复算为$1.000000000000002$。

再引入累计耗散$D'=\beta(U^2+V^2)^{3/2},D(0)=0$。记$w=\sqrt{U^2+V^2}$，机械能导数逐项为
\[\mathcal E'=UU'+VV'+Y'
=-\beta wU^2+V(-1-\beta wV)+V
=-\beta w^3.\]
因此
\[
\mathcal E=\tfrac12(U^2+V^2)+Y,\qquad
(\mathcal E+D)'=0.
\]
上述两组计算中，已存积分节点上的 $|\mathcal E+D-1/2|$ 最大值分别约 $1.6\times10^{-8}$、$1.1\times10^{-13}$。事件定位、极值搜索与耗散检查检验的是不同环节，应一起报告。

左图末点须满足 $Y=0,V<0$；最高点的条件 $V=0$ 不能用作落地事件。右图再核对整段运动的能量与累计耗散。
\end{solution}

\item \begin{minipage}[t]{\linewidth}
【实验】用配套数据分别拟合线性与含三次项模型。独立核对残差正交条件。若把所有力测量标准差加倍，拟合参数、归一化残差平方和及参数标准差怎样变化？

\end{minipage}
% 题设配图；与紧随其后的解答分别编号。
\ExpansionDiagram{fig:audit-m-q-117}{同一组伸长量与测力数据用于比较模型}{
\draw[->](-.4,0)--(6.8,0)node[right]{$x/(10^{-2}\,\mathrm m)$};
\draw[->](0,-.3)--(0,2.9)node[above]{$F/\mathrm N$};
\foreach \x/\force in {0/.20,1/1.01,2/1.83,3/2.70,4/3.65,5/4.71,6/5.87}{
 \draw(\x,{.4*(\force-.02)})--(\x,{.4*(\force+.02)});
 \fill(\x,{.4*\force})circle(1.5pt);
 \draw(\x,-.06)--(\x,.06)node[below,yshift=-1mm]{$\x$};
}
\foreach \f in {1,2,3,4,5,6}{\draw(-.06,{.4*\f})--(.06,{.4*\f});\node[left]at(-.08,{.4*\f}){$\f$};}
\node[right]at(3.6,.45){$\sigma_F=0.02\,\mathrm N$};
}

\begin{solution}
图\ref{fig:audit-m-q-117}的各误差棒对应相同的测量标准差，两模型采用同一组权重。伸长量和力的正方向见图\ref{fig:ex-4-9}，拟合残差画在图\ref{fig:orig-090}中。

配套数据为 $x=0,0.01,\ldots,0.06\,\mathrm m$，相应的 $F$ 为0.20、1.01、1.83、2.70、3.65、4.71、5.87 N，各标准差为0.02 N。线性模型写为 $F=b+kx$，三次模型写为 $F=b+kx+cx^3$，设计矩阵每行分别为 $(1,x_i)$ 和 $(1,x_i,x_i^3)$。

\begin{figure}[H]\centering
\begin{tikzpicture}[scale=.95]
\draw[thick](-1,2.5)--(1,2.5);\draw[exspring](0,2.5)--(0,.7);\draw[thick](-.45,-.1)rectangle(.45,.7);\node at(0,.3){物体};
\draw[bookguide](-.7,.7)--(1.7,.7);\draw[bookguide](-.7,1.6)--(1.7,1.6);\node[left]at(-.7,1.6){参考长度};
\draw(1.2,-.3)--(1.2,2.5);\foreach\h in{0,.2,...,2.4}{\draw(1.2,\h)--(1.35,\h);}
\draw[<->](1.7,.7)--(1.7,1.6)node[midway,right]{$x$};\draw[bookaxis,thick](0,-.1)--(0,-1)node[below]{$F$};
\end{tikzpicture}
\caption{题4.9：弹簧伸长与力的测量示意。参考长度、位移正方向和力的零点应在记录中说明。}\label{fig:ex-4-9}
\end{figure}

令 $W=\diag(1/\sigma_i^2)$，待求参数为 $\vct\theta$，残差为 $\vct r=\vct F-A\vct\theta$。目标函数及其对第 $j$ 个参数的偏导为
\[
S=\sum_i\frac{(F_i-\sum_\ell A_{i\ell}\theta_\ell)^2}{\sigma_i^2},
\qquad
\frac{\partial S}{\partial\theta_j}
=-2\sum_i\frac{A_{ij}(F_i-\sum_\ell A_{i\ell}\theta_\ell)}{\sigma_i^2}.
\]
令各偏导为零，得到 $A^{\mathsf T}W(A\hat{\vct\theta}-\vct F)=0$。本题各权重相同，公共因子可以约去，因此驻点条件也就是 $A^{\mathsf T}\vct r=0$。

为使各列的尺度接近，取 $u=x/(0.06\,\mathrm m)$，七个 $u_i$ 为 $i/6$，$i=0,\ldots,6$。把模型写为 $F=b+Ku$ 或 $F=b+Ku+Cu^3$。其中
\[
k=\frac{K}{0.06\,\mathrm m},\qquad
c=\frac{C}{(0.06\,\mathrm m)^3}.
\]
以下方程中的 $b,K,C$ 和 $F_i$ 均以牛顿为数值单位。线性拟合的正规方程为
\[
\begin{pmatrix}7&7/2\\7/2&91/36\end{pmatrix}
\begin{pmatrix}b\\K\end{pmatrix}
=\begin{pmatrix}1997/100\\4307/300\end{pmatrix}.
\]
第二行减去第一行的一半，得到
\[
\left(\frac{91}{36}-\frac74\right)K
=\frac{4307}{300}-\frac{1997}{200},
\quad\text{即}\quad
\frac79K=\frac{2623}{600}.
\]
所以 $K=7869/1400$，回代第一行，$b=1997/700-K/2=17/400$。

三次模型增加了 $u^3$ 这一列。逐项求和，正规方程为
\[
\begin{pmatrix}
7&7/2&49/24\\
7/2&91/36&2275/1296\\
49/24&2275/1296&67171/46656
\end{pmatrix}
\begin{pmatrix}b\\K\\C\end{pmatrix}
=\begin{pmatrix}1997/100\\4307/300\\108941/10800\end{pmatrix}.
\]
第二行减去第一行的一半，第三行减去第一行的 $7/24$ 倍，消去 $b$：
\[
\begin{aligned}
\frac79K+\frac{119}{162}C&=\frac{2623}{600},\\
\frac{119}{162}K+\frac{9847}{11664}C&=\frac{92071}{21600}.
\end{aligned}
\]
后一式再减去前一式的 $17/18$ 倍，得到
\[
\frac{65}{432}C=\frac{107}{800},\qquad C=\frac{2889}{3250}.
\]
依次回代，
\[
K=\frac97\left(\frac{2623}{600}-\frac{119}{162}C\right)
=\frac{435087}{91000},
\qquad
b=\frac{1997}{700}-\frac K2-\frac7{24}C=\frac{203}{1000}.
\]
这组小规模数据可用上述分数消元求解；数值程序则可对缩放后的设计矩阵直接使用 QR 或奇异值分解。换回原参数后，结果为
\begin{align*}
\text{线性：}\quad &b=0.04250\,\mathrm N,\quad k=93.67857\,\mathrm{N/m},\\
&\chi^2=297.63393,\quad \nu=7-2=5;\\
\text{三次：}\quad &b=0.20300\,\mathrm N,\quad k=79.68626\,\mathrm{N/m},\\
&c=4115.38462\,\mathrm{N/m^3},\\
&\chi^2=0.400275,\quad \nu=7-3=4.
\end{align*}
残差须用未四舍五入的参数回代。按 $x$ 递增顺序，两模型的残差（单位为牛顿，显示六位小数）分别为
\[
\begin{aligned}
\vct r_{\rm lin}^{\mathsf T}
&=(0.157500,\ 0.030714,\ {-0.086071},\ {-0.152857},\\
&\hspace{1.5em}{-0.139643},\ {-0.016429},\ 0.206786),\\
\vct r_{\rm cub}^{\mathsf T}
&=({-0.003000},\ 0.006022,\ 0.000352,\ {-0.004703},\\
&\hspace{1.5em}{-0.003835},\ 0.008264,\ {-0.003099}).
\end{aligned}
\]
用精确分数代入，线性模型满足 $\sum r_i=\sum u_ir_i=0$；三次模型还满足 $\sum u_i^3r_i=0$。这些等式分别对应设计矩阵的各列，说明残差正交于所有拟合方向。再计算 $\chi^2=\sum r_i^2/(0.02\,\mathrm N)^2$，便得到上述两个数值。图\ref{fig:orig-090}中线性模型两端为正、中间为负的残差也由这些数值给出。三次模型更贴合本组教学构造数据，但不能由此断言任意弹簧都遵循该模型。

全部标准差加倍，$W$ 整体变为 $W/4$。目标函数整体乘 $1/4$，极小位置不变，故参数估计不变、$\chi^2$ 除以4。把给定测量误差视为已知时，参数协方差为 $(A^{\mathsf T}WA)^{-1}$，随之乘4，参数标准差乘2；不能同时改用残差估计未知方差的约定。

线性模型的系统偏差无法仅靠调整零点消除；加入三次项后，各残差均落在给定的单次测量误差尺度内。
\end{solution}

\item \begin{minipage}[t]{\linewidth}
【综合】无阻力抛体从高度 $H>0$ 发射，初速为 $v_0$。自行推导最佳角度与最大射程，再请 AI 检验 $H=0$ 和 $H\gg v_0^2/g$ 两个极限。

\end{minipage}
% 题设配图；与紧随其后的解答分别编号。
\ExpansionDiagram{fig:audit-m-q-118}{从高于地面的位置发射，起落高度不同}{
\draw[->](-.4,0)--(5.8,0)node[right]{$ x $};\draw[->](0,-0.3)--(0,3.1)node[above]{$ y $};\draw[dashed](-.2,1)--(1.2,1);\node[left]at(0,1){$H$};\fill(0,1)circle(2pt);\draw[->,thick](0,1)--(1.3,2)node[above]{$v_0$};\draw[->](.6,1)arc(0:37.6:.6);\node[right]at(.65,1.2){$\alpha$};\draw[dashed](0,1)..controls(2,3.5)and(4,2)..(5,0);\draw[->](6.2,2.8)--(6.2,1.9)node[right]{$g$};
}

\begin{solution}
图\ref{fig:audit-m-q-118}中落地点在地面，发射点高$H$；可达边界应保留这项初始高度。
图\ref{fig:ex-4-10}把最佳轨迹与上包络画在同一坐标中，最远落地点正是包络与地面的交点。

以地面为 $y=0$、发射点为 $(0,H)$，令 $u=\tan\alpha$。消去飞行时间后，轨迹为
\[
y=H+xu-\frac{gx^2}{2v_0^2}(1+u^2).
\]
对给定正 $x$，将右侧视为 $u$ 的二次函数，其最大值在 $u=v_0^2/(gx)$。代回得所有轨迹的上包络
\[
y_{\rm env}(x)=H+\frac{v_0^2}{2g}-\frac{gx^2}{2v_0^2}.
\]
最大地面射程正是包络与地面相交处，因此
\[
R_{\max}=\frac{v_0}{g}\sqrt{v_0^2+2gH},\qquad
\tan\alpha_* =\frac{v_0^2}{gR_{\max}}
=\frac{v_0}{\sqrt{v_0^2+2gH}}.
\]
对应发射角在0至45度之间，飞行时间为正，因而符合正向发射条件。$H=0$ 时恢复 $R_{\max}=v_0^2/g$、$\alpha_*=45^\circ$。当 $H\gg v_0^2/g$，$R_{\max}\simeq v_0\sqrt{2H/g}$，$\alpha_*\simeq v_0/\sqrt{2gH}\to0$；飞行时间主要由高度决定，把初速度更多用于水平运动可得到较远射程。这两项也可直接代回轨迹式核对，无须依赖 AI 给出的结论。

\begin{figure}[H]\centering
\begin{tikzpicture}[x=3.2cm,y=1.9cm]
\draw[bookaxis](-.1,0)--(2.1,0)node[right]{$X=gx/v_0^2$};\draw[bookaxis](0,-.1)--(0,1.85)node[above]{$Y=gy/v_0^2$};
\draw[dashed,domain=0:1.732051,samples=80]plot(\x,{1.5-(\x)^2/2});
\draw[bookcurve,domain=0:1.732051,samples=80]plot(\x,{1+\x/sqrt(3)-2*(\x)^2/3});\fill(0,1)circle(1.6pt)node[left]{$gH/v_0^2$};
\fill(1.732051,0)circle(1.6pt)node[below]{$gR_{\max}/v_0^2$};\draw[bookaxis](0,1)--(.5,1.288675)node[right]{$v_0$};
\node[above right]at(.7,1.45){上包络};\node[below]at(.9,.8){最佳轨迹};
\end{tikzpicture}
\caption{题4.10解图：示取 $gH/v_0^2=1$，实线为最佳轨迹，虚线为可达区域上包络。二者在最远落地点相切，故由包络与地面的交点求最大射程。}\label{fig:ex-4-10}
\end{figure}

该点既是包络边界又可由实轨迹达到，所以判别式重根条件给出的确是最大可达射程。图示特例的最佳角为 $30^\circ$。
\end{solution}

\item \begin{minipage}[t]{\linewidth}
【综合】小珠沿第二章抛物线轨道运动，取 $a/R=0.1,0.3,0.6$。由能量积分求周期，并用运动方程积分核对；比较二阶周期修正，解释支持力为什么不做功。

\end{minipage}
% 题设配图；与紧随其后的解答分别编号。
\ExpansionDiagram{fig:audit-m-q-119}{受约束小珠的能量积分使用真实弧长}{
\draw[->](-.4,0)--(3,0)node[right]{$ x $};\draw[->](0,-0.4)--(0,2.4)node[above]{$ y $};\draw[thick,domain=-2:2,samples=80]plot(\x,{.35*\x*\x});\fill(1.7,{.35*1.7^2})circle(2pt)node[above right]{$A$};\draw[dashed](1.7,0)--(1.7,{.35*1.7^2});\node[below]at(1.7,0){$a$};\node[left]at(-1.7,1.8){$ y=x^2/(2R) $};\draw[->](2.6,1.7)--(2.6,.7)node[right]{$g$};
}

\begin{solution}
图\ref{fig:audit-m-q-119}中弧长微元与$\dd x$不同，周期积分的速度权重须保留斜率因子。
图\ref{fig:ex-4-11}的切线微元先给出速率与 $\dot x$ 的关系，再写能量积分。

轨道为 $y=x^2/(2R)$，故 $\dot y=(x/R)\dot x$。由初值 $x=a,\dot x=0$ 和能量守恒，
\[
\tfrac12m(1+x^2/R^2)\dot x^2+\frac{mgx^2}{2R}
=\frac{mga^2}{2R}.
\]
解出速度，再将从顶点至转折点的时间乘4：
\[
T=4\sqrt{\frac Rg}\int_0^a
\frac{\sqrt{1+x^2/R^2}}{\sqrt{a^2-x^2}}\dd x.
\]
令$x=a\sin u$，则$\dd x=a\cos u\dd u$、$\sqrt{a^2-x^2}=a\cos u$，两端变为$0,\pi/2$，端点根式因子被约去。以$T_0=2\pi\sqrt{R/g}$归一化，得到
\[
\frac T{T_0}=\frac2\pi\int_0^{\pi/2}
\sqrt{1+(a/R)^2\sin^2u}\dd u.
\]
运动方程沿用前文已由能量乘积求导和切向牛顿方程分别推导的式\eqref{eq:bead-ode}：
\[
(1+x^2/R^2)\ddot x+\frac{x\dot x^2}{R^2}+\frac{gx}{R}=0.
\]
以 $R$、$\sqrt{R/g}$ 为尺度，用与题4.5相同的积分容差和最大步长求第一次下降过零，周期为其时间的四倍。结果为
\begin{center}\small
\begin{tabular}{@{}c r r r@{}}\toprule
$a/R$&能量积分的 $T/T_0$&事件值减积分值&二阶近似\\\midrule
0.1&1.002495332&$-2.8\times10^{-13}$&1.0025\\
0.3&1.022133887&$-1.1\times10^{-13}$&1.0225\\
0.6&1.084689545&$-1.4\times10^{-13}$&1.0900\\\bottomrule
\end{tabular}
\end{center}
二阶式取$s=a/R$，用$\sqrt{1+s^2\sin^2u}=1+s^2\sin^2u/2+O(s^4)$。由于
\[\int_0^{\pi/2}\sin^2u\dd u
=\left[\frac u2-\frac{\sin2u}{4}\right]_0^{\pi/2}=\frac\pi4,\]
有$T/T_0=(2/\pi)[\pi/2+(s^2/2)(\pi/4)+O(s^4)]=1+s^2/4+O(s^4)$。下一项被积核为$-s^4\sin^4u/8$，其积分为负，因而表中二阶式略偏大，且振幅越大偏差越明显。

\begin{figure}[H]\centering
\begin{tikzpicture}[scale=1.15]
\draw[bookaxis](-2.1,0)--(2.6,0)node[right]{$x$};\draw[bookaxis](0,-.2)--(0,2.25)node[above]{$y$};
\draw[bookcurve,domain=-1.8:1.8,samples=80]plot(\x,{(\x)^2/2});\fill(1.7,1.445)circle(1.6pt)node[above]{静止释放};
\draw[bookguide](1.7,0)--(1.7,1.445);\node[below]at(1.7,0){$a$};
\coordinate(P)at(.8,.32);\fill(P)circle(1.6pt)node[above right,yshift=5pt]{$P$};
\draw[bookaxis,thick](P)--(.8,-.65)node[right]{$mg$};\draw[bookaxis,thick](P)--(.175,1.1)node[above left]{$\vct N$};
\draw[bookaxis](P)--(.175,-.18)node[below left]{$\vct v$};
\draw[bookguide](.8,.32)--(1.15,.32)--(1.15,.6);\draw(.8,.32)--(1.15,.6);\node[right]at(1.15,.47){$\dd y$};\node[below]at(.975,.32){$\dd x$};
\end{tikzpicture}
\caption{题4.11解图：中途运动点的速度沿切线、支持力沿法线。直角微元给出 $\dd s^2=(1+x^2/R^2)\dd x^2$，直接进入动能和周期积分。微元作放大示意。}\label{fig:ex-4-11}
\end{figure}

轨道静止且光滑，支持力沿法向、速度沿切向，逐时功率 $\vct N\cdot\vct v=0$。它可以改变速度方向而不改变机械能。题2.16的杆在转动，约束点的运动可使真实速度含沿支持力的分量，两种情形的做功条件不同。

图中支持力与中途速度垂直，所以没有支持力功；被积式的几何根号则来自同一条斜切线的路程。
\end{solution}

\item \begin{minipage}[t]{\linewidth}
【提高】以100 Hz采样70 Hz正弦，另以120 Hz采样同一信号，分别得到怎样的低频混叠？如果只知道信号为单频但不知道最高频率，两组记录是否必然唯一确定原频率？

\end{minipage}
% 题设配图；与紧随其后的解答分别编号。
\ExpansionDiagram{fig:audit-m-q-120}{70 Hz正弦的两种采样记录：实心点100 Hz，空心点120 Hz}{
\draw[->](-.4,0)--(6.8,0)node[right]{$ t/(10\,{\rm ms}) $};\draw[->](0,-1.8)--(0,1.8)node[above]{$ s $};\draw[thick,domain=0:6,samples=300]plot(\x,{sin(252*\x)});\foreach \x in {0,1,2,3,4,5,6}{\fill(\x,{sin(252*\x)})circle(1.8pt);}\foreach \x in {0,.833333,1.666667,2.5,3.333333,4.166667,5,5.833333}{\draw[fill=white](\x,{sin(252*\x)})circle(2.7pt);}\foreach \x in {1,2,3,4,5,6}{\draw(\x,-.1)--(\x,.1)node[below]{$\x$};}
}

\begin{solution}
图\ref{fig:audit-m-q-120}中离散记录只约束采样时刻的函数值，原频率仍可能相差共同的混叠周期。
图\ref{fig:ex-4-12}直接叠画连续信号、低频候选与采样点。下面在采样时刻验证这些点为何完全重合。

取信号 $s(t)=\sin(2\pi70t)$，两组采样都从 $t=0$ 开始。100 Hz采样在 $t_n=n/100$ 给出
\[
s_n=\sin(1.4\pi n)
=-\sin(0.6\pi n),
\]
即与负的30 Hz正弦完全相同。120 Hz采样则给出 $s_n=\sin(7\pi n/6)=-\sin(5\pi n/6)$，与负的50 Hz正弦相同。负号相当于相位相差 $\pi$。

\begin{figure}[H]\centering
\begin{tikzpicture}[x=.125cm,y=.9cm]
\foreach\yy/\rate/\alias in{0/100/30,-4.1/120/50}{\begin{scope}[shift={(0,\yy)}]
\draw[bookaxis](0,0)--(53,0)node[right]{$t/\mathrm{ms}$};\draw[bookaxis](0,-1.2)--(0,1.35);
\draw[bookcurve,domain=0:50,samples=220]plot(\x,{sin(25.2*\x)});\draw[dashed,domain=0:50,samples=220]plot(\x,{-sin(.36*\alias*\x)});
\pgfmathtruncatemacro{\last}{\rate/20}\foreach\n in{0,...,\last}{\pgfmathsetmacro{\xx}{1000*\n/\rate}\fill(\xx,{sin(25.2*\xx)})circle(1.6pt);}
\foreach\tick in{0,20,40}{\node[below]at(\tick,-1.15){$\tick$};}
\node[above]at(25,1.3){$f_s=\rate\,\mathrm{Hz}$，虚线为 $-\sin(2\pi\alias t)$};
\end{scope}}
\end{tikzpicture}
\caption{题4.12解图：实线均为70 Hz正弦；两组采样点分别同时落在负的30 Hz、50 Hz正弦上。连续曲线不同，样本相同。}\label{fig:ex-4-12}
\end{figure}

图~\ref{fig:ex-4-12}的采样点虽然重合，点间的连续曲线却不同，两组记录仍不能唯一确定原频率。600是100和120的最小公倍数，故频率 $70+600j$ Hz（$j$ 为正整数）在两组采样点多出的相位均为 $2\pi$ 的整数倍，也会产生相同记录。670 Hz就是另一个候选；要排除这类候选，还须给出频带等额外信息。
\end{solution}

\end{enumerate}

\chapter*{综合习题}\addcontentsline{toc}{chapter}{综合习题}
\markboth{综合习题}{综合习题}

\noindent\begin{minipage}{\textwidth}
\vspace{12pt}
\section*{一、带电约束系统的稳定性}
两个固定点电荷 $Q>0$ 位于 $(\pm a,0)$。质量 $m$、电荷 $q>0$ 的小珠只能沿 $y$ 轴运动，另受势能为 $ky^2/2$ 的弹簧作用，忽略重力。求总势能，在 $|y|\ll a$ 时展开至四阶；讨论原点稳定条件，并求存在时的非零小振动频率。在临界刚度处，说明原点是否稳定以及周期如何随小振幅变化。写 $k_e=1/(4\pi\varepsilon_0)$。

\end{minipage}\par
% 题设配图；与紧随其后的解答分别编号。
\ExpansionDiagram{fig:audit-m-q-121}{固定正电荷和沿对称轴运动的正小珠}{
\draw[->](-.4,0)--(3.5,0)node[right]{$ x $};\draw[->](0,-0.4)--(0,3)node[above]{$ y $};\fill(-1.5,0)circle(3pt)node[below]{$Q$};\fill(1.5,0)circle(3pt)node[below]{$Q$};\fill(0,1.8)circle(3pt)node[right]{$q,m$};\draw[dashed](-1.5,0)--(0,1.8)--(1.5,0);\node[below]at(-1.5,-.45){$-a$};\node[below]at(1.5,-.45){$a$};\node[left]at(0,1){$y$};\draw[->,thick](.4,1.8)--(.4,.8)node[right]{$F_s=-ky$};
}

\begin{solution}
图\ref{fig:audit-m-q-121}中小珠到两电荷的距离相同，电势能可先按标量相加，再与弹簧势能共同展开。
图\ref{fig:ex-5-1}中的两条斜线是到固定电荷的距离，沿导轨的合电场力与弹簧力决定局部稳定性。

小珠到两个固定电荷的距离相等，均为 $\sqrt{a^2+y^2}$，故总势能为
\[
U(y)=\frac12ky^2+\frac{2k_eQq}{\sqrt{a^2+y^2}}.
\]
两个电势都取正号，因为三个电荷同号。先提取距离因子 $a$，再用 $(1+u)^{-1/2}=1-u/2+3u^2/8+O(u^3)$，得到
\begin{align*}
U(y)&=\frac{2k_eQq}{a}
+\frac12\left(k-\frac{2k_eQq}{a^3}\right)y^2\\
&\quad+\frac{3k_eQq}{4a^5}y^4+O(y^6).
\end{align*}
记临界刚度 $k_c=2k_eQq/a^3$。若 $k>k_c$，二次项为正，原点稳定，线性方程为 $m\ddot y+(k-k_c)y=0$，所以 $\omega=\sqrt{(k-k_c)/m}$。若 $k<k_c$，二次项为负，原点不稳定。电荷的排斥作用使小珠离开原点后的力更偏向远离原点，因而需要足够大的弹簧刚度。

\begin{figure}[H]\centering
\begin{tikzpicture}[scale=.95]
\draw[bookaxis](-3,0)--(3.3,0)node[right]{$x$};\draw[bookaxis](0,-.7)--(0,2.8)node[above]{$y$};
\draw[line width=2pt,black!25](0,-.5)--(0,2.35);\fill(-2,0)circle(3pt)node[below]{$Q,\ (-a,0)$};\fill(2,0)circle(3pt)node[below]{$Q,\ (a,0)$};
\fill(0,1.6)circle(3pt)node[above right]{$m,q$};\draw[bookguide](-2,0)--(0,1.6)--(2,0);\node[below left]at(0,0){$O$};
\node[right]at(.5,2.35){弹性势能 $ky^2/2$};\node[left]at(-.1,2.5){导轨};
\end{tikzpicture}
\caption{综合习题一解图：带电小珠沿 $y$ 轴导轨运动，两个固定电荷位于 $x$ 轴上。}\label{fig:ex-5-1}
\end{figure}

若$k=k_c$，二次项恰好抵消。写$C=3k_eQq/(4a^5)>0$，局部势能首项为$Cy^4$，故原点仍为严格极小点，但线性频率为零。对小振幅$A\ll a$，忽略六阶及更高项，由能量守恒有$m\dot y^2/2=C(A^4-y^4)$。取$0$到$A$的正速度支计算四分之一周期，再令$y=Au$、$\dd y=A\dd u$、$\sqrt{A^4-y^4}=A^2\sqrt{1-u^4}$，两端变为$0,1$。周期因此为
\[
T\simeq4\sqrt{\frac m{2C}}\int_0^A\frac{\dd y}{\sqrt{A^4-y^4}}
=\frac4A\sqrt{\frac m{2C}}\int_0^1\frac{\dd u}{\sqrt{1-u^4}}.
\]
无量纲积分是常数，所以 $T\propto A^{-1}$，振幅越小，周期越长。回复力的首项为三次，振幅趋零时频率也趋零。

两电荷的横向力互相抵消、竖直力同向相加；这正是势能中因子2以及电荷作用削弱原点稳定性的来源。
\end{solution}

\noindent\begin{minipage}{\textwidth}
\vspace{12pt}
\section*{二、两种驱动频率的能量输入}
质量 $m$ 的物体满足 $m\ddot x+b\dot x+kx=F\cos\omega t+F\cos2\omega t$，$b>0$，且 $k/m=4\omega^2$。求纯稳态位移及在 $2\pi/\omega$ 上的平均耗散功率。当 $b\ll m\omega$ 时，比较两个位移谐波的幅值。解释为什么不能把两驱动力先相加成一个固定相量，再套单频公式。

\end{minipage}\par
% 题设配图；与紧随其后的解答分别编号。
\ExpansionDiagram{fig:audit-m-q-122}{两种驱动频率同时作用于一个位移自由度}{
\draw[very thick](0,-.8)--(0,.9);\draw[exspring](0,.35)--(2.4,.35);\node[above]at(1.2,.45){$ k $};\draw[thick](2.4,-.65)rectangle(3.5,.65);\node at(2.95,0){$ m $};\draw(0,-.35)--(.8,-.35);\draw(.8,-.6)--(.8,-.1)--(1.7,-.1);\draw(.8,-.6)--(1.7,-.6);\draw(1.35,-.55)--(1.35,-.15);\draw(1.35,-.35)--(2.4,-.35);\node[below]at(1.2,-.7){$ b $};\draw[->,thick](3.5,0)--(5,0)node[right]{$ F\cos\omega t+F\cos2\omega t $};\draw[->](2.4,-1.2)--(4,-1.2)node[right]{$x$};
}

\begin{solution}
图\ref{fig:audit-m-q-122}中两个力频率不同，先分别求单频响应，再以真实时间函数相加。
图\ref{fig:ex-5-2}两驱动力作用于同一位移，但频率不同。每条驱动箭头分别建立一项响应后再相加。

线性系统对不同频率的响应可以相加。每个驱动力先单独设同频特解，复刚度为 $D_n=k-mn^2\omega^2+\ii bn\omega$。用题设 $k=4m\omega^2$，得到
\[
X_1=\frac F{3m\omega^2+\ii b\omega},\qquad
X_2=\frac F{2\ii b\omega}.
\]
于是纯稳态位移为
\[
x_s=\Rea\big(X_1\ee^{\ii\omega t}+X_2\ee^{2\ii\omega t}\big).
\]
将第一项乘分母共轭、第二项用$1/\ii=-\ii$展开，得到实函数
\[
x_s(t)=\frac{F[3m\omega^2\cos\omega t+b\omega\sin\omega t]}
{9m^2\omega^4+b^2\omega^2}
+\frac{F}{2b\omega}\sin2\omega t.
\]
两个频率的位移幅值分别为
\[
A_1=\frac F{\sqrt{9m^2\omega^4+b^2\omega^2}},\qquad
A_2=\frac F{2b\omega}.
\]
瞬时耗散功率为$b(\dot x_1+\dot x_2)^2$。记各位移为$x_j=A_j\cos(j\omega t-\delta_j)$，则
$\dot x_j=-j\omega A_j\sin(j\omega t-\delta_j)$。交叉项用积化和差展开：
\[\dot x_1\dot x_2=\omega^2A_1A_2
\bigl[\cos(\omega t+\delta_1-\delta_2)
-\cos(3\omega t-\delta_1-\delta_2)\bigr].\]
在共同周期$T=2\pi/\omega$上，对任意相位$\phi$及非零整数$n$，
\[\frac1T\int_0^T\cos(n\omega t+\phi)\dd t
=\frac1T\left[\frac{\sin(n\omega t+\phi)}{n\omega}\right]_0^T=0.\]
所以交叉项平均为零。每项速度平方所需的积分为
\[\frac1T\int_0^T\sin^2(j\omega t-\delta_j)\dd t
=\frac1T\left[\frac t2-\frac{\sin(2j\omega t-2\delta_j)}{4j\omega}\right]_0^T
=\frac12.\]
因此
\begin{align*}
\overline P&=\frac b2\big(\omega^2A_1^2+4\omega^2A_2^2\big)\\
&=\frac b2\left[
\frac{\omega^2F^2}{9m^2\omega^4+b^2\omega^2}
+\frac{4\omega^2F^2}{4b^2\omega^2}\right].
\end{align*}
当 $b\ll m\omega$ 时，$A_1\simeq F/(3m\omega^2)$，而 $A_2=F/(2b\omega)$，故 $A_2/A_1\simeq3m\omega/(2b)\gg1$。第二谐波恰在系统的固有频率，较易获得大的响应。

\begin{figure}[H]\centering
\begin{tikzpicture}[scale=.83]
\draw[thick](0,-.8)--(0,.9);
\draw[exspring](0,.35)--(3,.35);\node[above]at(1.5,.55){$k$};
\draw(0,-.35)--(1.15,-.35);\draw(1.15,-.62)--(1.15,-.08)--(2.05,-.08);
\draw(1.15,-.62)--(2.05,-.62);\draw(1.7,-.56)--(1.7,-.14);\draw(1.7,-.35)--(3,-.35);
\node[below]at(1.55,-.65){$b$};\draw[thick](3,-.65)rectangle(4,.65);\node at(3.5,0){$m$};
\draw[bookaxis](3,-1.1)--(4.5,-1.1)node[right]{$x$};

\draw[bookaxis,thick](4,.35)--(5.8,.35)node[right]{$F\cos\omega t$};
\draw[bookaxis,thick](4,-.35)--(5.8,-.35)node[right]{$F\cos2\omega t$};
\end{tikzpicture}
\caption{综合习题二解图：两个不同频率的力同时作用于同一物块。箭头标明统一的力正方向。}\label{fig:ex-5-2}
\end{figure}

两个力的时间因子分别为 $\ee^{\ii\omega t}$ 和 $\ee^{2\ii\omega t}$，相对相位随时间变化；不存在一个恒定复幅值乘同一个时间因子来表示其和。因此需分别求响应后叠加。$b>0$ 保证自由分量最终衰减，这里计算的是纯稳态。

阻尼器耗散的是总速度平方；在共同周期平均后，两种频率的交叉项消失，所以平均功率才可分项相加。
\end{solution}

\section*{三、由能量表达式寻找独立运动}
某两自由度系统使用长度坐标 $x,y$，动能和势能为
\[
T=\frac m2(2\dot x^2+2\dot x\dot y+2\dot y^2),\qquad U=\frac k2(x^2+y^2).
\]
求质量与刚度矩阵、两个正常模及频率。初值 $x=a,y=0,\dot x=\dot y=0$，求完整运动与总能量。说明普通坐标中的动能交叉项为何不意味着系统不稳定。
% 题设配图；与紧随其后的解答分别编号。
\ExpansionDiagram{fig:audit-m-q-123}{位置初态与速度空间中的等动能曲线}{
\begin{scope}[xshift=-2cm]\draw[->](-.3,0)--(2,0)node[right]{$x$};\draw[->](0,-.3)--(0,1.8)node[above]{$y$};\fill(1.3,0)circle(2pt)node[above]{$(a,0)$};\end{scope}\begin{scope}[xshift=2cm]\draw[->](-1.7,0)--(1.8,0)node[right]{$\dot x$};\draw[->](0,-1.6)--(0,1.8)node[above]{$\dot y$};\draw[rotate=45](0,0)ellipse(.72 and 1.25);\node[below]at(0,-1.9){等动能曲线};\end{scope}
}

\begin{solution}
图\ref{fig:audit-m-q-123}分别画出位置初态与速度二次型，动能交叉项应计入质量矩阵。图\ref{fig:mixed-3}的两个转轴对应同相、反相运动，可先沿这两个方向求质量和刚度。

比较 $T=\dot{\vct x}^{\mathsf T}M\dot{\vct x}/2$、$U=\vct x^{\mathsf T}K\vct x/2$，其中 $\vct x=(x,y)^{\mathsf T}$，得到
\[
M=m\begin{pmatrix}2&1\\1&2\end{pmatrix},\qquad K=kI.
\]
动能交叉项包含对称矩阵的两个非对角项，因此 $M_{12}=M_{21}=m$。频率由广义特征方程确定：
\[
\det(K-\omega^2M)=(k-2m\omega^2)^2-m^2\omega^4
=(k-3m\omega^2)(k-m\omega^2)=0.
\]
同相、反相模态分别为 $(1,1)$、$(1,-1)$，角频率为 $\omega_1=\sqrt{k/(3m)}$、$\omega_2=\sqrt{k/m}$。取正常坐标 $q_1=(x+y)/\sqrt2,q_2=(x-y)/\sqrt2$，则
\[
T=\tfrac12(3m\dot q_1^2+m\dot q_2^2),\qquad
U=\tfrac12k(q_1^2+q_2^2).
\]
正常坐标的运动方程分别为
\[
3m\ddot q_1+kq_1=0,\qquad m\ddot q_2+kq_2=0.
\]
代入 $q_j\propto\ee^{\lambda t}$，得到 $\lambda^2+\omega_j^2=0$，故
\[
q_j=C_j\cos\omega_jt+D_j\sin\omega_jt,
\qquad q_j(0)=C_j,\quad \dot q_j(0)=\omega_jD_j.
\]
由 $x=a,y=0$、两速度为零，得到 $C_1=C_2=a/\sqrt2$、$D_1=D_2=0$。逆变换后，完整运动为
\begin{align*}
x(t)&=\tfrac a2(\cos\omega_1t+\cos\omega_2t),\\
y(t)&=\tfrac a2(\cos\omega_1t-\cos\omega_2t).
\end{align*}
每个模的初始能量均为 $ka^2/4$，相加得总能量 $ka^2/2$，与原坐标初始势能相同。对任意非零速度，变换后的动能均为正；势能也正定，所以原点稳定。交叉项来自坐标表示，不能脱离完整二次型单独判断正负。

\begin{figure}[H]\centering
\begin{tikzpicture}[scale=1]
\draw[bookaxis](-1.7,0)--(2,0)node[right]{$x$};\draw[bookaxis](0,-1.7)--(0,2)node[above]{$y$};
\draw[bookaxis,thick](-1.4,-1.4)--(1.5,1.5)node[above]{$q_1$};\draw[bookaxis,thick](-1.4,1.4)--(1.5,-1.5)node[below]{$q_2$};
\draw[bookguide](1.2,0)--(.6,.6);\draw[bookguide](1.2,0)--(.6,-.6);\fill(1.2,0)circle(1.7pt)node[above right,yshift=2pt]{$(a,0)$};
\node[left,align=center]at(-1.6,-1){同相\\惯性 $3m$};\node[left,align=center]at(-1.6,1.1){反相\\惯性 $m$};
\end{tikzpicture}
\caption{综合习题三解图：转过 $45^\circ$ 的两坐标分别对应同相和反相方向。初态在两方向上的投影相等，而惯性系数不同。}\label{fig:mixed-3}
\end{figure}

\end{solution}

\appendix
\begingroup\small\setlist[enumerate]{itemsep=2pt,topsep=3pt}
\chapter{习题提示}
\section{第一章}
\begin{enumerate}[label=1.\arabic*]
\item 用极形式写出四个辐角，再用多项式系数检查和与积。
\item 先画两个等长相量；最后分别代入 $t=0$ 核对函数与导数。
\item 使用倍角公式，注意取实部与相乘的先后不能交换。
\item 先写复刚度 $k-m\omega^2+\ii b\omega$；功率用速度幅值。
\item 对给定特解求两次导数；检查它的初值是否已经符合要求。
\item 先求电荷相量，电容电压为电荷除以电容。
\item 先检查柯西--黎曼条件，再验证其充分条件。
\item 分别列出围道内的简单极点和实积分所需的二阶极点。
\item 平方后整理 $x^2+y^2$ 的系数，注意该系数为零的情形。
\item 代入 $z=1/w$ 后乘以 $|w|$，再平方化简。
\item 同向峰相隔一个阻尼周期；先写对数衰减率。
\item 共同量是电压，应相加导纳。
\item 利用偶性，仅求余弦系数；在 $t=0$ 比较级数。
\item 参数不相等时先作部分分式，再对相等情况取极限。
\item 弹簧和阻尼器响应的是相对位移、相对速度。
\item 先说明 $\Omega-1$ 是否与 $\zeta$ 同阶。
\item 滞后 $\pi/2$ 确定复刚度实部为零；另一相位确定虚部。
\end{enumerate}
\section{第二章}
\begin{enumerate}[label=2.\arabic*]
\item 叉积方向用右手规则；混合积可作行列式。
\item 垂足沿法线寻找，直线与平面夹角由方向向量在法线上的投影求得。
\item 轨迹式中以 $\tan\alpha$ 为未知量，判断二次方程实根。
\item 极坐标的切向加速度含 $2\dot r\dot\theta$。
\item 先用纬度圆带积分，再将球体分成同心薄壳。
\item 由两端距离先求长半轴与偏心率。
\item 先看两端极限，再求导数符号。
\item 对电势求负梯度，远场先展开根号。
\item 由参数切向求单位切向，主法向指向曲线凹侧。
\item 参数积分会出现正弦、余弦平方和；格林定理路线只需面积。
\item 方向导数使用单位方向，不用未经归一化的向量。
\item 旋度为常向量，通量只取有向投影面积。
\item 球壳内边界的外法向朝向球心。
\item 展开 $z^3$，再用梯度点积判断等值线正交。
\item 势能首项为四次，先用能量估计速度尺度，再对周期作变量代换。
\item 惯性系极坐标中仍可直接写 $a_r,a_\theta$；支持力可能做功。
\item 先保留两电荷的相对位置，展开后再相减。
\end{enumerate}
\section{第三章}
\begin{enumerate}[label=3.\arabic*]
\item 行列式只依赖系数矩阵，不依赖右端。
\item 第二式减去第一式的两倍，检查剩下的是恒等式还是矛盾。
\item 第二向量先减去在第一方向上的投影。
\item 把初始向量拆成同相、反相两个分量。
\item 先取质心与相对坐标；零频模需包含匀速运动。
\item 固定端条件选择正弦形状，再归一化。
\item 转到同相与反相方向，高阶项只需在零方向上重点检查。
\item 到转轴的距离为杆上位置乘以夹角正弦。
\item 换基后长度由 $G=P^{\mathsf T}P$ 计算。
\item 先写弹簧连接点的位移，再写动能与弹簧势能。
\item 把矩阵拆成 $-I+N$，计算 $N^2$。
\item 联立两个相量方程，先观察哪个分子可能为零。
\item 对阻尼矩阵作与质量、刚度相同的模态变换。
\item 每个质点贡献 $m(r^2I-\vct r\vct r^{\mathsf T})$。
\item 比较失谐与耦合的比值；相位误差等于频差误差乘时间。
\item 先找全部驻点，再分别计算二阶导数；二阶为零时回到势能本身。
\item 若主物体稳态静止，辅助物体的弹簧力应抵消外驱动力。
\end{enumerate}
\section{第四章}
\begin{enumerate}[label=4.\arabic*]
\item 直接求两个函数的导数，检验原方程。
\item 先求周期相对修正，再乘累计周期数估计相位。
\item 对指数衰减条件取对数。
\item 用长度 $v_0^2/g$ 和时间 $v_0/g$ 逐项换元。
\item 一个方法求时间事件，另一个方法求能量积分，保留独立性。
\item 先用解析公式估计两个峰的间距，再决定频率网格。
\item 把包络上界写成 $1/\sqrt{1+(\delta/\epsilon)^2}$。
\item 先找落地，再找角度最大值；两种定位的精度都要检查。
\item 整体乘一个相同权重不会改变最小值位置。
\item 把落地条件改写为关于 $\tan\alpha$ 的二次方程。
\item 用 $x=a\sin u$ 消去转折点的积分奇性。
\item 寻找两采样率的共同整数倍，构造具有相同样本的另一频率。
\end{enumerate}

\endgroup
\backmatter
\chapter*{符号与计算约定}\addcontentsline{toc}{chapter}{符号与计算约定}
\begin{tabularx}{\linewidth}{@{}p{25mm}X@{}}\toprule
符号 & 含义\\\midrule
$\ii$ & 虚数单位，$\ii^2=-1$；程序中写作 \texttt{1j}。\\
$\bar z,|z|$ & 复共轭与复数模。\\
$\omega,f$ & 角频率与频率，$\omega=2\pi f$。\\
$\omega_0,\gamma$ & 无阻尼固有角频率与振幅衰减率，$\gamma=b/(2m)$。\\
$\zeta,\Omega$ & 阻尼比与驱动角频率比。\\
$\vct e_r,\vct e_\theta$ & 平面极坐标的径向和切向单位基矢。\\
$\theta,\phi$ & 球坐标极角与方位角；$\theta$ 从正 $z$ 轴量起。\\
$\mu,h$ & 引力参数与单位质量角动量；数值算法中 $h$ 另表示步长，按上下文区分。\\
$M,K$ & 质量矩阵与刚度矩阵；标量 $M$ 另可表示总质量。\\
$Q,P$ & 正交换基矩阵与一般换基矩阵。\\
$I$ & 单位矩阵或惯量张量，按公式上下文区分；$I_3$ 明指三维单位矩阵。\\
$\tau,q,p$ & 计算实验中的无量纲时间、位移与其导数。\\\bottomrule
\end{tabularx}

相量均采用 $\ee^{\ii\omega t}$，幅值为峰值。正弦信号有效值为峰值除以 $\sqrt2$。主辐角取 $(-\pi,\pi]$，零相量不赋予相位。矩阵特征向量写作列，整体变号代表同一方向。所有近似和数值误差结论都应结合正文给出的模型、初值、计算区间与光滑性条件理解。
\end{document}
